\quellseite{6}
\begin{titelseite}
{\Large Der Evangelische Geistliche.}\\[1.5em]
{\Large Erstes Bändchen.}
\end{titelseite}

\quellseite{8}
\begin{titelseite}
{\LARGE Der evangelische Geistliche.}\\[1.5em]
Dem nun folgenden Geschlechte evangelischer Geistlichen dargebracht\\[1.5em]
von Wilhelm Löhe,\\
weil. luth. Pfarrer zu Neuendettelsau in Franken.\\[1.5em]
{\Large Erstes Bändchen.}\\[1.5em]
Vierte, mit der dritten gleichlautende Auflage.\\[1.5em]
Gütersloh.\\
Verlag von C. Bertelsmann.\\
1872.
\end{titelseite}

\quellseite{10}
\chapter{Vorwort}

Die nachfolgenden Blätter sind eine Umarbeitung der in den Jahren 1847 und 1848 in der Erlanger Zeitschrift für Protestantismus und Kirche mitgetheilten „Beiträge zur Pastoraltheologie.“ Der Anhang für Pfarrersfrauen wurde zuerst in Brandt’s „Christlicher Jungfrau im Brautstande“ (Frankf. bei Schmerber, etwa 1838) abgedruckt. Diesem auf mehrfaches Verlangen veranstalteten neuen Druck ist Valentin Andreae’s vortreffliches Gedicht „Vom guten Leben eines rechtschaffenen Dieners Gottes“ vorangestellt, was gewiß keiner Entschuldigung bedarf.

Dies ganze Bändchen beschäftigt sich vorzugsweise mit den persönlichen Verhältnissen eines Pfarrers. Ist es dem Publicum, für welches der Verfasser schrieb, angenehm, so folgt einmal ein zweites Bändchen, in welchem Ansichten und Erfahrungen in Betreff der verschiedenen Arbeitsgebiete des
\quellseite{11}
evangelischen Pfarrers (Predigt, Katechese, Liturgie, besondere Seelsorge) mitgetheilt werden.

Der Standpunkt, von welchem aus geschrieben wurde, ist hauptsächlich der eines Landpfarrers in der Landeskirche.

Die confessionellen Fragen der Zeit treten in diesem Bändchen bei Seite, wie billig; weshalb auch der Titel so gewählt ist, wie außen zu sehen. Ein Titel wie „Der lutherische Geistliche“ hätte, abgesehen daß man beim Versuch, ihn zu gebrauchen, ganz besonders spürt, wie wenig ein Menschenname zur Bezeichnung des heiligen Amtes paßt, verheißen, was nicht in der Absicht lag zu geben. Daß der Verfasser deshalb seine confessionellen Ueberzeugungen nicht verleugnet, versteht sich von selbst.

Möge dies kleine Buch denen dienen, welchen es vermeint ist, und mögen sie es so herzlich auf- und annehmen, als es ihnen dargebracht wird! Gott segne den Leser und den Schreiber! Amen.

\begin{unterschrift}
Am Tage Marien Heimsuchung,\\
2. Juli 1852.
\end{unterschrift}

\quellseite{14}
\chapter[Das gute Leben eines rechtschaffenen Dieners Gottes]{Das gute Leben eines rechtschaffenen Dieners Gottes\\[1em]\normalsize Von\\ Johann Valentin Andreä,\\[1em] Würtemb. Hofprediger, Abt, Consistorialrath, Generalsuperintendent u.\,s.\,w.\\[1em] Geboren 1586, gestorben 1654.}

\quellseite{15}
\begin{quote}
„Ich hätte große Lust, Ihnen eine eigene vollständige Pastoraltheologie in Versen zu geben. „In Versen?“ Allerdings und dazu in Reimen, die trotz ihrer Rauhigkeit recht für ihren Gegenstand gemacht sind und ich gewiß nicht besser machen könnte. Dazu eine Pastoraltheologie, die nicht vollständiger, vielseitiger, wahrer, lehrender sein könnte. Sie glauben, ich scherze? Ich scherze nicht. Und dazu ist sie von einem der angesehensten, gelehrtesten, frömmsten, verdienstesten Theologen unserer Kirche. Er hat in ihr beinahe alle Erfahrungen seines Lebens (und in seinen Aemtern konnte er deren viel haben), den ganzen Schatz seines Herzens über das, was geistliches Amt, was dieses Standes Leid und Freude, Schimpf und Ehre ist, ausgeschüttet. Und in einer Sprache, die ich ihm beinahe in jedem abgebrochenen Artikel, in jeder verkürzten Sylbe, in jedem Reim und Nichtreim beneide. Und mit einem Salz, einer Wahrheit, wo es sein soll, mit einer Feinheit! wo es gerade heraus sein soll, mit einer Deutlichkeit! — Kurz, mein Freund, hier ist das Gedicht. Lesen Sie’s, auch wo es Ihnen wegen seiner abgebrochenen Form zuweilen etwas langweilig sein sollte, mit Ruhepunkten fort und ja zu Ende. Wo Sie Mitbrüder finden, die Stücke aus dieser Pastoraltheologie, in gutem und bösem Verstande, nöthig haben, leihen Sie damit nicht karg.“
\end{quote}

Herder\footnote{Herder gab das Gedicht in größten Theils glücklichen Kürzungen. Hier folgt es unverkürzt aus dem Original:} im 4. Theil der „Briefe, das Studium der Theologie betreffend“, Weimar 1786 p. 351 f.

\section{Geistliche Kurzweil}

\begin{center}
J.\,V.\,A.
\end{center}

\begin{center}
Zur Ergötzlichkeit einfältiger Christen mitgetheilt.
\end{center}

\begin{verse}
Die Gottes Furcht\\
Hält Widerpart\\
Der Welt Unart;\\
Das Herz bewart\\
So Gott nach Art.
\end{verse}

\quelle{Straßburg. In Verlegung Lazari Zetzners Erben. Anno 1619.}

\quellseite{16}
\subsection[H.\,M. Johanni Buttenlocher]{H.\,M. Johanni Buttenlocher, Pfarrer zu Illingen}

\begin{verse}
Als ich in meinen jungen Tagen\\
Oft hört von guten Pfründen sagen,\\
Wie daß nit feister Suppen wären,\\
Als die nun bey geistlichen Herren;\\
Die möchten mit geschmustem Mund\\
Umgaffen manche gute Stund:\\
Da dacht ich, hats die G’legenheit,\\
So muß ich auch in’s lange Kleid\\
Und sehen, wie ichs dahin bring,\\
Daß ich um lange Bratwürst sing.\\
Denn sollt ich viel umgehn mit Rechten,\\
So müßt ich erst mein Kopf aufstecken.\\
Sollt ich dann jeden Bauern salben,\\
So wär ich schmeckend allenthalben:\\
Hier will doch auch kein Feder glücken —\\
Mein Sach wird sich auf d’Rangel schicken.\\
Da red ich, muß ein andrer schweigen;\\
Da poch ich, muß ein andrer leiden;\\
Da geh ich vor, ein andrer nach;\\
Da schlaf ich zu\footnote{dormito.}, ein andrer wach.\\
\quellseite{17}
Hierzu war ich nun wohlgerüst,\\
Denn alle Kunst in mich genißt.\\
Ich hatt’ durchlernt der Logik Strick\\
Und der Rhetorik Büchlein dick.\\
Ich hatt erlernt des Himmels Sphär\\
Und was die Physik bringt für mehr,\\
Und was von Sitten Ethik sagt,\\
Und was Homerus einherjagt —\\
Das konnt ich gar, als wär’s nur Kraut,\\
Kein Baur hätt mir das zugetraut.

Drauf fiel ich in’s Compendium,\\
Und fehrt mich auch drin drei mal um,\\
Bis ich von Kunst ganz überging\\
Und mir der Witz zum Maul ausgieng.\\
Auch mir mein Röcklein rauscht daher,\\
Als ob ich schon Decanus wär.\\
Was ich nun sah, das konnt ich richten;\\
Was mir fürkam, das konnt ich schlichten;\\
Was mir aufgeben, ward vollendt;\\
Was d’Augen gsehen, machten die Händ.

Noch war kein Stell mir ausgeleert,\\
Wiewol ich wol der besten werth.\\
Jedoch dacht ich, nit jede Pfarr\\
Wird für dich sein die lange Harr.\\
Gleichwol muß sein diaconiret\\
Und dann bald drauf wol pastoriret;\\
So grätt’s dann auf das Decanat,\\
Bis daß du wirst „mein Herr Prälat.“\\
Will man dich dann zum Probst auch haben,\\
So mangelts dir nit an den Gaben.\\
Doch hüt mich Gott vorm Harzin Wald,\\
Den Bergen und den Klüften kalt;\\
\quellseite{18}
Denn mein Bauch ist an Wein gewöhnt,\\
Darum das Bachusgäu mir ziemt.\\
Da kann ich noch mein Glück erschleichen,\\
Inzwischen mich mit Wein bereichen:\\
Es geht doch so, wer wenig hätt,\\
Der kommt nit von seiner ringen Stätt.\\
Soll ich mein Karren weiter führen,\\
So muß nichts mangeln an dem — Schmieren.

Noch muß ein Paß ich thun quittieren,\\
Daß ich auch möcht die Kanzel zieren.\\
Es gschwand mir manchmal vor den Leuten,\\
So gieng mir aus die Red zum Zeiten:\\
Das Vest mir manchmal gar ausfiel\\
Und fällt am meisten auf den Stuhl.\\
Da muß ich andre zu mir bringen,\\
Die mehr umgangen mit den Dingen.\\
All die, so gut Postillen gmacht\\
Und sonst der’n Namen hochgeacht;\\
Die mußten mir wohl unter d’Preß,\\
Bis ich davon brächt alle Eß\\
Und käm in mich die Quintessenz,\\
Auch manch unaufgesucht Sentenz,\\
Damit ich wär für Groß und Klein\\
Gewürfelt wie ein Müllerstein,\\
Und ja kein Casus käm auf d’Welt,\\
Dem ich nit hätt sein Thema gestellt.

Also hatt’, ich mich ausgerüst,\\
Und fehlt nur, daß man es auch wüßt.\\
Drauf zog ich ins gelobte Land,\\
Da Wein wie Wasser, Korn wie Sand,\\
Und sucht mir aus ein feinen Platz,\\
Da ich mich einließ, wie ein Katz.\\
\quellseite{19}
Ich fragt’ die Leut, wo wär der Heerd,\\
Da man hätt, was man nur begehrt:\\
Da wär Wein, Korn, Obst, Holz und Weid.\\
Ich hört nit allweg guten Bescheid.\\
So wollt das Pflaster in den Flecken\\
Mich auch zuweilen lassen stecken;\\
Da gfiel mir nit der Kirchenthurm;\\
Dort warn nit recht gericht die Uhrn;\\
Bald wollt das Pfarrhaus mir nit ein,\\
Bei mir sollts wol noch anders sein.\\
In Summa, was ich contemplirt,\\
Das war von mir als reformiert,\\
Ich war der Mann, auf den gewart,\\
Was man so lange Zeit gespart,\\
Eim jeden Haus ein Stelz zu machen —\\
So gieng ich um mit Narrensachen.

In dem reist ich durchs grüne Gras,\\
Weil da ein schönes Wiesthal was.\\
Da traf ich an ein alt Person,\\
Von Haaren weiß, von Gsicht noch schön;\\
Die gieng mit einem Rechenstiel\\
Im Gras um, thät doch nit gar viel.\\
Ein Pfarrer sie sich wohl vergleicht;\\
Doch hätt ich gmeint, sie hätt sich gscheucht,\\
Mit grober Arbeit sich zu plagen,\\
Sie möcht doch wol ein Kunstbuch tragen,\\
Darin lesen wie mancher Mann\\
So meisterlich in Bann gethan.

Drauf mußt ich den Mann registrieren\\
Und in die Schul erst wieder führen.\\
Sprach: Bona dies, alter Herr,\\
Was habt ihr da für ein Gescherr?\\
\quellseite{20}
Er antwortet: „Semper quies!“ schnell;\\
„Mein Domine, das Gras ich zähl,\\
Daß mir kein Hälmlein komm davon.“\\
Ich dacht: „Mit dem Mann kriegst zu thun.“\\
Drauf mich räusper und so anfang:\\
„Ich weiß nit, ob ich irre gang,\\
Mich dünkt, Ihr seid des Dorfs Pastor.“\\
Er sprach: „Ich bins lang gewesen vor,\\
Eh dann der Herr die Welt erschaffen.\\
Vor vierzig Jahren ists geschehen\\
Und möcht nur wünschen, daß ein Junger\\
Auch unter meine Bauern kämmer,\\
Denn mir entgeht all Kraft und Saft.\\
Je matter Leib, je mehr man schafft.\\
Je weniger Kunst, je mehr mans treibt.\\
Je unwerther, je mehr man bleibt.“

Ich sprach: „Mein lieber alter Herr,\\
Ihr habt Euch nu gemäuset sehr\\
Und habt der alten Batzen viel,\\
Drum wollt ihr kehren um den Stiel.\\
Das möchten doch wir Junge leiden,\\
Die jetzund zehren auf die Kreiden,\\
Erwarten Glück, bei gesundem Leib\\
Ein guten Dienst und reiches Weib.“

Der alt Herr sprach: „Mein Studios,\\
Mich dünkt, Euer Kunst, die mach sich los.\\
Die Logik wird sich in Euch regen,\\
Daß Ihr mit mir recht so verwegen.\\
Wißt Ihr, was Luther in der Sach\\
Einsmals zu ein Nasweisen sprach?\\
Wir Alten, die mit Angst und Flehen\\
Dem Teufel in den H ...... gehen,\\
\quellseite{21}
Grüßen von euch Gnad, Docterlein,\\
Auf weichen Polstern gießen fein\footnote{In unsrer Ausg. ohne Interpunctiun. Herder: „Grüßen vor euch Gnad-Docterlein. Auf weichen Polstern gießen fein.“ Wir bezogen „fein“ auf Teufel.}.\\
Guckt vor so lang darein als wir,\\
Der Scherz wird euch gelegen schier\footnote{Bald aufhören, vergehen.}.

Der Fitz war mir sehr ungewohnt,\\
Ich wünscht, ich hätt des Manns geschont.\\
Drum zog ich bald ein andre Pfeifen,\\
Sprach: „Alter Herr, laßt das fürststreichen.\\
Es war mein Ernst ja nimmermehr,\\
Ich bin euch z’dienen gnogen sehr.\\
Nu will ich etwas Bescheidners sagen,\\
De illo tempore was fragen;\\
Ihr könnt mir geben guten Bescheid,\\
Was warn zu eurer Zeit für Leut,\\
Die selbst in Künsten wohl studiert,\\
Die Jugend löblich angeführt?“

Er sprach: „Ich denk der guten Tag.\\
Da war an Gelehrten wenig Klag.\\
Sollt ich die tapfern Leut all nennen,\\
Ich glaub, ich würde viel nit kennen.\\
Die sind nun todt und leben noch;\\
Nu leben viel und faulen doch.\\
Ich dank ihn’ ihrer guten Lehr;\\
Doch wie ich kommen bin hieher,\\
Hab ich viel anders müssen lernen,\\
Die Hülffen brechen und den Kernen\\
Mit bitterm Schweiß herfür gewinnen.\\
Das werd’t ihr auch noch einmal innen.“\\
\quellseite{22}
Ich sprach: „Ihr gebt aufs Geistlich acht\\
Und der Philosophie nichts acht’t.\\
Daher möcht es wohl kommen sein,\\
Daß euch die Welt nit wollt ein.“

Er lacht und sah mich höhnisch an.\\
„Was meint ihr denn, daß ich gethan,\\
Sprach er: Was möchte doch mein Hirn\\
Zu der Zeit gehabt han für Gestirn?\\
Ich war grammartig und was fein,\\
Und podet überzwerch herein,\\
Ich redt thörlich an manchem Ort\\
Und macht mich mausig immerfort,\\
Im Kopf hatt ich manches Gesperr,\\
Und sonst bisterlich Sachen mehr.\\
Ich log dick\footnote{Anspielungen auf die sieben freien Künste jener Zeit: Grammatik (grammartig), Poetik (podet), Rhetorik (redt thörlich), Musik (mausig), Logik (log dick) \&c.}, daß die Balken stoben\\
Und ecket aus, was krumm gebogen.\\
Meint ihr, daß man zu unsern Zeiten\\
Hab Meister gemacht aus Eselshäuten?\\
Oder hab einen graduiert\\
In dem, daß er gar nie studiert?\\
Oder hab einen heißen treiben,\\
Daß er sein Lebtag wird verschweigen?\\
Oder hab so grob numeriert,\\
Daß aus zwei über sieben wird?“

Der alt Herr hat mich wieder geschreckt\\
Und mir mein Meisterschrei bestekt.\\
Noch mehret ich mich mit aller Kunst,\\
Daß ich nit hätt studiert umsonst\\
\quellseite{23}
Und sprach: „Dürft ich ein Einzig fragen?\\
So ihr die Kunst habt all getragen,\\
Wie ist doch möglich, daß ein Bauer,\\
Der nur umgeht mit Arbeit saur,\\
Euch soll erst anders informieren?“

Er sprach: „Ja freilich deponieren,\\
Bis daß sich packt der hübsch Schulsack\\
Und nimmer quackt der Hackemack;\\
Bis daß verschwindt der Luft Gebäu,\\
Bis daß verdäunt der Pappenbrei,\\
Bis daß verräucht des Hirnes Dampf,\\
Bis daß vertobt der Witze Kampf,\\
Und nun die Praktik kommt zu Haus,\\
Die all Theorik treibet aus.\\
Da findt sich erst, was mir gethan,\\
Daß wir uns haben brauchen la’n.“

Die Ding mir spanisch Dörfer waren;\\
Ich hatt dergleichen nie erfahren.\\
„Wie, sagt ich, sollt der geistlich Stand\\
Von Bauern haben sein Verstand?\\
Sollt nit die hohe Schul uns weisen,\\
Wie wir bezähmen die Unweisen?\\
Was war dann die Theologie\\
Anders als ein Bauernkirchweih?“

Er sprach: „Ich muß euch das verzeihen,\\
Weil ihr noch lauft unter den Freien,\\
So ihr einstmals kommt in den Karren,\\
Da wird man mit euch anders narren.\\
Da werdt ihr sein Dorfstarr, Pfarrnarr,\\
Und alles Rufes Ofenscharr.\\
Da müßt ihr glauben, wissen, thun,\\
Leiden, lassen, fürchten und ho’n,
\end{verse}

\begin{quote}
\quellseite{24}
Was niemand darf, kann, mag, noch will,\\
Und dieses alles in der Still.\\
Denn wer sich dieses will beschweren,\\
Der mag sein Pfarr ein andern leeren.“
\end{quote}

\begin{verse}
Ich bat durch Gott den alten Herren,\\
Er wollt die Sachen nur erklären,\\
Denn ich fragt nit aus Uebermuth,\\
Sondern wie that ein junges Blut.\\
Könnt ich der Sachen ha’n Bericht,\\
Mein Tag wollt ichs vergessen nicht.
\end{verse}

\begin{quote}
„Gern, gern, gern, sprach mein alter Held.\\
Die Weis’ mir nun viel baß gefällt.\\
Weil ihr erst kommet von der Preß,\\
So seid ihr noch zu viel „zapffreß“\footnote{Gleichniß, vom jungen Wein hergenommen.}\\
So muß man euch ein wenig mischen;\\
Ich hoff, ihr sollt es noch erwischen\\
Und mit der Zeit den breiten Rucken\\
Lernen in engen Stand zu schmucken\footnote{Schmiegen.}.\\
So hört mit Fleiß, was ihr nit gwüßt,\\
Und büßet dann den Pfarrerslust.\\
Höret zuvor des Ordens Satz\\
Und zehret drauf die Einstands-Collatz\footnote{Collation, Gastmahl.}.\\
Höret zuvor meins Dorfs Beschwer,\\
Jückt euch die Haut, so kommet her.

I. „Ich hab gesagt, ein Pfarrer glaubt,\\
Daß kaum ein Mensch bringt in sein Haupt.\\
Er glaubt ein Gott, des niemand acht;\\
Ein jeder nach sein Götzen tracht.
\end{quote}

\begin{verse}
\quellseite{25}
Er glaubt ein Himmel, der verschmächt;\\
Ein jeder gern hie ewig zecht.\\
Er glaubt ein Höll, die niemand fleucht;\\
Ein jeder die breit Straße zeucht.\\
Er glaubt ein Gericht, das niemand besorgt;\\
Ein jeder auf die Rache borgt.\\
Er glaubt ein Lohn, den niemand will;\\
Ein jeder will die Hüll und Füll.\\
Er glaubt ein göttlich Regiment;\\
Ein jeder meint, daß Glück sei blind.\\
Er glaubt ein Tod, der alles scheidt,\\
Und jeder pocht auf lange Zeit.\\
So glaubt er, was die Welt verneint\\
Und ihren Augen ungereimt.\\
Damit zeucht er den schweren Narren\\
Und wird gehalten für ein Narren.
\end{verse}

\subsection{II.}

\begin{verse}
Darnach so weiß ein Seelenhirt,\\
Deß die Welt ungern inne wird.\\
Er weiß, daß großer Herren Pracht\\
Bei Gott aufs äußerst wird veracht.\\
Er weiß, daß großer Hirten Schlaf\\
Dem Wolf liefert manch armes Schaf.\\
Er weiß, daß große Leuteschinder\\
Verflucht sind auf Kindeskindern.\\
Er weiß, daß große Krapper-Mäuler\\
Endlich werden zu Höllenheuler.\\
Er weiß, daß große Federhahnen\\
Noch kommen in dem Pfuhl zusammen.\\
Er weiß, daß die groß Ueppigkeit\\
Der Welt gereicht zu Schmach und Leid.\\
Er weiß, daß jedes falsche Herz\\
Sich selbst noch stärkt in ewig-Schmerz.\\
\quellseite{26}
Das weiß er, wills schon niemand wissen\\
Und wird sehr oft darob geschmissen.\\
Damit zeucht er den schweren Narren\\
Und wird gehalten für ein Narren.

III. Drittens so muß ein Pastor thun,\\
Des jedermann will über stohn.\\
Er muß die Wahrheit jedem zeigen;\\
Darüber wird ihm zeigt die Feigen.\\
Er muß aufwischen jede Stund,\\
Darüber man ihm Uebels gunnt.\\
Er muß in d’Pest und Lasareth,\\
Da mancher weit zurüher geht.\\
Er muß zum Feur, Galgen und Rad\\
Ohn Gefängniß — und der Huren Bad.\\
Er muß verzweifelt Buben trösten,\\
Die Ruchlosen durchs Gesetze rösten,\\
Er muß jedermann helfen, bitten,\\
Rathen, warnen, traten und beschütten.\\
Er muß in alle Pfützen treten,\\
All Unlust putzen und ausjäten.\\
Das muß er thun ohn seinen Dank,\\
Bis er drob wird alt, krumm und krank.\\
Damit zeucht er den schweren Narren\\
Und wird gehalten für ein Narren.

IV. Viertens, ein Prediger muß leiden,\\
Da sonst der Thurm zu ist bescheiden.\\
Er leidet der Leut Abgötterei,\\
Aberglaub, Seg und Zauberei.\\
Er leidet der Haufen Sacrament,\\
Dadurch Gott und der Nächst geschändt.\\
Er leidet Verachtung Gottes Lehr.\\
Dafür Wollust wird trieben mehr.\\
\quellseite{27}
Er leidt Unghorsam und Gespött,\\
Daß mancher Pfaff für Ohren geht\footnote{„Das“ öfters, wo „des“ zu verstehen. Vielleicht: das oder um des willen „mancher Pfaff vor die Ohren (der Richter) geht.“}.\\
Er leidt Zorn, Neid, Rachgier und Grimm,\\
Zank, Hader, Schelten, Ungestüm.\\
Er leidt Ehrbruch, Unzucht und Schand,\\
So nur geacht für Narrentand.\\
Er leidt große und kleine Dieb,\\
Finanz und was ihm sonst nit lieb.\\
Er leidt Lügen und Afterreden\\
Praktik, Gelüst und viel Dup . . .\footnote{Nicht ausgedruckt. Von Herder sind die zwei Zeilen wie andere weggelaßen.}\\
Damit zeucht er den schweren Narren\\
Und wird gehalten für ein Narren.

V. Zum Fünften muß ein Priester laßen,\\
Daß die Welt liebt ohn alle Maßen.\\
Er läßt dem Hof sein reiches Kleid,\\
Und bleibt ihm die Kameelhaut bescheid.\\
Er läßt der Schul ihr große Witz\\
Und übt sich in der Liebe Sitz.\\
Er läßt der Reichen Silbergschirr\\
Und trinkt die Bächlein in der Irr.\\
Er läßt der Aufgeblasenen Wind\\
Und sich bei Christi Demuth findt.\\
Er läßt des Fleisches Lust und Geilheit\\
Und bindt sein Nücken jederzeit.\\
Er läßt sein Recht, sein Nutz, sein Fried\\
Und fügt sich, daß er Christi Glied.\\
Er läßt sein Haut, sein Fleisch, sein Wein,\\
Damit er mög bei Christo sein.\\
\quellseite{28}
Das alles muß er willig lassen\\
Und noch dabei sich selber hassen.\\
Damit zeucht er den schlimmeren Narren\\
Und wird gehalten für ein Narren.

VI. Zum Sechsten fürcht ein geistlich Mann,\\
Das sonst bei andern leicht gethan.\\
Er fürcht mit Scheu das End der Welt,\\
Dafür mancher sein Hauptgut zählt.\\
Er fürcht der Kirchen böse Feind,\\
Gewalt und Witz, die manches Freund.\\
Er fürcht der Ärgernis Gefahr,\\
Darin sich übt die größte Schaar.\\
Er fürcht des Glückes gute Wort,\\
Daß nit die Seele werd bethört.\\
Er fürcht seins eignen Gewissens Stimm,\\
Daß es nit schreie wider ihn.\\
Er fürcht der bösen Gesellschaft Schein,\\
Ohn welche mancher nit kann sein.\\
Er fürcht der hohen Gaben Glanz,\\
Die sonst auch Guts verbleiben ganz.\\
Das ist sein Sorg, sein Furcht, sein Angst,\\
Welchs alls die Welt verlacht vorlangst.\\
Damit zeucht er den schlimmeren Narren\\
Und wird gehalten für ein Narren.

VII. Zum Siebenten ein Clericus,\\
Was niemand will, wohl nehmen muß.\\
Er nimmt wenig als niemand glaubt,\\
Denn der thut wohl, der Pfründen b’raubt.\\
Er nimmt das Schlechst vom Pfleger sein,\\
Die schlechste Frucht, den sauren Wein.\\
Er nimmt mit Müh, das saur verdient;\\
Noch hält man als für Ehren die Pfründ.
\quellseite{29}
Er nimmt mit Schmerz von seinen Bauern,\\
Die ihn bezahlen wie die Lauern.\\
Er nimmt alls Faul von falscher Hand\\
Der gilt als er den Tod empfand\footnote{Wie wenn er den Tod empfände.}.\\
Er nimmt mit Dank, was ungern geht.\\
Und bitt ein Dieb um Seinigs stät(s).\\
Er nimmt, daß er niemals genüßt,\\
Dem jedermann ihn drum beschmeißt.\\
Also muß er im Bettel reisen\\
Und endlich lassen arme Waisen.\\
Damit zeucht er den schweren Karren\\
Und wird gehalten für ein Narren.
\end{verse}

\begin{quote}
Wie dünkt euch nu, mein junger Hach?\\
Ist euch zur Pfarr nochmal so gach?\\
Dünkt euch nochmal, ihr seid gefaßt\\
Zu dem Stand, den so mancher haßt?\\
Gelüst euch noch der Pfarrer Braten,\\
Oder wollt ihr der gern entrathen?“
\end{quote}

\begin{verse}
Ich sprach: „O liebster Vater mein,\\
Euer Red, die gehn in’s Herz hinein.\\
Ich bin erschlagen und verstummt\\
Und dank doch Gott für diese Stund,\\
Daß ich durch euern weisen Mund\\
Erfahren soll den rechten Grund.\\
Doch bitt ich, wollt mich weiter lehren,\\
Wo ich mich nu hinaus soll kehren,\\
Denn ich einmal bin Gott verbunden.“

Er sprach: Der Weg ist längst gefunden.\\
Ihr habt gewählt den höchsten Stand,\\
Der hat mehr Gefahr denn Meeres Sand
\end{verse}

\quellseite{30}
\begin{verse}
Und wird durch d’ Welt stets angerannt,\\
Darum bedürft ihr Gottes Hand.\\
Kein Stand auf Erd je werther war,\\
Als der durch Gott berufen dar,\\
Sein Wort und Willen zu verkünden\\
Und dadurch pflegen Gottes Kinden, —\\
Sein Wahrheit und Gerechtigkeit,\\
Sein Wahrheit und Barmherzigkeit,\\
Sein Langmuth und auch großen Zorn,\\
Sein Wunder und des Heiles Horn\\
Fürtragen durch des Geistes Sprach\\
Den Frommen z’Gut, der Welt zu Nach, —\\
Da Gott eins Menschen Zung und Hand\\
Gebracht\footnote{Gebraucht?} gleichsam zu seinem Beistand.\\
Sein Geist und Pfand zu dispensieren,\\
Damit in sein Reich einzuführen.\\
Dem wird vertraut Gotts liebstes Gut\\
Und Jesu Christi Fleisch und Blut,\\
Als auch des Geistes Freudenöl,\\
Damit beseeligt manche Seel.\\
Den Stand laßt euch kein Mensch verleiden,\\
Vor dem all andre Ständ sich neigen.\\
Ist nun der Stand so hoch und werth,\\
So hat er billig sein Beschwerd.\\
Der Teufel ist kein Ding so feind,\\
Als wo Christi Pferch wohl verzäunt.\\
Die Welt braucht immer mehr Betrug,\\
Als daß der Pfaff werd abgeschmeigt mit Fug.\\
Das eigen Fleisch läßt nit sein Tück,\\
Daß er ein fromm treu Herz berück.\\
So bringt der Baalspfaffen Schaar\\
Der Kirchen erst die größt Gefahr.
\end{verse}

\quellseite{31}
\begin{verse}
Denn nie kein Blutvergießen hat\\
Wie Heuchelei der Kirch geschadt,\\
Da man sich selbst, nit Christum sucht,\\
Und mangelt stets an guter Frucht,\\
Da man mehr Witz und Klügeln will,\\
Als Christi Einfalt steckt das Ziel,\\
Oder sonst geht in großen Haufen,\\
Den Leithämmeln all nach hinlaufen.\\
In Summa: wer nit fleißig wacht,\\
Der ist in manch Gefahr gebracht.\\
Je mehr Gefahr, je minder Sold,\\
Eim Gottesdiener soll kein Gold.\\
Wer hie sein Besoldung will einnehmen,\\
Den wird der HErr einmal nit kennen.\\
Hie solls sein g’arbeit, ghüt und gewacht!\\
Dort wirds sein b’lohnt und hoch geacht.\\
Hie solls sein mühsam und unwerth,\\
Dort wirds sein ruhsam und hochgeehrt.\\
Hie solls sein arm, schlecht und bethört,\\
Dort wirds sein warm, recht und gelehrt.\\
Kein Frommer legt hie Gülden an,\\
Wie der aus U X machen kann.\\
Fromm Geld wird hie nicht augmentirt,\\
Wie dem\footnote{Das „dann“ im Orig. ist wohl Druckfehler.}, der die Schreibfeder führt.\\
Fromm Geld mußt\footnote{Wächst.} nit so wunderbar\\
Wie dem, so feist wird in eim Jahr.\\
Fromm Geld läßt sich nit zFuß ereilen.\\
Wie Bösgeld von den auf den Gäulen.\\
Frommgeld vergnügt, wie Gott es fügt.\\
Bösgeld verstäubt, wie viel man treugt\footnote{„Treigt.“}.
\end{verse}

\quellseite{32}
\begin{verse}
„Wollt ihr nu weiden Christi Heerd,\\
So seht, daß ihr berufen werdt\\
Durch Christi Ordnung, nit oblique\\
Durch Geschlecht, Weib, Geld und sonst inique.\\
Gott b’ruft recht durch den obern Mund,\\
Er b’ruft auch in des Herzens Grund\\
Und wie der fromme Luther gmeint,\\
So stünd auch sehr viel bei der Gemeind.\\
Eilt nit zu sehr, Gott weiß euch wohl,\\
Euer Theil euch noch wohl werden soll.\\
Laß laufen, was nit bleiben will;\\
Gott findet die Seinen in der Still.\\
Wahrlich, daß man viel Miethling duldt,\\
Das ist des losen Laufens Schuld.\\
Kein Wurm dem Körper ist so gefähr,\\
Als der gern an sein Stelle wär.\\
Den Leichnam läßt man kaum erhalten,\\
So will schon ein’r sein Dienst verwalten.\\
O wenn Verfolgung reget sich,\\
Wie mancher schrie nit: „Hie bin ich!“

„Seid ihr dann zu der Kirchen kommen,\\
Den schweren Eid auf euch genommen;\\
Da rüst euch nu mit Herz und Muth,\\
Daß ihr alls nehmen wollt für gut.\\
Ja, wie jener uns thät bescheiden,\\
Müßt ihr auch lernen Henken leiden.\\
Weh euch, so man euch zu viel lobt!\\
Wohl euch, wenn die Welt heftig tobt.\\
Weh euch, so euch der Dienst wird süß.\\
Wohl euch, so ihr findt viel Verdruß.\\
Weh euch, so euch die Welt gefällt,\\
Wohl euch, so sie euch Fallen stellt.\\
Weh euch, so ihr nach Ehren strebt.
\end{verse}

\quellseite{33}
\begin{verse}
Wohl euch, so ihr im Niedern lebt.\\
Weh euch, so ihr auf Titel schaut.\\
Wohl euch, so euchs wenig vertraut.\\
Weh euch, so ihr hie haltet mit.\\
Wohl euch, so euch die Welt ausschütt.\\
So könnt ihr Gotts Haushalter sein,\\
Der Welt ein Dorn, ein Ruth und Pein.

Noch müssen wir das Hauskreuz tragen,\\
Wie jeder Ehmann wird beladen.\\
Was jedem gschicht, das kann uns werden,\\
All täglich Fäll gehören auf d’ Erden.\\
Wollt ihr doch hie den kürzsten Weg,\\
Daß euch genug göttlicher Segen;\\
So laßt nit z’viel auf Erden gahn,\\
Der Himmel steht euch besser an.\\
Gewöhnt eur Leut zu schlechter Art.\\
Nicht ehers lernt sich als Hoffahrt.\\
Laßt Arbeit thun, was essen will;\\
Zur Ruh bleibt Zeit noch überviel.\\
Traut nit zu wohl ein jedem Maul;\\
Das Bös ist frisch, das Gut geht faul.\\
Veracht nit leichtlich armer G’stalt,\\
Gott viel Geheimnis dabei vorbehält.\\
Wißt nit zu viel, das sag ich z’vor,\\
Daß ihr nit seid des Dorfs Doctor.\\
Glaubt auch nit alles, was man leugt;\\
Unzeitig Eifer manchen treugt.\\
Ich gib euch auch das noch zu Bericht,\\
Verlaßt euch auf kein Menschen nicht.\\
Gott sei euch einig euer Scopus,\\
Dazu der Mensch euch helfen muß.\\
Sonst wo ohn Gott der Mensch soll helfen,\\
Da gilts laufen, schmieren und gelsen,
\end{verse}

\quellseite{34}
\begin{verse}
Und ist doch nichts als Wort und Schein;\\
Der gewinnt’s, der über euch kann sein;\\
Damit hat euch euer Götz gelassen;\\
Wer nimmer hat, der mag fort passen.\\
O kurze Zeit und schnöde Freud,\\
Wie manchem hast du Gold gezeigt\\
Und ihn gesetzt ins tiefe Grab.\footnote{Roth?}\\
Der glaubt es, ders versuchet hat.\\
Ich bin, mit Züchten z’reden, auch\\
Der Leut, die nit gehangt im Rauch.\\
Kömmt man mich gar in Ofen stecken,\\
Man wird mit brauchen andre Stecken.\\
— Han meine Leut so ghalten d’Leut,\\
So ist es Zeit, daß ich mich leid.“

Hiermit hält er sich schier erzürmt,\\
Vielleicht viel Beschwerliches aufgezürmt.\\
Ich fiel in Dreck und sprach: „Ich Thor\\
Nenn euch mit Ehren Präceptor.\\
Mein lieber, frommer, weiser Herr,\\
Wär ich vorlängsten kommen her,\\
Mein Ohren sollten kürzer sein,\\
Mein Rüssel reiner als beim Schwein.\\
Ich hab gefolgt der Narren Zunft,\\
Da oberherrscht die Unvernunft.\\
Ich meint, ein jedes Dorf hätt Schätz,\\
Die man nur fieng ohn Strick und Netz.\\
Nu gib ich mich in eurer „Lehr“,\\
Daß ihr meinen Wurm recht metzt\\
Und legt mir ab mein Ring und Hut\footnote{Die academische Magisterzierde.},\\
Das Röcklein und das Sträußlein gut,
\end{verse}

\quellseite{35}
\begin{quote}
Und stoßt mir d’Nasen in das Buch,\\
Daß ich solch neue Leges such,\\
Damit, wenn ich komm unter d’Leut,\\
Ich nit umgeh als der nit g’scheid.“
\end{quote}

\begin{verse}
Das schlug mein alten Herren zu.\\
Er sprach:
\end{verse}

\begin{quote}
„Ich nicht lieber thu,\\
Als jungen Leuten, die noch jähren,\\
Was ihnen noch weit fehlt, zu lehren.\\
Es mag es aber, was noch glitzt\\
Und noch wohl hintern Ohren schwitzt,\\
Nit allweg leiden, daß wir Gecken\\
Ihn’ wollen ihr groß Kunst erschrecken,\\
Die sie im großen Buch erquckt.\\
Darum sich mancher Alter duckt\\
Und denkt: Laß vor die klugen Nasen\\
Anlaufen, daß die rothen Masen\\
Ihn’ geben Lehr, wie in der Welt\\
Es manchem Frechen hoch gefällt.\\
Doch muß ich leider auch bekennen,\\
Und werd es mit meinem Schmerzen innen,\\
Daß nit alles, was schwarz, geistlich ist;\\
Daß nit all Geistliches lauter Christ,\\
Daß nit all Lauteres ist gesund,\\
Daß nit all Gesundes ist fürn Mund.“
\end{quote}

\begin{verse}
Hierauf hat mich der ehrlich Mann,\\
Ich wollt mit ihm zu Hause gahn,\\
Daselbst ein Süpplein helfen essen:\\
Das Schwätzen wird sich nicht vergessen.\\
Er muß heimtragen an der Stangen\\
Den hübschen Vogel, den er gefangen,\\
Und ihn seiner alten Mutter bringen.\\
Die weiß doch auch von diesen Dingen,\\
\quellseite{36}
Darum sie auch den jungen Tropfen\\
Die Gauchfedern weiß auszurupfen\\
Und sagt ihm umsonst ihren Text.\\
Das Haus, das sei doch allernächst,\\
Da er mit seinem Holderstock\\
Oft spalten manchen dicken Block,\\
Lieb und Leid williglich „gelait“,\\
Manch’ tiefe Hauswunden geheilt,\\
Vor manchem Sturmwind sich geduckt,\\
In manchem Wetter sich geschmuckt,\\
Vor manchem Unglück sich entzuckt,\\
Durch manches Löchlein durchhin guckt.

Also gieng ich mit Scham und Freud.\\
Mein Herz war eng und sich ausbreit.\\
Mein Kunst war klein und hört doch viel.\\
Mein Reu war groß. Eilt doch zum Ziel.\\
Ich wollt nit, daß ich wälsche Land\\
Dafür hätt gesehen allesammt.\\
Denn ein deutsch Herz, so man das findt,\\
Ist werther als viel fremdes Gsind.\\
Der sagt, was fehlt, und räth hierzu.\\
Hiemit kommt man mit Gott zur Ruh.\\
Was aber nur schwätzt mum, mum, mum,\\
Und wirft den Brei im Maul herum,\\
Das braucht viel Zeit, Geld, Müh und Sorg,\\
Daß man im Eiteln gar erworg,\\
Nun wünsch ich, daß all meine Geselln\\
Ihm auch abtrennen lan die Schelln\\
Und geben sich in Christi Orden,\\
Der nie kein Frommen süß ist worden.\\
Hiemit folg ich meinm Alten nach.\\
Wer bessers weiß, der besser die Sach.
\end{verse}

\quellseite{37}

\quellseite{38}
\chapter{Einleitendes}

\quellseite{39}
\quellseite{40}
\section[I. Vorbereitung zum geistlichen Amte]{I. Einige Bemerkungen über die Vorbereitung zum geistlichen Amte}

1. Erst auf den Universitäten beginnt bei uns die eigentliche Vorbereitung zum geistlichen Amte. Auf den Gymnasien unterscheiden sich die zukünftigen Diener der Kirche von den übrigen Schülern durch nichts, wenn nicht etwa durch sorgfältige Erlernung der hebräischen Sprache. Uebrigens genießen sie dieselbe allgemeine Vorbildung wie alle, die auf Universitäten zu irgend einem besonderen Lebensberufe vorbereitet werden sollen. Man muß so gerecht sein zuzugestehen, daß diese allgemeine Vorbildung auch auf anderem als dem gewöhnlichen Wege erreicht werden könnte; aber eine andre Frage ist, ob sie dem Geistlichen überhaupt entbehrlich sei? Wollte man sie umgehen, so würde die Stellung des Geistlichen im öffentlichen Leben eine ganz andere und sein Gesichtskreis mehr als gut wäre eingeengt werden. So wie es nun einmal geworden ist, verlangt man von einem Diener der Kirche Kenntniß und Ueberblick aller menschlichen Verhältnisse und muß sie auch verlangen. Gerade dazu bedarf er aber einen Grad der allgemeinen Bildung, vermöge dessen er allen Ständen der Gesellschaft
\quellseite{41}
gleich stehen und offenen Zugang zu jeder Höhe derselben in Anspruch nehmen kann. Es ist daher auch jene allgemeine Vorbildung für den Geistlichen, so wie er unter uns gestellt ist, nicht wohl entbehrlich. Sie ist es um so weniger als man dem Alter eines gewöhnlichen Gymnasialschülers eine reuelose, besonnene Wahl des Lebensberufes nicht mit voller Seelenruhe zutrauen kann. Zwar gibt es eine Art von innerlicher Berufung Gottes zum geistlichen Amte\footnote{Vocatio Dei interior. So nennen unsere Theologen die innerliche, gleichsam eingeborne Neigung zum h. Amte. Mit Recht, denn Gott gibt auch solche Neigungen. Doch muß man die Umstände prüfen, durch welche oft eine solche Neigung recht und schlecht werden kann. So ist z.\,B. kein wahrer innerer Beruf da, wenn bei allem Drang doch keine Gabe da ist. Cf. Tarnov. de minist. p. 63. Joh. Fecht Instructio pastoralis. edit. II. Rost. 1728. p. 10 ff. Mayeri Museum ministri ecclesiae. 1703. p. 7.}, welche sich oftmals von Jugend auf durch eine unüberwindliche Liebe und Neigung zu diesem Amte kund gibt. Aber diese Neigung wird durch eine Hinausschiebung der endlichen Entscheidung für den geistlichen Beruf bis ans Ende der Gymnasialjahre weder ausgetilgt noch vermindert; sondern sie wird im Gegenteil, wenn sie rechter Art ist, gerade dadurch bewährt, geläutert und gestärkt und geht dann um so sicherer und zuversichtlicher ihrem Ziel entgegen. Es würde ja auch die ausgesprochenste Neigung bemißtraut werden müssen, wenn sie die herkömmliche Vorbildung nicht überwinden und sich die Früchte derselben aneignen könnte.

2. Zuweilen hat man (S. z.\,B. Fecht p. 10) geklagt, daß sich in der protestantischen Kirche selten bessere Köpfe zum
\quellseite{42}
Studium der Theologie wenden. Hätte diese Klage Recht, so wären beide, die besseren Köpfe und das Studium der Gottesgelehrtheit zu bedauern. Auf keinen Fall läge aber die Schuld an der Gottesgelehrtheit selbst. Diese, sowie das geistliche Amt, zu welchem sie befähigen soll, machen gerade die größten Ansprüche an die Kräfte eines Menschen. Die Theologie steigt in die höchsten Höhen, in die tiefsten Tiefen, hat das weiteste Feld der Erkenntniß, und an sie grenzen ringsum alle andern Gebiete des menschlichen Wissens an, empfangen von ihr Licht, Leben und Bedeutung. Das geistliche Amt erfordert unter allen menschlichen Berufsarten die meiste geistige Productionskraft. Wer täglich zu geben hat, der muß eine reiche Quelle haben, aus welcher er schöpfen kann; wo nicht, so wird er bald in quälende, peinliche Armuth versinken. Man sollte deshalb denken, daß gerade die reichsten Geister zur Theologie gezogen werden müßten, wie das Eisen zum Magnet. Geschieht nun das nicht, so liegt die Ursache, wie gesagt, nicht in der Theologie. Oft darf ein Jüngling nicht wählen wie er will. Manchen verhindert das sündige Getriebe eines bösen Herzens, die vorhandene Gabe und den Trieb des h. Geistes zu erkennen. Bei den sogenannten höheren Ständen, die aber an guten Köpfen nicht gerade fruchtbarer als die niedrigeren zu sein pflegen, erstickt oft Reichthum und Verweltlichung die tief im Herzen vorhandene Lust zu einem Studium und Beruf, welcher das Sichtbare gering und klein, die unsichtbaren Güter groß, das Gewissen unruhig machen, und die eingeborne Sehnsucht nach Bleibendem mächtig anfachen könnte.

\quellseite{43}
3. Eine andere, mit dem vorigen keineswegs verneinte Frage ist die, ob nicht auch mittelmäßige und geringe Gaben in der Kirche unter gewissen Umständen mit Segen verwendet werden können? Die Antwort ist nicht schwer. Nur zur Abwehr des Vorwurfs, als wäre Theologie und geistliches Amt der besseren Köpfe nicht werth, ist oben von hohen Anforderungen die Rede gewesen. Andererseits wird niemand leugnen, daß Theologie und geistliches Amt zweierlei sind, daß man das Amt führen kann, ohne ein in wissenschaftlicher Erkenntniß weit geförderter Theologe zu sein, daß die Productionskraft, welche der Geistliche bedarf, keineswegs allein von gelehrten Kenntnissen abhängt, sondern sich oft gerade dann am reichsten und fruchtbarsten erweist, wenn es gilt, allgemeine, jedermann zu Gebote stehende Erkenntniß auf die mancherlei Fälle des innern und äußern Lebens anzuwenden. Hier wird oft ein Tropfen Wassers zur Quelle, die ins ewige Leben springt. Woher käme es sonst, daß so mancher gering begabte Pfarrer mit dem segensreichsten Erfolge arbeitet, während umgekehrt oft hochbegabte und gelehrte Männer in ihren Gemeinden wie hochragende, aber unfruchtbare Bäume stehen? — Die Wahrheit ist, daß es im Reiche Gottes mancherlei Gaben gibt, deren keine zu feiern braucht; — daß zum Bau Zions mancherlei Gaben nöthig sind; daß eine jede ihre Arbeit und ihren Segen finden kann, den man sich nur nicht durch eigene Schuld, durch hochmüthiges, lähmendes Beneiden fremder Gaben und Erfolge verkümmern muß.

4. Bezieht der zukünftige Diener der Kirche nach vollbrachten Gymnasialstudien die Universität, so wird die
\quellseite{44}
Forderung an ihn gestellt, einen Theil seiner Zeit noch dem Studium der allgemeinen Wissenschaften zu widmen. Es ist gegen diese Forderung nichts einzuwenden; da man ja die Nützlichkeit und Nothwendigkeit allgemeiner Bildung für den Geistlichen anerkennen muß und das Gymnasium die Aufgabe, allgemeine Bildung zu geben, vollständig weder löst noch lösen kann. Aber ob nicht die Art und Weise, wie der Jüngling jener Forderung genügen möge, frei zu geben oder wenigstens eine gewisse Rücksicht auf Anlage und Individualität, sowie auf das erwählte Berufsstudium zu nehmen wäre? Wenn nur allgemeine Bildung erreicht wird, so kann es am Ende gleich sein, wie viele Zeit auf einmal oder nach und nach von dem oder jenem Jüngling dazu angewendet wird. Man läßt gerne die allgemeinen Studien dem Berufsstudium vorangehen, und gewiß manche Studien, z.\,B. Logik und Aehnliches, passen ganz wohl auf die Schwelle des Gymnasial- und academischen Studiums. Aber ob dies von allen allgemeinen Studien gilt? Ob nicht vielfach der Fortschritt vom Besondern zum Allgemeinen, vom Berufsstudium zum allgemeinen eben so der natürliche wäre, wie in der Liebe der von der Bruderliebe, als der früheren, zur allgemeinen Liebe, als zu der aus jener folgenden? Ob nicht für die Weitschaft vieler Köpfe manches Studium zu ferne, manches zu hoch liegt und überdies männliche Sinnen erfordert? Vielleicht gäbe es eine gute Anzahl Raisonneurs unter den Studierenden weniger, wenn sie sich nicht einbildeten, vor allem philosophisch im Sinne der neueren deutschen Wissenschaft gebildet sein zu müssen. Sie lernen den
\quellseite{45}
philosophischen Sprachgebrauch mit leichter Mühe, wenden ihn jugendlich muthwillig auf alles an, auch auf ihre Lehrer; sonderlich auf diejenigen, welche sie in ihre Berufsstudien einführen sollen, und der Spaß wird kein Ende, bis etwa die Pforten eines Examinationssaales sich ganz in der Nähe zeigen. Dann wohl heißt es: Philosophandum est, sed paucis, nam omnino nocet.

5. In Anbetracht des eigentlich theologischen Fachstudiums wäre sehr zu wünschen, daß zunächst das gelehrt und gelernt würde, was dereinst im Amte frommt. Ein gewisses sorgloses Lernen, ein freies Sichbewegen in unbeengten Räumen der Wissenschaft bietet allerdings seine Vortheile; aber es bleibt doch immer das Vortheilhafteste das Nöthige zu lernen, und ob es gleich seine Grenzen hat, gibt es doch auch den meisten Studierenden Weitschaft genug. Was hilfts, wenn der Lehrer lehrt, nicht, was sich im dereinstigen Beruf bewähren und segensreich erweisen wird, sondern was ihn selbst, den Lehrer, den Mann der Wissenschaft, in seinem wissenschaftlichen Berufe drängt und bewegt? Was hilfts, den Schüler zum Standpunkt des wissenschaftlichen Theologen emporzuheben und in Untersuchungen einzuführen, die sein unreifer Geist nicht würdigen kann? Was hilft es, in erhabenen Sphären der Wissenschaftlichkeit zu schweben oder vielmehr durch des Lehrers Flügel getragen zu werden, da man doch bald in das nüchterne Land des Berufes als ein Fremdling herunter fallen wird? Weder Lehrer noch Schüler genügen damit ihrem Berufe, jene nicht dem gegenwärtigen, welcher doch zunächst die Ausrüstung zukünftiger Pfarrer
\quellseite{46}
gebeut; diese nicht dem zukünftigen des geistlichen Amtes. Wie sehr ist bei solchem Lernen namentlich die Schaar von mittelmäßig oder gering begabten Köpfen zu bedauern, welche noch überdies die Mehrzahl der Studierenden ausmachen, da sie auf diesem Wege nicht bloß nicht zum Ziel gefördert, sondern verkehrt und untüchtig werden! Gewiß, zur Vermeidung so manches Müßiggangs des Geistes, so manches Traumes, mancher Täuschung, so mancher schädlichen Irrfahrt sollte man immer und immer wieder die Studierenden erinnern, daß nicht der Katheder ihrer Lehrer, sondern das bescheidene Loos einer Dorfkanzel auf die meisten unter ihnen warte. — Ein Jüngling, der sich an der Hand eines Lehrers der Philosophie in unabsehbare Räume der Wissenschaft verflogen hatte, fand zu seinem Befremden, daß er dem Lehrer und Vater seiner Jugend fremd geworden war. Da er die Ursache nicht in sich selbst suchte, fragte er den Lehrer und bekam die Antwort: „Du bist Student, ich Pfarrer, dein Weg ist überall, meiner ist schmal und eng, der Weg des geschlossenen Berufes; du hast weite Aussichten, ich sehe wie ein Bergansteigender nichts als das Stück Weges vor meinen Füßen. Wenn auch deine Wissenschaft sich engert, du zum Berufe einlenkst, den ich habe, wirst du mir wieder kindlich nahe werden.“ — Das geschah, und es geschieht öfter. Aber oft ist der Uebergang schwer, und mancher bleibt in den Seilen seiner Universitätsjahre so hängen, daß er nie ins Geleise eines ernsten Berufslebens sich zu fügen vermag. Manchem kommt sein bißchen Wissenschaftlichkeit so theuer zu stehen, daß er völlig unbrauchbar ins Amt kommt und
\quellseite{47}
eben deshalb ein unglücklicher Mensch für Lebenszeit bleibt. Da muß er denn seinem academischen quadriennium oder quinquennium ein „perdidi“ nachrufen, ohne den Fehl so leicht wieder gut machen zu können.

6. Daß mit alle dem nicht dem Lernen und wahrer Gelehrsamkeit Schmach nachgesagt werden soll, versteht sich von selbst und wird hier noch ausdrücklich versichert. Es wäre thöricht, so ohne allen Unterschied Lernen und Gelehrsamkeit zu schelten, da sie ohnehin, recht genommen, sich selten genug finden. Man studiere nur das Rechte und das allen Ernstes. Man höre nur auf treue academische Lehrer mit aller Aufmerksamkeit; es lernt sich von ihnen mehr und besser als aus dem Buche. Vom Mund zum Ohre geht ein lebendigerer und segensreicherer Weg als vom Buch zum Auge. Habet nescio quid latentis energiae viva vox et in aures discipuli de ore autoris transfusa fortius sonat, sagt Hieronymus ad Paulinum. Man sammle und speichere auf. „Hat einer nur fein viel gelernt,“ sagt A. H. Francke, „es wird nichts so gering sein, er wird es zu seiner Zeit wohl brauchen können, so gar, daß bei diesem verworrenen Zustand der Kirche auch die unnützen Dinge ihm oft werden müssen zu Nutzen kommen, daß andern damit gedienet werde.“ Man kann unbedenklich mit Theophylactus sagen: „Doctrina est virtus et character episcopi“ und mit Hilarius: „Summa omnium virtutum episcopalium est scientia et doctrina.“

Alles, was mit dem Vorausgehenden gesagt werden sollte, läßt sich in die Worte zusammenfassen: „Lerne fürs Amt und werde fürs Amt so gelehrt, als nur immer
\quellseite{48}
möglich.“ Wir können es nicht besser als mit den Worten unseres ersten Pastorallehrers Hartmann sagen: „Medicus corporis coram aegroto multa erudite de discrimine, qualitate, de causis item et medela morborum disserens suam audientibus ostentare et verbis depraedicare atque adserere potest peritiam: infirmo autem aliter prodesse non poterit, quam si pharmaca sua praesenti malo applicet et ea dicat potissimum, non quae in arte medica sunt acutissima et subtilissima, aut quae ad fallendum tempus dicis causa aegroto narranda nuper legit et excogitavit, sed quae unice faciunt ad praesentis morbi depulsionem et aegroti informationem et consolationem. Ita pastor ecclesiae, qui longe pretiosissimarum animarum est medicus, non quam multa (quibus non raro obtunduntur auditores) et quam erudite excogitata sibique arridentia et occurrentia, sed quam accommodata et pro plerorumque captu maximeque ad auditorum suorum emendationem, informationem et consolationem facientia dicere queat, sollicitam navabit operam. v. Act. 20, 19.“

7. Ist es nun mit dem zum heiligen Amte vorbereitenden Studium hie und da ein übles Ding, so wäre das doch immer noch erträglich, wenn nur nicht an diesem einen Stücke so viel anderes hinge. Oratio, meditatio, tentatio faciunt theologum, sagt Luther mit vollem Rechte, und dieser Ausspruch umfaßt jedenfalls in übersichtlicher Kürze alles, was zur Ausbildung für das heilige Amt nöthig ist. Aber wer sagt heut zu Tage diesen Spruch sich und andern im Ernste? Wer betont jetzt noch
\quellseite{49}
in vollem Ernste oratio und tentatio in gleichem Maße wie die meditatio? Es könnte ja, denkt man, der wissenschaftlichen Entwicklung Eintrag thun, wenn ein Student oratio und tentatio eben so angelegentlich ins Auge fassen wollte als sein Studieren. Und doch ist’s wahr, was Eberlin von Günzburg sagt: „Zwei Dinge sind, welchen alle andere Arbeit des Studierens weichen soll: Andacht und eures Bruders Noth!“—Und nicht mit Unrecht warnt ein trotz allen üblen Gerüchtes achtungswerther Lehrer (Fecht): „Frigida est plerumque per totam vitam pietas, quae non statim cum ipsa studiorum origine conjuncta fuit!“ Und doch sind die Versuchungen der Studierenden so groß, daß diese ohne Unterlaß beten sollten, um ihnen mit heißer Seele zu entrinnen! Man denke doch nur an die Gefahren, welche überhaupt, besonders aber in unsern Tagen im academischen Studium liegen! Nur ein betender Geist ist zur Erfahrung, Erduldung und Ueberwindung von Versuchungen geschickt; nur betende Studenten können die Schlingen und Netze wahrnehmen und vermeiden, welche der Satan allen denen zu legen pflegt, die sich vornehmlich mit Verstandeswerk befassen! Aber leider, gerade Gebet und Anfechtung pflegen auf Universitäten völlig übersehen und vernachlässigt zu werden. Daher sind aber auch die meisten von Universitäten abgehenden Candidaten so gar unerfahren und ungeschickt zu Gebet und Versuchung, so unpraktisch und zaghaft in den Geschäften des heiligen Amtes, in welchen sie sich allenfalls versuchen, daß man erstaunen muß. Selbst den Neophyten gleich, machen sie an andern, theuer erkauften, oft
\quellseite{50}
leicht verderbten Seelen Experimente, deren Werth oder Unwerth sie aus eigner Erfahrung nicht kennen. Ja, es werden von jungen Geistlichen viele Seelen, die sich ihnen vertrauend anschließen, irre geführt oder doch aus Noth und Verzweiflung nicht errettet, weil den nunmehrigen Seelenhirten während ihrer Vorbereitung die Zeit nur zwischen Buch und Krug verrann.

Wenn ein A. H. Francke und andere in herzlich guter Meinung, ohne es zu wollen, durch populäre Vorlesungen, weil sie nicht eben so gründlich als wohlgemeint waren, hie und da dem echten Studium der Theologie Abbruch thaten und der Oberflächlichkeit zu Entschuldigungen halfen, so muß das nicht gelobt werden. Aber es gibt ein wahres, tiefes, umfassendes Lernen (man denke an Thomas Aquinas, an Luther und Johann Gerhard etc.), welches nie zum Zweck wird, sich nie überhebt, aber gerade darum desto gewaltiger und heilsamer wirkt und zugleich in sich selbst desto vollkommener ist und mehr und mehr wird. Es gibt eine Gelehrsamkeit, die kräftig zum Leben vorbereitet und den Schatten der Francke’schen Richtung nicht wirft. Sie kann in Versuchung die Fackel tragen und dem betenden Geiste zur Läuterung und Zuversicht dienen. Diese ist es, welche der Kirche Noth thut, die ja keine bloße Schule ist; nach ihr sollten alle Studierenden streben. — Ein Dic cur hic\footnote{Cujuslibet theologiae studiosi museo inscripta esse debet memorabilis Pauli sententia: „Pietas ad omnia utilis est“ 1. Tim. 4, 8. et illa Psalmis: „Impiorum via peribit!“ Ps. 1. (Fecht p. 16.)} — dürfte daher über
\quellseite{51}
dem Studierzimmer jedes jungen Studenten stehen, und keiner sollte vergessen, daß er einem Berufe entgegengeht, nach welchem er sich selbst selig machen soll und die ihn hören werden. Das heilige Amt soll ja am Ende nichts anders, als im Thale dieses tödtlichen Lebens und lebendigen Todes den Weg zum ewigen Leben erleuchten und die Seelen mit starker Hand auf demselben führen. Was soll also die Vorbereitungszeit, wenn nicht dazu befähigen in allerlei Weise?

\section[II. Erprobung der Tüchtigkeit]{II. Erprobung der Tüchtigkeit zum heiligen Amte}

\subsection{Examen}

8. Von alten Zeiten her dient das Examen zur Erprobung derjenigen Tüchtigkeit, ohne welche niemand eine Vokation oder die Ordination bekommen soll. Durch diese enge Beziehung des Examens zur Vokation und Ordination wurde man früherhin veranlaßt, demselben die Stelle zwischen Vokation und Ordination einzuräumen. Ohne Vokation glaubte man zu einem Examen keinen Anlaß zu haben. Indes ergab sich doch bei dieser Stellung des Examens eine Unbequemlichkeit, welche sich immer fühlbarer machte. Indem man nämlich nur bereits Vocierte zum Examen zuließ, kam man in den Fall, zuweilen schon vocierte Pfarrer des Amtes und der Ordination für unwerth oder für untüchtig zu erklären und ihnen unmittelbar vor der Thüre des Amtes und einer damit gegebenen zeitlichen Versorgung den Eingang zu verwehren. Die üble Lage
\quellseite{52}
solcher examinirten Pfarrer versetzte natürlich die Examinatoren selber in eine mißliche Lage; ja, es konnte bei sich erweisender Untüchtigkeit schon vor Schluß des Examens aus derselben eine gewisse Befangenheit und Ungerechtigkeit des Urtheils entspringen, wie sie sich überhaupt leicht einzustellen pflegt, wenn man Rücksichten nehmen muß. Viel erträglicher für den Abzuweisenden, und deshalb viel leichter für den Abweisenden ist eine Abweisung ohne Beziehung auf eine bestimmte Vocation. Darum hat man es schon in früheren Zeiten nicht bloß räthlich befunden, die Candidaten vor erfolgtem Rufe zu examiniren; sondern hie und da, z.\,B. im ehemaligen schwedischen Herzogthum Bremen wurde auch der gute Rath befolgt\footnote{S. J. Gerhards Loc. XXIV. S. XIII. p. 171. Hartmann p. 116 V. Fecht p. 49.}. Unsere jetzige Sitte, das examen pro ministerio gleich am Schlusse der Studienzeit zu gewähren, vermiede also wenigstens die Inconvenienz der obenbenannten Examina.

9. Ein Examen pro ministerio ist in Gottes Wort nicht bloß gerechtfertigt, sondern gefordert. 1. Tim. 3, 10 befiehlt der Apostel: „Die Diakonen lasse man zuerst versuchen; darnach lasse man sie dienen, wenn sie unsträflich sind“; und 1. Tim. 5, 22 heißt es: „Die Hände lege Niemand bald auf.“ Beide Stellen befehlen nicht nur eine δοκιμασία der Aspiranten zu kirchlichen Aemtern, sondern auch Strenge der δοκιμασία. Zwar könnte man in ersterer Stelle nur eine Prüfung fürs Diakonat begründet finden; allein das Diakonat erscheint dieses Falls nur als
\quellseite{53}
Eingang ins geistliche Amt, und es muß 1. Tim. 3, 10 die Prüfung der Aeltesten um so mehr mit eingeschlossen sein, als die Erfordernisse, welche beides fürs Diakonat und fürs Presbyterat 1. Tim. 3 gestellt werden, ohne Prüfung gar nicht erkannt werden können.

10. Prüfungsgegenstand im Allgemeinen ist die Tüchtigkeit (ἱκανότης) zum Lehramte. (2. Tim. 2, 2. „Was du von mir gehört hast durch viele Zeugen, das befiehl treuen Menschen, die da tüchtig sind, auch Andere zu lehren.“ Vgl. 2. Kor. 2, 16. 3, 5.) Zu dieser Tüchtigkeit rechnen unsre Väter gemäß der heiligen Schrift mit Recht:

1. confessionis sinceritatem, Reinheit des Bekenntnisses;\\
2. docendi dexteritatem, Lehrhaftigkeit;\\
3. morum integritatem, Sittenreinheit.

11. Von den drei Erfordernissen ist das erste und dritte in gleichem Maße von allen Aspiranten des geistlichen Amtes zu fordern. Eines von diesen beiden ist so nötig, als das andere; jedes ist unerläßlich; beide zusammen sollten den character indelebilis des evangelischen Geistlichen bilden. In neueren Zeiten hat man aber gerade auf diese beiden Erfordernisse nicht immer das scharfe Auge gerichtet, welches man auf sie hätte richten sollen. Namentlich hat man Bekenntnis-Reinheit und Treue mit Mißtrauen angesehen, ja gar wie eine unheilvolle Beschränktheit angefahren. Dagegen hat man ungebührlicher Maßen aus dem Examen pro ministerio hauptsächlich ein „wissenschaftliches“ Examen gemacht. Da man die Kirche
\quellseite{54}
zur Schule erniedrigt hatte, mußte man nichts Höheres, als das oft so leicht zu erringende und sehr täuschende Lob der Wissenschaftlichkeit. Recht ist es, mit Theophylakt zu behaupten: „Doctrina est virtus et character episcopi“ und mit Hilarius: „Summa omnium virtutum episcopalium est scientia et doctrina“; auch ist es wahr, daß man die Studenten und Candidaten zum Lernen anzutreiben alle Ursache habe. Aber auf der andern Seite sind Theophylakt und Hilarius nicht Lobredner dessen, was man heut zu Tage für Wissenschaftlichkeit ausgibt, und jeden Falls ist es unstatthaft und schädlich, das zweite Erforderniß vor 1 und 3 hervorzuheben, zumal in dem hie und da geübten Maße, nach welchem ausgezeichnete Wissenschaftlichkeit etwa auch die baarsten Mängel in den beiden andern Beziehungen zu bedecken vermochte.

12. Völlig richtig, nur mit zu wenig innerer Zustimmung sagt Hüffell: „Schon die Apostel gaben sich viele Mühe, Uebereinstimmung in der reinen Lehre zu sichern (Act. 15, 22. 29. Röm. 14, 20. 2. Tim. 2, 15. Tit. 3, 10. 1. Cor. 5, 1–5); sie warnten vor falscher Lehre und fürchteten untergeschobene Schriften (2. Thess. 2, 2). Rechtgläubigkeit war daher der Mittelpunkt aller Tugenden der neuen Kirche, und Origenes konnte schon denjenigen, welcher von den Hauptlehren der Kirche abgefallen war, härter beurtheilen, als den, der eine böse That begangen hatte.“ Gewiß hatte Origenes in diesem Stück vollkommen Recht, eben so Irenäus, Tertullian, Cyprian, Lactantius in ähnlichen Stellen. Nur wenn man verkennt,
\quellseite{55}
daß heilige Gedanken Ursprung alles heiligen Lebens sind, und daß jeder Mangel des Gedankens einen Mangel des Lebens nach sich zieht, wenigstens konsequenter Maßen nach sich ziehen muß, kann man den Aposteln und den Vätern Unrecht geben. Will doch der Jünger der Liebe diejenigen, welche eine andere Lehre bringen, nicht einmal beherberget und gegrüßt wissen, 2. Joh. 10. 11. So klar war er sich selbst, so zuversichtlich die hohe Meinung, welche er von dem Werthe völliger, reiner Lehre hatte! — Pure, wirkungslose Orthodoxie ist scheinheilige Bosheit, aus welcher kein Schluß auf wahre Orthodoxie zu machen ist. Man fordere nur Orthodoxie und zwar genaue, mit dem Bekenntnisse der Kirche scharf zusammentreffende! Je ernster und strenger die Forderung, desto erschwerter und kenntlicher wird die Heuchelei.

13. Ueber Vernachlässigung des Lebens und Wandels bei der δοκιμασία der Candidaten klagt schon A. H. Francke in seinen pastoralen Aphorismen. Und doch liegen die Anforderungen der heiligen Apostel gerade in diesen Stücken 1. Tim. 3, 2 ff. und Tit. 1, 7 ff. so klar vor Augen. „Ein Lehrer soll sein: unsträflich, eines Weibes Mann, nüchtern, mäßig, sittig, gastfrei, nicht ein Weinsäufer, nicht pochen, nicht unehrliche Handthierung treiben, sondern gelind, nicht habsüchtig, nicht geizig, der seinem eigenen Hause wohl vorstehe, der gehorsame Kinder habe mit aller Ehrbarkeit, ehrbar, nicht zweizüngig, der das Geheimniß des Glaubens in reinem Gewissen habe, und Summa: ein Lehrer
\quellseite{56}%
, Priester oder Prediger, Bischof oder Kaplan soll also beschaffen sein, daß er ein gutes Zeugniß habe von denen, die draußen sind, auf daß er nicht falle dem Lästerer in die Schmach oder Stricke.“ So faßte Hutter die apostolischen Anforderungen zusammen, und gewiß, nach diesem Maße sollte ein jeder Candidat gemessen werden, so weit es nämlich bei dem jugendlichen Alter, in welchem Candidaten zu stehen pflegen, angewendet werden kann. — Es gibt zwar eine höhere Anforderung, welche alle andern in sich schließt, aus welcher alle andern fließen, nämlich die Liebe zu Christo. Als der Herr den gefallenen Jünger Petrus in seinem Apostolate bestätigte, ihm seine Lämmer und Schafe aufs neue anbefehlen wollte, fragte er ihn zu dreien Malen um die Liebe. Πέτρε, φιλεῖς με; Εἰ φιλεῖς με, ποίμαινε τὰ πρόβατά μου. Die Pastorallehrer unserer Kirche machen auf dies hohe Erforderniß zum heiligen Amte aufmerksam (S. z.\,B. Balduin Instit. past. p. 65.) und ein großer Hirte des Alterthums, Johannes Chrysostomus, sagt mit Bezug auf die Liebe zu Christo (Περὶ ἱερωσύνης ed. Tauchn. Λογ. II. Κεφ. II. 93.): Ἀγέσθωσαν εἰς μέσον οἱ πολλῷ τῷ μέτρῳ πλεονεκτοῦντες ἁπάντων, καὶ τοσούτῳ ὑψηλότεροι τῶν ἄλλων κατὰ τὴν τῆς ψυχῆς ἀρετήν, ὅσον τοῦ παντὸς ἔθνους Ἑβραίων κατὰ τὸ τοῦ σώματος μέγεθος ὁ Σαούλ, μᾶλλον δὲ καὶ πολλῷ πλέον. Indes ist der Blick ins Herz und armen Menschen nicht gegeben, wir sind an die kenntlichen Früchte gewiesen und thun darum wohl, wenn wir keine Liebeserscheinung, keine Liebenswürdigkeit eines Candidaten
\quellseite{57}
für Wahrheit nehmen, wenn er sich, am Maße von 1. Tim. 3. und Tit. 1. gemessen, nicht bewährt.

14. Was endlich die Lehrhaftigkeit anlangt, welche der Apostel 1. Tim. 3, 2 mit dem Worte διδακτικός ausdrücklich fordert, so rechnet man nach Luc. 21, 15 dazu στόμα καὶ σοφία d.\,i. δύναμιν ἐρμηνευτικὴν und σύνεσιν πνευματικήν. Zu letzterer gehören alle gelehrten Kenntnisse, welche man von einem Diener des Wortes mit Fug und Recht verlangen kann, zu ersterer das praktische Geschick im Reden, Lehren und im Umgang. Hutter faßt es schön zusammen, wenn er auf Grund der h. Schrift (1. Tim. 4, 13. 2. Tim. 2, 2. 15. 3, 14. 16. 4, 1. 2. 5. Tit. 1, 9) sagt: „Evangelische Prediger müssen die Eigenschaft haben: 1) daß sie nützliche Kenntnisse und Sprachen studiert und gelernt, 2) in der Schrift wohl belesen, 3) die Artikel der christlichen Religion aus Gottes Wort wohl verstehen und allerhand Irrthum widerlegen können, 4) daß sie mit solchen Gaben von Gott gezieret, daß sie, was sie studiert und gelernt haben, von sich an Tag geben, und also auch Andere lehren und unterweisen können.“

15. So ernst und gebieterisch die erste und dritte Anforderung an alle Examinanden in gleichem Maße zu machen sind; so offenbar ist es, daß rücksichtlich der Lehrhaftigkeit nicht jede Pfarrstelle gleiche Forderungen und denselben Grad von Ausbildung erheischt, daß also auch im Examen nicht an alle gleiche Forderungen gestellt werden müssen. Das ganze Notenwesen, welches heut zu Tage allerdings oft zur Profanation der Kirche gemißbraucht wird, hat dennoch eine Seite, welche jeden Falls anerkannt
\quellseite{58}
zu werden verdient. Es liegt nämlich in demselben unleugbar noch die Anerkennung, daß verschieden begabte, und in verschiedenem Grade ausgebildete Leute im h. Amte Raum finden können. (Vgl. S. 3.) Tum συνέσεως πνευματικῆς, tum δυνάμεως ἐρμηνευτικῆς sunt certi gradus, nec summus illorum donorum gradus exigendus ab eo, qui ad inferius aliquod ecclesiae ministerium promovetur, — sagt einer, dem ein kompetentes Urtheil zuzutrauen ist, Johannes Gerhard.

16. Ganz offenbar hat man in neuerer Zeit bei den Prüfungen protestantischer Candidaten eher zu viele, als zu wenige Anforderungen an die σύνεσις πνευματική der Examinanden gemacht. Die ganze verkehrte, übermäßig wissenschaftliche Richtung der protestantischen Theologie brachte das mit sich. Dennoch hört man oft aus dem Munde römisch-katholischer Laien, daß von ihren Geistlichen mehr gefordert werde. Dieselbe Sprache führten öfters auch römische Schriftsteller, besonders von der Partei der Jesuiten, die es den früheren Lutheranern so gerne zum Vorwurf machten, „daß bei ihnen zu Predigern berufen und aufgestellt würden gemeine Idioten und Laien, welche durchaus nichts wissen, oder verstehen, auch wohl weder lesen noch schreiben können, als Fleischer, Weber, Köche, Schuster, Schneider \&c., welche auch, nachdem sie eine Zeit lang gepredigt und es ihnen nicht länger bei solchem Thun zu verbleiben gefallen will, sich wiederum zu ihren Handwerken und Handthierungen begeben.“\footnote{Wie gar nicht tadelhaft unter Umständen das sein kann, beweist folgender amerik. Fall. Der erste Sendling nach Amerika, der von Neuendettelsau ausgingen, wär zuvor Schuhmacher gewesen. In Amerika wurde er Pastor. Ein Katholik nöthigte ihn, sein Amt niederzulegen. Er wurde nun Kirchenvorsteher unter einem andern Pfarrer und nährte sich von seinem alten Handwerk. Als sich sein Leiden wieder verlor, wählte ihn die Gemeinde wieder zu ihrem Pfarrer. Weit entfernt, daß deshalb ein Schatten auf den Mann oder die Gemeinde fallen muß, kann der Fall vielmehr beiden zur Ehre gereichen.} Dergleichen Vorwürfe
\quellseite{59}
könnte man den Römischen dadurch vergelten, daß man auf dasjenige Minimum von Anforderungen verwiese, welches sich im Pontificale Romanum findet. Nach diesem ist erforderlich sufficiens litteratura grammatica, examen natalium, ne ordinentur illegitimi, exploratio membrorum corporis, ne sit in eis debilitas, deformitas vel impotentia, quam ob causam palpare et numerare jubent ordinandi faciem, oculos, aures, nasum, manus, digitos, pedes, si opus fuerit discalceatos, ne forte sint lignei. Wir können ihnen jedoch eine bessere Antwort geben.

Die Römischen sahen bei ihren Vorwürfen auf die Reformationszeit, in welcher allerdings hie und da weniger ausgebildete Prediger auf die Kanzeln traten. Heutiges Tages könnten sie dieselben Vorwürfe in Bezug auf viele unserer Heidenboten und viele lutherische Prediger Nordamerikas wiederholen. Wir könnten ihnen darauf mit Hutter antworten: „Gesetzt, daß zur Zeit der Verfolgung oder bei andern dergleichen unvermeidlichen Nothfällen in Mangel gelehrter Leute und ordentlich berufener Prediger einer oder mehr aus den Laien oder Handwerksleuten zu Lehrern wären aufgestellt worden; so wollen uns die Feinde
\quellseite{60}
berichten, ob denn solches für eine Sünde zu rechnen und so hoch zu tadeln sei, und ob es nicht viel besser und fürträglicher sei, bei solchem Zustand etwas von Gottes Wort hören, denn durchaus alles Berichtes und Trostes beraubt sein. Und wäre sich durchaus nicht zu verwundern, wenn gleich bei vorigen Zeiten im finstern Papstthum solche gemeine Leute aufgetreten und öffentlich gelehrt hätten an den Orten, da der Papst und seine Meßpfaffen das heilige Evangelium nicht allein nicht gelehrt, sondern auf allerlei Weise und Wege dasselbe verfolgt, ja gänzlich zu verdunkeln sich höchstes Fleißes bemüht haben. Und gehört freilich hieher, was dort Christus den Schriftgelehrten und Pharisäern, als solche dem Volk und Kindern das Hosianna verbieten wollten, geantwortet hat: „Ich sage euch, wo diese schweigen, werden auch die Steine schreien.“ — Gott kann Ihm ein Lob zurichten aus dem Munde der jungen Kinder und Säuglinge, ja aus dem Munde der gemeinen Laien, und durch solche Leute Seine Kirche erbauen, besonders wenn die vermeinten Geistlichen zu Wölfen und zu Tyrannen werden und fromme, reine Lehrer verfolgen und aus dem Mittel hinwegräumen.“ Uebrigens wäre es ja doch auch Beschränktheit, nur denen Amtsfähigkeit zuzutrauen, welche unsern gewohnten Schul- und Bildungsweg durchlaufen haben. Gar mancher unstudierte Prediger in Nordamerika steht an Tüchtigkeit und Segen hinter seinem studierten Collegen keineswegs zurück, und Aehnliches könnte man aus den verschiedensten Zeiten der Kirche zum Beweis anführen. Seckendorf stellt in seinem Christenstaat, 3. B. 7. C., §. 5., eine bündige und schöne Vergleichung der ersten
\quellseite{61}
Zeit nach der Reformation mit den ersten Zeiten der Kirche nach Christo an. Er zeigt, daß man nicht alles nach den Umständen der eigenen Zeit und Lage beurtheilen müsse, daß der h. Paulus nicht viel Edle, nicht viel Weise gefunden habe, welche sich Christo hingegeben hätten, daß er seines Gleichen gemeine Bürger und Handwerker durch Auferlegung seiner Hände zu Presbytern gemacht und von ebendenselben ausdrücklich bezeugt habe: „Der heilige Geist habe sie zu Aufsehern und Bischöfen gesetzt, zu weiden die Gemeine Gottes“ Act. 20, 8. Diese seien dann nicht bloß selbst „die allerbesten Bischöfe, Pfarrer und Lehrer“ gewesen, „sondern hätten auch ihre Söhne wieder zu gleichem Amt unterrichtet und manchmal in viel Gliedern nach einander einer Gemeine vorgestanden“ (S. 430). Das sagt Seckendorf nicht zur Schmach der Gelehrsamkeit, sondern nur „damit niemand den Wahn faßte, es wäre die Art des Studierens, des academischen Lebens, der graduum u.\,dgl. (so jetzt gebraucht, auch nicht verworfen wird) ein Substantialwerk der Auferziehung der Prediger und Lehrer.“ Uebereinstimmend urtheilt auch Joh. Gerhard L. XXIV. K. III, S. 1, §. LXXIV, (T. XII, p. 70 f.). Er weist es als die Ansicht des Alterthums nach, daß unter den Heiden, da, wo die Kirche und das Predigtamt noch nicht geordnet, oder die Lehrer Wölfe geworden sind, jeder Christ die Verpflichtung habe, sich seiner Nebenmenschen durch Lehre anzunehmen. Er weist auf die Wirksamkeit der christlichen Jünglinge Frumentius und Aedesius unter den Abyssiniern, auf die Wirksamkeit gefangener christlicher Frauen unter den Heiden, und bestätigt
\quellseite{62}
diese alten Beispiele durch die Urtheile des Origenes und Ambrosius. Der Letztere sagt zu Ephes. 4: „Ut cresceret plebs et multiplicaretur, omnibus inter initia concessum est et evangelizare et baptizare, et scripturas in ecclesia explanare.“ Ganz konform mit Ambrosius sagt auch Luther: „Ja, wie? Wenn ein Christ nicht dazu berufen, so darf er nicht predigen? Antwort: Hie sollst du den Christen in zweierlei Ort stellen. Aufs erste, wenn er ist an dem Ort, da keine Christen sind, da bedarf er keines andern Berufes, denn daß er ein Christ ist, inwendig von Gott berufen und gesalbt; da ist er schuldig, den irrenden Heiden oder Unchristen zu predigen und zu lehren das Evangelium aus brüderlicher Liebe, ob ihn kein Mensch dazu beruft etc., denn Noth bricht alle Gesetze und hat kein Gesetz. Aufs andere, wenn er aber ist, da Christen an dem Ort sind, die mit ihm gleiche Macht und Recht haben, da soll er sich selbst nicht hervorthun, sondern sich berufen und hervorziehen lassen, daß er an Statt und Befehl der Andern predige und lehre.“ G. tom. II. Jen. germ. f. 244.

17. Da das erste und dritte Erforderniß beim Examen der Candidaten durchaus unerläßlich und mit gleicher Strenge an alle zu richten sind, hingegen innerhalb des zweiten Erfordernisses Stufen Statt finden, so wird insonderheit dies zweite der Spielraum eitler und hochmüthiger Bestrebungen. Man bemerkt dabei nicht, daß man unvermerkt am dritten Erforderniß einbüßt, so wie man sich dem Hochmuth in Bezug auf das zweite ergibt. Indem man von hohem Fluge träumt, wird man untüchtig und sinkt
\quellseite{63}
tief herunter. Denn es liegt wahrlich für die innere Tüchtigkeit des Pfarrers mehr an Demuth und Heiligkeit des Lebens, als an Gabe und Gelehrsamkeit. Man soll auch von denen lernen, die uns nicht wohlwollen. Darum sei es erlaubt, aus den in vieler Hinsicht vortrefflichen „Bemerkungen über die Seelsorge besonders auf dem Lande von P. Aegidius Jais, 5. Aufl. 1843“ einige starke Sprüche für die Nothwendigkeit eines reinen Lebens zu entlehnen:

\begin{quote}
p. 65. Vita clerici evangelium est populi.
\end{quote}

\begin{quote}
„ 63. „Non bene auditur, qui non bene diligitur, sagt Gregor der Große. Diligo will aber mehr sagen, als amo. Non tantum te amo, schrieb Tullius an seinen Freund, verum etiam diligo.“
\end{quote}

\begin{quote}
p. 17. „Das Evangelium der meisten Weltmenschen ist das Leben der Priester, wovon sie Zeuge sind.“ Massillon.
\end{quote}

\begin{quote}
p. 65. „Die beste Vorbereitung, das schönste Exordium einer Predigt ist der fromme Wandel des Predigers.“
\end{quote}

\section{III. Die Candidatenjahre}

18. Wenn ein Studiosus die Universität verlassen und das Examen glücklich überstanden hat, so hat er insgemein eine gute Reihe von Jahren zu warten, bis ihm eine Pfarrstelle übertragen wird. Es fragt sich deshalb, wie diese Wartezeit am besten angewendet werden könne. Jedenfalls
\quellseite{64}
sollte bei dieser Berathung nicht das bessere oder größere Stück Brot den Ausschlag geben, da ohne Zweifel eines jeden werdenden, wie jedes gewordenen Kirchendieners beständiges Symbolum das Wort des Propheten sein muß (Jerem. 17, 16): „Menschentage habe ich nicht begehrt, das weißest Du.“ Die Herrlichkeit des heiligen Amtes soll billig jeden Candidaten ganz erfüllen, — und alle seine Wünsche sollen in dem einzigen aufgehen, daß er für dasselbe immer mehr reifen und tüchtig werden möge. Darum soll er auch alle seine Zeit so anwenden, daß er die Erfüllung seines innigsten und größten Wunsches befördere und nicht selbst Ursache sei, wenn er um das Werk seines Lebens betrogen wird.

19. Viele Candidaten bringen ihre Wartezeit als Lehrer an gelehrten Schulen, als Hofmeister in einzelnen Familien, als Vikarien zu; einzelne treten in öffentliche Predigerseminarien ein. Es fragt sich nun, was für einen Werth für den dereinstigen Beruf diese verschiedenen Arten, die Wartezeit anzuwenden, haben können, und welche zu erwählen das weiseste sei?

20. Was zuerst das Lehren in Schulen und einzelnen Familien anlangt, so hat es seinen gewissen, unbestreitbaren Nutzen. Die Studienzeit gewährt in der Regel keine Gelegenheit, das Schulwesen kennen zu lernen und durch Uebung Lehrweisheit zu erlangen; trägt nun auch die Wartezeit in dieser Beziehung nichts aus, so ist nicht abzusehen, mit welchem Fug man dereinst die Inspection einer Schule übernehmen könne. Wie oft geschieht es, daß auch sehr tüchtige Geistliche wegen ihrer Unerfahrenheit und ihres
\quellseite{65}
Ungeschicks im Schulwesen keinen Einfluß auf Schulen und Schullehrer gewinnen können, vielmehr ihren Schullehrern selbst die Veranlassung zu hochmüthiger Selbsterhebung und Ungehorsam geben! Ist daher auch nicht gerade mit Luther zu rathen, zehn Jahre lang Schule zu halten, ehe man eine Pfarrei übernimmt; so ist doch jeden Falls zu rathen, daß man sich in der Wartezeit mit dem Schulwesen und Schulhalten befasse und durch Lehren lehren lerne. Darum ist gewiß die Zeit nicht verloren, welche ein Candidat als Lehrer an öffentlichen und Privatschulen zubringt. Und zwar wird es seinem zukünftigen Amte am meisten frommen, wenn er in ganz gewöhnlichen Schulen in Städten und auf dem Lande als Lehrer eintritt und in Gottes Namen die Elemente alles Wissens der heranwachsenden Jugend beibringt. Wollte Gott, es gäbe nicht bloß einzelne Gemeinden, in welchen Candidaten zugleich die Stelle des Pfarrers und des Schullehrers vertreten (Exposituren, Schulexpositi); sondern es würden die Candidaten der Theologie regelmäßig eine Zeit lang für die Schulen verwendet. Dies würde nöthigen, schon auf den Universitäten mehr Rücksicht auf das Schulwesen zu nehmen, es würde den Schulen und den Candidaten nützlich sein, und endlich würde dadurch auch auf die gewöhnlichen, nicht studierten Schullehrer eine zugleich demüthigende und erhebende Wirkung ausgehen. Gar nicht davon zu reden, daß eben damit auch ein Einigungsmittel mehr zwischen Schule und Kirche gegeben wäre.

21. Weniger zu empfehlen, als das Schulhalten, ist die Ausfüllung der Wartezeit durch Hofmeisterei. Zu geschweigen, daß sich der Candidat in den Häusern der höheren
\quellseite{66}
Stände leicht an Bedürfnisse und eine Lebensweise gewöhnt, von denen er sich späterhin wieder entwöhnen muß, daß dies ohne Schmerz und Verdruß nicht abgeht, — daß er sich leicht einen Ton des Lebens und Benehmens aneignen kann, der für einen zukünftigen Vater des Pfarrvolkes nicht paßt; so bietet auch der Privatunterricht eines oder einiger Kinder der vornehmeren Stände, so viel Noth, Uebung und Versuchung er einem jungen Manne bereiten kann, doch nicht die Frucht, welche man ihm für sein zukünftiges Amt wünschen muß. Seine Beziehungen sind in Betreff der Lehrgegenstände zu manchfaltig, in Betreff der Lehrlinge zu einförmig, als daß er lernen könnte, was er durch Schulhalten für Schule, Kinderlehre und Amt lernen soll. Doch kommt hiebei allerdings auch Lust und Gabe in Betracht.

22. Fragt man weiter, was ein junger Candidat anzufangen habe, um sich zum eigentlich geistlichen Amte zu befähigen; so möchte die Wahl zwischen Predigerseminarien und Vicariaten gegeben sein. Vielleicht ist kein einziges Predigerseminar so gestellt, daß es die Candidaten in’s heilige Amt einführen könnte. Was hilft es, wenn man noch so viele Predigten schreibt, so lange diese Predigten zu keiner Gemeinde in Beziehung stehen, so lange sie Exercitien sind, bei denen man sich Gemeindezustände denken muß, welche man am Ende doch nur durch ein reges, theilnehmendes Leben innerhalb einer Gemeinde recht erkennen kann. So wenig ein Maler dadurch ein Portraitteur wird, daß er Gesichter nach Phantasieen malt: so wenig wird einer ein Prediger durch fingirte Predigten. So wenig ein
\quellseite{67}
Hirn- und Herzpredigen in bald dieser, bald jener Gemeinde, so wenig sogenannte Studentenpredigten nützen können; mehr nützen sie am Ende doch als das Ausarbeiten und Aufsagen solcher Predigten\footnote{Damit soll nicht geleugnet werden, daß auch Uebungen der Art für den allerersten Anfang und zur Uebung in der Form ersprießlich sein können.}. Eben so ist es mit den Katechesen. Wenn man gleich Kinder aus den Schulen beruft, damit die Candidaten sich an ihnen üben, so sind doch dergleichen Uebungen für die Kinder oft nur katechetische Späße; die Kinder verweigern dem Katecheten dasjenige Herz, welches die Sache verlangt, und mit Einem Wort: man predige oder katechisire, so fehlt dem jungen Manne die amtliche Stellung, die Alles und Alles ändert\footnote{Man hat die Erfahrung gemacht, daß amerikanische Nothhelfer, denen man in Deutschland für Predigt und Katechese äußerst wenig Geschick zutraute, in Nordamerika bedeutendes in beiderlei Hinsicht leisteten, so wie sie ins Amt kamen. Solch ein Unterschied ist zwischen Vorbildung und Leben, und so gewiß bildet das Leben am meisten.}. — Was man aber außer Predigen und Katechisiren in einem Predigerseminare Praktisches vornehmen wolle, ist gar nicht abzusehen. Es müßte denn ein Predigerseminar zur Bedienung einer Gemeinde verwendet werden. Dann müßte es aber einen Rector und moderator studiorum haben, damit nicht die Candidaten experimentieren, ohne zu wissen, nach welchem Ziele sie ringen. Dann müßte der Rektor Pfarrer und der Seminaristen dürften nicht zu viele sein. Das käme dann aber, richtig gefaßt, auf das hinaus, was wir sofort sagen werden.

\quellseite{68}
23. Wir glauben nämlich, daß Vicariate bei weitem die beste Gelegenheit gewähren, einen Candidaten ins heilige Amt einzuführen. Doch muß man uns hier wohl verstehen. Wir reden nicht von Vicariaten bei alten, abgelebeten, bereits kindisch werdenden Pfarrern; diese erfordern geübte, erfahrene Candidaten. Wir reden auch nicht von Vicariaten bei überbürdeten Pfarrern, nicht von Vicariaten zur Bedienung von übergroßen Gemeinden, nicht einmal im allgemeinen von Vicariaten in Stadtgemeinden. Im ersten Fall fehlt es an Unterweisung, in den beiden letzteren fehlt etwas, was zu allseitiger Orientirung und Erfassung des heiligen Amtes so ersprießlich scheint, wir möchten es das Compendiöse nennen. Namentlich bieten Stadtgemeinden kein übersichtliches Bild einer Gemeinde und gemeindlicher Verhältnisse, und auch der Pfarrer selber pflegt in Städten in seiner wahren Eigenschaft und Stellung weniger kenntlich zu sein, als auf dem Lande. Wir reden von Vicariaten bei tüchtigen Landpfarrern, die alt genug sind, um Erfahrung gemacht zu haben, jung genug, um die Mühe des Lehrens und Anlernens zu ertragen, lehrhaft genug, um ihre Erfahrungen mittheilen zu können, neidlos genug, um aufgehenden Sternen in ihren Gemeinden den nöthigen Raum zu gönnen, damit sie sich in Gemeinde, Gemeindeleben und Amt einleben können. Man wird freilich sagen: solche noch kräftige, übersichtlichen Gemeinden vorstehende Landpfarrer nehmen keine Vicarien. Und man hat Recht, sie nehmen keine, d.\,i. sie haben keinen Anlaß, Vicarien auf eigene Kosten zu nehmen. Aber was hinderts denn, dasjenige Geld, welches man auf so manche unnütze
\quellseite{69}
Anstalt verwendet, zur Unterstützung von Candidaten anzuwenden, die eine Art von Biennium bei erfahrenen Pfarrern antreten wollten oder auch sollten? Müßten sich die Candidaten in äußerlichen Dingen auch sehr einschränken, was wäre es denn? Dürfen sie sich etwa nicht vorüber zu dem, was ihnen doch zu Theil werden wird, wenn sie in Amt und Ehe treten, — zu einem armen Leben? Gewährt ihnen Amt und Kirche nicht Hoheit genug, um sich für zeitliche Armuth zu trösten? — Freilich müßte man einerseits den Candidaten die Wahl des Pfarrers überlassen, der ihnen den Uebergang aus der Theorie zur Praxis zeigen soll, denn Vertrauen zum Lehrer ist nöthig, und gerade bei einem Lehrer dieser Art am meisten; anderer Seits müßte man vorsehen, daß nicht eine falsche Wahl getroffen würde. Oder was wäre für Nutzen und Segen zu hoffen, wenn z.\,B. sogenannte erfahrene Geistliche gewählt würden, welche gelernt haben, spielend ihr Amt, wie ein Handwerk zu verrichten, die höhere Idee und das höhere Ziel der Kirche und des Amtes aber verlernt und verloren haben? Es müßten an Geist und Gemüth gehobene, höher begabte Männer sein, denen man Candidaten vertrauen sollte. Deren könnte man aber auch ganz wohl nicht bloß einen, sondern auch mehrere Candidaten anvertrauen. Die Pfarrereien solcher Landpfarrer könnten die edelsten Pflanzschulen der Kirche, rechte Seminarien sein, — und solche Landpfarrer eigneten sich hinwiederum am besten zum Amte eines Seminarregens.

24. Zu einem solchen Biennium bedürfte es der Ordination nicht, da gerade die Geschäfte, welche die
\quellseite{70}
Ordination erfordern, eines Theils durchaus dem Pfarrer zukommen und deshalb von keinem Andern übernommen werden können, andern Theils ganz wohl durch Zusehen und Mithelfen gelernt werden können\footnote{Concionari quidem licet non ordinatis, si fiat exercitii causa; sacramenta vero administrare, confessiones audire et similia soli debentur ordinario ministro: in his enim exercitio non est opus. Balduin p. 77 f. Wiewohl eine Theilnahme an der Seelsorge, so weit sie außerhalb des Beichtstuhls geübt wird, dem jungen Candidaten unter Aufsicht und Verantwortung des Pfarrers gar wohl gestattet werden kann. Gibt es doch sogar viele Fälle, in denen der Pfarrer die Mittelsperson eines Gehülfen ganz wohl und zu großem Segen gebrauchen kann.}. Der Mangel der Ordination brächte gerade einen sehr heilsamen Unterschied zwischen Pfarrer und Candidaten vor die Augen der Gemeinde.\footnote{Eine solche Art der Amtseinweisung würde sich auch sehr gut mit der Stellung eines Candidaten vertragen, der die Schule eines Pfarrers übernommen hätte. Ein Candidat könnte Schule halten und vicariieren. Müßten die Candidaten durch die Schule ins Amt gehen, so fände sich Gelegenheit zur Vorbildung am leichtesten. Es würde auch sehr dienlich sein, wenn sie Cantorei und Messnerei eine Zeit lang versähen.}

\section{IV. Vicariate}

25. Wir haben bereits von Vicariaten geredet, welche Candidaten bei tüchtigen, lebenskräftigen Pfarrern suchen sollten, um von ihnen in die Geschäfte des heiligen Amtes eingeführt zu werden. Man sollte diese Stellung zu einem
\quellseite{71}
Pfarrer nicht einmal Vicariat nennen, weil die Eigenschaft eines Vicars oder Stellvertreters dabei so wenig hervortritt und der Candidat nicht um des Pfarrers, sondern um sein selbst willen die Stellung einnimmt, ja gesucht und als Wohltat angenommen hat. Ganz anders ist es mit den eigentlichen Vicariaten, die ein Candidat nicht zu seiner Amtseinführung, sondern zur nöthig gewordenen Unterstützung eines überbürdeten oder kranken oder alten Pfarrers übernimmt, in welchen er wirklich pastor vicarius und Stellvertreter des Pfarrers ist, für welche deshalb auch Ordination erfordert wird. Man kann sagen, daß sich bereits mit dem Eintritt in solche Vicariate die Wartezeit eines Candidaten schließe. Hat er doch von da an alle geistlichen Amtsbefugnisse, und gehört doch alles, was ein ordinierter Vicar nicht kann oder darf, sammt und sonders nicht zum heiligen Amte. Das größere Einkommen und die etwaigen Befugnisse, Pflichten und Rechte, welche ein Pfarrer vor dem Vicare voraus hat, bringen in dem geistlichen Amte eines Ordinirten eben so wenig eine Aenderung hervor, als dies durch eine Verheirathung und die Ehe geschieht. Dennoch aber ist ein ordinierter Vicar kein Pfarrer, und es ist hauptsächlich der Unterschied zwischen ihm und einem Pfarrer, welchen wir nun besprechen möchten.

26. Das natürliche Verhältnis eines Vicars zu seinem Pfarrer ist jeden Falls das der Unterordnung. Der Pfarrer, welcher durch Umstände veranlaßt wird, einen Vicar anzunehmen, hört deshalb nicht auf Pfarrer zu sein. Er behält alle seine Pflichten, alle seine Rechte, all’ seine
\quellseite{72}
Würde, all’ seine Verantwortung und ist auch verantwortlich für seinen Vicar, so lange und in so weit er selbst zurechnungsfähig ist. Der Vicar hingegen ist und bleibt nur Hilfsgeistlicher (Cooperator), Stellvertreter und zwar immerhin nur in denjenigen Geschäften, in welchen der Pfarrer eine Vertretung braucht und begehrt. Der Pfarrer kann sich nach seinem eigenen Ermessen gewisse Geschäfte vorbehalten; ja man wird es ihm gar nicht verdenken können, wenn er diejenigen Amtshandlungen, welche ihn eigentlich zum Pfarrer machen, z.\,B. Beichte und Sacramentsverwaltung, schwer und langsam dem Vicar überträgt. Möglich, daß er manche Geschäfte schlechter verrichtet, als sein Vicar; daß er seiner Gemeinde die größte Wohlthat erzeigte, wenn er alle seine amtlichen Befugnisse dem Vicar übertrüge. Wenn ers aber nicht thut, wenn er sich gegen die Abtretung jedes Geschäfts, jeder Befugniß wie gegen den Tod wehrt, wer kann ihn, so lange er nicht durch rechtlichen Spruch des Amtes für unwerth oder für untüchtig erklärt ist, dazu zwingen? Er thut und läßt immer nur ihm Zuständiges, er thue und lasse, was es auch sei. In dieses Verhältnis muß sich jeder Vicar von vornherein ergeben. Er gehorche, diene und thue, was ihm aufgetragen wird. Auch er kann alt werden und in den Fall kommen, einen Vicar haben zu müssen, er schaue sich selbst in seinem alten Pfarrer an und schone sich selbst in ihm. Es ist nur Achtung vor seiner eigenen Zukunft, wenn er den Unterschied zwischen Pfarrer und Vicar wohl faßt und fest hält. Edler und schöner ist nichts an einem Jüngling und Mann, als Bescheidenheit im alten Sinne des Wortes
\quellseite{73}
und Zufriedenheit mit der Lage, in welcher man ist und sein soll. Kein liebenswürdigerer und geliebterer Vicar, als der sich gern und freudig und eifrig in seinen Schranken hält und bewegt; ihm wird auch gewiß je länger, je mehr vertraut werden.

27. Da ein junger Mann, auch wenn er großen Amtsverstand hat, doch jeden Falls keine große Amtserfahrung haben kann; so liegt ihm beim Eintritt in ein wirkliches Vicariat nichts näher, als Mißgriffe und Fehltritte. Er sollte drum froh sein, wenn ihm nicht gleich beim Anfang viele Verantwortlichkeit auferlegt wird. Auch wenn er schon ein Vicariat der ersten Art, ein Biennium der Praxis unter einem tüchtigen Pfarrer hinter sich hat, sollte er doch nicht zu viel Vertrauen auf seine Reife setzen, den Schritt zu größerer Selbständigkeit nicht leichtsinnig thun, vielmehr die Leitung seines Pfarrers, wenn er sie nur irgend haben und brauchen kann, suchen und annehmen. Er sollte gerne dessen Amanuensis und Schüler sein, insofern es möglich ist, zumal da es ohne alle Aufopferung von Selbständigkeit und Originalität geschehen kann und vielmehr das Zeichen eines hohen Grades von Selbständigkeit und Herrschaft über das eigene Herz und die eigenen Kräfte ist. Und doch wird gerade dies den meisten Vicarien so schwer! Die Jugendluft der Selbständigkeit und Ungebundenheit macht es ihnen so schwer. Will nun der Pfarrer, wie es auch oft der Fall ist, seiner Rechte Herr bleiben, verhehlt ers nicht, gibt ers offen kund: so ist der Vicar unglücklich und beginnt bald alles Thun und Lassen seines Pfarrers im schlimmen Lichte zu sehen. Kommt es vollends, daß der
\quellseite{74}
Pfarrer, welcher vielleicht auch einmal eine Zeit höheren Schwungs, größerer Wirksamkeit und Anerkennung gehabt hat, in Erinnerung jener Zeit seine gegenwärtigen Leistungen überschätzt, etwa gerade in dem Maß überschätzt, wie es seinem Vicar mit den eigenen Leistungen begegnet; schlägt er etwa die Erfolge seines Vicars, den Beifall, welchen er findet, die Kraft, welche er entwickelt, zu gering an; durchschaut er vielleicht den Hochmuth, die Eitelkeit, die Selbstsucht des Vicars oder frißt Neid und Ungunst das Herz „des Alten“ an; dann gute Nacht, Pfarrer, gute Nacht, Friede und Freude! Die Stellung eines solchen Jünglings zu einem solchen Manne oder Greise ist dann eben so schmerzig, wie die einer schlimmen Schnur zu einer schlimmen Schwiegermutter, und das aus ähnlichen Gründen. Wer will dann aber solchen Leuten helfen?

28. Wie gut wäre es, wenn Pfarrer, die Vicarien haben müssen, gegenüber diesen diejenige Demuth beweisen, welche alte Jünger so herrlich ziert und den Jünglingen Achtung und Ehrerbietung abgewinnt! Nichts jämmerlicher als wenn die Alten nicht alt, nicht schwach sein, nicht aufhören, nicht zurücktreten, nicht sterben wollen, und eben damit beweisen, daß sie die Lection nicht gelernt haben, die sie im langen Leben lernen sollten. Ihr ergrauendes Haar, ihr matter Leib sollte sie mächtig aufschrecken und ihnen zu Gemüth führen, daß es nun Zeit sei demüthig zu werden und in Demuth die Aufgabe des Lebens zu vollenden. Je länger Leben, je schwerere Last der Sünden, je demüthigendere Erkenntniß der alten Sünden und Gebrechen, je weniger Ansprüche, je mehr Genügen an der
\quellseite{75}
Gnade, je mehr Sehnsucht nach Verklärung und ewiger Vollendung, je größere Mildigkeit gegen jedermann, je mehr Hoffnung für das zukünftige Geschlecht, je willigere und völligere Anerkennung jugendlicher Thätigkeit und Tüchtigkeit, je größere Nachsicht gegen jugendliche Eitelkeit, gegen jugendlichen Uebermuth! So sollte es sein! Und daß es so sein sollte, sollte man auch erkennen. Und wenn man sich zu solcher Erkenntniß und zu solchem Verhalten unwillig fände, sollte man vor eigenem Hochmuth erschrecken und um so tiefer aus der Seele um Demuth, um Licht und Gnade schreien, damit man in dem Verhältniß zum Vicar eine Schule täglichen Sterbens und Wiederauferstehens fände und sich fröhlich und getrost in sie und ihre gute Lehre begäbe. Ach, wie viele Alte übersehen das, daß ihnen im Verhältniß zu den Jungen die letzte Gelegenheit und Schule der schönsten Vollendung vor ihrem Hingang dargeboten wird! Ach wie viele alte Pfarrer benehmen sich gegen ihre Vicarien so, daß diese, falls sie fromme Jünglinge sind, sich daraus nichts Besseres abnehmen können, als die Lehre und den Entschluß: „So machst du’s dermaleins in deinem Alter nicht.“ — Gott schenke doch seiner Kirche und ihrem Lehrstande demüthige Greise!

29. Aber auch andererseits, wie viel, wie viel versehen die Vicarien! Wie gut wäre es, wenn diese beherzigten, was geschrieben steht: „Es ist ein köstlich Ding geduldig sein und auf die Hülfe des HErrn hoffen. Es ist ein köstlich Ding einem Manne, daß er das Joch in der Jugend trage!“ Klagl. 3, 26. 27. Wie gut wäre es, wenn sie muthig das Joch der Unterordnung auf sich nehmen
\quellseite{76}
und die Erfahrung machten, daß es nicht schwer, sondern sanft ist! So oft ich des Verhältnisses gedachte, in welchem der Diaconus Laurentius gegen seinen Bischof Sixtus stand, und des herrlichen Gesprächs, welches beide auf dem Todesgang des Bischofs miteinander hatten, in welchem sich kindliche Liebe und Vaterliebe so schön ausgesprochen hat; gedenke ich auch des Verhältnisses, das zwischen Vicar und Pfarrer ist, und wünsche nichts mehr, als daß zwischen allen Vicarien und Pfarrern ein gleiches Zwiegespräch der Herzen und der Lippen möge sein können und wirklich sein. Quo progrederis sine filio pater? rief Laurentius dem Bischof zu, welcher zum Tode gieng. Quo, sacerdos sancte, sine ministro properas? Noli me derelinquere, pater sancte, quia thesauros tuos jam expendi, quos tradidisti mihi. Quo progrederis? Tu nunquam sine ministro sacrificium offerre consueveras. Quid in me displicuit paternitati tuae? Numquid degenerem me probasti? Experire, utrum idoneum ministrum elegeris, cui commisisti dominici sanguinis dispensationem etc. Solche innige Kindesliebe zu dem Greise fand auch einen herrlichen Wiederhall aus dem Herzen des Bischofs. Non ego te desero, fili, neque derelinquo, sed majora tibi debentur pro Christi fide certamina. Nos quasi senes levioris pugnae cursum recipimus: te autem quasi juvenem manet gloriosior de tyranno triumphus. Post triduum me sequeris sacerdotem, Levita! — Ach wie ganz anders klingt oft Rede und Gegenrede unter Pfarrern und Vicarien! Da ist keine Achtung vor dem Alter im Herzen, keine Ehrerbietung, keine Zartheit, keine Rücksicht noch Schonung,
\quellseite{77}
keine Liebenswürdigkeit, keine Freundlichkeit, keine Kindlichkeit im Benehmen der Vicarien. Nirgends ist es ihnen langweiliger als in der Gesellschaft des Pfarrers: sonst gesprächig, verstummen sie hier; da findet sich trotz alles Besinnens kein Gegenstand der Unterhaltung, da quält man jedes Wort hervor; es ist viel, wenn sich nach und nach eine sein sollende Höflichkeit und kalte Förmlichkeit einführt; oft ist auch das nicht; Unwille, Widerwille und Entfremdung nehmen oftmals zwischen beiden Platz. Und wie im Hause, so im Amte. Da vergißt der Vicar, daß er Vicar ist, und redet, als wäre die Gemeinde seine Gemeinde, als wäre er ihr Vater, als amtierte er im eigenen, nicht in fremdem Kreise, als wäre er der Anfänger alles Glaubens in der Gemeinde, der einzige, der erste Evangelist, der seinen „Alten“ in allen Stücken weit übersähe, in ihm nur Hinderniß alles Gedeihens fände und alles und jedes ganz anders und besser machen würde, wenn er nur dürfte, wenn er nur nicht immer gehemmt würde \&c. Wie oft, wie oft kann man im Leben zu dieser Schilderung die Belege finden! Wie schwere und große Vicariatsünden lädt sich mancher Jüngling auf! Wie groß Unrecht geschieht oft den Alten! Welch eine Saat der Rache säet damit die Jugend für das eigene Alter! Leider findet dann der unbescheidene, unverständige Uebermuth des Vicars oft noch Beifall bei der Gemeinde und bei jenen Vorlauten, die immer dem aufgehenden Gestirn sich zuwenden! — Ach, es ist so viel Trübes und Betrübendes im Leben, und es könnte doch so gar viel schöner sein, wenn nur jeder seinen Beruf erkennen und innerhalb desselben
\quellseite{78}
ganz einfach sein Christentum beweisen und bewähren möchte. Kein Christ, geschweige Gottes Diener ist der Vicar oder Pfarrer, der nicht nach innigster Eintracht mit seinem Nächsten aus voller Seele strebt! — Gott schenke uns doch wahre Christen für das heilige Amt, Christen zu Pfarrern und Vicarien!

\section{V. Pfarrverwesung}

30. In manchen Ländern werden vacante Pfarreien bis zur Wiederbesetzung von benachbarten Pfarrern versehen. Ohne allen Zweifel kann auf diese Weise eine vacante Pfarrei nur sehr dürftig versehen werden, während doch zugleich auch die Pfarrei, welche sich mit der vacanten in ihren Pfarrer theilen soll, zu kurz kommt. In andern Ländern werden vacante Pfarreien durch eigens aufgestellte Pfarrverweser versehen, und man wählt, wie sich das von selbst ergibt, solche Pfarrverweser aus der beweglicheren Classe der Geistlichkeit, also aus der Zahl der Candidaten. Es wird sich mit nichten leugnen lassen, daß besonders aufgestellte Pfarrverweser den Gemeinden mehr Segen bringen, als angestellte benachbarte Pfarrer, denen ihre eigene Pfarrei schon Arbeit genug darbietet. Aber eben so gewiß ist es, daß ein Verweser doch den Pfarrer nicht ersetzt, daß in dem Alter des Verwesers und in der kurzen Dauer der Verwesung, sowie im Worte „Pfarrverwesung“ wesentliche Bedenken gegen eine zu große Begünstigung der Verweserei als Institut gegründet sind.

\quellseite{79}
31. Wenn der Pfarrverweser ein kräftiger Jüngling und namentlich ein lebhafter und angenehmer Prediger ist,\footnote{Vergleiche, von lebhaften und angenehmen Predigern, Dr. Heinr. Müllers Geistliche Erquickstunden. Nr. 197.} so pflegt er viele und vergleichsweise ungleich mehr Gunst und Beifall zu finden, als ein gleich kräftiger und begabter Pfarrer. In Stadtgemeinden zieht ein begabter Predigerjüngling die weibliche Jugend und überhaupt die weibliche Bevölkerung auf eine Weise und in einem Maße an, welche die Schicklichkeit beleidigen und um deren willen man zuweilen erröthen möchte. In Landgemeinden gewinnt er überhaupt leicht die Mehrzahl der Bevölkerung, und zwar aus guten, oder vielmehr schlechten Gründen und Ursachen.

Dem Pfarrer reicht der Landmann in der Regel seine Gebühren mit einer gewissen invidia, welche bei dem Verweser wegfällt. Ja dieser darf, ohne Ungunst zu erregen, auf Entrichtung der Gebühren dringen, weil er sie nicht für sich einnimmt, sondern verrechnen muß. Dagegen ladet der geringe Sold des jungen Mannes zu Geschenken ein; — wem man freiwillig gibt, den nimmt man eben damit in Schutz und Protektion, und hat man gegeben, so ist man Parthei dessen, dem man „Verehrung“ gethan hat. Je mehr man gibt, desto mehr hängt man dem Manne an. (S. Micha’s Priester. Richt. 17, 7–13.) — Sodann bleibt ein Pfarrer länger und gibt weniger Aussicht auf Veränderung; er tritt in viel engere Berührungen mit der Gemeinde, lernt Leute, Verhältnisse, Schäden und Sünden viel genauer kennen; er kann weniger betrogen und leichter
\quellseite{80}
enttäuscht werden; seine genauere Kenntniß der Dinge macht ihn ungeneigt zu loben und gibt seiner Bußpredigt mehr Schein der Absichtlichkeit: sein ganzes Benehmen trifft mehr, schneidet eher und tiefer in die geheimen Wunden und thut leichter weh 2c. 2c. Ganz anders ist es bei einem Verweser. Er setzt sich nirgends fest, er ist ein Gast, und je mehr seine Talente in unbefangener jugendlicher Anmuth auftreten, desto freundlicher nimmt man seinen Eifer auf; die Zeit seines Bleibens und der Mangel an Urtheil und Erfahrung verhindern ihn, die Leute, die Verhältnisse, die Schäden und Sünden derselben kennen zu lernen; man fürchtet ihn nicht, wie den Pfarrer, weil er in die geheimen Gebrechen und Sünden der Gemeinde nicht eingeweiht ist, nicht eindringen kann; er kann auch bei gutem Willen, schon aus Mangel an Zeit manches nicht anfangen, was die Gemeinde verderben, ihr etwa ein fälschlich angemäßtes Recht oder Gut gefährden könnte; er kann leichter über pfarrliche Verhältnisse ausgeholt, leicht ins Interesse gegen den vorigen, auch wohl nachfolgenden Pfarrer gezogen, während seiner Amtsführung kann etwa leichter etwas erreicht werden, was dem Eigennutz der Gemeinde dient; er hat weniger Interesse, vorsichtig gegen die Gemeinde zu sein. Kurz, die Wirksamkeit eines Verwesers kann, mag sie auch mächtig eingreifen, doch so tief nicht gehen, als die eines bleibenden Hirten. So ermunternd, erfrischend, erquickend und weckend sie auch sein, und so hoch sie etwa deshalb geschätzt werden kann; eine gewisse Oberflächlichkeit bleibt von dem Leben und Wirken des Pfarrverwesers doch immer unzertrennlich. Und gerade das ist’s, was viele,
\quellseite{81}
oft die schlimmsten, Gemeinden so zufrieden mit Verwesern macht. Hat ein Verweser nun überdies eine gute Stimme, hübschen Vortrag, lebhafte, konkrete, eindringende Gedanken, persönlich gewinnendes Benehmen; so kann sich die Gemeinde im eigentlichen Sinn in einen solchen jungen Geistlichen „verlieben“, alle Leidenschaften gehen zu Gunsten des Lieblings empor, und man sieht mit Schmerz der Zeit seines Abzugs entgegen. Man wünscht ihn statt des designierten Pfarrers behalten zu dürfen; es fehlt nicht an Bitten bei Konsistorien und Patronen, ihn zum Pfarrer zu bekommen, — und schlägt aller Wunsch und alle Bemühung fehl, so thut man sich wenigstens selbst auf alle mögliche Weise Genüge. Man bedeckt den jungen Mann mit Geschenken und Verheißungen, mit Lobeserhebungen und Liebeserweisungen. Und daß dann der Zeitungsabschied nicht fehlen darf, versteht sich von selbst. Eine Zeitung ist ein öffentlicher Ort für alle Rendezvous, die gerne bemerkt sein wollen. — So gibt’s denn freilich einen harten Anfang für den neuen Pfarrer, der zuweilen in seinem Schicksal keinen bessern Trost, als den der Gerechtigkeit finden kann. Denn auch er hatte als Pfarrverweser seinen nachfolgenden Pfarrern, vielleicht ganz unschuldig, dieselben Leiden bereitet und leidet nun billig ein Gleiches. Jedoch gibt es auch noch einen (schlechten) Trost für den neu eintretenden Pfarrer. Empfängt ihn seine Gemeinde ziemlich kühl, so freut er sich eines desto herzlicheren Entgegenkommens seiner Kollegen und Amtsnachbarn, für welche ein junger, kräftiger Verweser oft eine harte Geduldsprobe ist und, zumal wenn er das Unglück hat, sie durch leuchtende Gaben zu übertreffen,
\quellseite{82}
auf sie wie eine Rakete wirkt, die man unter einen Haufen Leute wirft.

32. Was die Verweser selbst anlangt, so hat die Pfarrverwesung auch für sie zweierlei Seiten. Es läßt sich nicht leugnen, daß Pfarrverwesungen vielfache Gelegenheit verschaffen, Erfahrungen zu sammeln, und daß junge Männer, die zuerst in einer Schule, dann in einem Biennium, wie es oben beschrieben ist, und zuletzt in einem eigentlichen Vicariate gearbeitet hätten, keine angemessenere Stufe für ihre weitere Ausbildung finden könnten, als eine Verwesung. Nicht selten feiern junge Männer in Verwesungen die eigentliche Blüthezeit ihres geistlichen Lebens und Wirkens, — und die Schule, welche sie in ihnen durchmachen, ist um so reizender, weil sie mit einer Selbstständigkeit durchgemacht wird, die andern Schulen abgeht. Es ist eine Schule, in welcher man sich als Meister geriren kann. Und hier eben liegt auch alle Gefahr. Ein Pfarrverweser hat dieselben Amtsrechte und Amtspflichten, welche der Pfarrer selbst hat, nur nicht auf so lange Zeit. Zwar geht ihm Erfahrung des Lebens- und Amtes ab, aber da ein Pfarrverweser selten in die schwereren und mühsameren Regionen des Amtes kommt; so langt er bei sonstiger Gabe mit seiner Wenigkeit doch aus, und der Erfahrungsmangel macht sich ihm in der Regel nicht so fühlbar, daß er deshalb in Gewissensnöthen käme. So pflegt dann ein Pfarrverweser des Amtes irdische Wohlfahrt und zeitliches Gelingen zu genießen, aber auch unter starker Versuchung zum Hochmuth zu wandeln. Wie mancher junge Geistliche hat als Verweser den Grund zu seinem selbstgenügsamen Leben und
\quellseite{83}
trotzigen Dünkel gelegt, der ihn hernachmals untüchtig und zugleich unglücklich gemacht hat, wenn er den Applaus verstummen hörte und den rechten Grund in sich nicht fand!

33. Pfarrverwesern kann daher nichts Besseres gerathen werden, als daß sie sich vor allen Dingen mit dem genügen lassen, was ihr Beruf mit sich bringt, und keine Schranke übertreten, die ihnen Amt oder Tugend steckt. Sie füllen eine Lücke aus zwischen einem Vorgänger und Nachfolger und werden selbst nicht Nachfolger der Pfarrer genannt. Beschränkung auf das Nothwendige, Arbeit in Wort und Lehre in der jeder Gemeinde gewohnten Weise, in dem jeder Gemeinde eignenden Maße, Enthaltung von folgenreichen Unternehmungen 2c. gehört zum Charakter einer Verwesung. Der heilige Respekt, welchen ein Verweser vor dem innigen Verhältnis eines Hirten zu seiner Heerde hat, soll in all seinem Thun und Lassen hervortreten. Dazu gehört eine heilige Pietät gegen den abgegangenen und gegen den nachfolgenden Pfarrer. Diese Pietät im eigenen Herzen tragen und sie der Gemeinde einprägen, ist seine Schuldigkeit und außer der reinen Lehre die beste Frucht, die er bringen kann und bringen soll, wo nur immer höhere Gründe und Rücksichten nicht anderes erheischen. Deshalb hat ein Verweser in allen amtlichen Verhältnissen genaue Erkundigung einzuziehen, wie der abgegangene Pfarrer es gehalten, und ist es nicht gegen Wort und Amtsseid, so schließe er sich derselben Art und Weise an. Jedenfalls hat er Mißtrauen in die eigene Meinung zu setzen, mit aller Vorsicht über den abgegangenen oder künftigen Pfarrer sich zu äußern und die Achtung
\quellseite{84}
vor der Amtsführung desselben für Tugend zu erkennen. Neues aufbringen, ziemt einem Verweser nicht; kann es immer sein, so lasse er alles beim Alten; ist es nicht möglich, so treffe er solche Einrichtungen, welche dem neuen Pfarrer nicht zu nahe treten, ihn nicht beschränken, es sei denn, daß er irgend etwas, was in der Kompetenz eines Pfarrers steht, nach Uebereinkunft mit dem designierten neuen Pfarrer thäte und ordnete. Sei er in seinem Amte gesegnet, wie immer es sei; so vergesse er nie, daß es seine eigenthümliche Zier und Aufgabe ist, als ein Gast, als ein Davoneilender, der sich nicht einbürgern darf, mit aufgeschürztem Kleid, mit dem Wanderstab in der Hand zu wirken — und die Gemeinde und deren Herzen dem neuen Pfarrer im möglich besten Zustande zu überliefern. — Ist eines Vicars Demuth die Demuth eines Sohnes gegen den Vater; so ist die Demuth eines Verwesers Johannesdemuth und ihr Spruch heißt: „Ich muß abnehmen.“ Damit setzt er sich selbst das schönste Denkmaal.

\section[VI. An welche Pfarrstelle]{VI. An welche Pfarrstelle bewirbst du dich?}

34. Gregor der Große sagt in seiner Regula pastoralis (P. I. C. 7.): Moyses suadente domino trepidat, et infirmus quisque, ut honoris onus percipiat, anhelat; et qui ad casum valde urgetur ex propriis, humerum libenter opprimendus ponderibus submittit alienis; quae egit ferre non valet, et auget, quod portat. Wie er klagten so viele andere auch in der protestantischen Kirche, und
\quellseite{85}
was sollte man erst heut zu Tage zu klagen haben, wenn man sich aufs Klagen einließe! Es ist ein ungeheurer und unverantwortlicher Mißbrauch, welcher mit den Meldungen getrieben wird.\footnote{Zu den traurigsten Seiten unserer luth. Landeskirchen gehören ohne Zweifel die Systeme der Anstellung nach Anciennetät der sich Meldenden. Bequem sind sie, aber eine Quelle unsäglicher Uebel! Eine starke Behauptung, zu deren Aufrechthaltung man aber Bücher voll Beleg und Beweis schreiben könnte! Möchten diese Gebäude dahinfällen! — Nach alter Anschauung ist der Pfarrer Bräutigam, die Gemeinde Braut, der Amtsantritt eine Hochzeit, bei welcher die Braut dem Bräutigam angetraut wird. Da wechselt man nun die Ehe und Frau so oft, um so verhältnismäßig geringer Ursachen willen, und so wenig männliche Liebe und Treue ist vorhanden, daß es kein Wunder ist, wenn auch die Gemeinden nichts mehr von weiblicher, treuer Liebe übrig haben. Wer sein Weib so leicht und oft verläßt, kann nicht bloß wegen Kürze der Zeit, sondern auch wegen wandelbarer Gesinnung keine rechte Wirksamkeit finden. — O daß Gott Buße den Hirten gäbe — in allen Dingen — sonderlich in diesem! — Weide dich, aber vergiß nicht, daß der Erzhirte Rechenschaft fordern wird.} Dennoch aber hebt er den rechten Gebrauch der Meldungen, welcher in 1. Tim. 3, 1. fest gegründet ist, keineswegs auf, zumal in einer Zeit, in welcher es vor lauter Meldungen dahin gekommen ist, daß kaum jemand mehr ohne Meldung zu einer Berufung kommt. De clerico renitente zu schreiben, ist wohl für unsre Verhältnisse ziemlich überflüssig (Dr. Ziegler schrieb darüber); wohl aber wäre es gut, wenn der Eifer, sich um eine Pfarrei, und, wenn man schon eine besitzt, weiter zu melden, durch das ἐπιλαμβάνων εἶναι, welches 1. Tim. 3, 2. auf das ὀρέγεται v. 1. folgt, ein wenig gemäßigt und abgefühlt würde. (S. in Gregors Regula das hieher schön passende Capitel.)

\quellseite{86}
Jeden Falls soll hier diesem Mißbrauch der Meldungen kein Lob gesprochen werden, indem wir nun einige Fragen in Anbetracht der Pfarrstellen, um die man sich melden kann, zu lösen suchen.

35. Soll man sich um eine Stelle melden, bei welcher man nicht allein im Amte ist, bei welcher man Collegen hat? Collegen im Amte sollten zwar immer die besten Collegen sein, und, wenn irgend unter Menschen, sollte unter ihnen die Bruderliebe blühen und Früchte tragen. Allein es steht nun einmal nicht wie es sollte, und Psalm 133 pflegt selten von Collegen in Gemeinschaft gesungen zu werden. Die Liebe muß, wenn Friede sein soll, von zwei Seiten kommen, und eine Liebe von zwei Seiten her, wie selten ist diese! (Es ist eine Schmach, daß man so sagen muß.) Darum sorgt ein Weiser, daß seine Liebe nicht zu sehr in Versuchung komme, und ein kluger Mann traut sich selbst nicht zu viel. Wer sich um eine Stelle melden will, bei welcher er Collegen hat, der wird darum wohl thun, sich nicht bloß um die Beschaffenheit der vorhandenen Collegen zu erkundigen, sondern sich vor allem zu überzeugen, ob die den Collegen zustehenden Rechte und Pflichten ausgeschieden und genau bestimmt sind oder nicht. Wenn nicht, so ist es mißlich, eine solche Stelle anzunehmen, da Ungewißheit der Rechte Streit zu erregen und die Liebe zu ertöden pflegt. Sind aber die gegenseitigen Rechte und Pflichten wohl abgegrenzt, oder kann man sicher erkennen, daß sich alles im Frieden ordnen lasse, bevor man die Stelle antritt; so kann ein Mann, der sich der Liebe und Friedfertigkeit bewußt ist, der gerne
\quellseite{87}
demüthig und der andern Last zu tragen bereit ist, es wohl wagen, College zu werden. Der HErr kann ihn ja behüten vor dem Uebel und Unglück der Zwietracht, welche so viele edle Seelen um alle Freudigkeit und Heiterkeit bringt, sie vor der Zeit alt und grämlich macht! Er kann auch mehr als das, er kann wahre Bruderliebe und tiefen Frieden, völlige Eintracht schenken. „Wenn jemands Wege dem HErrn wohlgefallen, macht er auch seine Feinde mit ihm zufrieden,“ so lesen wir Spr. 16, 7.; warum sollte der HErr Collegen die richtig wandeln, nicht miteinander zufrieden machen können?

\subsection[36. Stadt oder Land]{36. „Soll man sich in die Stadt oder aufs Land melden?“}

Auch das ist keine unwichtige Frage. Wer die Wirksamkeit in einer Stadt vorzieht, muß gewöhnlich auch Collegen „mit in den Kauf nehmen“ und sich gefaßt machen, die Pein jener Rivalität auszustehen, von der wir bereits andeutende Worte gesprochen haben. Schon um des willen fürchtet sich mancher vor „Stadtstellen.“ Aber es ist das nicht der einzige Grund der Furcht; man pflegt auch sonst die Wirksamkeit in der Stadt für schwieriger zu halten als die auf dem Lande. Während das Leben eines Landpfarrers fast idyllisch geschildert wird, erkennt man die Aufgabe eines Stadtgeistlichen für besonders mühevoll und schwer. Und doch ist sehr die Frage, wessen Aufgabe die schwierigere sei, und ob sich nicht am Ende das Zünglein der Waage mehr auf die Seite des Landpfarrers neigen wird? Wer beides versucht hat, in der Stadt und auf dem Lande das Amt zu führen, wird beistimmen. Man kann dies zu behaupten wagen, auch wenn andere
\quellseite{88}
die gegentheilige Meinung haben sollten, auch wenn sie dieselbe, wie z.\,B. Hüffell, mit Stellen aus Göthe und Rousseau belegen. Ein träger Pfarrer oder ein rosiges Gemüth, das alles gern rosig ansieht, sich ungern aus dem Traume jugendlichen Schwärmens reißen läßt, kann wohl einmal einen Sommer auf dem Lande geruhlich und gemächlich leben, aber es kommt auch der Winter — und wen auch der nicht von den idyllischen Ansichten vom Landpfarrerleben heilt, den werden über kurz oder lang die Leute heilen, wenn sie in den Fall kommen, sich geben zu wollen, wie sie sind. Insgemein sieht es auf dem Lande aus, wie es etwa Pestalozzi im ersten Theil von „Lienhard und Gertrud“ beschreibt und mans aus dem schon älteren Büchlein A. Hartmanns, dem kleinen Pastorale für Landpfarrer,\footnote{„Unvorgreifflich-Einfältig- und wohlgemeinter Entwurf, wie ein Dorfpfarrer seiner anvertrauten Gemeinde erbaulich vorstehen möge, Christo Jesu, dem Erzhirten unsrer Seelen zu Lob, Ehr- und Preiß, und redlichen Unterhirten zu Nutz und Dienst in der Furcht des HErrn, ehemals verfaßet, anjetzo aber an das Licht gegeben v. M. Andreas Hartmann, hochfürstl. württemb. Pfarrer in Truchtelfingen. Ulm, im Verlag Geo. Wilh. Kühns. 1710.“} ersehen kann. (Vgl. A. Hartmanns „Gespräch vom Landleben.“)

Wer eine Pfarrstelle in der Stadt antritt, muß sich seine ganze Stellung erst erobern, nichts ist fertig. Er kommt, er predigt, er wird gehört, sein Publikum bildet und sammelt sich völlig frei, sein Beichtstuhl wird gesucht oder gemieden, je nachdem ein jeder will. Da gilt ein jeder Pfarrer so viel er kann, und es kann freilich recht traurig und betrübt hergehen, wenn man keinen oder wenig
\quellseite{89}
Anklang findet. Aber es ist doch immer leichter, sich in einen kleinen Kreis von Beichtkindern zu finden und zu fügen, als ein Beichtvater von Unwilligen, Widerwärtigen und Feinden werden zu sollen. Das aber ist auf dem Lande der Fall und kann nicht wohl ausbleiben. Der neue Landpfarrer muß so und so viel ihm zuvor ganz unbekannte Pfarr- und Beichtkinder annehmen, und sie müssen ihn nehmen, mögen sie ihm geneigt sein oder nicht, mögen sie Vertrauen zu ihm haben oder keines. Da gilt es dann eine schwere Aufgabe gegenseitiger, oft höchst widerwärtiger Gewöhnung, über deren Lösung viel Zeit vergeht, manch übles Wort geredet, manch unlauterer Schritt hin und her gethan wird.

Die elende Köderung und Verführung der Beichtkinder, welche ein junger Stadtgeistlicher von seinen Collegen zu erfahren hat, und all die Pein und Reiberei, die er da zu genießen bekommt, sammt allem andern städtischen Ungemach dürfte von geringerem Belang gegen die Schwierigkeit sein, das Vertrauen einer Gemeinde zu gewinnen, die aus oft so überaus verschiedenartigen Elementen zusammengesetzt ist.

Wer eine Pfarrstelle in der Stadt übernimmt, muß sich auf die Stadtsitte verstehen, und wissen, wie man mit den Leuten umzugehen hat; allein er kann sich doch die Stadtsitte, welche sich fast allenthalben gleich oder doch ähnlich findet, leichter aneignen. Lebte er doch auf dem Gymnasium, auf der Universität in ihr heran, so daß er sie, auch wenn er vom Lande gebürtig wäre, doch nicht erst bei Uebernahme der Pfarrstelle als etwas ganz Neues kennen lernen
\quellseite{90}
muß. Dagegen hat der Landmann auch „seinen Comment“ und zwar einen sehr hartnäckigen und in verschiedenen Gegenden sehr verschiedenen. Wie soll ein Pfarrer, der nicht gerade in seiner Pfarrei aufgewachsen ist, die Sitte und den Brauch seiner Gemeinde kennen oder kennen lernen? Vor ihm sucht sich Sitte und Brauch vielfach unsichtbar zu machen; schämisch tritt der Landmann vor dem Pfarrer zurück, expliciert sich nicht, kann’s auch oft nicht, weiß manchmal nicht, warum er vom Pfarrer nicht erkannt noch verstanden wird, da er doch sich und seines Gleichen von Jugend auf auch ohne Rede und Erklärung verständlich ist. Wahrlich ist es nichts Leichtes, nur z.\,B. den Dialekt und Redebau des Landmanns zu verstehen; und noch weniger ist es leicht, sich in’s ganze Wesen und Leben einer Landgemeinde einzuleben, sie zu verstehen, und, wohl gemerkt, ihr wieder verständlich zu werden, ohne daß man Ton oder Art verläßt, wie sie nun einmal dem Pfarrer gegenüber der Gemeinde gebühren. So stehen sich Sitte und Heerde fremd und ferne gegenüber und es erwachsen dem Amte aus diesem Verhältniß gar viele Hindernisse.

Wenn man gewissen Leuten glauben sollte, so wäre die Bildung des Städters ein Hinderniß, die Einfalt des Landmanns ein Förderniß des heiligen Amtes. Allein die Bildung, welche der Städter vor dem Landmann voraus hat, ist eine formale, welche im Gegentheil dazu hilft, daß des Pfarrers Sprache und Ausdruck leichter verstanden wird. Gerade dies aber, nicht mehr und nicht weniger, geht dem Landmann ab im Vergleich mit dem Städter, — also ein Vortheil. Uebrigens ist der Landmann mit seinem
\quellseite{91}
Lebenskreise ganz außerordentlich vertraut, klug und listig, nimmt es in schlauer Wahrung seines Vortheils mit vielen Städtern auf, bleibt hartnäckig dabei, nicht verstehen zu wollen, was er nicht versteht, wird darum nicht leicht betrogen, und verhält sich dem Städter gegenüber ungefähr wie der Jude gegenüber dem Christen, nämlich vermögen und oft überlegen in der Selbstsucht. Es ist nicht zu sagen, wie sehr der Landmann die niederen Regionen des Lebens kennt und beherrscht, wie er alles zu benützen weiß, und wie gleichgültig es ihm insgemein ist, ob er Sünde und Schuld auf sich häuft oder nicht. Betrug, Diebstahl, Lüge, Meineid — ach, und welche grobe Sünde müßte man ausnehmen? — findest du bei dem Landmann geschickt oder ungeschickt verdeckt. Verschmitzt und pfiffig, daß er nicht ertappt werde, geht er die bäuerliche Bahn, — und seine ganze Einfalt wird, manche tröstliche Ausnahme zugestanden, meist nur zur Außenseite. Und nun da durchdringen, bei einem Publikum durchbringen, das sich den Pfarrer nicht gewählt, das ihn mit Mißtrauen aufgenommen hat, das ihn mit scharfem Aug’ und lauschendem Ohr betrachtet! Wahrlich, das ist eine sehr schwierige, große Aufgabe, zu deren Lösung es großer Gaben des Geistes und Gemüthes und einer mächtigen Geduld bedarf, die man ohne die Kräfte der zukünftigen Welt, die im Worte wirken, unmöglich nennen müßte. Weit leichter scheint ein einseitig begabter Mensch in Städten zum Ziele kommen zu können, als auf dem Lande. Dort leichter wie hier gewinnt man wahre, ungeheuchelte Anerkennung. Etwa bedarf es dort mehr äußerliche Gewandtheit in Rede und Benehmen,
\quellseite{92}
aber Gewandtheit des Geistes, Bewußtsein des Zieles und der Mittel, Consequenz und Beharrlichkeit, Stärke, Geduld und Langmuth, kurz alle Mannestugenden bedarf der Landpfarrer in ausgezeichnetem Maße; dazu auch, schon um seiner selbst willen, viel Freude und Lust an heiligen Studien, da ein immerwährendes Wirken unter einem rohen Volke nothwendig eines starken Gegenmittels bedarf, wenn man nicht von demselben unbewußt lernen und annehmen soll, was sich nicht schickt.

Gewiß könnte man noch gar manches über die schwierige Stellung eines Landpfarrers sagen; es sei aber an dem genug, und ein Verständiger nehme sich daraus ab, was er thut und wählt, indem er sich um eine Landpfarrei bewirbt. — Ueberhaupt, welche Pfarrstelle man auch erwähle, man beachte, überlege, erwäge alles. Es ist unter den Pfarrstellen von gleicher Art oft doch noch große Verschiedenheit. Da prüfe ein jeder vor allem sich selbst, ob er auch für die Pfarrstelle passe, um die er sich melden will? ob er in die Verhältnisse tauge? ob er nicht unglücklich werden und unglücklich machen wird? Denn daß jede Stelle den rechten Mann und jeder Mann die rechte Stelle bekomme, hat jeder seines Orts und Antheils gewiß zu sorgen.

\quellseite{93}
\trenner

\begin{quote}
Nulla ars doceri praesumitur, nisi intenta prius meditatione discatur. Ab imperitis ergo pastorale magisterium qua temeritate suscipitur, quando ars est artium regimen animarum? Quis autem cogitationum vulnera occultiora esse nesciat vulneribus viscerum? Et tamen saepe qui nequaquam spiritalia praecepta cognoverunt, cordis se medicos profiteri non metuunt; dum qui pigmentorum vim nesciunt, videri medici carnis erubescunt etc. etc.
\end{quote}

\quelle{Gregorii M. Regula pastoralis.\\P. I. C. I.}

\trenner

\quellseite{94}
\chapter{Erstes Auftreten im Amte}

\quellseite{95}

\quellseite{96}
\section[I. Du gehörst allen Pfarrkindern]{I. Du gehörst allen deinen Pfarrkindern}

37. Die Gemeinden, wie sie heut zu Tage zu sein pflegen, sind aus mancherlei widersprechenden Bestandtheilen zusammengesetzt. Sollte man sie classificieren, so könnte man sie vielleicht in zwey Hauptklassen theilen: in die Klasse derjenigen, welche Religion und Kirche achten, und in die Klasse der Verächter und Feinde. Jede Hauptklasse zerfällt dann aber wieder in mancherlei Unterabtheilungen. So viele Parteien sich nun aber auch in einer Gemeinde finden mögen, die Pflicht eines Pastors, welcher einmal eine solche Gemeinde zu übernehmen den Muth hat, ist es, alle für seine Pfarrkinder zu erkennen. Sein Auge unterscheide und erkenne jede Partei, ja, wenn es möglich ist, jede einzelne Seele; sein Herz, seine ungetheilte Liebe, sein großer Fleiß gehöre allen und jeden. Gleichwie ein Hirte Böcke und Schafe, Schafmütter und Lämmer, kranke und gebrechliche Schafe, allerlei Schafe zu seiner Heerde zählt und ihnen seine Hirtenliebe zuwendet; so wende auch ein Seelenhirte sein Aug und Herz voll
\quellseite{97}
Liebe den verschiedenartigen Gliedern seiner Gemeinde zu; seine heilige Absicht, welcher er auf das Consequenteste nachjage, sei das Heil aller und jeder. Jedes Schaf sei ihm Zweck, — jedes werde nach Bedürfniß, alle mit gleicher Liebe behandelt. Der Herr, der Erzhirte, läßt neun und neunzig Gefundene und sucht ein Verlorenes; Er liebt nicht bloß die Liebenswürdigen, sondern auch die armen Sünder, die am Weg im Staube liegen. Ihm nach liebt und handelt jeder fromme Seelenhirte. Nichts vermag ihn, von irgend einem seine Liebe abzuwenden, so lange nicht Gottes Wort selbst ein Anderes gebietet. Keine Beleidigung überwindet ihn zum Zorn, sie erregt nur seine Liebe und sein Erbarmen aufs Neue. Von ihm muß man sagen können: „Beleidige ihn, so liebt er dich um so mehr.“ Eine unüberwindliche Liebe muß ihn erfüllen und treiben. — Man könnte dagegen einwenden, daß ein solches Verfahren gegen die in der heiligen Schrift begründete Scheidung zwischen der Bruderliebe und allgemeinen Liebe streite, daß Indifferentisten und Feinde Christi keine Brüder seien, daß ihnen Zucht und Strenge zuzuwenden sei. Allein das alles hebt ja das Gesagte doch nicht auf. Wir haben nicht gesagt, daß man die Feinde Christi als Brüder lieben solle, daß man Bruderliebe und allgemeine Liebe vermengen, daß man die Zucht und Strenge unterlassen solle, wo sie hingehört. Im Gegentheil, wir glauben, daß die Zucht und Strenge selbst nur Formen der Liebe seien, daß die Steigerung der Zucht durch ihre Vermahnungsgrade auch als Steigerung der Liebe zu fassen sei, daß man alle lieben müsse, aber jeden nach Nothdurft, zu seinem Heil, so wie
\quellseite{98}
alle gerettet werden sollen, wenn gleich einem jeden die Rettung auf anderem Wege begegnet. — Die gegenwärtigen Gemeinden ihrer großen Mehrzahl nach kennen keine Zucht; sie wissen von der heiligen Pflicht strafender, züchtigender und eben damit rettender Liebe noch nichts; die herrlichen Stellen in Matth. 18. und 1. Cor. 5. 2c. sind ihren Augen, so klar und deutlich sie reden, dennoch verborgen und dunkel; alles, was Christ heißt, gilt ihnen auch dafür; soll es zu einem Unterscheiden kommen, soll züchtigende, rettende, erziehende Liebe eintreten, so muß erst von den Pfarrern durch Wort und That darauf hingewirkt werden.

Am neuen Pfarrer liegt es, unter dem gnädigen Beistand des h. Geistes eine neue bessere Zeit herbeizuführen. Das aber kann wieder nur langsam geschehen. Der Pfarrer kommt in die ihm neu übertragene Gemeinde; er kennt die Einzelnen nicht; er weiß noch nicht, wo er allgemeine, wo Bruder-Liebe, wo Zucht zu üben hat. So viel aber weiß er, alle einzelnen sind ihm zugewiesen: so umfaßt er denn auch alle mit Liebe, bis sich die Liebe zu bestimmten und mannigfaltigen Liebeserweisungen abklärt. Er ist für alle berufen; er hat den Ruf angenommen; wie sein Beruf, so sein Wille, seine Absicht; wo möglich, alle — wenn das nicht, so viele als immer möglich, zu retten, das ist’s, wozu er in der Gemeinde steht: dahin geht alle seine Kraft, seine Treue, seine Lust, seine Liebe, — und so weiß er von keinem Widerspruch, wenn er sich selbst sagt: „Du gehörst allen deinen Pfarrkindern.“

38. So wie der neue Pastor das Amt in seiner Gemeinde angetreten hat, nahen sich ihm seine Leute. Es naht
\quellseite{99}
die fromme Einfalt, welche das Hirtenamt lieb hat, und jedem neuen Pfarrer in Liebe und Hoffnung das Herz aufschließt: darum hüte sich ein neuer Pfarrer abzustoßen. Es machen aber auch unter den vordersten diejenigen ihre Erscheinung, welche, Pietisten von Partei, anspruchsvoll erwarten, von jedem Pfarrer zuerst bemerkt, bevorzugt und an’s Herz gezogen zu werden. Unter diesen sind Leute, die zu läutern, zu wenden, zu ändern sind, von denen man mit Unrecht die Hoffnung wegwürfe. Darum müssen auch sie ein freundliches Auge und den Wiederklang eines väterlichen Herzens finden. Unter ihnen sind aber auch zuweilen unerträgliche, verzweifelte, unmögliche Menschen, die je länger, je weniger gefallen, — dem Pfarrer, welcher sich ihnen vertraut, je länger je mehr Verlegenheiten bereiten, gern seine Vertrauten spielen, für ihn Partie nehmen, so lang ers für sie thun kann, und giftig stechen, wenn er es nicht mehr kann. Und um ihretwillen würde der neue Pfarrer seine Freundlichkeit und Liebe gegen die Annahenden mit dem würdigen Ernste einer aufrichtigen, offenen Zurückhaltung, — einer „aufrichtigen, offenen Zurückhaltung,“ sage ich, weil man am Ende auf die freundlichste, unverholenste Weise Grund und Ursach der Zurückhaltung offenbaren kann. — Zuweilen nahen sich gerade die besten Gemeindeglieder dem Pfarrer zuletzt. Es ist nicht recht, daß es also ist; die „besten“ fehlen hierin; aber es ist so, und so geschieht es. Trotz ihres Zurückhaltens müssen sie von dem Pfarrer berücksichtigt werden, noch ehe er sie kennt; und weil es ihnen nicht bloß zum Aerger, sondern zum Aergerniß zu gereichen pflegt, wenn ihr Pfarrer sich an die ihnen
\quellseite{100}
wohlbekannten Pfäffischen von Fach anschließt, von denen wir bereits redeten; so muß er schon um ihretwillen mit seinem persönlichen Vertrauen langsam thun. Neben den erwähnten „Besten“ nahen zuletzt auch die Schlimmsten. So schlimm sie sind, wissen sie doch, ja sie vor andern, daß es mit dem Christenthume der Parteileute nichts ist. Gibt sich der Pfarrer diesen hin, so verlieren sie von vornherein auch das persönliche Vertrauen zu ihm, da sie zu seinem Amte ohnehin schon längst kein Vertrauen hatten. Darum auch um der Schlimmsten willen thue der Pfarrer langsam mit seinem Vertrauen. Vielleicht könnte man die Hauptgrundsätze eines richtigen Benehmens in folgenden Worten zusammenfassen:

\begin{quote}
„Mache nicht leicht und am wenigsten gleich anfangs mit einem Pfarrkind Gemeinschaft oder gar Kameradschaft. Vergiß auch in deinem Privatleben nicht, daß du Hirte aller deiner Schafe sein sollst. Du bist ein Mensch, wie andere, bedarfst und brauchst persönliche Liebe, wirst sie auch finden, verlaß dich darauf; du wirst sie am sichersten und süßesten finden, je weniger du darnach jagst, je mehr du ganz Pfarrer, ganz Pastor bist. Nahe in der Regel allen nur amtlich, bei den Gelegenheiten des Amtes, mit den Gnaden des Amtes; alle persönliche Liebe ergieße sich in dein amtliches Annahen, so daß du als der herzlichste Freund erscheinst, wenn du in Jesu und der Seelen Geschäften kommst. Stehe als Bischof, als Hirte, Hüter, Wächter, über allen und verlaß deinen Standpunkt nicht, damit du allen, einem jeden zu seiner Zeit nahen, allen allerlei sein könnest. Benimm dich so, daß keiner um der Vorliebe willen, welche du dem
\quellseite{101}
oder jenen erzeigst, dir abhold werde und dir in Zeiten der Noth das Herz verschließe. So benimm dich, daß alle je länger je mehr an deine Liebe, an deine liebevolle Weisheit, an deine liebevolle, weise, starke Männlichkeit und Hirtenwürde glauben können.“
\end{quote}

39. Aus diesen allgemeinen Regeln geht alle Weisheit des Benehmens in speciellen Fällen für einen neuen Pfarrer hervor. Wer sie beobachtet, findet das rechte Benehmen gegen Wohlhabende und Arme, welches so gar oft mangelt. Die Anmahnung des „Nothigen“ und Armen ist unbequem und unbehaglich für den (ach auch für den neuen und besonders für den jungen Pfarrer), welcher selbst gern Gemach und Ruhe eines glücklichen Lebens genösse. Dagegen ist die Anmahnung des Reichen und Glücklichen wohlthuend dem, der gerne glücklich wäre, und um des wohlthuenden, behaglichen Gefühles willen, das mit dem Glücklichen kommt, hält man ihn gern auch für einen Inhaber inneren Glückes, Friedens und eines gerechten Geistes. Davon hüte sich ein junger Seelsorger besonders in seinem Anfang, wo alle Eindrücke überwältigender wirken. Thu langsam mit deinem Vertrauen gegen Reiche wie gegen Arme! — Ebenso ists mit dem Benehmen eines Seelsorgers gegen die, welchen das gute Gerücht wahrer Frömmigkeit vorangeht, wie gegen renommirte, grobe Sünder. Keiner werde als ein Heiliger oder Vollkommener aufgenommen. Aus dem Benehmen des Pfarrers soll sich jedermann noch ein „Vorwärts“ abnehmen und einen Reiz empfangen können, dem noch nicht erreichten Ziele nachzujagen. Freundliche,
\quellseite{102}
brüderliche Ehrerbietung vereinigt sich damit ganz wohl. Anerkennung und der Ruf zu einem heiligen Vorwärts können zusammengehen — beide, zusammen ertönend, werden von denjenigen in ihrem Zusammenhang leicht verstanden, welcher selbst von dem apostolischen Worte: „ich jage ihm aber nach“ durchdrungen ist. — Eben so sehr aufs Ziel gerichtet sei auch das Benehmen des Pfarrers gegen grobe Sünder und Feinde. Ein Seelsorger sieht in einem jeden Menschen, der ihm naht, einen Kandidaten des ewigen Lebens, zwischen welchem und einem Heiligen kein Unterschied ist, als die wunderbare Zeit und Stunde, da es anfieng besser zu werden. Diese Zeit und Stunde kann aber für einen jeden Lebendigen noch kommen. Da ein Seelsorger das weiß, so behandelt er einen jeden Sünder, der ihm naht, mit der Fülle jener Liebe, welche keine Sünde gutheißt, rechtfertigt oder entschuldigt, aber für jeden Sünder noch hofft und betet. Ein unverhofftes Schweigen, wo man eine Strafpredigt erwartete, — ein unerwartet freundliches Begegnen dessen, welchem man gewissenshalber auch eine herbe Strenge zu gute gehalten hätte, — Güte von Seiten eines Mannes, in dessen ganzem Leben und Benehmen man geneigt ist eine verkörperte Predigt zu sehen, haben oftmals noch auf solche gewirkt, von denen man aufgehört hatte zu hoffen. Das ist überhaupt in seinem Maße wahr; es ist doppelt wahr für Anfangstage eines Pastors, in welchen viel zu gewinnen ist, was späterhin nicht mehr gewonnen werden kann.

40. Wohlthuend und ermuntend für Alle wirkt bei einem neuen Pfarrer die in Wort und Benehmen ausgesprochene Voraussetzung, daß sein Pfarrkind das
\quellseite{103}
Gute und den rechten Weg verschmähe, daß nicht böser Wille, sondern Ursachen von minderer Verschuldung den Fortschritt bis dahin aufgehalten haben. Es ist deprimierend und hemmend, wenn vornherein Mißtrauen herrschend erscheint. Es wird sich freilich böser Wille genug überall finden, aber da man nicht weiß, bei wem er ist, und zur Erkenntnis eine kurze Begegnung nicht hinreicht; so ist es keine List, sondern menschliche Wahrheit und Klugheit, bis zu mehrerer und etwa widerlegender Erfahrung alle mit Hoffnung zu behandeln. Für den Anfang jedes Wirkens ist Vertrauen nötig, und der Spruch: „Die Liebe hofft Alles“ hat am Anfang einer Laufbahn seine natürlichste Stelle. Versteht sich, ist dies Hoffen nicht ungefesselt, — nicht ohne Unterschied, — nicht sanguinisch, — nicht auf Menschliches, sondern auf Gottes Gnade bauend; es besteht mit der Erwartung des Menschenkenners, auch in der neuen Gemeinde den breiten Weg Betretener zu finden, ganz wohl. Je vollendeter der Mann und Seelsorger, desto friedlicher und leichter geht Hoffen und Fürchten, Trauen und Mißtrauen zusammen.

\section{II. Richte dich auf ein langes Bleiben}

41. Alle menschlichen Verhältnisse sind unvollkommen, und nichts macht sich in neuen Verhältnissen leichter sichtbar, als diese Unvollkommenheit. Weil sich nun neue Pastoren darauf nicht gerichtet haben, so werden sie durch die Unvollkommenheiten und Mängel der neuen Pfarrei, der
\quellseite{104}
Gemeinde, des Hauses, des Ertrags oft sehr überrascht und niedergeschlagen. Ehe sie sich Zeit gelassen haben, sich innerlich mit ihrer Lage auszusöhnen und sie zu überwinden, — während sie fühlen, daß sie das Heimweh nach den früheren Verhältnissen beherrscht, daß sie noch gar nicht eingewöhnt sind u.\,s.\,w., mit einem Wort schon in den allerersten Anfängen geht ihnen die Geduld aus. Sie versichern dann oftmals, was Erfahreneren so lächerlich klingt, daß sie auf der neuen Stelle nicht lange bleiben wollen. Sie überlegen gar nicht, daß ihr Bleiben und Gehen nicht einmal in ihrer Macht steht, sonst würden sie ja gleich wieder gehen. Sie bedenken nicht, daß mit dem unwahrscheinlichen Geschwätz nichts geändert wird, so wenig als einer aus seiner Haut fahren kann, der es gern möchte. Es ist ein unmännliches, kindisches Guthabenwollen, welches Zeugnis gibt, daß man nicht will, was man soll, — daß man den hohen Gedanken der Hirtenliebe nicht erfaßt, — daß man sich selbst, nicht das Amt und Reich Gottes gesucht hat, wenn man sogar unter dem Heimweh nach irdischem Gemach seufzt und schmachtet. Man vertreibe sich den Jammer durch Erkenntnis seiner Sündhaftigkeit und bitte Gott und im Stillen auch der neuen Gemeinde die Lieblosigkeit eines solchen Sehnens ab. Ohne Zweifel würde es die Gemeinde sehr herunterstimmen und die Wirksamkeit eines solchen Hirten sehr hindern, wenn es offenbar würde, wie unbefriedigt und mißvergnügt er mit der neuen Stellung sei.

42. Statt sich der Betrachtung vorhandener Mängel und Uebelstände zu überlassen, erkenne man lieber das sich
\quellseite{105}
einstellende Mißvergnügen als eine Anfechtung. Hat man diese überwunden, so wird man Stärke angezogen haben zur Lösung der Aufgabe, die man übernommen hat. Um sie aber zu überwinden, stelle man sich recht lebhaft vor, daß man nicht lebt, um es irdisch gut zu haben, sondern um dem Herrn und Seinem Reiche zu dienen. Man erinnere sich, daß man dies in seinem Lebensberufe so ganz und in so hohem Maße vermöge, als in dem des Geistlichen. Man vergesse nicht, was die Gemeinden sind, was ihnen werden soll und was aus ihnen werden soll, — und daß ihnen alles, daß sie selber alles durch den Dienst ihres Hirten werden sollen. Der hehre Beruf, die schwere Verantwortung erwecke aus dem quälenden Traume selbstsüchtigen Berufes, sporne zur Amtsarbeit, über der man das Eigene vergißt! — Man mache sich recht klar, was man soll; man erwäge recht genau, was man vermag; man entschließe sich, dem scharf erkannten Ziele treulich und geduldig nachzujagen. Nur zu erreichen, was man soll, sei man bedacht, nicht aber quäle man sich, zu rechnen, wie lange Zeit dazu nötig sei. Man verspreche unter keiner Bedingung, eine bestimmte Zahl von Jahren zu bleiben; aber man verweigere das Versprechen nicht aus Lust, zu gehen. Die Zeit komme gar nicht in Anschlag. So lange Gott will, thue man stille und treulich, was Gott will. Seinen Willen erkenne man im Beruf, ihm opfere man freudig den eigenen Willen, ihm widme man alle Kräfte. Unsre Zeit steht in seinen Händen. Er weiß, wann es für uns an jedem Orte Abend ist. Er gibt Geduld, bis daß es Abend ist.

\quellseite{106}
\section[III. Arbeite, was vorliegt]{III. Arbeite, was vorliegt; das Andere laß an dich kommen}

43. Nichts gewöhnlicher, als daß neue Pastoren auf einmal alles, oder doch zu viel angreifen und bekehren wollen. Wenn das Versagen des Anfangs vorüber, so meldet sich bei manchen Pfarrern etwa der Muth, ja wohl Uebermuth des Anfängers. Weil man es mit der Aufgabe zu leicht nimmt und die Kräfte überschätzt; so schickt man sich an, im Sturm zu erobern. Man greift zu hart an, — und die Folge ist, daß man schnell ermüdet und nachlassen muß. Man betreibt zu viel und vielerlei auf einmal und zugleich, da muß man bald einen Haufen angefangener Dinge liegen lassen. So viel man zu viel gethan, so viel muß man dann zu wenig thun. Das erweckt inwendig ein lähmendes, beschämendes Gefühl der Ohnmacht, und in der Gemeinde das gerechte Urtheil, welches aber sehr hinderlich ist, daß der Pfarrer ein Anfänger sei und die Sache noch nicht verstehe. Damit ist denn der Drang nach großer Wirksamkeit theuer genug bezahlt. Die Lorbeeren, die man schon hatte — in Gedanken, sind verdorrt, ehe sie reiften, und das Beste von der Sache ist, daß die todeswürdige Eitelkeit einen tüchtigen Todesstoß erhalten hat.

44. Wer nicht den Eindruck eines Anfängers machen will, der sei ein Anfänger. Obwohl vollbrünstigen Verlangens, das Ziel zu erreichen, vergeße man
\quellseite{107}
nicht, daß die besten Fußgänger im Anfang einer Reise kurze Tagereisen machen und langsam gehen. So nach und nach erstarkt die Kraft zu größerer Anstrengung und Ausdauer. Wer darum weise ist, der legt es so an, daß er, statt bald aufhören zu müssen, je länger, je mehr thun kann. Ein solches Beginnen erkennt und würdigt das Volk bald und traut einem solchen Anfänger schnell Geschick und Meisterschaft zu, wie es denn auch in seinem Urtheil desfalls nicht irrt.

45. Gemäß diesen Grundsätzen thut ein neuer Pastor anfangs nur das Nöthige und wendet allen Fleiß an, das nach Maßgabe der Verhältnisse recht und völlig zu thun. Ehe er des Nöthigen völlig Meister ist, greift er nicht zu anderem. Die Arbeit, die er thut, überwindet er; aber er läßt sich mehr von ihr suchen, als er sie sucht. Er läßt es an sich kommen. Bei solchem Langsamthun lernt er sein Terrain immer besser kennen, — und während er Schritt für Schritt vorwärts geht, ist sein Gang richtig und gewiß und er braucht nicht leicht einen Schritt rückwärts zu gehen. Ein Uebel nach dem andern wird abgethan, ein Steg nach dem andern gewonnen, — und je unbemerkter alles gehen kann, je natur- und sachgemäßer jeder Fortschritt ist, je weniger er herbeigezwungen erscheint; desto ungestörter geht es weiter. Allmählich merkt dann das Volk am Werk den Meister. Noch ehe der Widerstand der Feinde rechte Gelegenheit gefunden, wird man für bewährt erkannt. Wie viel aber solche Bewährung hilft, das erfahre man, um es zu loben.

46. Diesen vorsichtigen, aber segensreichen Gang beachte
\quellseite{108}
man namentlich in der Seelsorge. Man glaube ja nicht, daß irgend ein Mittel sei, gleich anfangs sicher und schnell die Gemeinde kennen zu lernen. Methodistische Besuche\footnote{Sehr gut weist S. J. Baumgarten (kasuist. Pastoraltheol. S. 847 ff.) die Hausbesuche in Schranken. In „Anleitung der Erziehung der Kinder, Versorgung des Hausgottesdienstes, Vorsorge für das Seelenheil des Gesindes, Eintracht in Ansehung des gesellschaftlichen Lebens, so wie das Vertrauen der Zuhörer auf erlaubte Weise zu wecken“, statuirt er sie. Dagegen mißbilligt er methodische Hausbesuche als unanständig und nachtheilig. Es sei „keine Obliegenheit dazu da. Auch werde keine entfernte Zubereitung der Zuhörer erreicht.“} fordern auch die Leute zu einer Methode heraus. Willst du sie methodisch kennen lernen, so werden sie sich methodisch zu verstecken und zu verdecken wissen. Sie werden damit leicht, du aber desto schwerer das Ziel erreichen. Darum stehe ab von Methoden und halte von vornherein den Grundsatz fest, daß die ganze Seelsorge kasuell ist. Warte auf die Gelegenheiten und nimm dieselben wahr. Die Menschen pflegen sich zwar auch bei diesen nicht herauszugeben. Aber die Gelegenheiten, bei welchen ein Pastor seinen Pfarrkindern im Amte zu nahen hat, pflegen doch von der Art zu sein, daß das Herz mehr als sonst aufgeht. Dazu kann ein Pfarrer gar vieles durch ganz einfache Betrachtung und Beobachtung der allgemeinen Verhältnisse seiner Gemeinde lernen und sich in sein Arbeitsfeld und dessen Eigenschaften einstudieren. Ohne viel zu laufen und zu rennen, während er sich selten macht, wird er auf diese Weise bald genug
\quellseite{109}
seine Gemeinde und deren einzelne Glieder kennen lernen. Bald genug wird er mit allen in irgend ein Verhältnis kommen, bald genug wird er der geistige Mittelpunkt seines Volkes werden, — und die Beziehungen werden in aller Stille immer häufiger und inniger werden, die Bande sich immer zahlreicher und inniger knüpfen. Es wird sich offenbaren, daß kein Mensch stiller und allgemeiner, kein Mensch verborgener und doch öffentlicher wirkt, als ein Pastor, der in den Schranken seines Berufes geht.

47. Diese Art und Weise der Wirksamkeit bringt es auch mit sich, daß man im Predigen nicht gleich vornherein seelsorgerlich und kasuell rede.\footnote{Der beste Anfangsseelsorger ist der, welcher allen Fleiß anwendet, die Seelen durch Predigt und Katechese und Liturgie zu gewinnen. Das Gewinnen durch allgemeiner und objectiver gehaltene Vorträge eignet sich mehr als alles andre zum Anfangsleben.} Ueberhaupt ist die Unwissenheit in Anfangszuständen ein Vorrecht, welches ein Pfarrer mit bestem Wissen und Gewissen gebrauchen darf. Er hat kein Urtheil über seine Gemeinde, kann keins haben, will sich auch jedenfalls kein schiefes Urtheil bilden, und behandelt sie deshalb so lang als möglich so, wie man eben nach der heiligen Schrift alle Menschen, auch die man nicht kennt, ansehen und behandeln muß, nämlich als theuer erkaufte, des Hirtenamtes bedürftige Seelen. So predigt er denn auch Buße und Erlösung im allgemeinen. Er steht auf den Sprüchen, daß Gott allen Menschen an allen Enden der Welt gebiete, Buße zu thun, daß er wolle, daß allen Menschen geholfen werde etc. Kommt Zeit, kommt Rath und Veranlassung genug, kasuell und
\quellseite{110}
speciell zu werden. Es gibt, versteht sich, Fälle, in denen man gleich anfangs klar sieht, sehen muß, ja sich gedrungen und gezwungen fühlt, Zeugniß abzulegen, trotz dem, daß man kaum in die Gemeinde eingetreten ist. Solche Fälle rechtfertigen sich aber vor dem Verständigen von selbst, und der andern Urtheil wiegt nicht. Aber insgemein verdirbt sich der für spätere Zeiten Achtung und Wirkung, der zu früh casuell und speciell wurde. Er konnte natürlich das Rechte nicht treffen, wurde lächerlich, — und damit verdirbt es ein Jeder für lange Zeit. — — Es ist freilich so, daß Predigten erst recht interessant, ansprechend und wirksam werden, wenn sie so recht für die Verhältnisse der treffenden Gemeinde gemacht sind. Aber das ist nun eben so, daß man vornherein so nicht predigen kann. So spare man auch den guten Wein auf später und freue sich, wenn nur je länger je mehr Wasser in Wein, nicht umgekehrt Wein in Wasser verwandelt wird!

48. Es gibt einen Weg, gleich anfangs wenigstens eine äußerliche Kenntniß der Gemeinde zu bekommen, — einen Weg, den man zwar durchaus nicht überschätzen darf, der aber nichts desto weniger bemerkens- und beachtenswerth genug ist, um hier genannt zu werden. Er besteht nämlich im Studium oder in der Anlegung eines genauen Familienregisters. Ein Familienregister enthält eine tabellarische Uebersicht der Namen, Geburts-, Tauf-, Confirmations-, Trauungs- und Sterbetage, so wie der hauptsächlichsten andern Ereignisse im bürgerlichen und sittlichen Leben einer Familie und ihrer Glieder. Es wird nach Ortschaften und Hausnummern angelegt, und eng anschließend
\quellseite{111}
kann man alsdann auch das Beichtregister gleichfalls familienweise anfertigen, so daß es zur Ergänzung des Familienregisters dienen kann. Für den, welcher nicht Hirte der Gemeinde ist, sind solche Register von keinem Werthe; aber der Pastor lernt sie sehr hoch schätzen. Nicht bloß lernt man beim Lesen und noch mehr beim Anlegen der Register die Gemeinde in einem gewissen Maße kennen; sondern dieselben sind dann auch eine große Unterstützung für das Gedächtnis und belebte Andenken des Pfarrers, anderer Vortheile, die man bei Aufstellung von Geburts-, Confirmations-, Trau- und Todtenscheinen, Familienschematen \&c. davon genießt, gar nicht zu gedenken. Ist es nöthig, solche Register ganz frisch anzulegen (und in wie vielen Gemeinden mag es nöthig sein!), so gewährt die allerdings mühevolle, später aber sehr fruchtbare Arbeit eine vortreffliche Gelegenheit die Hausväter vorzurufen und mit ihnen ihre Familien, deren Namen \&c., durchzugehen. Da schließt man dann Bekanntschaften mit den Gemeindegliedern auf die einfachste Weise. Es sei erlaubt, im Anhang\footnote{Siehe den Anhang am Schluß.} einige Seiten aus einem Familien- und Beichtregister der nun bezeichneten Art mitzutheilen. Aus ihnen wird man leicht erkennen, wie sehr zur Kenntniß der Gemeinde ein Familienregister dienen kann. Vielleicht segnet manch junger Pfarrer, der diesen Rath — Familienregister und darauf gegründete Beichtregister anzulegen — befolgt, diesen scheinbar geringhaltigen Paragraph später einmal von Herzensgrund.

49. Vertrauen von Seiten der Gemeinde ist einem
\quellseite{112}
Hirten nöthig, ja, sehr nöthig! Aber eben weil man nur dem Bewährten vertrauen kann, ist es natürlich, daß man dies Vertrauen nicht fordern, noch erwarten darf, ehe man sich bewährt hat. Man lege es deshalb nur darauf an, sich zu bewähren. Das geschieht aber, wenn man sein redlich und einfältig thut, was eben vorliegt, und in allen Lagen beweist, daß man nichts wolle, als recht thun zum Heil der Gemeinde. Vertrauen ist unmittelbar niemals zu erreichen, sondern nur mittelbar. Mittel dazu ist Arbeit, Ackerbearbeit, treue, geduldige Liebesarbeit. Nicht das Vertrauen ist Augen- und Zielpunkt, sondern das Reich Gottes und seine Gerechtigkeit. Vertrauen fällt dem zu, der nach dem Reich Gottes und seiner Gerechtigkeit trachtet. Es ist darum eine unwürdige Buhlerei, wenn ein neuer Pastor um jeden Preis Vertrauen zu gewinnen sucht. Gottesmänner wünschen sich Vertrauen, thun, was Vertrauen erwecken kann, sie laufen ihm aber nicht nach, sondern warten sein. Es kommt auf diese Weise in der Regel auch eher und bleibt länger, wenn es einmal da ist. Auch hier heißt es: „Arbeite, was vorliegt; das Andere laß an dich kommen.“

50. Jeder Pastor wird und muß wünschen, in seiner Gemeinde möglichst viel zu erreichen. Aber da er weiß, daß er selbst zu diesem Produkte nur ein Faktor ist; da er nicht weiß, ob sich die andern Faktoren dazu finden werden; da er weiß, daß Gott sein nicht bedarf zu seines Reiches Mehrung, und nicht weiß, ob ihn Gott in dem von ihm gewünschten Maße brauchen und segnen werde; da er weiß, daß sein erstes Ziel die eigene Seligkeit ist, und nicht weiß, ob für seine Seele nicht viel Gelingen
\quellseite{113}
schädlich wäre; so verspricht er sich nicht, was er wünscht, sondern er arbeitet nur so, daß er seinen Platz ausfülle, daß es an ihm nicht fehle, zu erreichen, was er wünscht. Das übrige stellt er betend und flehend dem Herrn, welcher auch in seinen verborgenen Rathschlägen und dunkeln Führungen die Liebe ist und bleibt. „Er thut, was vorliegt; das andere läßt er an sich kommen.“

\section[IV. Schone deinen Vorfahr und deinen Nachfolger]{IV. Schone deinen Vorfahr und deinen Nachfolger\footnote{Schone deinen Kollegen, dürfte man dazu setzen. Denn wie oft übersehen namentlich jüngere Pfarrer diese Regel und laden durch Unbedacht, z.\,B. durch methodistische und unberufene Hausbesuche, einen Verdacht des Ehr- oder Geldgeizes auf sich!}}

51. Wer Gelegenheit gehabt hat, mit neueingetretenen Pastoren umzugehen, der wird wissen, wie mißgestimmt sie in der Regel gegen den Vorfahr zu sein pflegen. Der hat nichts geleistet, viel gefehlt und muß an allem Schuld sein, was sich nur übles in der Gemeinde und im Verhältnis zu ihr findet. Jede Anhänglichkeit an ihn, die sich etwa hie und da, wenn auch nur ganz schüchtern und leise kund gibt, hat er sich durch seine fehlerhafte, menschengefällige Amtsführung erworben; jede Mißstimmung, die sich bemerklich macht, hat er verdient. Und so sehr quält die innwendige Abneigung das arme Herz, daß es bei jeder gegebenen Gelegenheit in Worte übergeht. Es kann wohl auch kommen, daß man die Gelegenheit, wenn sie nicht kommen will, mit Gewalt herbeizieht. Dabei ists dann gleich viel, zu
\quellseite{114}
wem man spricht, ob zu vertrauten Freunden, oder zu Pfarrkindern, die doch billig eine Pietät gegen ihren ehemaligen Hirten behalten sollen. — Nun ist es freilich wahr, daß mancher einen Vorfahr gehabt hat, welchen zu loben Sünde wäre, dessen Achtung bei der Gemeinde auf sündlichem Grunde ruht, gegen welchen nur der ihm gleiche unbußfertige Sünder Pietät und Neigung bewahren kann. Aber oft ist der Vorfahr ein redlicher Mann gewesen, nur kein vollkommener. Einen solchen sollte man schonen. Und oft ist es nur der Neid auf des Vorfahr’s entschiedene und überwiegende Trefflichkeit, und auf die sich kundgebende Anerkennung derselben, was das Herz und den Mund mit Klagen füllt. Da sollte man es wagen, den Neid durch Liebe auszurotten, denn es giebt kein besseres Mittel gegen den Neid als die Liebe (ein Paradoxon, aber ein wahres), und man sollte sich deshalb vielmehr dem an den Hals hängen, der mit nichts, als mit Gab und Tugend, seinem Nachfolger ein wenig Leid bereitet hat. — Ach, es wirkt auf den neuen Pfarrer ein so schlimmes Licht, wenn, so viel aus seinen Worten geschlossen werden kann, die bessere Zeit erst mit seinem Kommen begonnen hat, vor ihm nur Finsterniß die Gemeinde bedeckte. Dagegen ist eine unaffektirte, redliche Anerkennung des Vorfahr’s schön und herzgewinnend! Es scheint ein edles und gerechtes Herz zu verrathen, dem man daher auch gerne mit Vertrauen entgegenkommt, wenn einem Pfarrer Lob und Anerkennung seines Vorfahr’s leicht und ohne Mühe von den Lippen fließt. Wenn mancher neue Pfarrer wüßte und erkennete, wie viel weiter er durch Anerkennung des Guten, als
\quellseite{115}
durch Mißbilligung der Mängel des Vorfahr’s kommen kann: er würde gegen den neidischen Hochmuth seiner Seele ganz anders zu Felde ziehen, und der Herr würde ihm dann auch zu Hilfe kommen und ihn von Sünd’ und Sündenschmutz erretten. — Schone den Vorfahr, du schonst damit dein Gewissen, dein eigenes Ansehen bei der Gemeinde, deine Wirksamkeit! Das bleibe unvergessliche Regel eines neuen Pfarrers.

52. Wenn nun die erste Hälfte der Ueberschrift dieses Paragraphen „Schone den Vorfahr“ auch leicht erkannt wird; so ist es doch mit der zweiten anders. Das „Schone den Nachfolger“ scheint keinen Sinn zu haben. Aber es scheint doch nur. Kennt man gleich den Nachfolger nicht mit Namen, hat man gleich kein Werk von ihm, keine von ihm geschlossene Wirksamkeit vor Augen, an der man sich versündigen und dadurch ihm selbst wehtun könnte, falls er noch leben sollte; so kann man doch durch ein rücksichtsloses Vorwärtsgehen den Nachfolger, er sei auch, wer er sei, sehr beschränken und einengen. Vieles liegt in der Befugniß eines Pfarrers, bei dessen Ausübung ihn die Rücksicht auf den Nachfolger und das Wort des Apostels leiten sollte: „Ich habe es zwar alles Macht, aber es frommt nicht alles.“ Wer z.\,B. in irgend einem Stück mehr leisten kann, als ein anderer, als vielleicht sein Nachfolger leisten kann, der wird wohl zu überlegen haben, ob es der Gemeinde nicht heilsamer ist, ihre Pastoren in Einigkeit handeln zu sehen, als Zeuge zu sein, wie der Vorfahr den Nachfolger übertrifft. Wenn ein geringeres Maß von Arbeit mit verschiedener Gabe und in verschiedener Weise, aber mit gleichem Fleiß
\quellseite{116}
und Ernst geleistet wird, wird im Ganzen mehr Segen zu gewinnen sein, als wenn sich einer vor dem andern geltend macht. Dort hat die verschiedene Gabe doch ihre Wirkung, ohne einen andern zu stören; hier wird die Gemeinde zu Wählerei und menschlicher, meist nicht leidenschaftsloser Bevorzugung eines vor dem andern verleitet. Um ein Beispiel zu geben, wollen wir einmal setzen, ein Pfarrer hätte ein großes Maß von körperlicher Kraft, vermöge dessen er wöchentlich sechsmal zu predigen tüchtig wäre. Kraft macht Lust, und so gienge er denn ans Werk, unbekümmert, ob auch sein Nachfolger ein gleiches leisten könnte oder nicht. Würde man dies Verfahren nicht tadeln müssen? — Etwas ganz anderes wäre es natürlich, wenn etwa ein täglich Predigen durch besondere Umstände eine Zeit lang erheischt würde. Dies würde sich aus den Umständen rechtfertigen. Und wenn diese Umstände sich in die Zeit des Nachfolgers verzögen, so würde dem Vorfahr kein Tadel gesprochen werden können.

Fälle dieser Art gibt es viele. Ein frommer Mann handelt also, daß er, so viel an ihm liegt, weder anwesend, noch abwesend jemandem anstößig wird. Aufs wahre Heil der Gemeinde bedacht, sucht er sich in keinerlei Weise zu erheben. Er denkt sich immer als Glied einer Reihe, — und zwar, wie das meistens der Fall sein wird, als ein Glied, das zwischen andern steht. Sein Sinn ist, dem Vorfahr nach-, dem Nachfolger aufs ersprießlichste vorzuarbeiten. Gleichwie er selbst sich freut, wenn ihm sein Vorfahr Anknüpfungspunkte gelassen hat; so sucht er es auch seinem Nachfolger leicht zu machen, in seine Fußstapfen
\quellseite{117}
zu treten. So viel an ihm liegt, arbeite er dahin, daß man ihn persönlich möglichst wenig vermisse. Er wünscht Verborgenheit seiner Wohlthaten schon wegen des Zusammenhangs mit Vorfahren und Nachfolgern, öffentlich zu werden begehrt er nicht; es ist ihm genug, daß der Herr sieht. — Gienge es, wie er wünschte, so würde keine Versetzung, kein Tod eines Pfarrers Lücken lassen, er würde es dahin führen, daß die Leute seiner Gemeinde sagen müßten: „Bei uns hat man seit Menschengedenken nur an Auf- und Abzug oder an Leichen gemerkt, daß das Amt in andere Hände gekommen ist.“

Dies äußern wir hier, weil Anfänger am meisten dagegen sündigen, und doch in ihrem Lebensalter so wenig als in einem andern Umstand die Versicherung besitzen, daß sie so bald keinen Nachfolger haben werden.

\section[V. Rücksichtlich deines äußern Lebens]{V. Rücksichtlich deines äußern Lebens unter Andern ordne alles gleich anfangs}

53. Ein Pfarrer übernimmt in unserer Zeit gleich beim Antritt seines Amtes eine ziemliche Anzahl äußerlicher Rechte und Pflichten. Das weiß er. Drum ist es gut, wenn er auch gleich anfangs sich die vollständigste und genaueste Kenntniß derselben verschafft. Es ist etwas Unbequemes und dem Gefühle nach fast Unschickliches dabei, wenn man gleich vorn herein nach seinen Grenzen, nach aller seiner Wirksamkeit, nach seinen Rechten und Pflichten fragen soll. Und doch ist es natürlich, daß man bei
\quellseite{118}
Uebernahme, ja, wo möglich, schon vor Uebernahme\footnote{Ein wohl gemeinter Rath in Betreff der Aufzugskosten: Glaube nicht, dem Edelmuthe der Gemeinde trauen zu dürfen, sondern mach getrost die Sache des Kostenpunkts schon vor deinem Aufzug fix und fertig. Versäumst du es, so wirst du’s büßen, und die Freude an der neuen Stellung wird schnell zu Wasser werden!} eines Wirkungskreises denselben so genau kennen lerne, als es sein kann. Auch findet man eine genaue Erforschung, so unbequem und drückend sie auch sei, doch immerhin am Anfang noch am schicklichsten; und erst nach angetretenem Amte, bei vorkommenden Fällen Nachfrage halten, macht den Eindruck der Unbesonnenheit und Unüberlegtheit bei der Meldung. Und doch ist es gerade eines Mannes Zeichen, seines Berufes Meister zu sein, was nun einmal ohne Kenntniß seiner Grenzen und Weitschaft unmöglich ist. — Das kann leicht sein, daß ein genaues Erforschen der Rechte und Pflichten zu manch unangenehmer Untersuchung, zu mancher Verwickelung, mancher Reibung, manchem Anstoß, manchem Verdruß führt. Allein es ist allewege gerathener, wenn ja Verdruß entstehen muß, ihn vorweg hin zu nehmen, weil dann muthmaßlich zum Vergeben und Ausgleichen mehr Zeit übrig bleibt, als wenn man das Unangenehme aufspart. Hat man auch vornherein angestoßen und den und jenen von sich entfernt, so kommt er doch meist wieder, wenn es offenbar wird, daß man dennoch in der Ruhe und Unbefangenheit des guten Gewissens bleibt und sich in Ausübung der Liebe nicht stören läßt, auch nicht von denen, die sich beleidigt wähnen. Ja, der den Muth hatte, anfangs anzustoßen, um hernachmals
\quellseite{119}
desto ungestörterem Frieden zu bewahren, wird je länger je mehr in der Trefflichkeit seiner Gesinnung anerkannt. — Freilich aber muß erst eine treffliche Gesinnung da sein, und es muß nicht Ehrgeiz und Habsucht die Triebfeder sein, um deren willen man gleich anfangs mit den Interessen anderer anbindet.

54. Mit dem bereits Gesagten steht es im Zusammenhang, daß man gleich anfangs sich über diejenigen schriftlichen Arbeiten ins Reine bringe, welche nach einmal bestehender Ordnung zur Unterstützung der weltlichen Behörden den Pfarrern auferlegt sind, oder welche zu geordnetem Regimente der Kirche nöthig erfunden werden. Ein genauer Kalender der periodisch wiederkehrenden Arbeiten, eine pünktliche Bezeichnung der Termine, an welchen sie abgeliefert werden müssen, hänge zu immerwährendem Merkzeichen und zur Erinnerung an einem Orte, dem man oft nahe kommt. Eben weil diese Dinge großen Theils fade schmecken und unangenehm sind, eben weil jeder Schreiber sie leicht eben so gut, oder gar besser, als ein Pfarrer liefern könnte; soll man sich mit ihnen recht vertraut machen und sie sich möglichst schnell vom Halse schaffen.\footnote{Doch thue man sie auch nicht ab, um sie abgethan zu haben und sie vergessen zu können. Geflissentliche Vergessenheit könnte sich durch die Nothwendigkeit, öfter wieder auf sie eingehen zu müssen, rächen. Sie schnell abthun, aber ihnen auch die gebührende Beachtung und ein treues Gedächtniß widmen!} Eulenspiegel’s Art, bei Bergen zu laufen bis zum Gipfel und dann den Nachkommenden zuzusehen, ist in diesen Dingen leicht nachzuahmen und sehr räthlich. Man lerne gleich vornherein
\quellseite{120}
die rechte Art dieses sogenannten „Amtierens“ vollständig kennen, damit man es hinter sich habe. Thut man das nicht, so kann man sich viel unnütze Mühe machen, viel Unmuth und Verdruß. Denn es ist erfahrungsmäßig wahr, daß diese Dinge desto schwerer gelingen, je weniger man sie kennen lernt, daß manche, die ein geschickter Arbeiter in halben Stunden vollendet, halbe, ja ganze Tage wegfressen können, wenn man sie nicht recht angreift, — daß sie, einmal verunglückt, gleich kleinen Rechenfehlern den unfleißigen und ungründlichen Arbeiter auf eine erbärmliche Weise verzehren können. Wie eine Neckerei eines bösen Geistes verfolgen sie manchen Pfarrer sein langes Leben, weil er nicht von Anfang gelernt hat, sie leicht und lustig zu thun.

Einem angehenden Pfarrer ist sehr zu rathen, seine Registratur kennen zu lernen, und wofern sie’s bedarf, zu ordnen. Es ist eine gewaltige Arbeit, eine auch nur mäßige Registratur zu ordnen. Manch schöner Tag geht drauf, Fascikeln zu ordnen, Roteln zu schreiben, Repertorium und alphabetisches Register der einzelnen Aktenstücke herzustellen. Aber, man wendet diese Mühe nur einmal an und genießt sie dann während einer ganzen Amtsführung. Denn eine geordnete Registratur fortführen und gebrauchen, ist leichte Mühe. Dagegen eine ungeordnete fortführen und gebrauchen, ist eine unmuthvolle Pein. — Nichts spart mehr Zeit, als einmal seine Registratur genau und völlig kennen und ihrer Meister sein.

55. Manche schriftliche Arbeiten, z.\,B. Rechnungen für Armenpflege, Kirchenstiftungen, Schulen \&c. liegen dem
\quellseite{121}
Pfarrer nicht als unerläßliche Pflicht ob, sondern man sieht es nur gerne, wenn er sich ihnen unterzieht. — Unerfahrene Pfarrer, welche an den Mitteln nicht genug haben, die ihnen das heilige Amt zur Einwirkung auf die Leute darreicht, weil sie dieselben nicht kennen, glauben sich durch Uebernahme jener Arbeiten Einfluß auf die Gemeinde zu sichern.\footnote{Es gibt jedoch auch sehr gewandte Pfarrer, denen Rechnungen \&c. nicht schwer werden, die schnell damit fertig werden und weniger Zeit verlieren, wenn sie die Arbeiten selbst machen, als wenn sie fremde Arbeiten überwachen. Bei denen mag ein anderes Urtheil gelten.} Und doch ist dieser Einfluß, auch wenn er erreicht wird, nicht groß, nicht nachhaltig, und wirkt nicht einmal Dank. — Der Schullehrer fertigt jene Arbeiten vielleicht eben so gut, als der Pfarrer. Die Wohlfahrt, welche der Pfarrer der Gemeinde damit erweist, ist nur sehr gering. Der Entgang an Zeit für das heilige Amt ist aber oft groß. Gar nicht angeschlagen, in wie vielfache Berührung mit weltlichen Behörden er sich begibt, wie oft dergleichen Arbeiten, um den Pfarrern, die sie gemacht haben, die weltliche Ueberlegenheit fühlen zu lassen, verworfen und bekrittelt werden, wie vielfach sie Anlaß geben, daß man sich an Pfarrern reibe und dadurch auch diese in eine unwürdige leidenschaftliche Aufregung versetzt. Schon um sich selber sündliche Erregung zu ersparen, soll ein Pfarrer dergleichen Arbeiten denen überlassen, welchen sie gehören, und sich so unabhängig als möglich von zeitlichen Geschäften zu stellen suchen. Ein Geistlicher thue sein Amt, andere Aemter überlasse er andern und sei nicht hungrig
\quellseite{122}
nach einem Lobe, wie das eines guten Geschäftsmannes ist.\footnote{Ist’s ein besserer Ruhm, wenn der Pfarrer ein guter Parlamentarius ist? — „Lassen Sie diejenigen zum Landtag gehen, die Zeit haben,“ sagte eine arme Kranke zu einem Pfarrer, der die Kranken andern überlassen und zum Parlamente gehen wollte.} Denn dies Lob ist gering: kein Studium der Theologie, keine Weihe, keine geistliche Gabe ist dazu nöthig: es geht die Kirche und ihre Diener im Grunde nichts an und bleibt billig denen, welche zu dergleichen in der Welt bestellt sind.

Recht schön ist, was Balduin hievon sagt:

\begin{quote}
Omne ministerium suum habet locum et quasi circulum, intra quem continetur, ne vagae sint ejus expeditiones, sed certo hominum generi prosint. — Intra cancellos muneris sui se contineant ministri verbi! Ἀλλοτριοεπισκοπὴν prohibuit Dominus! Luc. 22, 25. Si ministri verbi in domo domini, quae est Ecclesia, versari debent, non versentur in republica, seu in foro politico, alterum pedem in suggestu, alterum in curia habentes; non dicent magistratui leges forenses, non cudant consilia politica, non procurent bella, non gubernent tota regna et provincias. — In genere et in thesi de causis quoque criminalibus judicium ferre possunt, ut quod homicidae sint occidendi, immorigeri liberi puniendi, adulterii tollendi; idque non fit praeter officium ministri, cum verbum Dei ad omnia, quae sub sole fieri possunt, se extendat, ideoque olim sacerdotibus in specie praeceptum fuit Deut. 17, 8. Sed facti praesentis examinationem et theseos ad hypothesin applicationem
\quellseite{123}
merito committunt magistratui. Sic litigare in foro etiam ob alienationem bonorum ecclesiasticorum a ministrorum officio alienum est; ideo praefectis aerarii committi debet, ne ministri a laboribus sacris avocentur et aliorum in se concitare odia cogantur. Idcirco legibus Imperatorum etiam clerici ab actionibus publicis in curia exempti sunt et absurdum esse dicitur clericis, immo etiam opprobriosum, si peritos se velint etiam ostendere disceptationum forensium\footnote{Merkwürdige Stellen über Verweltlichung des bischöfl. Amtes aus Gregor und Gehler s. bei Hutter l. c. S. 41. Ueber Unabhängigkeit von weltlicher Gewalt, welche gar nicht begehrt wird, s. Luther W. A. Bd. XV, S. 2380. Hüssel, S. 37.} etc. etc.
\end{quote}

Führt ein Pfarrer solche Grundsätze von vornherein aus, so gewöhnt sich die Gemeinde bald daran und findet hernachmals selbst manches Gute daran. Hat auch der Vorfahr hierin eine andere Praxis befolgt, so braucht man ihm deshalb nicht nachzufolgen, weil man weder ihm, noch dem vielleicht anders handelnden Nachfolger in diesem Stücke Mißgunst erregt und das Volk doch einmal auch an einem Beispiel merkt, daß es seinen Pfarrer zu anderem, als zu Schreiberdiensten hat, und daß es ihn nach denselben nicht messen, noch schätzen müsse. — Ueberwachen soll ein Pfarrer jene Arbeiter, aber machen soll er sie nicht.

56. Zum Ordnen von Anfang her gehört auch das, daß sich ein Pfarrer sogleich klar mache, wem er unter-, wem gleich-, wem übergeordnet sei, und welche Befugnisse ihm nach allen diesen Seiten hin zustehen. So wie
\quellseite{124}
er das erkannt hat, richte er auch sein Benehmen darnach ein. Er sei ein Genosse weder seiner Ueber-, noch seiner Unter-Geordneten. Er sei jenen in demüthiger Einfalt gehorsam, so weit es das Wort des Herrn zuläßt, — und sei ein freundlicher, milder Vorgesetzter diesen. Er sei gegen seine Untergebenen so, wie er wünscht, daß seine Uebergeordneten gegen ihn sein möchten. Gefällt es ihm nicht und am wenigsten von geistlichen Uebergeordneten, immer nur angeherrscht zu werden, — wünscht er sehnlich, statt des leider der Stellung wegen fast unumgänglichen Kanzleistils der judicia ecclesiastica magistratuum territorialium\footnote{Man suche sich hiezu in A. H. Francke’s Collegium pastorale alle diejenigen Stellen auf, welche im Register bei dem Namen „Konsistorium“ beziffert sind.} den väterlichen Ton frommer und weiser Väter und Bischöfe zu vernehmen; so verdiene er auch sein Schicksal nicht durch den Ton, den er gegen seine Untergeordneten anstimmt. Er lasse sich das Ruder nicht aus den Händen winden, wo es ihm gebührt; er mache nicht Kameradschaft mit Kantor und Meßner; er duze sie nicht und trinke nicht Brüderschaft mit ihnen; aber er sei väterlich, freundlich, brüderlich, zuvorkommend, hilfreich und verschmähe es nicht, durch eine Bitte zu erlangen, was er auf diese Weise gern, mit strengen Worten aber nur ungern bekäme und erreichen würde. — Kurz, er sei Pfarrer und niemals Herr und lasse den weltlichen Behörden den weltlichen Imperativ, da ja Gottes Befehl und Seines Wortes Kraft ihm zur Seite stehen.

\quellseite{125}
\section[VI. Stelle dich im gesellschaftlichen Leben richtig]{VI. Stelle dich auch im gesellschaftlichen Leben gleich anfangs richtig und so, daß du immer vorwärts gehen kannst}

57. So, wie die Sachen stehen, erleidet es keinen Zweifel, daß der Stand des Geistlichen am meisten Anlaß zu allgemeiner Fortbildung gibt. Während der Jurist, der Arzt 2c. engbegrenzte Lebenssphären behandeln und durch ihre Geschäfte zu Einseitigkeit gewissermaßen gedrungen werden, ist der Geistliche durch sein Amt zur Betrachtung der allgemeinen Interessen der Menschheit hingeleitet und es bleibt ihm die Möglichkeit, an jeder Bewegung des geistigen Lebens Theil zu nehmen, mehr als andern Ständen gemäß. Daher kommt es auch, daß sich fast in jedem Fache menschlichen Wissens Dilettanten aus dem Stande der Geistlichen hervorgethan haben. Und eben daher kommt es auch, daß man den Geistlichen in Gesellschaften gerne sieht, von ihm Belebung und Erfrischung erwartet. Es ist fast eine — namentlich in den kleineren Städten, wo sich Familien von höherer Bildung seltener finden, — allgemeine Anforderung an den und Erwartung von dem eintretenden Geistlichen, daß er gesellschaftlich sei. Ist er’s, so hat er Achtung und Applaus der sogenannten Honoratioren und Gebildeten, die dann vielleicht dem guten Gesellschafter zu Gefallen des Jahres ein oder etliche Male auch zur Kirche kommen (Sylvesterabend!). Ist er’s nicht, so hat er in der Regel
\quellseite{126}
— wenigstens Anfangs — den Unmuth der Honoratioren zu tragen.

58. Nun erleidet keinen Zweifel, daß es in der Welt noch Gesellschaften geben kann, in denen Ehrbarkeit herrscht, in welchen der Name des Hochgelobten mit Ehrerbietung genannt wird, in welchen ein Geistlicher erscheinen kann, ohne die Stimmung der Andacht und sein, niemals von ihm zu trennendes heiliges Amt zu Hause lassen zu müssen. Es leidet auch keinen Zweifel, daß ein Geistlicher in solchen Gesellschaften erscheinen darf, wenn er darüber keine Amtspflicht, keinen Kranken, keinen Unterricht, keine Vorbereitung zur Predigt etc. versäumt, wenn ihn die Liebe zum Herrn und Seiner Gemeinde zu nichts Besserem leitet — wenn er nicht — zu üblem Exempel namentlich in unsern Tagen! — das Familienleben und die Erziehung seiner Kinder hintansetzt.\footnote{Die natürlichste Erholung, die lieblichste unter allen findet man im Umgang mit Weib und Kind. Wem die nicht mundet, der macht sein Weib nicht glücklich und ist nicht glücklich. Er wird auch schuldig an der hartnäckigen Unart seiner Kinder.} Aber — dieser Art Gesellschaften sind selten, sehr selten. Gesellschaften entgegengesetzter Art sind häufiger, sehr häufig. — Da nun ein Geistlicher die öffentlichen Gesellschaften seines neuen Wohnsitzes nicht gleich anfangs kennen kann, und er in der Regel darüber die Wahrheit nicht erfährt; so ist es das Gerathenste, sich anfangs wenigstens der Theilnahme zu entschlagen, bis man seine Umgebung in anderer Weise genauer kennen gelernt hat.

Man stößt freilich an, wenn man die in diesem Stücke
\quellseite{127}
gewöhnliche Zuvorkommenheit nicht erwidert, Einladungen von sich weist etc. Allein es ist eben kein Zeichen von einem Leben in göttlichen Dingen, wenn man schon durch den Gedanken, eine üble Meinung zu erregen, von derjenigen Weisheit und ihrer Konsequenz abgebracht wird, welche nun einmal alles als Mittel zum Zweck des Amtes ansieht. Dem Amte genügen; so leben, daß man sich nicht selbst im Laufe aufhalte! Nicht leicht einen Schritt in eine Gesellschaft thun, ohne überzeugt zu sein, daß man ihn nicht werde rückwärts thun müssen! Sich nicht unfreundlich, aber jeden Falls so benehmen, daß man frei bleibt von gesellschaftlichen, oft so störenden Zumuthungen, daß man frei bleibt vom Zwang der Verhältnisse, daß jede Annäherung in ihrer freien Liebe erkannt werde! — — Wollte Gott, man hätte anfangs so viel Muth! Man würde dennoch Gleichgesinnte finden! Man würde nicht vereinsamen! Im Anfang ist der Lockung am schwersten zu widerstehen! Im Anfang wird am meisten gesündigt! Gesündigt auf den guten Namen der christlichen Freiheit! Gesündigt — ja leider, und aus Liebe zu den gesellschaftlichen Kreisen das Vertrauen des Volks versehrt, das Takt genug hat, seine Hirten nicht in Gesellschaften zu suchen! —\footnote{Etwas ganz anderes ist es mit den herkömmlichen Einladungen eines Pastors zu Hochzeiten, Taufschmaus etc. Ist sich der Pastor eines redlichen Willens bewußt, hat er eine mannhafte Seele, die leicht Einfluß gewinnt, sieht er nicht Unordnung und allerlei Verlegenheit vorher; so wird er wohlthun, wenn er solche Einladungen nicht ausschlägt. Nicht eine große Wirkung hoffe er von seinem Dortsein, aber er beachte auch die kleine nicht, welche seine Gegenwart vielleicht haben kann, nämlich Verhinderung des Bösen. Wohl, wenn ihm gegeben ist, die ganze Gesellschaft an sein Gespräch zu fesseln und auf eine freundliche Weise bei Gutem zu erhalten! — Ist der Prediger nicht da, so benützt oft ein boshafter Schullehrer oder sonst ein Uebelgesinnter dergleichen Gelegenheiten, bösen Samen auszustreuen, die Gemeinde aufzureizen etc. Das alles ist nicht der Fall, wenn der Pfarrer die Gespräche leitet. (Cf. S. J. Baumgartens Casuist. Pastoraltheologie S. 863.)} Wären wir doch einmal weise genug, zu erkennen,
\quellseite{128}
daß es sich eben so sehr um die eigene Seele, als um das Heil des schwachen (?) Volkes handelt! — — Möchten wir uns doch enthalten von gesellschaftlicher Freude, wenn sie nicht in dem Herrn und zum Heile der Gemeinde genossen werden kann.

\quellseite{129}
\section{Die Mannestugenden}

Prudentia, Justitia, — Fortitudo, Temperantia (Buch der Weisheit 8, 7.) sind große Pfarrertugenden.

Ist der ein Mann und rechter Pfarrer, dem die letztgenannte fehlt? Gehen ihm nicht auch die drei ersten ab?

\begin{quote}
Δεῖ οὖν τὸν ἐπίσκοπον\\
εἶναι\\
σώφρονα.
\end{quote}

\quelle{1. Tim. 3, 2.}

\quellseite{130}
\chapter{Wandel}

\quellseite{131}
\begin{quote}
— μήπως ἄλλοις κηρύξας\\
αὐτὸς ἀδόκιμος γένωμαι!
\end{quote}

\quelle{1. Cor. 9, 27.}

\quellseite{132}
\section{I. Eingang zum Nachfolgenden}

59. Alles amtliche Verhalten eines Geistlichen, so wohlgethan es sei, erwirbt ihm die Achtung seiner Gemeinde nicht, wenn ihm ein wahrhaft geistlicher Wandel fehlt. Gott Lob, daß im Bewußtsein unseres Volkes Lehre und Leben noch so unzertrennlich verbunden sind, daß es am Beispiel des Geistlichen eine lebendige Predigt zu schauen begehrt. Darin liegt für den Geistlichen eine mächtige Anreizung zu einem würdigen Leben. Die Forderung des Volkes ist für ihn der Wiederhall seiner eigenen Predigt, und ein Beweis, daß nicht umsonst gepredigt wird. In ihr hat der Prediger, der so selten versteht, für sich selbst zu predigen, auch ein wenig Predigt, für welche er Gott danken darf.

60. Der Wandel ist ein gedoppelter, er ist entweder verborgen oder öffentlich. Der verborgene Wandel hat den Namen, weil er vor den Menschen verborgen ist, und sie ihn entweder nicht sehen können, da er im Innern, allein vor den Augen des Allwissenden geführt wird, oder doch insgemein nicht bemerken, indem er zwar auch äußerlich sichtbar ist, aber doch nicht zu Tage liegt und
\quellseite{133}
sich dem Auge des Nächsten aufdrängt. Der öffentliche Wandel ist schaal und leer, wenn ihm der verborgene Wandel vor dem Angesicht Gottes fehlt. Wie sich ein Leichnam von dem beseelten Leibe unterscheidet, so unterscheidet sich die pur äußerliche Ehrbarkeit von dem heiligen Leben eines Mannes, der inwendig und auswendig, in seinem heimlichen und öffentlichen Wandel vom Geiste der Heiligung durchdrungen ist.

61. Wenn der äußere und innere Wandel zusammenstimmen, so geht in Erfüllung, was der Apostel sagt: „Wenn uns unser Herz nicht verdammt, so haben wir eine Freudigkeit zu Gott“ 1. Joh. 3, 21. Freudigkeit und gutes Gewissen vor dem Herrn ist dem nöthig, welcher Israel seine Fehle und dem Volke Gottes sein Uebertreten verkündigen soll. Im Gegenteil ist es eine heimliche Schande, die Gott offenbaren und rächen wird, wenn man andere lehrt und selbst verwerflich wird. Ein gutes Gewissen macht eine stille, klare, zuversichtliche, muthige, mannhafte Seele. Und gerade eine solche bedarfst du zum heiligen Amte.

62. Ein reiner Wandel vor Gott und Menschen und das daraus hervorgehende gute Gewissen ist einem Geistlichen um so nöthiger, weil er in einer bösen Welt und im Haße der Teufel wandeln muß und darum der Verläumdung nicht entgehen kann. Der hat wenig Erfahrung, der bei einem jeden entstehenden Gerüchte mit einem semper aliquid haeret daher kommt. Es beliebt manchmal dem Lästerer, ein wenig Wahrheit zum Haltpunkt großer Lügen zu benützen. Oft aber erlügt und erdichtet er geradezu,
\quellseite{134}
und seine Kinder reden oft, wie sich der Landmann ausdrückt, „was sie mögen.“ Daher bleibt keiner länger unbescholten, als bis es irgend einem Buben einfällt, ihn zu schelten. Das kann aber bei einem Pfarrer, der ohne Ansehen der Person sein Amt verwaltet, gar leicht geschehen. Und ist nur einmal ein Scheltwort recht teufelisch unverschämt ausgesprochen, so kommt es auch bald in Umlauf; jeder glaubt dann auch, „was er mag.“ Und kurz, es gibt auf Erden und im Himmel keinen, welcher der Anfechtung böser Zungen entgehen kann. Das thut aber einem armen Menschen oft so weh, zumal wenn er einsieht, wie viel auf ein anerkanntes Beispiel für das Amt ankommt. Wohl dann, wenn „ihn sein Herz nicht verdammet!“ Mit der Zeit gewöhnt sich alles, und hat man sich nur erst gewöhnt, in der Anfechtung böser Zungen zu leben, so kommt das gute Gewissen auch zu Recht und Macht. Die Freudigkeit eines treuen Pfarrers ist alsdann, wie eines Einhorns, und seine Ruhe wie die eines Felsens.

\section{II. Verborgener Wandel}

63. Wir wollen nicht von Allem reden, was zum verborgenen Wandel gehört, können es ja auch nicht. Aber einiges sei besprochen. — Alles geistigen Lebens Quelle ist das Wort des Herrn. Darum vermahnt auch der Apostel seinen Timotheus, „anzuhalten mit Lesen.“ Das exegetische Studium führt zuweilen auch zu einer schönen Stunde der Erbauung; aber in der Regel ist es ein Theil der
\quellseite{135}
Arbeit eines Pfarrers. Ganz anders ist es schon mit dem kursorischen Lesen, bei welchem man nicht auf die einzelnen Schwierigkeiten achtet, sondern den Zusammenhang des Ganzen ins Auge faßt und sich seinem Eindruck hingibt. Aeltere Theologen haben das kursorische Lesen der heiligen Schrift sehr gerühmt und darauf gesehen, daß sie täglich etwas in fortlaufender Ordnung läsen und in gewisser Frist die heilige Schrift „hinausbrächten.“ In neuerer Zeit, wo man die Beständigkeit und Stärke des Willens nicht mehr zu haben scheint, wie früher, fand man in dem Vorschlag, die Bibel kursorisch und so zu lesen, daß man in bestimmter Frist immer ein gewisses Pensum, und so nach und nach die ganze heilige Schrift durchliefe, eine Beschränkung der christlichen Freiheit. Wie wenn man jemand zwingen könnte, es zu thun! Wie wenn es nicht möglich wäre, christlich frei die heilige Schrift in fortlaufender Reihe bis zu Ende zu lesen! Wie wenn ein frei erwähltes tägliches Lesen des göttlichen Wortes eine Gefahr mit sich führen könnte! Wie wenn es den Vätern gefährlich gewesen wäre! Man hat es nie versucht, nie geübt und hält es doch für gefährlich! Unkenntniß, Mangel an Erfahrung, Trägheit, Voraussicht des eigenen Unbestandes — das ist es, was von dem großen Segen eines regelmäßigen Lesens der heiligen Schrift abhält. Aber man gehe nur in die Stille und übe eine Weile ganz einfältig, was die Alten übten, und man wird aufhören, die Einwendungen für weise zu halten, die doch der Weisheit so offenbar ermangeln.

64. Ein Geistlicher soll die heilige Schrift lesen; aber er dürfte auch wohl menschliche Erbauungsschriften zur
\quellseite{136}
eigenen Seelenweide gebrauchen. Leider gibt es bei uns noch wenige, fast keine Erbauungsbücher für evangelische Geistliche, obwohl die Arbeiten und Bedürfnisse des Geistlichen besonders für ihn eingerichtete Schriften dieser Art viel mehr rechtfertigen würden, als wenn man z.\,B. Erbauungsschriften für das weibliche Geschlecht oder den jungfräulichen Stand zurichtet. Hätten wir nur ein evangelisches Buch, welches in seiner Weise so vortrefflich wäre, wie das römische Brevier in seiner Weise wirklich ist, so würde einem großen Mangel abgeholfen und eine reiche Segensquelle geöffnet sein. Da nun aber so etwas nicht vorhanden ist und auch schwerlich zu erwarten sein dürfte; so gebrauche der evangelische Geistliche um so eifriger die mannigfaltige ascetische Literatur der lutherischen Kirche und verachte diese süßen, noch lange nicht genug erkannten Früchte des heiligen Geistes nicht. Zwar ist die Ausdrucksweise manches berühmten lutherischen Erbauungsbuches manieriert und deshalb denen nicht zu empfehlen, welche gern in jeder absonderlichen Form der Rede sich so verfangen, daß sie nachahmen, oder aber davon abgestoßen werden; aber wir haben ja Schriftsteller, die nach Form und Inhalt vortrefflich sind, denen in der Geschichte deutscher Prosa ihr Platz eben so gewiß anzuweisen ist, als die Lieder der lutherischen Kirche in der Geschichte deutscher Poesie ihren Platz und ihren Lorbeer bereits gefunden haben. Wir wollen von Einem nicht reden, dem auf diesem Gebiete alle Parteien huldigen, von Luther. Aber man erinnere sich an einen Scriver, an Philipp Nicolai, an Meyfart \&c. Wahrlich, wir besitzen in diesen und vielen andern Schriftstellern reiche, anmuthige Paradiesesgärten voll
\quellseite{137}
hoher Bäume und schöner, duftender Blumen, voll köstlichen Obstes und frischer Brunnen. Der selbst reich begabte Geistliche, wie der von geringeren Fähigkeiten kann hier Segen, Fortschritt, Reichthum, Stärkung, Tröstung und Ermunterung finden. Es sollte überhaupt keinen Geistlichen geben, welcher in diesen unsern reichen Gärten nicht bekannt und bewandert wäre, und es sollte längst dafür gesorgt sein, daß diejenigen, welche in sie eintreten wollen, Anleitung und Rath fänden. Es sollte längst eine Geschichte der geistlichen Literatur unserer Kirche geben und eine Sammlung des Herrlichsten und Besten. Das wäre auch in der That der beste Anfang, die beste Vorarbeit zu einem evangelischen Brevier. Aber freilich, selbst diese Vorarbeiten herbeizuschaffen, bedarf es betender, andächtiger, feiernder Seelen, die zu eigener Erbauung lesen, die an einer in Gottes Wort lebenden, anbetenden Seele die Kraft alles dessen zu erproben und zu messen verstehen, was sie lesen. — Möchten sich viele jüngere Geistliche aus dem Gesagten die nöthige Vermahnung und Ermunterung nehmen und in ihren noch jugendlichen Jahren ein heiliges, süßes, andächtiges Studium beginnen, welchem man die stillen Stunden eines ganzen Lebens widmen dürfte, ohne Mangel oder Langeweile zu finden.\footnote{Es dürfte bei dieser Gelegenheit wohl auch der Umgang mit den Schriften der Kirchenväter älterer Zeiten zum Zwecke der Erbauung empfohlen werden. Wer Erfahrung von dem Segen dieses Theils der christlichen Literatur machen wollte, dürfte nur z.\,B. folgendes Buch eine Weile betend gebrauchen: „Precationes ex veteribus orthodoxis doctoribus, ex ecclesiae hymnis et canticis, ex psalmis denique Davidis collectae, et nunc recens recognitae et auctae per Andream Musculum D. Lipsiae anno 1573.“ Frühere Auflagen waren schnell vergriffen worden — unter den Lutheranern, denen es gewidmet war, — und zwar in den ersten Jahrzehnten nach der Reformation. Leider dürften wir dem jetzigen Geschlechte ein lateinisches Betbuch kaum aufs neue bieten, so ausgezeichnet es wäre.}

\quellseite{138}
65. Eine Uebung und Aeußerung des inwendigen Lebens ist bei uns ganz verloren gegangen, nämlich die Meditation, die Erwägung göttlicher Worte oder Wahrheiten vor dem Angesichte Gottes. Was man bei der Betrachtung des Textes zum Behufe der Predigt jeden Falls verlangt um der Invention willen, das verlange man doch von jedem Bibelleser, zumal von einem Pastor. Die Meditation fördert vielen Reichthum in die Seele, und ein Geistlicher würde sich größtentheils die Mühe der Invention ersparen können, wenn er es nicht verschmähte, um sein selbst willen in täglicher Meditation zu leben. Er würde erfahren, daß die Gottseeligkeit zu allen Dingen, also auch zum heiligen Amte nützlich ist. Dem mangelt es an Gedanken zum Predigen gewiß nicht, welcher in der Meditation die Höhe und die Breite und die Länge und die Tiefe der göttlichen Liebe erwägt. — Der selige Calvör\footnote{Ein noch immer sehr brauchbares, früher auch viel gebrauchtes Buch der älteren Zeit dürfte hier ins Andenken zu rufen sein. Es ist in einer der vielen Ausgaben betitelt: „Praxis Pietatis, das ist: Uebung der Gottseligkeit, anfänglich in englischer Sprache beschrieben durch Herrn D. Ludwig Baili, Bischofen zu Bangor, und über dreißig mal gedruckt. Jetzt seiner Würde halben wiederum hervorgegeben. Lüneburg, gedruckt und verlegt durch Johann und Heinrich Stern S. Erben. Anno 1670.“ Der andre Theil „begreift in sich die fürtreffliche Uebung gottseliger und andächtiger Meditation, wie nämlich ein frommer Christ in Betrachtung göttlicher und himmlischer Sachen zu gewisser Zeit sich üben und hierdurch sowohl in diesem verbrüchlichen, als folgendem ewigen Leben herzerquickende Freud und beständige Ergetzung erlangen könne. Sammt einer nützlichen Vorrede von Prüfung des Gewißens und etlichen schönen Exempeln solcher Betrachtungen. Aus dem Englischen und Französischen in unsre deutsche Sprache gebracht.“ — Das ganze Buch wäre einer neuen Ausgabe werth.} gab 1691, hernach 1711 unter dem Titel „Christliches
\quellseite{139}
Kleeblatt“ drei erbauliche Schriften heraus, deren drittes die Aufschrift führt: „Gebahnter Weg zu der Ruhe in Gott.“ In diesem handelt er 1. von dem inwendigen Herzensgebete, 2. von der sogenannten Andachtsleiter, 3. von der Seufzerandacht. Zuletzt folgt als Anhang eine „kleine Beischule.“ Von den drei Hauptabschnitten des „Gebahnten Wegs“ legte der Verfasser dieser Blätter die Hauptgedanken in einem Tractate nieder, dem er den Titel „Sabbath und Vorschabbath“ gab. Er ist 1851 zum zweiten Male als Nr. 6 der v. Abth. II. der Gesellschaft für innere Mission nach dem Sinne der lutherischen Kirche herausgegebenen Tractate in Nördlingen bei Beck gedruckt worden. Möge es erlaubt sein, hier aus dem sehr einfachen Büchlein (p. 36–41) denjenigen Theil einzurücken, welcher von der Gebets- oder Andachtsleiter handelt oder, was einerlei ist, von der Meditation. Die Sache ist so unbekannt geworden, daß eine einfache, eingehende Belehrung am besten zeigen kann, was und wie es gemeint ist.

\subsection{Von der Andachts- und Himmelsleiter}

1. Was versteht man unter der Andachts- oder Himmelsleiter?

\quellseite{140}
Den rechten, seligen, zur Andacht und zum Himmel leitenden Gebrauch des göttlichen Wortes, wie es in der heiligen Schrift oder auch in menschlichen Büchern, dem Katechismus 2c. vorkommt.

2. Ich gestehe dir, daß ich nun doch nicht weiß, was du meinst?

Sieh, dieser selige Gebrauch der göttlichen Wahrheit besteht in drei Theilen: 1. in der Betrachtung, 2. in der Prüfung, 3. im Gebet. Man kann die Schrift schnell lesen, und das hat seinen eigenen Vortheil und ist nöthig. Aber man soll auch jene drei Theile des seligen Gebrauchs der Schrift nicht verabsäumen. Es liegt viel daran! —

3. Nun wohlan! Also wie übst du die Betrachtung?

a) Ich nehme vor mich einen bekannten Spruch, einen Liedervers, ein Stück aus dem Katechismus, ein Gebet 2c. Ich lasse kein Wort unerwogen vorüber gehen, sondern bedenke und beherzige eines nach dem andern.

b) Damit ich desto tiefer und gründlicher in den Sinn und Verstand der vorgenommenen Worte hineindringen möge, so wiederhole und überlege ich sie langsam, mit andächtiger Aufmerksamkeit etliche Male, theile auch den Spruch in gewisse Fragen, frage nach allen Umständen (wer? was? wo? warum? 2c.)

c) Ich erinnere mich dabei aus der Schrift, dem Gesangbuch, dem Katechismus anderer Stellen, welche mit der vorgenommenen übereinstimmen, und bewege auch diese in meinem Herzen.

d) Ich bedenke, was ich aus meinem Spruche, Verse 2c. für eine Lehre, Vermahnung, Trost nehmen könne.

4. Wie übst du die Prüfung.

Ich frage, ob der Spruch auch mich angehe. Das sehe ich sonderlich daraus, wenn ich die Personen\footnote{Ebenso kann man auch die Tugenden und Laster 2c. erforschen, von denen die Rede ist, und sich prüfen, ob man sie an sich habe oder nicht 2c.} ansehe, von
\quellseite{141}
welchen oder zu welchen der Spruch gesagt ist. Gehöre auch ich zu ihnen, so gehört der Spruch auch mir. Ich wende und deute ihn dann ganz auf mich, wie wenn er zu mir insonderheit gesprochen wäre. Alle Fragen, die ich bei der Betrachtung gethan habe, beantworte ich noch einmal mit Beziehung auf mich. Bin ich aber unter den Personen nicht, auf welche sich der Spruch bezieht, so ziehe ich ihn auch nicht auf mich. Und zwar, ist’s ein Trostspruch, so deute ich ihn nicht eher auf mich, als bis ich ein solcher Mensch geworden bin, welcher sich des Trostes annehmen kann. Ich fange deshalb sofort meine Besserung an, indem ich zu Gott um Gnade und Bekehrung seufze. Ist der Spruch, der mich nicht angeht, ein Strafspruch, so danke ich zwar Gott, daß Er mich vor der gestraften Sünde behütet hat, bitte Ihn aber auch, Er wolle mich ferner behüten, meine andern Sünden und Schwachheiten mir gnädig vergeben. —

\subsection{5. Wie übst du das Gebet?}

Was ich bei Betrachtung und Prüfung gefunden habe, das fasse ich vor Gott in ein Gebet zusammen.

\subsection[6. Wie schreitest du zur Herzensandacht?]{6. Wie schreitest du dann aus dieser Vorbereitung zur Herzensandacht weiter?}

Die Tugenden des inwendigen Lebens, an welche der Spruch erinnert, benütze ich sogleich zur Herzensandacht, zu dem friedlichen, seligen Werke meines inwendigen Sabbath’s. Die Sünden und Laster, welche bezeichnet werden, fasse ich ins Auge, vermaledeie, verwerfe sie und bitte Gott, daß ich dies mehr und mehr von ganzem Herzen thun könne, und daß meine Seele von Lust und Liebe zu ihnen abgeschieden sein und ruhen möge. So bin ich über die drei Stufen der Himmelsleiter zur vierten, d.\,i. durch die Vorbereitung zur wahren Herzensandacht gekommen. Und wenn ich nun durch geistliches Entfernen vom Bösen, durch Nahen zum Guten, durch die selige Bewegung meiner Seele im Herzensgebete gestärkt
\quellseite{142}
worden bin, so gehe ich hin an mein Tagwerk und übe auch äußerlich das Gute und hasse und lasse das Böse. —

7. Eigentlich sagst du mir mit dieser von dir sogenannten Himmelsleiter nichts Neues. So sollte man immer die heilige Schrift gebrauchen und sich durch sie zu Feierstunden des inwendigen Lebens, zum Vorschmack des ewigen Lebens leiten lassen.

Freilich ist es nichts Neues. Aber wer gebraucht die heilige Schrift so? Was ist leichter, als dies Betrachten, Prüfen, Beten — und was macht weiser, stärker, seliger im Guten? Wie wird man da die Kräfte der zukünftigen Welt inne, welche im Worte verborgen liegen! Wie wird da das Wort süßer denn Honig und Honigseim! — Die Sache ist nicht neu. Aber für viele Christen mag die Ausübung ganz neu sein. Wollte Gott, daß diese Himmelsleiter, dieser Gebrauch des göttlichen Wortes allgemeiner, geübter und geliebter wäre! Man würde es an allen Orten und in allen Stücken spüren.

8. Ich wünschte, du zeigest mir an einem oder einigen Beispielen den Gebrauch der Himmelsleiter.

Ich möchte dich doch fürs Erste zu eigenen Versuchen ermahnen. Vielleicht könnte ich dir später eine oder einige Betrachtungen, Prüfungen, Gebete vorlegen. Man lernt hier aus fremdem Beispiel mehr, wenn man durch Versuch und Erfahrung diejenigen Punkte gefunden hat, welche einem besonders schwer werden. — Die Uebung ist übrigens nach der dir gegebenen Anleitung ganz leicht.

9. Nur um ein einziges Beispiel bitte ich dich. Wie wendest du etwa Betrachtung und Prüfung bei dem Spruche Joh. 3, 16: „Also hat Gott die Welt geliebt“ etc. an?

Fragen zur Betrachtung sind diese:

Wer hat die Welt geliebt? — Was hat Er geliebt? — Wie, wie sehr hat Er sie geliebt? — Was heißt: „Er gab Seinen Sohn?“ — In welcher Absicht hat er Seinen Sohn gegeben? — Was heißt: „verloren werden?“ —
\quellseite{143}
Was heißt: „das ewige Leben haben?“ — Wer soll nicht verloren werden? — Was heißt: „an den Sohn glauben?“ — Wer unter denen, die an Jesum glauben, ist ausgenommen von der Verheißung des ewigen Lebens? — — Welche anderen Sprüche stimmen mit diesem Spruche? \&c.

Bei der Prüfung kann man etwa sprechen:

Geht der Spruch auch dich an, o meine Seele? Er geht dich eines Theiles an, weil du zur Welt gehörst, welche Gott so sehr geliebt hat. Also auch dich geht das herrliche „Also“ an. Also auch für dich ist der Sohn, — der Eingeborene — gegeben. Also auch du sollst nicht verloren gehen \&c. — Aber du würdest doch verloren gehen, wenn du nicht zu den Gläubigen gehörtest. Die selige Verheißung des ewigen Lebens geht dich nicht an, wenn du nicht an Jesum Christum glaubst. Geht dich also der Spruch mit seiner seligen Verheißung wirklich an? Glaubst du? — Ist dein Glaube rechtschaffen?

10. Aber ist das nicht eine pure Katechese?

Ja, aber eine solche, bei welcher Lehrer und Schüler eine und dieselbe Person ist. Da klingen die Fragen ganz anders und die Antworten desgleichen. Das drängt mächtig vom Herzen und geht zum Herzen.

11. Es kann aber doch auch pure Kopfsache sein.

Es kann, aber das ist dann Schuld der Menschen selbst, und dann ist’s freilich eine pure Schulmeisterei. Aber treib es nur mit rechtem Ernste, in abgeschiedener Stille, als Einer, der Gott sucht. Du wirst einen großen Segen finden. Auch muß ja nicht gerade Alles gefragt sein. Noch viel weniger muß Alles in Fragen geschehen. Ich setze nur Kürze halber die Fragen. Uebe dich und du wirst bald merken, daß bei der ganzen Sache eine selige Freiheit herrscht, alles ängstliche, sklavische Wesen verschwindet bald.

12. Ich habe dich noch nicht gebeten, mir zu zeigen, wie
\quellseite{144}
man nach Betrachtung und Prüfung über Joh. 3, 16. Alles in Gebet zusammenfasse.

Du hast wohl gethan. Denn das braucht kein Muster, zumal im Beten noch weniger als bei Betrachtung und Prüfung Aengstlichkeit obwalten soll.

13. Du verweist mich nun auf Uebung und Erfahrung?

Ganz und gar. Ein anderes Urtheil als das, welches aus Erfahrung hervorgeht, gilt hier nicht. — Geh hin, prüfe, erprobe, was Du gelernt hast!

66. In dem vorstehenden Auszug über die Meditation ist bereits davon die Rede, daß man das zu betrachtende Wort zur Selbstprüfung anwenden solle. Es gibt aber auch eine von der Betrachtung des Wortes unabhängige Selbstprüfung, welche man Sündenprüfung nennen könnte. Im Herzen eines Christenmenschen soll, auch wenn er das Wort nicht aufgeschlagen vor sich liegen hat, immer so viel Licht und Kraft des Wortes sein, daß er beim Blick ins Herz und in den Wandel seine Abweichung und Entfernung von dem Wort und Willen Gottes findet. Wohin er geht, soll eine Bestrafung des göttlichen Geistes mit ihm gehen, er soll sich in Zucht und Strafe des Geistes wissen und fühlen. Er soll aber nicht bloß diese unausbleibliche Folge eines dem Worte zugewendeten Lebens in sich tragen und gleichsam dulden; sondern er soll der demüthigenden Wirkung des Wortes und Geistes durch Fleiß und Treue in Erforschung seiner Sünden entgegenkommen. Wenn er die Bestrafung des Geistes in sich spürt, bete er Ps. 139, 23. 24: „Erforsche mich, Gott, und erfahre mein Herz; prüfe mich und erfahre, wie ichs
\quellseite{145}
meine, — und siehe, ob ich auf bösem Wege bin und leite mich auf ewigem Wege.“ Insbesondere geziemt es einem Pfarrer, in beständiger Selbstprüfung, Reue und Buße zu leben — um der Amtssünden willen. Der nicht allezeit in ernster Prüfung lebt, kann, ohne es zu wissen, viele Sünden bei sich führen; hingegen ist große Demüthigung und ein ernster Drang zu Dem, der die Gottlosen gerecht und heilig macht, in einem Herzen, welches die Pflicht der Selbstprüfung nie versäumt. Das gilt besonders für den Pfarrer. Lebt er in Prüfung seiner selbst, so wird ihm der ganze Complex seiner Pflichten immer neu und immer mehr vor die Seele treten, und er wird manches nicht für immer und nicht so oft übersehen, was außer ihm jedermann an ihm sieht und mißbilligt, was ihm aber in menschlicher Schwachheit und Vergeßlichkeit leicht und oft entgehen könnte.\footnote{Heilsame Uebung! Prüfungsfragen über Amtsverirrte und Amtstugenden aufsetzen. Nach diesen sich oftmals oder nach einer Prüfungstafel erforschen! Jeder Beruf hat seine specielle Prüfungstafel.}

67. Die Selbstprüfung eines Christen, besonders aber eines Geistlichen, sollte sich jedoch nicht bloß auf seine Sünden beziehen, sondern auch auf die Gaben und deren Maß. Nichts ist gewöhnlicher als Ueberschätzung der eigenen Gaben, und die eitelhafte Anmaßung, welche damit zusammenhängt, begegnet einem oft gerade bei geringbegabten Männern mit der zuversichtlichsten Gebärde. Jeder Pfarrer kommt sich selbst (und natürlich auch seinem Fleische und Blute, seinem Weibe) als ein keineswegs unbedeutender Redner, als ein tüchtiger Katechet und
\quellseite{146}
besser Seelsorger 2c. vor. Es wird daher geringe Einwirkung auf die Gemeinde und geringer Amtserfolg aus allen möglichen Ursachen hergeleitet, und nicht daher, woher es oft so richtig wäre. Man sollte denken, es müßte Trost gewähren, wenn man die geringe Leistung aus der geringen Gabe herleiten dürfte. Aber freilich, wie kann denn das trösten, was die bitterste Enttäuschung und die schmählichste „Demüthigung“ mit sich führt?! Nein, da hätte man schier umsonst gelebt! Man ist durch die Schulen, durch die Universitäts- und Candidatenjahre gegangen, hat so viel „durchgemacht“ und ertragen, und nun wäre das eigentlich alles nichts gewesen, und man hätte besser einen andern Beruf gewählt! Das kann man nicht annehmen. Nein, nein, nicht an der Begabung liegt der geringe Erfolg, sondern die Gemeinde hat nur kein Auge für die vorhandenen Gaben ihres Pfarrers, — es mangelt an Anerkennung dessen, was da ist, — die Schuld ist an den Leuten. So denkt das stolze Herz des Pfarrers; die Seinigen und seine persönlichen Freunde stimmen bei, und in der ganzen Gemeinde hat niemand „das Herz,“ die allgemein bekannte Wahrheit dem zu sagen, den sie betrifft. So geht dann der arme Pfarrer in der Gemeinde mit einem Anspruch und unter Voraussetzungen herum, um deren willen er hinter seinem Rücken verlacht wird. Der Hochmuth des Selbstgerechten ist empörend, der des Eingebildeten aber ist lächerlich: sträflich aber sind beide in gleichem Maße und zu beseitigen ist dennoch der letztere schwerer als der erstere. Man findet eher einen bußfertigen Sünder, als einen Menschen, welcher für seine kleine
\quellseite{147}
Gabe und Bedeutung ein offenes Auge hat und dabei heiter ist und sichs wohl gefallen läßt in der Gemeinschaft, in welcher er nur ein kleines Licht ist und dafür gilt. So ist der Mensch! Er ist noch stolzer auf die Gabe, die er zu haben meint, als auf seine Gerechtigkeit. Er will lieber ein Dieb, als dumm heißen; er verzeiht lieber jene, als diese Nachrede. Und doch, wie steht sich namentlich ein Pfarrer durch Gabenhochmuth selbst im Lichte! Würde er erkennen, in welchem Stücke er Gabe hat, in welchem nicht: so würde Wahrheit und Demuth die Gemeinde geneigt machen, die vorhandene Gabe sich desto lieber dienen zu lassen; er selbst würde sich theils vergebliche Arbeit ersparen, indem er nicht mehr, wie zuvor, mühselig Zeigen und Trauben von seinem Distel- und Dornstrauch erzwingen wollte, theils aber auch glücklicher arbeiten, indem er die wirklich vorhandene Gabe bildete und pflegte. Auch die geringe Gabe würde eher wachsen, wenn sie richtiger erkannt und ihr alsdann die Hülfe gebracht würde; die ihr förderlich sein könnte. Man soll ja freilich auch die geringe Gabe erziehen, bilden, stärken; aber wie man ein kleines Kind nicht erzieht, wenn man es wie einen Mann behandelt, so wird auch aus der kleinen Gabe nichts, wenn man ihr zumuthet und ihr Ziele steckt, wie ihr nicht ziemt. — Außer dem ist die Demüthigung, welche man durch Herabsteigen von erträumten Höhen empfindet, nur einen Augenblick schmerzlich; sie führt zur Wahrheit, und wer sich zur Wahrheit führen läßt, kommt zu einem Glücke. Die wahre Ansicht, welche ein Mensch von sich selbst bekommt, macht bescheiden, still, glücklich und heiter
\quellseite{148}
im vorhandenen Maße, gerecht gegen andere, kurz, sie ist gesegnet, wie jede Wahrheit, wenn man den Muth hat, sie zu umfassen. Darum wage es doch, und öffne die Augen für dich selbst und deine Gabe! Prüfe dich im Lichte Gottes! Lies und studiere die apostolischen Stellen, welche von den Gaben und ihrem Maße handeln, und bitte Gott, daß er dir verleihe, nicht mehr sein zu wollen, als du von Gottes Gnaden bist; das aber sei dann auch mit Kraft. Füll deinen Platz aus, und sei damit zufrieden. Man verlangt vom Haushalter nicht mehr, als daß er treu sei. — Predige manchmal für andere und für dich vom Glück des Kleinwerdens und Kleinseins. Du weißt, daß dies Glück den Größten unter den von Weibern Geborenen zum Liebling aller macht. Sei glücklich in deinem kleinen Maße, so wird man dich mehr lieben und mehr von dir lernen, als wenn du dich aufbläsest, so du doch nichts bist.

68. Vom Herzensgebet ist wenig zu sagen, weil ohnehin jedermann von dessen Nothwendigkeit und Wichtigkeit überzeugt ist, und es eine Schande wäre, denen erst Anweisung dazu geben zu wollen, welche Christi Schafe und Lämmer in dasselbe einführen sollen.\footnote{Vielleicht möchte aber doch gut sein, aus der oben Nro. 66 angegebenen Schrift von Calvör oder dem Tractat „Sabbath und Vorschabbath“ die Lehren unsrer frommen Väter vom Herzensgebet kennen zu lernen. Man kann ohne diese Handleitungen aus dem Herzen beten, und wie schlimm wäre es, wenn es anders wäre! Aber warum soll man eine auf gründlicher Erfahrung beruhende Anweisung verachten?} Doch erlauben wir uns:
\quellseite{149}
A. zu erinnern an die Orte, welche sich einem Geistlichen zur Uebung stillen Gebets am meisten empfehlen:

a) Eine Betkammer sollte jedes Haus haben, sonderlich ein Pfarrhaus. Da es nicht so ist, so muß die Studierstube zur Betkammer des Pfarrers werden.

b) Die Sacristei der Kirche wird oft zu allerlei Dingen gemißbraucht. Oft ist sie ein Ort des Geschwätzes für Pfarrer und Mesner und Kirchenpfleger. Oft muß sie noch zur Studierstube dienen, und der Pfarrer memoriert in ihr, bis er auf die Kanzel geht. Wie viel besser wäre ihre stille Abgelegenheit zum Herzensgebet, zur betenden Vorbereitung, zur Fürbitte für die Gemeinde und alle Nöthen der Kirche und der Welt angewendet! Eine Sacristei sollte zu dem Ende immer einsam, verschließbar, mit einem Betstuhl oder Betaltare versehen sein. Man kann ohne dergleichen beten; dennoch ist die durch einen Betstuhl verschaffte, zum Gebet einladende, knieende — und wohlgemerkt durch den Betstuhl und die daran befindliche Armbank bequem gemachte — Gebärde des Leibes und eine stille feierliche Ruhe des Ortes mit nichten zu verachten.

c) Besonders bequeme Gelegenheit zum Herzensgebet gewähren die Amtsgänge. Nirgend ist man einsamer, als auf dem Wege; nirgend dem Lauscher weniger ausgesetzt, wenn man beim Beten in vernehmlichen Worten zu reden gewohnt ist. Auch hat
\quellseite{150}
Ambulando-Beten eine große Lieblichkeit. Man versuche es nur!

B. Zweitens erinnern wir an den Segen der Fürbitte. Martin Boos, von dem wir vieles lernen können, wenn er gleich nicht in allen Stücken zum Muster aufzustellen ist, behauptet, er habe seine früheren Amtserfolge erbetet. „Hinter dem Choraltar in Wiggensbach,“ schrieb er öfter in seinen Briefen an vertraute Freunde, „habe ich mir meinen lebendigen Glauben und all die Gnaden und Gaben, die mir der Herr schenkte, erfleht.“ (S. Goßners Martin Boos, p. 34.) Was waren seine Gebete um Gnaden und Gaben anders als Fürbitten für die Gemeinde, der zum Heil und Segen alle Gnaden und Gaben, die er wollte, dienen sollten? Aber ganz abgesehen von dem Erfolge wirklicher Erhörung ist doch das Gebet um Amtssegen und die Seligkeit der Gemeindeglieder eine nothwendige Folge wahrer Liebe zu Gott, der solches Gebet will, und wahrer Liebe zur Gemeinde, die des Gebetes bedarf. Wer nicht betet für die Heerde, ist kein rechter Hirte; Fürbitte ist Amts- und Liebespflicht zugleich.\footnote{Dem frommen Pfarrer wird sein Seelenregister zum mahnenden Register nöthiger Fürbitten. Er wird auch wohl thun, für seine täglichen stillen Halbenstunden einen Auszug aus dem Seelenregister zu fertigen und immer neue herzustellen, in welchem alle Personen und Sachen eingezeichnet sind, welche in der Fürbitte nicht vergessen werden dürfen. Unser Liebesgedächtnis ist untreu und trügerisch, bedarf Mahnung und Unterstützung.} Ja, beten ist unüberwindlicher und unabweisbarer Drang eines rechten Hirtenherzens. Ein rechter Hirte
\quellseite{151}
betet, auch wenn er keine Erhörung spürt. Er betet auch nicht bloß um eigene, sichtliche Wirksamkeit. Als Boos Pfarrer am Rhein war, im Jahre 1820, besuchte ihn ein Freund. Den führte er in die Umgegend und sagte: „Sieh da, an diesen Bergen ist fast kein Plätzchen, wo ich nicht schon oft auf meinem Angesichte lag und weinte und flehte, daß mir der HErr wieder die Gnade gebe, mit Freudigkeit meinen Mund aufzuthun und mit Segen zur Erweckung der Herzen sein Wort zu verkündigen, aber ich finde kein Gehör.“ Er konnte nicht anders als beten. Er war auch erhört mehr als er wußte. Und gewiß ist wahr, was sein Lebensbeschreiber sagt (p. 11): „Ohne Zweifel hat der HErr hier — in Sayn am Rhein — mehr als anderswo durch ihn gethan — wenigstens an ihm — er wollte es ihm aber nicht wissen lassen. Es sollte still geschehen.“ — — Merken wir uns das zu unserer Beschämung! Der HErr aber bessere Seine Knechte und lehre sie, bußfertig ihrer Sünden Reinigung zu suchen! Sein heiliger Geist lehre sie beten!

\section{III. Stille und Öffentlichkeit}

69. Strauß sagt einmal mit Recht, daß das Leben des Pfarrers zugleich das öffentlichste und einsamste sei und sein müsse. In der rechten Abwechselung der Einsamkeit und Öffentlichkeit steht und geht nicht nur nach Thomas a Kempis das heilsame Maß eines jeden Christenlebens%
\quellseite{152}%
;\footnote{Nemo secure apparet, nisi qui libenter latet. De Imitat. Christi. L. I. Cap. XX. 2.} sondern sie ist insonderheit ein Geheimniß der oft wunderbar wirkenden Erscheinung eines Pfarrers. Immer still und zurückgezogen zu leben, gestattet das Amt nicht, welches den Geistlichen zum Mittelpunkt, zum Versammlungspunkt der Gemeinde macht und ihn bei jeder Amtshandlung und in Ausübung aller seiner Pflichten mitten unter die Gemeinde stellt. Umgekehrt aber erfordert das Amt nicht bloß nicht ein immerwährendes Leben unter der Gemeinde, sondern es legt Verbot dagegen ein. Sieht man auf den Pfarrer selbst, so wird er durch ein immerwährendes Leben in Mitte seiner Pfarrkinder in Gefahr kommen, das Auge für sie zu verlieren. Wer eine Sache nur aus der Nähe sieht, sieht am Ende bloß sie, verliert die Vergleichung mit andern Dingen, gewahrt nicht mehr ihren Zusammenhang mit anderem, nicht mehr ihre wahre Gestalt. So erkennt der Pfarrer die wahre Gestalt der Gemeinde nicht mehr, wenn er ganz und gar unter und mit ihr lebt. Das ist der Vorzug, welchen der Pfarrer vor dem Schullehrer hat. Der letztere steht näher bei und mitten unter der Gemeinde, kann in Einzelheiten mehr Einsicht in die Verhältnisse ihrer Glieder haben, in solchen das Urtheil des Pfarrers oft berichtigen. Aber das Gesammturtheil über alle und über die einzelnen Pfarrkinder kann der Pfarrer richtiger und besser haben, und dadurch leichter die richtige Behandlung seiner Gemeinde ausfindig machen. Und das ist wichtig! — Was die Gemeinde
\quellseite{153}
anlangt, so wird sie, wenn sie den Pfarrer immer in nächster Nähe sieht, die Achtung gegen ihn schwerer bewahren. Es wird ihr mit dem Pfarrer ergehen, wie dem Pfarrer mit ihr. Gleichwie dieser durch zu große Nähe die wahre Gestalt seiner Gemeinde verlernt, so auch die Gemeinde die wahre Gestalt des Hirten. Alle Beziehungen hören auf, wenn man zu nahe steht. — Was die Amtshandlungen eines in Mitte seiner Gemeinde zu viel lebenden Pfarrers anlangt, so werden sie nicht bloß in Folge der vorbemerkten Wahrheiten an Eindruck verlieren, sondern auch darum, weil ihnen die rechte Vorbereitung und zu dieser die rechte Ruhe und Stille abgeht.

70. Die Erfahrung beweist, daß der Pfarrer bei einem stillen nur durch die amtliche Thätigkeit in der Stille unterbrochenen Leben die richtigste Ansicht von seiner Gemeinde gewinnt. Dieses stille, abgeschiedene, aber dabei immer wache, aufmerksame, beobachtende Leben gibt aber dem Pfarrer auch den Höhenpunkt eines Aufsehers (Episcopus) und Wächters, und dieser Höhenpunkt darf ihm nicht fehlen. Er muß nothwendig über der Gemeinde leben, in deren Mitte er wohnt, wenn er sie und ihre Bedürfnisse überschauen, die ihr nahenden Gefahren bemerken und gegen dieselbe Fürsorge treffen soll. — Auf das Volk macht ein stilles Leben des Pfarrers einen Eindruck des Schicklichen, einen durchaus günstigen Eindruck. Es verlangt von seinen Pfarrern keine Kameradschaft, sondern es spricht übel von dergleichen. Ja, man kann wohl sagen, daß nur der Pfarrer populär im guten Sinne des Wortes sein kann, der den
\quellseite{154}
rechten Punkt der Weite und des Abstands von seiner Gemeinde gefunden hat. Wir wollen zur Würdigung des rechten Benehmens im allgemeinen auf einen hierhergehörigen besondern Fall hinweisen. Ein Pfarrer geht viel aus, sei es auch nur z.\,B. um Hausbesuche zu machen. Ein anderer geht, so oft es seine Amtspflicht erfordert, und macht sich sonst selten. Vorausgesetzt, daß beide freundlich und leutselig sind, behauptet der letztere den entschiedenen Vorzug in der Achtung der Gemeinde. Jener wird in dem Maße weniger wirken, als er öfter kommt; dieser wird, eben weil er seltener kommt, mehr Aufmerksamkeit, mehr Willigkeit, mehr Nachdruck finden. Die selteneren Gänge des letzteren, weil sie immer in bestimmter Beziehung auf Bedürfnisse der Gemeinde stehen und von denselben erheischt sind, erinnern von selbst auch die, welche ihn nur gehen sehen an die himmlische Berufung, welche der Pfarrer repräsentirt. Des ersteren wird man wenig achten, weil mehr Willkür, mehr bloß menschliches Gutmeinen als Amtspflicht seinen Gängen zu Grunde liegt. Er wird bei seinen Pfarrkindern wie einer angesehen werden, der zu geschäftig ist, um anderes als Geschäfte dieser Welt zu treiben. Man frage nur die Erfahrung, ob es nicht so ist.\footnote{„Man muß sich zur Unzeit nicht sehen lassen, um zu rechter Zeit gesehen zu werden.“ Das war der Grundsatz von Martin Boos. Wenn ein Freund ihn einlud, mit ihm im Orte umherzugehen, sprach er: „Laß uns dieses nicht thun; wir müssen uns vor dem Volke nie als in unserem Berufe sehen lassen. Wenn sie uns außer dem Berufsgeschäft zu viel und zu oft sehen, sehen sie uns nicht in unserem Berufe — als solche, wie sie uns sehen sollen. Wir sind Boten des HErrn, darum wollen wir den Leuten aus den Augen bleiben, wenn wir ihnen gerade keine Botschaften zu bringen haben, wollen uns nie zeigen und sehen lassen, als wenn wir Aufträge vom Herrn und Gelegenheit haben, sie auszubringen; damit, wenn uns die Leute auf der Kanzel oder am Krankenbette oder sonst in unserm Beruf sehen, sie glauben, wir seien vom Himmel gefallen, d.\,i. wir seien die übrige Zeit beim HErrn, im Umgange mit Ihm, wir kommen von Ihm her, haben uns das, was wir ihnen sagen und bringen, bei Ihm geholt, gehört, und seien also wahre Zeugen des Herrn.“ Goßners M. Boos p. 33. Man muß Boos verstehen und ihm nichts Thörichtes aus eitler Consequenzmacherei zutrauen, dann hat er gewiß Recht.}

\quellseite{155}
71. Was endlich die Amtstüchtigkeit anlangt, so wird sie gewaltig unterstützt, wenn ein Pfarrer gerne in der Stille ist, um zum Heile seiner Gemeinde immerzu zu lernen. Wer immer geben soll, muß immer haben, und da er, was er zu geben hat, nicht aus sich selbst schöpfen kann, so muß er immer an der Quelle sein, um zu schöpfen. Manche Wahrheit entwickelt sich, wenn man sie einmal hat, schon durch das bloße Leben immer mehr; aber auch ihre Entwickelung ist doppelt und dreifach reich, wenn das Leben still und öffentlich, innerlich und äußerlich zugleich ist. Das göttliche Wort und die theologischen Wissenschaften können ein stilles Studium gewähren, welches, weit entfernt von jener Wissenschaftelerei, die dem Amt entfremdet und für dasselbe untüchtig macht, vielmehr den Pfarrer für das Amt vorbereitet, stärkt, kräftigt und gründet, seiner Erfahrung klärend und verständigend zur Seite geht, immer neue Blicke in die Herrlichkeit der gepredigten
\quellseite{156}
Einen Wahrheit eröffnet und vor dem Unglück bewahrt, daß viele träge Pfarrer haben, daß ihnen ihr Stand, wenn sie seiner einmal mechanisch mächtig geworden sind, verleidet wird, und sie, während von ihren Lippen Honig des ewigen Lebens fließen sollte, vor den Ohren guter Freunde bekennen, daß Christenthum löse ihnen selber die höchsten Fragen nicht.

Ein angehender Pfarrer versehe daher seine Amtsgeschäfte mit dem möglichsten Fleiß. Die übrige Zeit wende er treulich auf Studium und Vorbereitung. Diese Regel dürfte um so leichter zu befolgen sein, als ohnehin, je vertrauter ein Pfarrer mit seiner Gemeinde wird, die Beziehungen sich mehren und immer weniger Zeit und Muße übrig bleibt, der Wahrheit in aller Stille nachzugehen.\footnote{Wer im Volke nicht erscheint, wird nicht wirken, weil er es, — es ihn nicht kennt. Wer überall zu treffen ist, wo es Leute gibt, dem wird das Vertrauen mangeln. Vom Amte zur Studierstube, von dieser ins Amt — das ist der Weg des Pfarrers, und zwar eben so nothwendig, als der Wasser schöpfen muß, der gießen will. Des Lehrens Bedingung ist Lernen. Des Tröstens Bedingung ist Beachten und Ueberwinden der eigenen Anfechtung. Und des Lebens Schwung und Weihe ist ein bethend Herz. Einsamkeit ist der Quellort aller Ströme — und nichts Herrliches wird publice geboren. Nicht immer lernen, nicht immer lehren! Lehren und lernen, so gehts ein- und vorwärts.}

\section[IV. Einrichtung des Pfarrhauses]{IV. Einrichtung des Pfarrhauses, des Tisches, der Kleidung}
\quellseite{157}

72. Es ist hier nicht die Rede von Vorschriften über Diät und Kleidung etc. des Pfarrers, wie sie etwa die Pastoralmedicin, wenn es ja eine besondere Wissenschaft dieses Namens geben sollte, aufstellen dürfte. Wir reden bloß von dem Schicklichen und von der Pfarrersweisheit, welche auch im gewöhnlichen Leben alles auch das Aeußerlichste in Beziehung auf das heilige Amt zu setzen weiß.

Diät und Kleidung und alle Einrichtung des äußern Lebens hängen zunächst vom Vermögen eines Mannes ab. Da nun die meisten Pfarrer aus den armen und mittleren Ständen entsproßen zu sein pflegen, und im Amte nicht leicht einer wohlhabend wird, der seinem Berufe treu lebt, Liebe und Erbarmen ausübt; so sind auch den meisten Pfarrern die rechten Schranken schon durch ihre Umstände gesteckt. Es scheint aber auch im Amte und seiner Beschaffenheit selbst die Weisung zu liegen, daß sich ein Pfarrer in dem Aufwande seines täglichen Lebens den mittleren Ständen gleich stelle. Ein reiches und ein allzuarmes Leben sind beide fürs Amt hinderlich und beschwerlich. Wer alle Tage herrlich und in Freuden lebt, taugt schwerlich zum Seelsorger Lazari und seiner Freunde; und ob er auch mit dem Geiste über seiner Herrlichkeit stände, es würde doch nicht gut sein, wenn seine gesammte Einrichtung und
\quellseite{158}
sein Leben, Reichthum und Fülle zur Schau stellte. Der Reichthum eines Pfarrers darf groß sein, wenn er durch Wohlthun die Armen in einer Gemeinde zufrieden stellen soll. Ein reicher Pfarrer wird schwerlich dem Vorwurf des Geizes entgehen können. So lang er nicht arm geworden ist durch Geben, wird er meistens für geizig gehalten werden, und wenn er selbst über das Maß seines Vermögens hinaus gäbe. Armuth aber ist fast noch schlimmer als Reichthum. Es ist schwer, für die geistlichen Bedürfnisse einer ganzen Gemeinde zu sorgen, wenn sich die Sorge ums liebe Brot beständig an die Fersen hängt. Dazu wird man auch für Geschenke und den moralischen („unmoralischen“) Eindruck derselben viel empfänglicher werden, wenn man arm ist. Es bleibt daher der Befehl des HErrn, daß, die das Evangelium predigen, sich vom Evangelium nähren sollen (Luc. 10, 7 ff. 1. Cor. 9, 14. Gal. 6, 6 ff.), die rechte Anweisung für die Gemeinden, die den Namen „christlich“ haben wollen, — und für Pfarrer ist es das beste Gebet, mit dem reichen König Sprichw. 30, 7 ff. zu sprechen: „Armuth und Reichthum gib mir nicht; laß mich aber mein bescheiden Theil Speise dahin nehmen.“

Gemäß dieser Ansicht von der Lage eines Pfarrers im Allgemeinen möchte auch Diät und Kleidung einzurichten sein. Nicht arm, nicht reich, bescheiden sei alles. Nicht Ueberfluß, nicht Mangel werde zur Schau getragen, wenn anders das letztere zu vermeiden ist. Nicht Ueberfluß, auch wenn er da wäre, denn es reizt die Gemeinde und andere zu Erwartungen und Forderungen,
\quellseite{159}
welche nicht befriedigt werden können, — es macht einen dem Amte widersprechenden Eindruck. Nicht Mangel, wenn es sein kann; denn es erweckt eine Theilnahme der Gemeinde und eine Mildthätigkeit, von welcher man gar zu leicht abhängig wird. Man wird so leicht zum Bettler, wenn man arm ist.\footnote{Eine Bürgschaft für Einhaltung der rechten Mitte gewährt eine einfache Gewöhnung, ein von Bedürfnissen der Weltkinder freigewordenes, genügsames, in der Uebung eines armen Lebens genügsames Herz. Wenn nichts am Herzen liegt, was die Welt bedarf, der sieht die wahren Lebensbedürfnisse und sucht die Befriedigung derselben ohne unreines, selbstsüchtiges Wesen.}

73. Vielleicht würde man die ganze Einrichtung eines Pfarrers am besten mit dem Namen Einfalt bezeichnen können. Von dem Modetone der Welt und altväterischem, die Zeit nicht achtenden Tone gleichweit entfernt sei Diät, Kleidung und Umgebung eines Pfarrers. Weder der Reiche, noch der Arme werde abgeschreckt und finde des Pfarrers Lebensart von der seinigen allzuferne. Jenem, wie diesem, wenn jener nicht verzogen, dieser nicht zu tief heruntergedrückt ist, müsse es in der Wohnung und im Familienkreis des Hirten wohl werden können. Das geistige Interesse des Lebens walte vor und mache das vergehen, was an Gemächlichkeit und Fülle des zeitlichen Lebens mangelt. Die Fülle an ewigen Gütern, welche vorhanden ist, ersetze alles. Alles in einem Pfarrhause predige Unabhängigkeit vom Zeitlichen, Hingabe an das Ewige. „Schlecht und recht, das behüte mich, denn ich harre dein“. Ps. 25, 21. sei eines Pfarrers Wahlspruch wie in allem persönlichen
\quellseite{160}
Benehmen, so in dem, was die Aufschrift dieses Paragraphs besagt.

74. Eine Ermahnung zur Ordnung im Pfarrhaus möchte eben so wohl passend als überflüssig genannt werden können. Es versteht sich von selbst, daß die Diener Dessen, welcher selbst ein Gott der Ordnung ist, die Ordnung lieb haben. Wer inwendig in Erkenntniß und Begehren geordnet ist, sucht auch auswendig und um sich her Ordnung herzustellen. Die Erfahrung zeigt, daß ein gewisser Grad von innerer Verwirrung und Unordnung durch Ordnen in äußerlichen Dingen gehoben werden könne. Ein Mann, der viel mit Correspondenz zu schaffen hatte, beseitigte Unruhe und Unklarheit der Seele oftmals durch Ordnen seiner Briefe. Daraus erkennt man die Wechselwirkung, welche zwischen Aeußerem und Innerem besteht, und eben daraus erklärt und rechtfertigt sich das Mißtrauen, welches man gegen unordentliche Männer und Frauen und sonderlich gegen unordentliche Pfarrersfamilien hat.

75. Was insbesondere die Kleidung anlangt, so ist es wenigstens Sitte geworden, daß Geistliche sich dunkle Farben, namentlich die schwarze beilegen. In Ländern, wo man vorzugsweise schwarz zu gehen pflegt, ist freilich der Geistliche damit nicht ausgezeichnet. Es fragt sich, ob er in diesem Falle, oder auch sonst, eine Auszeichnung durch die Form des Kleides suchen solle? Man verwechsele diese Frage ja nicht mit der andern: „ob der Geistliche in der Auszeichnung durchs Kleid etwas suchen solle?“ Die letztere ist leicht verneint; die erstere kommt in Anregung
\quellseite{161}
so oft ein Geistlicher ein neues Kleid bedarf. Die Frage ist unwichtig, wenn gleich im practischen Leben vielleicht doch nicht so gar unwichtig, als es scheinen könnte. Jeden Falls will sie, wenn sie einmal ihres Ortes gethan ist, recht beantwortet sein.

In der neueren Zeit hat man es unpassend gefunden, daß Geistliche im Chorrock „Aufwartungen“ machen, und man hat darauf gedacht, ihnen für feierliche Gelegenheiten, wo sie nicht im Chorrock erscheinen sollen, ein eigenes Kleid zu geben, welches ihnen aber doch fürs gewöhnliche Leben nicht gestattet sein sollte. Man hat aber vergessen, daß man damit den großentheils armen Pfarrern eine empfindliche Ausgabe verursacht, welche vermieden würde, wenn man ihnen entweder gestattete, im Chorrock „aufzuwarten“, oder im Schnitte des Visitengewandes für gewöhnlich zu erscheinen. Das letztere möchte jeden Falls das bessere sein. Der Chorrock bezeichnet die amtliche Function, nicht den Stand; eine „Aufwartung“ hat mit amtlichen Functionen nichts zu schaffen. Sofern nun aber billig das ganze zeitliche Leben des Pfarrers in Harmonie mit dem amtlichen Leben geführt werden und wie aus einem Guß erscheinen soll, und dies der allgemeine Anspruch der Gemeinde eben sowohl als die innere Forderung des Gewissens und das Gebot des göttlichen Wortes ist; dürfte wohl auch für die Kleidung des Geistlichen mehr, als bloß die negative Bestimmung, „Unanständiges und Modisches zu vermeiden“ gegeben werden. Es ist nur ein Schritt weiter, eine positive Bestimmung zu geben, — und man sollte sie, ohne viel Wesens
\quellseite{162}
von einer so äußerlichen Sache zu machen, ganz einfach geben.

Freilich hat man bei uns für dergleichen vielfach ein so auffallendes Ungeschick, daß man am Ende am besten thäte, dem Geistlichen den breitkrämpigen Hut und den langen schwarzen Ueberrock für das gewöhnliche Leben und für Aufwartungen zu gestatten. Diese Kleidung wäre eben so würdig und dem Beruf entsprechend, als sie eine allzu auffallende Unterscheidung von andern Ständen vermiede. Soll der Geistliche als solcher erkannt werden, so würde ein stehender oder doppelter Kragen und ein längeres Herabgehen der beiden Reihen großer Knöpfe Kennzeichen genug sein. — Für die schmutzigen Wege auf dem Lande würden sich lange Stiefel jedenfalls besser eignen, als lange Hosen. Ein besseres Aussehen geben sie ohnehin. — Für Aufwartungen könnten ja allenfalls Schuhe vorbehalten sein.

Jeden Falls verdient der eben so unschickliche als abgeschmackte Frack ein Todesurtheil.

\section{V. Liebhabereien}

76. Häufig findet man, daß sich Pfarrer Liebhabereien hingeben. Kaum wird man irgend ein Gebiet menschlichen Wissens oder Könnens erwähnen, ohne daß einem auch bei nur mäßiger Bekanntschaft unter Geistlichen irgend ein Pfarrer, der sich damit befaßt und dadurch bemerklich
\quellseite{163}
gemacht hätte, einfiele. Philosophie, Philologie, Jurisprudenz, Kameralia, Politik, Naturwissenschaften aller Art; Botanik, Mineralogie, Geologie, Pomologie, Blumisterei, Oekonomie, Jagd, Fischerei 2c. 2c.: kurz, was man nur nehmen kann, ein Pfarrer hat sich darin versucht und hervorgethan. — Es fragt sich nun, ob man das loben oder tadeln, — ob man einen angehenden Pfarrer davor warnen oder zur Nachfolge ermuntern soll?

77. Gewiß ist, daß keine Wissenschaft an und für sich selbst ein so weites Gebiet hat und für die übrigen Wissenschaften so viele Anknüpfungspunkte darbietet, als die Theologie. Keine ist an und für sich so hinreichend, auch den reichsten Geist zu beschäftigen; aber auch keine so einladend, fremde Gebiete zu betreten, als gerade sie. Bei keiner verstünde sich’s so von selbst, daß man auf dem eigenen Gebiete bliebe, als bei ihr; — und doch ist sie es wiederum, die am ersten einen Seitenweg und Seitenblick entschuldigte. Ein theologischer Sinn findet überall Bestätigung der theologischen Wahrheit; aus allen Wissenschaften kann Gewinn für die Theologie gezogen werden. Aber auch umgekehrt: die Theologie bedarf keiner andern Wissenschaft, um fest zu stehen.

Bei so gestalteten Umständen wird man nun allerdings keinen Pfarrer tadeln können, der irgend ein, zumal ein vorzugsweise verwandtes Gebiet betritt, um von seinem eigentlichen Berufsstudium daselbst auszuruhen und neue Spannung zu gewinnen. Jeder angespannte Bogen, der tüchtig bleiben soll, bedarf der Abspannung. Es geht eben nicht, daß man immer einerlei treibe; am wenigsten geht
\quellseite{164}
es bei den höchsten Regionen des menschlichen Wissens. St. Johannes bedurfte des Rebhuhns, und selbst Adam im Paradiese bebaute das Land bei seinem Leben voll himmlischer Beschauung. — Aber tadelhaft ist und bleibt es ohne Zweifel, wenn die Nebenbeschäftigung zur Hauptbeschäftigung, zur Liebhaberei und Leidenschaft wird. Ein Pfarrer hat wol über sich zu wachen, daß nicht irgend Etwas ihn von der Liebe zum heiligen Amte und dem theologischen Studium abwende. Große Sünde, großer Undank ist es, wenn ein Mann, dessen Aug und Herz Befehl hat, sich an Betrachtung der hochgelobten Schönheit des Dreieinigen zu weiden, diesem Berufe Geringeres vorzieht. Er warte des Seinen und laße andern christlichen Gelehrten die Aufgabe, ihre Wissenschaften zum Heile der Kirche, zur Bestätigung ihrer ewigen Wahrheit auszubauen. Der im Amte trägste Pfarrer ist dennoch vom Amte so gehalten, daß er sich nicht mit voller Kraft auf eine andere Wissenschaft hinwenden kann. So viele Beispiele von Pfarrern es daher gibt, die sich mit fremder Wissenschaft etc. befaßt haben, so selten findet man doch Beispiele, daß einer auf einem fremden Gebiete etwas Großes geleistet hat. Wer Allotria treibt, ist von seinem Berufe nicht durchdrungen, leistet in ihm nicht, was er soll, und in den Allotriis am Ende auch nicht viel. So verschwendet er das Leben und setzt sich der schrecklichsten Verantwortung aus.

78. Besonders mögen sich unverheirathete oder kinderlose Pfarrer vor Liebhabereien hüten. Denn die Erfahrung beweiset es, daß diejenigen, welche nicht auf
\quellseite{165}
den von Gott gesegneten Pfaden geordneter, wahrhaft menschlicher Liebe zu Weib und Kindern wandeln oder wandeln können, geneigt sind, ungeordnete Liebe zu erwählen. Der unedlere Mensch verfällt alsdann auf die Liebe zu großen Sünden; der edlere auf scheinbar subtilere Irrthümer, auf geistige Liebhabereien: beide aber sündigen dennoch vor Gott ohne großen Unterschied. Darum hat der ehe- oder kinderlose Pfarrer seine Kräfte desto mehr dem Dienste der Kirche zu widmen, 1. Cor. 7, 32. Alle aber haben zu bedenken, daß sich zwar keinem Amte so viel Schlendrian und Trägheit anhängt, als dem geistlichen, daß aber auch keines so große Forderungen an seine Diener macht. Je höher das Amt, desto leichter kann es gemißbraucht, aber auch desto schwerer recht gebraucht werden. Das geistliche Amt bedarf der vollen, ungetheilten Liebe und Kraft derer, die es tragen, wenn es nur einigermaßen recht geübt sein soll.

Es ist darum nichts natürlicher, nichts schöner, als daß ein Pfarrer seinem Amte und seinem Amtsstudium lebe; und sonst keiner Sache! Das Rebhuhn Johannis ungetadelt!

\section[VI. Geschenke, Sammlungen, Beichtgeld]{VI. Geschenke, Sammlungen, Beichtgeld, Accidentien}

79. Dem Pfarrer werden fast überall mehr oder weniger freiwillige Gaben verehrt. Manche freiwillige Gaben sind zum Brauch, zur Observanz geworden, z.\,B. Flachssammlungen, Eiersammlungen und das Beichtgeld. Rücksichtlich dieser Dinge eben sowohl, als rücksichtlich seines
\quellseite{166}
Ausgebens muß ein Pfarrer jeden Falls beim Eintritt ins Amt Bescheid wissen, wenn er nicht gleich vornherein sich seine Stellung zur Gemeinde verderben will.

80. Was die Geschenke anlangt, so werden sie entweder aus purer Dankbarkeit und Liebe oder mit der Absicht, etwas zu erreichen, gegeben. Der Pfarrer kann als Schulinspektor vielleicht einmal eine Gunst erweisen, bei Kirchstuhlverleihungen etc. eine Rücksicht nehmen etc.; dergleichen sucht dann das Volk mit Geschenken. Gut ists, wenn ein Pfarrer den Unterschied macht, daß er in allen Fällen, in welchen sich eine Nebenabsicht indicirt, keine Geschenke annimmt, weder wenn die Absicht, etwas zu erreichen, zu vermuthen steht, noch wenn sie als Danksagung für erfüllte seelsorgerische Pflichten erklärt werden. Der Ruhm der innocentia kommt einem Pfarrer oft so sehr zu Statten, daß er mit einer gewissen Härtigkeit gesucht zu werden verdient. St. Paulus weiß und verehrt des Herrn Gebot, daß vom Evangelium leben soll, wer das Evangelium predigt, dennoch erklärt er es für seinen Ruhm, den er sich von niemand nehmen lasse, daß er den Gemeinden umsonst und ohne Lohn gedient, sich seinen Lebensunterhalt selbst verdient habe. Die Umstände, welche ihn dazu nöthigten, finden sich auch bei andern oftmals vor, aber nicht ebenso die uneigennützige armuthsfrohe Gesinnung des heiligen Apostels. — Ganz etwas anderes ist es jedoch mit jenen arglosen Geschenken, welche man dem Pfarrer bei besondern Gelegenheiten und zu bestimmten Zeiten von Alters her zu geben pflegt. Diese soll und darf man eben so wenig zurückweisen, als man ein
\quellseite{167}
Wohlgefallen daran äußern darf, daß nach mehr schmeckt. Nicht bloß ist oft das Geschenk eine wirkliche Liebesweisung, der man sich nicht roher Maßen entziehen darf; sondern es ist zuweilen in der That eine Art von Stolz, sich nicht beschenken lassen zu wollen. Die Verehrung, welche dem Pfarrer durch Geschenke geschieht, deutet auf das schöne, heilige, patriarchalische Verhältnis hin, in welchem ein Pfarrer zu seiner Gemeinde stehen soll. Indem man sie in ihren Zeichen und Aeußerungen zurückweist, tödtet man sie selber, — und zwar dann nicht bloß für die Zeit der eigenen Amtsverwaltung, sondern auch für die der Nachfolger, welchen das rechte Verhältnis zur Gemeinde zu zerstören und überdies zeitlichen Verlust an gerechtem Einkommen zuzufügen, man überall kein Recht hat. Es gibt eine Weise, Geschenke anzunehmen, bei welcher man in schönster Dankbarkeit zu seiner Gemeinde stehen kann, ohne im mindesten abhängig zu werden. Diese suche man!

81. Aehnlich ist es mit den §. 79 erwähnten observanzmäßigen Sammlungen von Flachs, Eiern u.\,dgl., welche noch jetzt in vielen Gegenden bestehen. Man kann sie auf eine geizige und habsüchtige Weise vornehmen, durch welche das Volk oft so sehr beschwert wird, daß man es in seiner großen Geduld und Langmuth oft wirklich bewundern muß. Man kann sie aber auch auf eine stolze und übermüthige Weise vornehmen, bei welcher man die Gabe mit großem Unwillen und ohne allen Dank dahin nimmt und zu verstehen gibt, daß man sich, wenn es nicht um des Nachfolgers willen geschähe, gar nicht die Mühe nehmen würde. Viele haben auch gesucht, dergleichen Sammlungen
\quellseite{168}
in eine Geldabgabe zu verwandeln, welche sie dann im Ganzen aus dem Gemeindesäckel erhoben oder durch den Schultheiß einsammeln ließen. Alle diese Arten und Weisen, Sammlungen vorzunehmen oder zu umgehen, sind gewiß nicht zu billigen. Die erste habsüchtige Weise nicht, das begreift jeder; aber auch die andern nicht. Ein demüthiger Pfarrer verachtet es nicht, seine Sammlungen selbst vorzunehmen, oder doch sie durch die Pfarrerin, die hier ganz an ihrem Platze ist, vornehmen zu lassen. Er erscheint gerne als ein Bittender und Dankender und freut sich, durch sein Beispiel Bitte und Dank zu lehren und Liebe im Nehmen zu erweisen, wie er sie gerne auch im Geben erweist. Ja, gerade diese Sammlungen geben Anlaß zu den freundlichsten persönlichen Begegnungen mit den Gemeindegliedern. Da steht und spricht ein jeder den Pfarrer oder die Pfarrerin in einem ganz lieblichen, rein menschlichen Verhältnis. Da kommt man vermöge einer durch das Herkommen geheiligten Gewohnheit in alle Häuser und findet offenere Herzen — mit Ausnahme derer freilich, die den Pfarrer fliehen, um nicht geben zu müssen. Oft wird auf diesem Wege ein geistliches, oft ein leibliches Bedürfnis der Gemeindeglieder erkundet. Der Arme, der nichts geben kann, aber bei dieser Gelegenheit doch besucht wird, spürt die Freundlichkeit, gewinnt Zuneigung und Liebe zum Pfarrer, eröffnet manchmal sein Herz. Die Kinderzucht und überhaupt das Verhältnis der Eltern zu den Kindern kann erforscht und dem Seelsorger zu Hilf und Tröstung bekannt werden. Kurz alle Tugenden, welche Hausbesuchen u.\,dgl. persönlichen Berührungen des Pfarrers mit seinen Beichtkindern
\quellseite{169}
zugeschrieben werden können, haben die bei Gelegenheit der Sammlungen angestellten Besuche, während sie manche Nachtheile nicht haben, die sich bei weniger berechtigten und motivirten Besuchen finden. — Man kann bei diesen Gelegenheiten sogar viel Barmherzigkeit üben, indem man die Gaben nicht nimmt,\footnote{Es ist Schwachheit, dergleichen Geschenke oder auch Accidentien und Stolgebühren insgemein von sich zu weisen. Es wird nicht einmal gut angesehen, für Ehrgeiz, ja gar für berechnenden Geldgeiz gehalten und schadet dem Nachfolger. Dagegen kann man sich in einzelnen Fällen seines Rechtes mit desto größerem Nachdruck begeben. Das wird mit Dank angenommen und man bindet damit weder sich selbst, noch andern die Hände. Cf. S. J. Baumgarten l.\,c. S. 341. — Häufiger als eine zu weit gehende Großmuth der Pfarrer im Erlassen der Gaben und Accidentien findet sich eine zu große Strenge in deren Beitreibung. Wenn auch kein Mensch ableugnet, daß ein Pfarrer in gewissen Fällen wegen seiner Gebühren vor Gericht klagen dürfe; so werden diese Fälle doch nur ganz selten eintreten dürfen. Es ziemt sich, möglichst selbstverleugnend und mild zu verfahren. Cf. Deyling Prud. past. p. 242.} wohl auch selbst noch allerlei Gaben gibt. Indem man als ein Nehmender von Haus zu Haus geht, kann man ein Engel Gottes sein, der im Verborgenen, ohne daß es vermuthet oder erwartet wird, wohl thut. Schon um deswillen soll man dergleichen Sammlungen nicht von sich weisen. Am allerwenigsten aber soll man sie in eine Geldsumme verwandeln; bequem ist das wohl, aber es ist damit auch aus einer Gelegenheit, die von Segen triefen konnte, für immer weiter nichts als eine kleine Nutznießung des Pfarrers geworden. Herzlos und lieblos erscheint es, nur um ein wenig Mühe und Demüthigung zu ersparen,
\quellseite{170}
sich außer jenen freundlichen Verkehr mit der Gemeinde zu setzen. Für sich selbst sorgt ein solcher Pfarrer, nur das Seine sucht er, und, wenn er es für immer einführt, Geld statt Gaben zu nehmen, so verhindert er seine Nachfolger auf eine unverantwortliche Weise, das Schönere und Bessere zu thun. — Es hat ja ohnehin das Leben eines jetzigen Pfarrers nur noch wenige Spuren des ihm gebührenden herzlichen, patriarchalischen Wesens. Warum ihm denn auch diese letzten Spuren nehmen, die in der Hand rechter Pfarrfamilien sich so schön und segensreich gestalten können?

82. Was das Beichtgeld anlangt, so ist es etwas ganz anderes, es einzuführen, als es abzuschaffen. Es einzuführen, wo es nicht ist, wäre gewiß nicht weise. Aber eben so wenig weise ist es, es abschaffen zu wollen, wo es noch besteht. Man führt in der Regel dagegen an, daß es allzuleicht als Bezahlung der Absolution genommen werden könnte. Allein der Irrthum ist dem Volke durch Belehrung leicht zu nehmen. Er ist zu abgeschmackt, als daß er nicht leicht eingesehen werden sollte. Ist er nun erkannt oder ist ihm durch Belehrung vorgebeugt, so ist nicht abzusehen, was am Beichtgeld so gar Schlimmes sein sollte? Es ist ein Geschenk, welches man dem Beichtvater zum Dank für gehabte oder zu übernehmende Mühe macht. Und was ist da Uebles dran? Sind wir so gar verschämt und zart — Beichtväter und Beichtkinder —, daß uns ein kleines Geschenk (denn was ists ganz und gar, ein Groschen oder wie viel mehr an den meisten Orten?) schon weh thut, wenn es nur in der Nähe
\quellseite{171}
eines großen unschätzbaren Geschenkes gegeben wird? Man sollte ja vielmehr denken, daß in der Nähe ewigen Gutes aller Schein einer Bezahlung durch eine irdische Kleinigkeit verschwinden müßte!

82. Eben so ist es mit den sogenannten Accidentien oder Stolgebühren.\footnote{Cf. Deyling P. II. Cap. III. §. XVII. Hartmann Past. evang. L. IV. C. I. L. III. C. XXIII. p. 641. sq. Seckendorffs Christenstaat III. B. VI. Cap. §. 3. Speners theol. Bedd. I. * 71. II. 438. IV. 6, 6. Letzte Bedenken I. 601. 606. Böhmers Jus paroch. Sect. VII. C. II. Jus eccl. T. I. P. VII. 2. 5. 610. II. 3. 5. 176. 300. I. P. VII. 2. 11. Brunnemann Jus eccles. L. II. C. V. §. 7. Balduini cas. consc. C. X. Cas. V. d. 1120. Dunte Cas. consc. p. 432. 433. 513. Roque’s Geist eines evang. Lehrers 1. Thl. 1. Versuch. §. LX–LXIV. — Cf. Baumgartens casuist. Pastoraltheologie S. 333 ff.} Man belehre die Leute, und das Anstößige fällt weg. So wenig das apostolische Wort auffallend sein kann, daß der vom Evangelium leben soll, der es lehrt, so wenig kann in einer Gabe, welche nach bestehenden Ordnungen bei Gelegenheit gewisser amtlicher Handlungen zum Unterhalt des Pfarrers gereicht wird, etwas Auffallendes liegen. Quando igitur pro laboribus sacris salarium in genere recte datur, quidni etiam pro actu speciali aliquid dari posset? (Deyling Instit. Prud. past. ed. III. p. 240.)

84. Was die Strenge der Forderung anlangt, so haben unsere Väter zwischen den Accidentien und den Stolgebühren einen einflußreichen Unterschied gemacht.

\begin{quote}
„Latius aliquando patet nomen accidentalium quam jurium stolae. Illa enim vox in genere de omni
\quellseite{172}
emolumento, quod ratione officii extraordinarie percipitur, accipi solet, sine respectu ad officium aliquod speciale praestandum, e.\,g. quod a parochianis ex mera liberalitate tempore novi anni ministris ecclesiae datur: quod nullo jure perfecto exigi potest. E contrario jura stolae jure perfecto exiguntur ob speciale aliquod officium praestitum.“ S. Deyling p. 239 f.
\end{quote}

Das ist: der Begriff der Accidentien fällt zwar gewisser Maßen mit dem der Stolgebühren zusammen und steht alsdann dem des salarium substantiale, welches certas pensiones et praestationes continet, gegenüber; sofern aber Accidentien und Stolgebühren unterschieden werden, ist der Begriff jener der weitere, indem er auf die freiwilligen Geschenke geht, die ohne Bezug auf eine bestimmte Amtshandlung zu bestimmter Zeit gegeben werden, während die Stolgebühren für bestimmte Amtshandlungen verabreicht werden. Daß man Accidentien nicht fordern könne, leuchtet demnach ein; die Stolgebühren aber kann man rechtlich fordern, wenn sie verweigert werden. Doch wird unter ihnen selbst wieder ein Unterschied gemacht. Je nöthiger und wichtiger für das ewige Leben eine Amtshandlung ist, desto weniger darf der Pfarrer die Verrichtung derselben von der Zahlung abhängig machen oder hernachmals sich zu Klage und Streit wegen Verweigerung der Gebühr verleiten lassen. Die Taufe, die Absolution, das heilige Abendmahl kann unter keiner Bedingung um Gebührenverweigerung willen aufgeschoben werden. Andere Gebühren mag der Pfarrer eher fordern.

\begin{quote}
„Ex juribus stolae parochus quaedam suo jure exigit virtute vel legis vel consuetudinis rite introductae.
\quellseite{173}
Quo pertinent proclamatio, copulatio sacerdotalis, preces publicae, deductio funerum, parentatio vel concio funebris. Alia autem accidentia, e.\,g. pro baptismi collatione ac S. Coenae administratione, quin ob nummum confessionalem sponte oblatum accipere quidem, sed invitis extorquere non potest.“ S. Deyl. l.\,c. p. 240 f.
\end{quote}

Vielleicht wäre es gut, das „invitis extorquere non potest“ auf alle eigentlichen Stolgebühren auszudehnen.

85. Jeden Falls sollte in unserer Zeit eine Stolgebühr ausnahmsweise vor den andern entweder abgethan oder auf irgend eine die Armen nicht drückende Weise ersetzt werden. Wir meinen die Gebühr für Communionen der Kranken und Sterbenden. Weil eine Gebühr zu entrichten ist, muß der Kranke den Genuß des Sacraments entbehren, bis alle Lebenshoffnung ausgelöscht ist und eben damit die Hoffnung, das Abendmahl in der Kirche mit weniger Kosten zu genießen. Wie manchmal ist aber dann der Kranke nicht mehr fähig, das Sacrament zu empfangen! Und wie viele gehen um der Gebühr willen, ohne den Segen des Sacraments neu empfangen zu haben, aus der Welt! Das und dergleichen sind unerträgliche Fälle. Es liegt dabei wenig Besserung darin, daß den Armen die Gebühr erlassen wird. Wer ist arm, wer nicht? Diese Frage läßt sich verschieden beantworten; es kann auch verschiedene Praxis geben. Dadurch bleibt denn doch für die meisten die Furcht vor der Gebühr, wie die Erfahrung beweist. Auf dem Kranken-, dem Sterbebette aber sollten solche Fesseln und Bande nicht mehr binden; da sollte der Genuß des Sacraments so frei gegeben und so leicht gemacht
\quellseite{174}
werden als möglich. Und darum sollte die unleidlichste unter allen Gebühren, die für die Krankencommunion, wo nur immer möglich, fallen.

86. Vielleicht könnte überhaupt für den Unterhalt des Pfarrers auf andre Weise gesorgt werden, als durch Accidentien und Stolgebühren. Zwar stammen sie von den Oblationen der Alten, welche sie mit Verschonung der Gäste und Zuschauer bei einzelnen, sie allein betreffenden Amtshandlungen der Pfarrer zum Besten der Armen, zur Besoldung der Pfarrer etc. auf den Altar niederlegten; — und Oblationen sollte man lieber herstellen, als vollends sammt der letzten Erinnerung ausrotten. Wir könnten versucht sein, mit Chrysostomus (Homil. 86 in Matth.) die Rückkehr zur alten Sitte zu empfehlen, auf den Unterhalt der Geistlichen durch freie Liebe anzutragen. Nunquam ecclesia earumque ministris melius fuit prospectum, quam ubi sustentationem suam ex oblationibus habuerunt, quae admodum fuerunt largae et copiosae. Deyling p. 224. Allein leider ist, was die Mehrzahl unsrer Gemeindeglieder anlangt, Liebe und Gemeinschaft (κοινωνία) erstorben und ganz dahin. Die Pfarrer würden jetzt vielleicht darben müssen, wenn sie von Liebesgaben leben sollten. Müßten sie doch um ihren armen Unterhalt oft genug trotz bestehender Rechte mit den Gemeinden erst hadern!\footnote{Doch darf man nicht vergessen, daß die Gemeinden gegenwärtig nicht im Fall sind, ihre Pfarrer versorgen zu sollen. Wäre der Fall gegeben, so würde dadurch auch der Sinn zum Geben und zum Unterhalt des Pfarrers beizutragen, ins Leben gerufen werden. Unsere Dotationen und Rechte lassen den Sinn nicht aufkommen, der, wie Beispiele zeigen, sich lebendig erweisen würde, wenn Noth an Mann ginge.} Darum
\quellseite{175}
wird auch wohl in der Regel kein Segen verschüttet, wenn Accidentien und Stolgebühren dahin fallen und ein anderer Weg der Versorgung von Pfarrern eingeschlagen wird, wofern es möglich ist. Hat es gleich bei rechten Pfarrern nicht bloß nichts Anstößiges, sondern sogar etwas recht Herzliches und Inniges, Liebe und Gaben der Gemeindeglieder bei solchen Gelegenheiten auffallender zu erfahren und dahinzunehmen, in denen auch der Pfarrer den Einzelnen auffallender dient, so könnten doch vielleicht auch andere Gelegenheiten, dem Pfarrer Liebe und Wohlthat zu erweisen, aufgefunden werden, die gleichfalls ein inniges und herzliches Anmahnen zuließen. Die hohen Feste des Jahres, der Kirchweihtag, der Einsetzungstag des Pfarrers \&c. z.\,B. gäben auch Gelegenheit, wenn schon immerhin so herzlich, wie die Gelegenheiten, bei denen gegenwärtig die Gaben gespendet werden, die erwähnten nimmermehr sind. — Man möchte zuweilen wünschen, daß Beichtgeld und Accidentien vermieden werden könnten, wenn man das Jagen nach Beichtkindern, die man doch verwahrlost, und den unwürdigen Handel mit Accidentien bedenkt, der hie und da an der Tagesordnung ist. Wenn nur auch ein menschliches Verhältnis denkbar wäre, in welchem gar kein Mißbrauch, gar kein schändliches, eigennütziges Benehmen denkbar wäre! Aber das ist nicht einmal bei einer durchaus fixen Geldbesoldung möglich. Auch wenn man um des abscheulichen Mißbrauchs willen jeden
\quellseite{176}
möglichen heilsamen Gebrauch der bestehenden Einrichtung daran geben wollte, würden doch Habsucht und die Feilheit immer neue Wege finden. Es ist darum doch am besten, bei dem alten Brauch zu bleiben und den Mißbrauch zu brandmarken.

\section{Anhang zum vorigen Paragraph}

87. Da die Sache mit den Accidentien und deren Aufhebung, namentlich in unserer Zeit, eine keineswegs unwichtige ist; so erlauben wir uns, aus S. J. Baumgartens „Kurzgefaßter casuistischer Pastoraltheologie“ (Halle 1752) die treffende Frage und Antwort hier einzurücken. Wir lesen p. 332 ff.:

\begin{quote}
„Können Accidentien oder Bezahlung gewisser Amtsverrichtungen rechtmäßig sein und gefordert werden?“
\end{quote}

„Diese Frage ist mit Ja zu beantworten. Bei der Entscheidung aber sind drei Stücke zu beobachten:\\
1. die Einschränkungen der zu entscheidenden Frage;\\
2. die Entscheidungsgründe;\\
3. die Einwürfe und deren Beantwortung.

Die Einschränkungen der zu entscheidenden Frage sind folgende. Die erste enthält eine Erklärung der Sache selbst, was Accidentien seien? Wir wollen zwei Merkmale angeben,
\quellseite{177}
dadurch diese Einkünfte eines Lehrers von den übrigen ordentlichen unterschieden werden.

Das erste Merkmal ist: daß diese Einkünfte von einzelnen Zuhörern entweder zu einer gewissen verordneten Zeit, oder bei gottesdienstlichen Handlungen, so an einzelnen Personen geschehen, den Lehrern gezahlet werden.

Diese Art der zufälligen Einkünfte gottesdienstlicher Lehrer sind aufkommen bei Gelegenheit der öffentlichen Oblationen oder gemeinschaftlichen Beitrags in der ersten Kirche, so alle Glieder beim ordentlichen Gottesdienst und Haltung des Abendmahls zu thun gewohnt gewesen. Denn, weil diese Oblation gleich vom Anfange zum Unterhalt der Lehrer und Armen bestimmt worden; so ist es nach und nach geschehen, daß bei Handlungen einzelner Personen, dieselben, um andere gebetene Gäste mit Unkosten zu verschonen, diese Oblation allein überkommen und, da sie gering gewesen, den Lehrern ganz allein überlassen, ihre Willigkeit zu dergleichen Handlungen zu befördern. Worauf hinlänglich erhellet, daß dergleichen dem ersten Ursprung nach keine Bezahlung göttlicher Gnadenmittel gewesen.

Das zweite Merkmal ist, daß solche Einkünfte durch Kirchengesetze entweder befohlen oder doch bestätigt sein müssen, um selbige von selbst eingeführten zu unterscheiden.

Die zweite Einschränkung betrifft eine Vorschrift, so bei Förderung derselben zu beobachten ist, nemlich: es muß alle Verletzung der allgemeinen Pflichten, auch der Schein des Geizes vermieden werden; folglich Keiner gottesdienstliche Verrichtungen vervielfältigen und den Zuhörern zu seinem Nutzen aufbürden.

Die dritte Einschränkung ist, daß man die Willigkeit, den Zuhörern mit seinem Amte zu dienen, nicht darnach abmesse, folglich weder Alles bezahlt haben wolle, wodurch Manchem eine Furcht vor Besuchung des Predigers ankommt, noch auch die Bemühung selbst und den Grund derselben nach den Vortheilen,
\quellseite{178}
die man erhält oder wenigstens hofft, abmesse und einrichte.

Unter diesen Einschränkungen ist es rechtmäßig, dergleichen Accidentien zu geben und zu nehmen, aus folgenden Gründen:

1. Weil es ein Stück und Theil des Unterhalts der Lehrer ist, folglich die Verbindlichkeit der Zuhörer, ihre Lehrer zu erhalten, sich mit darauf erstrecket, sonderlich wenn der übrige ordentliche Unterhalt nicht hinlänglich ist.

2. Weil dergleichen Verordnungen mit zur rechtmäßigen Befugniß einer gottesdienstlichen Gesellschaft oder Obrigkeit gehören, folglich auch von Lehrern ohne Versündigung beobachtet werden können.

3. Weil GOtt selbst im alten Testament beim israelitischen Gottesdienst verordnet, daß außer denen gewöhnlichen Zehenten auch ein Theil der Opfer an die Priester abgegeben werden mußte, woraus erhellet, theils, daß dergleichen an sich nicht unrechtmäßig sein könne; theils, daß es nach der Weisheit GOttes für eins der bequemsten Mittel angesehen worden, diesen Unterhalt den Lehrern zu verschaffen.

Damit aber folgende Einwürfe nicht Platz finden.

Erstlich ist es nicht eine Verkaufung göttlicher Dinge, wie man vorgibt, und daher Gelegenheit nimmt, dergleichen mit allerhand gehäßigen Benennungen zu belegen, sondern nur eine Erkenntlichkeit für die Bemühung der Lehrer in einzelnen Fällen, welche nicht unbillig ist, indem keine allgemeine Verbindlichkeit gezeiget werden kann, einzelnen Personen gottesdienstliche Handlungen nach ihrem Gefallen zu verrichten. Ueberhaupt beweiset dieser Einwurf zu viel, nemlich die Unrechtmäßigkeit des ganzen Unterhalts der Lehrer von der Gemeinde, welcher doch in ausdrücklichen Befehlen der heiligen Schrift gegründet ist.

Zweitens streiten auch die Worte Christi Matth. 10, 8. damit nicht, denn 1. zeiget der Zusammenhang, daß der Heiland nur von wunderthätigen Gaben rede; 2. beantwortet der
\quellseite{179}
Heiland selbst gleich im Folgenden diesen Einwurf, da er seinen Jüngern befiehlt, von Jedermann, wo sie hinkommen würden, den Unterhalt anzunehmen, ohne Ausnahme der Gutwilligkeit solcher Leute, bei denen sie sich befinden würden.

Drittens erweisen die Worte Petri Apg. 8, 20. nicht, daß es sündlich sei, für gottesdienstliche Handlungen Etwas anzunehmen, sondern wie der ganze Zusammenhang dieser Geschichte zeigt, nur die Sündlichkeit und Unrechtmäßigkeit der Handlung und Erkaufung eines gottesdienstlichen Amts.

Viertens kann die Stelle 1. Tim. 6, 5. gar wohl damit bestehen. Denn 1. zeiget der Zusammenhang, daß nicht von Lehrern, sondern von gemeinen Christen die Rede sei, die durch Aufbürdung sowohl neuer Meinungen als anderer gottesdienstlichen Uebungen Anderer Gemüther an sich zu ziehen und dadurch einen Gewinn zu erhalten gesucht. 2. Ist ein großer Unterschied unter der Treibung eines Gewerbes mit der Gottseeligkeit und der Belohnung für besondere Bemühungen beim Gottesdienst, welches letztere in eben diesem Briefe verordnet ist.

Damit aber hiebei aller Mißbrauch vermieden werde, so dienen dazu folgende Lehrsätze:

Erstlich, es wäre besser, auch für Lehrer und Zuhörer bequemer, wenn dergleichen Einkünfte gänzlich aufgehoben werden könnten. Entweder durch eine reichliche Gutthätigkeit der Glieder der Gemeinde und willigen Beitrag zur Erhaltung der Lehrer, die mit Verrichtung gottesdienstlicher Handlungen nicht verknüpfet wäre; oder auch, welches noch besser wäre, durch anderweitige Bestimmung eines reichlicheren Unterhalts; theils, weil dabei weniger Anstoß zu besorgen in Ansehung Solcher, die die Prediger als eine Last bürgerlicher Gesellschaften ansehen; theils, weil in diesem Fall manche einzelne Arten des Gottesdienstes den Leuten nicht aufgedrungen werden dürften.

Zweitens, ein Lehrer thut nicht unrecht, wenn
\quellseite{180}
er entweder aus zweifelhaftem Gewissen, oder zur Vermeidung des Anstoßes bei Zuhörern, oder bei besorglicher Unvermögenheit derselben und eigener anderweitigen hinlänglichen Versorgung, oder auch zur Beförderung des guten Vertrauens der Zuhörer gegen sich, und der Fruchtbarkeit seines Amts an denselben sich dieses seines Rechts begiebet. Doch ist es besser, wenn es nur in einzelnen Fällen und bei einzelnen Personen, als durch eine gänzliche Aufhebung oder öffentliche Bekanntmachung vor der Gemeinde geschieht. Weil

1. dadurch sowohl der nachfolgenden als auch anderer Lehrer Nachtheil verhütet wird, die es weder so halten, noch auch halten können.

2. Weil dergleichen die Zuhörer an ihrer Obliegenheit auch in andern Fällen irre macht; inmaßen sie leicht weiter gehen und den ganzen Unterhalt der Lehrer für unrechtmäßig halten können.

3. Weil ein Prediger weit eher den Zweck erreichen kann, Gutthätigkeit auszuüben, wenn er sich seines Rechts nur in einzelnen Fällen nicht bedienet.

4. Weil es weit schwerer und mit einer größern Collision der Pflichten verbunden ist, im Fall der Noth und bei merklicher Verschlimmerung der vormaligen beglückten Umstände, sich eines Rechts wieder zu bedienen, dem man feierlich und öffentlich entsagt hat, als wenn es nur in einzelnen Fällen und bei einzelnen Personen geschiehet.

5. Weil, wie mir aus fremder Erfahrung bekannt ist, auch die beste Absicht eines Lehrers bei gänzlicher Aufhebung rechtmäßiger Accidentien ungleich beurtheilt und als ein verborgener Ehrgeiz angesehen wird, folglich der gesegneten Führung seines göttlichen Amts viel Hindernisse im Wege gelegt werden können, welches bei dem Gegentheil leichter zu vermeiden ist.

Drittens ist es auch aus obigen Gründen unrecht,
\quellseite{181}
diese Obliegenheit entweder ganz oder zum Theil der Gemeinde nachzulassen, wenn es gleich zu Erbauung öffentlicher Gebäude bestimmt würde. Weil eine Gemeinde dadurch in die größte Versuchung gesetzt werden kann, ihre Predigerwahl darnach einzurichten und solches mit zu ihrem Beruf zu rechnen.

Hieraus können auch folgende Fälle entschieden werden:

1. Ob ein Prediger mit gutem Gewissen über Einziehung seiner Einkünfte bei der Obrigkeit klagen könne? welches in Dingen, so ihm durch obrigkeitliche Gesetze gebühren, erlaubt und billig, und die Obrigkeit verbunden ist, die Beibehaltung solcher Einkünfte durch richterliche Mittel zu bestimmen.

2. Ob ein Prediger bei Erhaltung des Berufs zu einer Gemeine, in welcher sein Vorgänger dergleichen Einkünfte öffentlich aufgehoben, berechtiget sei, selbige wieder zu fordern und den Berufenden darüber gehörige Vorstellung zu thun? welches zu bejahen ist.

3. Ob eine solche Gemeinde aus dem besondern Verhalten eines Lehrers ein beständiges Recht erhalte, allen ihren künftigen Lehrern keine Accidentien zu geben; welches verneint wird.“

\section{VII. Leben}

88. Wenn das Einkommen eines Pfarrers Geldbesoldung wäre, die ihm in gleichmäßigen Fristen, monatweise, quartalweise \&c. ausgezahlt würde; so würde, auch wenn das Ganze eine geringe Summe betrüge, eine gewisse Gleichmäßigkeit des Pfarrhaushalts erreicht werden können. Wenn einer weiß, daß er täglich einen halben
\quellseite{182}
Gulden auf seinen Haushalt verwenden darf und er denselben alltäglich wirklich in der Hand hat; so kann er sich nach seiner Decke strecken und, zwar immer nur so gut es gehen will, aber doch in einem gewissen Maße sorglos unter ihr sich bergen. Ganz anders und sehr erschwert hingegen ist der Haushalt der namentlich noch jüngeren, also gering besoldeten und zugleich unerfahreneren Pfarrer dadurch, daß die Besoldung nicht nur in Geld, sondern größtentheils in Naturprodukten besteht, deren man nicht eher, als zur Zeit der Aernte habhaft werden kann. Dadurch kann es kommen, daß man zu einer bestimmten Jahrszeit die Fülle, während der übrigen Zeit des Jahres aber Mangel hat. So lange der Mangel währt, muß sich dann der Pfarrer behelfen, — borgen, leihen; wenn seine Fülle kommt, hat er zu zahlen, um hernachmals wieder zu borgen und Mangel zu leiden. Allerdings gibt es nun Charaktere, welche diese Noth schnell begreifen und denen es ein leichtes ist, zur Zeit der Fülle eine solche Anordnung zu treffen, daß die Gesammteinnahme auf die sonst eintretende Hungerzeit richtig vertheilt wird. Mancher, der einmal aufmerksam gemacht ist (und das ist die Absicht dieser Sätze!), bedenkt das schon vor Antritt der Pfarrei, so wie er die von dem Naturlauf abhängige Beschaffenheit seiner neuen Einkünfte aus der Fassion erkennt. Mancher hat Wirthschaftstalent genug, nur einmal in Noth zu kommen, vorausgesetzt, daß ihn nicht besondere, geldraubende Ausnahms- und Unglücksfälle treffen. Viele aber kommen ihr Leben lang nicht in die Reihe; sie müssen immer sorgen, irgend einmal nicht solvent zu sein, ohne daß sie sich
\quellseite{183}
bis zu dem muthigen Entschluß erheben können, beim Eintritt eines solchen Falls in Gottes Namen mit Darlegung der Umstände für den Augenblick zufrieden zu sein. Diese Lage veranlaßt manche zu einer immerwährenden Unzufriedenheit mit ihrer Pfarrei, welche aus der vielleicht nicht eben unansehnlichen Gesammtsumme ihres Einkommens nach Fasson nicht zu erklären ist.

Hier gilt es nun Klugheit und Rechnen und feste Konsequenz, rücksichtlich alles Gebens. Wo nur immer möglich, darf eben so wenig in der Zeit der Fülle, als in der Zeit des Mangels das überschritten werden, was einmal die Rate der fraglichen Zeit und Frist ist. Will dann trotz aller Treue eines Haushalters das nicht zureichen, was ordnungsmäßig da ist, so kann man die Augen in ruhigem Hoffen zu Dem emporheben, der da reich ist über alle.

89. Am schmerzlichsten müßte es einem Pfarrer sein, wenn er seinem Taglöhner den bedungenen Lohn nicht geben könnte und denjenigen über den Sonntag müßte warten lassen, der vielleicht gerade vom Tagelohn seines Pfarrers die sonntäglichen Bedürfnisse der Seinigen bestreiten wollte. Dem Armen nicht geben können, was sein ist, — nicht zu rechter Zeit geben können, ist ein bitterer Tropfen im Lebenskelch eines Mannes, der andere zum Hausstande mit den Worten einsegnet: „Du wirst dich nähren deiner Hand Arbeit, wohl dir, du hast’s gut!“ — Erfahrungen wie diese lehren erst recht, was das heißt:

\begin{quote}
„Im Schweiß deines Angesichts sollst du dein Brod essen. Mit Kummer sollst du dich nähren, bis daß du zur Erden
\quellseite{184}
werdest; denn du bist Erde und zu Erden sollst du werden.“
\end{quote}

— Einem Pfarrer läuft freilich der Schweiß und die Thräne nicht wie dem Landmann übers Angesicht, aber es ist nur desto mehrthuender, wenn Schweiß und Thräne sich ins Innere ergießen.

90. Weniger kränkend aber voll Schmerzen bitterer Reue ist es, wenn ein Pfarrer sich sagen muß, daß seine Insolvenz, sein Mangel daher rühren, daß er in Ehrengaben und im Schenken nicht Maß gehalten, — von Ausgaben für Luxusartikel nicht zu reden, die einem Manne in mittleren oder geringeren Vermögensumständen nicht gebühren.

91. Ehrengaben kommen häufig an einen Pfarrer. Da stelle er sich vornherein unerschütterlich fest auf den Satz, daß sein Stand nimmermehr verlangen könne, was seine Vermögensumstände vermehren. Dagegen andere thun, was sie wollen; das irre einen treuen Haushalter nicht! Du darfst und sollst nicht, was du nicht kannst, dabei bleibt es. Und dabei bleibe es auch, so schwer es auch gehe, so selten auch einer dem völlig gemäß handelt, so selten einer hierin ohne Schuld, ohne Strafe, ohne bittre Reue bleibt. Eben so ist es mit den Geschenken. Es geht schwer, abzuschlagen, wenn man oft so flehentlich gebeten wird. Aber es ist doch keine Tugend, sondern eine Schwachheit, ja mehr als das, es ist fast Bosheit, wenn man über Vermögen gibt. Weib und Kind muß es dann entgelten, wenn ein Pfarrer sich den Ruhm der Barmherzigkeit bald hie, bald da erwirbt. Wenn Weib und Kind versorgt sind, so thue mit deinem Uebrigen, was
\quellseite{185}
du willst. Aber sieh wohl zu, daß du nicht unbarmherzig gegen die Hausgenossen werdest, indem du andern Barmherzigkeit erweisest. — Hier öffnet sich einem armen Pfarrer eine Schule heiliger Entsagung, welche mit Heiterkeit durchgemacht, den Mann bewährt und ohne Zweifel zur geheimen, aber vor Gott und Seinen Engeln anerkannten Vollendung des Armen führt. Menschen haben hiefür kein Auge; was in dieser Rücksicht im Herzen eines armen Pfarrers von edlem Gemüthe vorgeht, taugt auch nicht für die Oeffentlichkeit; aber jener Tag wird’s klar machen! — Ein Pfarrer, der nicht gibt, weil er nicht geben kann, erscheint seinem Volke immer oder doch meist nur geizig zu sein, so lange seine Armuth nicht bis zum Mangel des Nothwendigsten ausschlägt. Das weiß ein Pfarrer. Wenn er nun ungeirrt von Vorurtheil und Argwohn seines Volkes mit aller Macht andere zur Barmherzigkeit ermuntern soll, so gehört dazu eine hohe Seele, vorausgesetzt namentlich, daß die Seele nichts lieber thäte, als mit mächtigem Beispiele der Barmherzigkeit voranzuwandeln. — Dergleichen Feuerproben der Heiligung gibt es für einen armen Pfarrer im Amte gar viele.

92. Hier ist auch wohl ein Wort vom Leihen und Borgen zu reden. Nichts gewöhnlicher, als daß man vom Pfarrer Geld leihen und Waaren (Getreide und andere Früchte) auf Borg nehmen will. Den Pfarrer ermuntern Stellen, wie: „Leihet, daß ihr nichts dafür hoffet!“ Auf der andern Seite kann er voraussehen, daß er mit den Seinigen bei solchem Leihen Mangel leiden werde. Er kennt sein Einkommen genau vorher; er weiß, daß ihm
\quellseite{186}
nicht mehr zu Theil werden wird. Er weiß auch, daß die, welche vom Pfarrer leihen oder borgen wollen, in der Regel die sind, welchen niemand sonst leiht oder borgt, — die Armen, die keine Sicherheit der Zahlung oder Rückerstattung geben können. Diese Armen sollten nicht leihen, nicht borgen, sondern ehrlich um Geschenke bitten; aber sie hoffen, wieder geben zu können, wo nichts zu hoffen ist, — sie leihen mit dem Willen, aber ohne alle Wahrscheinlichkeit der Wiedererstattung. Was ist da zu thun? Ein wohlhabender Pfarrer leihe, und das nach des Herrn Gebot und ohne Zinsen, wenn er arme Leute vor sich hat. Ein armer Pfarrer handle mit offenem Bekenntnis seiner Armuth nach dem Grundsatz, daß er nicht mehr ausgeben dürfe, als zu seinem Unterhalt durchaus nöthig sei, daß er seine Einnahme nach und nach, oft kümmerlich bekomme, daß er nicht ausleihen könne, was er nicht habe, daß er aber gerne schenken wolle, wenn er etwas übrig habe. Denn es kann kommen, daß er zwar nicht die größere Summe, die er leihen soll, in Gefahr setzen, wohl aber eine kleinere schenken darf. Auf diese Weise konsequent verfahren, hat einen großen Segen. Man wird, wenn auch nicht alsbald, doch je länger je mehr erkannt. Je länger, je mehr wird der Pfarrer mit den Anforderungen seiner Armen verschont; je länger, je mehr wird seine Handlungsweise in ihrer Harmonie mit der Predigt, je länger, je mehr wird er als wohlthätig nach dem Maße seiner Umstände erkannt; je länger, je mehr kommt er in den Stand, ungebeten, nach eigenem Ermessen, still und ohne Geräusch das wenige, was er etwa erübrigt, zu
\quellseite{187}
Wohlthaten anzuwenden. Ein Pfarrer kann ja auch diese Grundsätze, nach denen er handelt, auf und unter der Kanzel offen besprechen und in ihrer schriftmäßigen Richtigkeit und Weisheit darlegen. Damit hilft er dem Verständnis der Gemeinde nach und bringt sich selbst schneller aus dem Mißverständnis. Es geht, wie auch sonst in schwierigen Fällen und Lagen: Gott läßt es den Aufrichtigen gelingen, und die ehrliche Darlegung der Umstände bringt gerechtere und billigere Beurtheilung derselben.

\subsection[93. Unsre Zeit ist eine Zeit der Vereine]{93. Unsre Zeit ist eine Zeit der Vereine\footnote{§. 93 soll keine Demonstration gegen die Vereine sein. Als vor einigen Jahrzehnten das protestantische Vereinswesen einen größeren Aufschwung nahm, war der Schreiber dieses Buches allerdings nicht gleich im Klaren, ob es sich bei der Gefahr der Vereinserei, in welche die Vereine so leicht gerathen, für einen lutherischen Pfarrer und Christen schicke, auf dem Weg der Vereine zu wirken. Allein bald zeigte sich’s ihm doch, daß es unweise sein würde, vereinzelt, bloß unter Benützung persönlicher Bekanntschaften und ohne alle Form kirchliche Werke fördern zu wollen, während sich die Christen aller Confessionen zu Vereinen zusammen schließen. Die Form ist nöthig für alles Wirken, zu welchem die eigenen persönlichen Kräfte nicht ausreichen; kein rechter Mann liebt es ohne Controle, das ist ohne Form zu wirken; formloses Handeln auf alleinige eigene Gefahr hin kommt so schwer zum Bewußtsein eines guten Gewissens. Aber auch abgesehen von der Förderung kirchlicher Werke haben die Vereine noch einen seelsorgerischen Nutzen, um dessen willen sie verdienen, nicht bloß wie Auswüchse des kirchlichen Lebens geduldet, sondern mit aller Treue gepflegt und mit großem Fleiße gefördert zu werden. Allerdings ist die Kirche der von Gott gewollte Verein, in welchem alle andern aufgehen sollen; allein auch wenn die Kirche wäre, was sie sein soll, würde sie doch zu ihrem Zweck nur dadurch gelangen, daß sie ihre Kinder je nach den verschiedenen Gaben und Kräften, welche sie haben, für gewisse einzelne Liebeszwecke zusammen treten und zusammen arbeiten ließe. Da sie nun aber in dieser Welt nicht ist, was sie soll, und namentlich in unserer Zeit die Bösen an Zahl- und Einfluß überwiegend sind, die Kirche allenthalben in der größten Gefahr schwebt, daß der Süßteig von dem Sauerteig ganz und gar durchdrungen werde; so ist jedes Mittel zu ergreifen, durch welches die Besseren zu Hause und in Verbindung gebracht und gestärkt werden können. Da leuchtet es denn ein, daß die christliche Vereinsthätigkeit zu diesem Zwecke benützt werden kann. Die kirchlichen Versammlungen und die Erbauungsstunden pflegen die Gemeinschaft des Glaubens, die christlichen Vereine aber lehren die Kinder der Kirche mit vereinter Kraft Liebe üben und füllen damit eine Lücke aus, welche ohne sie das christliche Leben zu haben pflegt. Das gewöhnliche Leben in der Gemeinde und Familie öffnet allerdings für alle Tugenden und deren Uebungen weites Feld, auch für die Werke der Barmherzigkeit zeigt es Gelegenheit genug; die Vereinsthätigkeit aber gestattet doch mehr Blick in das menschliche Elend, weckt den Eifer, für dasselbe zu wirken, bringt die Gleichgesinnten zusammen, vereinigt nicht bloß ihre Kräfte, sondern auch ihre Herzen und hilft dazu, daß die Christen das Wörtchen „Wir“ im Sinne der kirchlichen Liebe und Einigkeit sagen lernen. Weit entfernt, daß also den Schreiber dieses die längere Erfahrung von der Vereinsthätigkeit abgebracht hätte, ist er im Gegentheil zu der gründlichen Ueberzeugung gelangt, daß ein Seelsorger schuldig sei, seine Pfarrkinder zum Anschluß an Vereine aufzufordern und anzuleiten. Es versteht sich aber von selbst, daß man sich nicht allen möglichen Vereinen anschließen kann, weil man dadurch den seelsorgerischen Zwecken der Vereine entgegen handelte. Dagegen aber sollte ein jeder Christ einem oder etlichen Vereinen mit allem Ernste angeschlossen sein.}}

Ein
\quellseite{188}
Verein nach dem andern entsteht, und es entstehe einer, von welcher Art es sei, so wird auf die Pfarrer, als auf die Vorgänger in guten Werken, gerechnet. Es kann leicht kommen, daß man zur Theilnahme an etwa 10 bis 12 Vereinen aufgefordert wird. Man rechne nur zusammen, es wird für manche Pfarrer kaum übertrieben erscheinen.
\quellseite{189}
Gäbe nun ein Pfarrer zu jedem solchen Verein auch nur jährlich 2 fl., so würde eine Summe herauskommen, welche nach dem Maße seiner Einnahme sehr bedeutend wäre. Es ist keine Kleinigkeit, wenn ein Pfarrer jährlich etwa 20 bis 25 fl., also den zwanzigsten, dreißigsten Theil seines Einkommens an Vereine abgibt. Etwas Sonderliches für jene Vereine thut er damit nicht, und doch raubt er sich vollends das, was er nach Amtspflicht seinen armen Pfarrkindern schuldig ist. Er wage es darum, nirgends einen regelmäßigen Beitrag zu versprechen, der ihm die Kraft, für seinen Wirkungskreis zu sorgen, schmälern könnte. Er fasse vor allem seine armen, kranken Pfarrkinder ins Auge, dann erst thue er mit dem, was übrig bleibt, nach freiem Willen und gutem Gewissen. Die Unterstützung bedürftiger Glaubensgenossen und die Mission unter den Heiden wird ein frommer Pfarrer immer im Auge behalten, und es wird ihm die größte Freude sein, wenn er kann, etwas und vieles zu thun, was seine Freude an jenen heiligen Liebeswerken bezeugen und seine öffentliche Ermunterung zu denselben bekräftigen kann. Aber auch hierin wird er vor allen Dingen nach dem Maße seines Vermögens thun und sich mögliche Miskennung nicht irren laßen. Bei diesem Verfahren wird dem Pfarrer viel böses Gewissen erspart; seine Seele möglichst sorgenfrei erhalten; die Wohlthat, die er bei solchen Grundsätzen thut, ist wirklich Wohlthat, nicht aber wie sonst oft Qual.

Mit alle dem sehen wir natürlich nicht auf reiche und dabei geizige Pfarrer. Ihnen soll keine Entschuldigung
\quellseite{190}
für ihre Sünde geliefert werden. Es ist allein davon die Rede, wie sich die große Anzahl armer Pfarrer gleich beim Eintritt ihres Amtes nothgedrungen zu stellen habe, um nicht aus mißverstandener Barmherzigkeit in unerträgliche, das Amt hemmende Armuth, in immer neue Verlegenheiten, wohl gar auf und über die Grenzen der Ehrlichkeit zu kommen.

94. Hier möchte es auch wohl an der Stelle sein, kürzlich auf zwei Dinge hinzuweisen, welche nicht selten im Pfarrerleben vorkommen.

Wir meinen fürs erste den Kauf und Verkauf, welcher bei allen Leuten, so auch bei Pfarrern vorkommt. Es ist schon ärgerlich und dem Amte nachtheilig, wenn ein Pfarrer oder dessen Angehörige bei gewöhnlichen Einkäufen jüdisch feilschen und handeln; so sehr man auch durch die Unredlichkeit wuchernder Kaufleute dazu herausgefordert wird, bleibt es eben doch immerhin dem Amte zuträglicher, etwas zu nachgiebig, als zu vorsichtig zu sein. Noch ärgerlicher und hinderlicher ist es aber, wenn ein Pfarrer beim Verkauf seiner Naturalien melkförmig verfährt. Es ist, wenn auch nicht immer, doch meistens ein zweideutiger Ruhm für einen Pfarrer, wenn er seine Feldfrüchte immer für höhere Preise absetzt, als seine Pfarrkinder. Es kann zwar hie und da einmal ganz unschuldige Ursachen haben; aber dem sei, wie es will, ein Pfarrer soll es mit seiner Waare nicht aufs Höchste treiben. Auch die Noth soll ihn nicht vermögen, den Wucherern gleich zu erscheinen. Er kann nicht alles, was andere Leute können; er hat es zwar alles Macht, aber es frommt
\quellseite{191}
nicht alles; und zum Frommen anderer muß ein Pfarrer durchaus leben, auch mit eigenem Schaden. Es ist daher ein altes Sprichwort: „Ein Pfarrer muß Haare lassen!“ — Mögen daher neuangehende Pfarrer (deren manchen vielleicht öfter die umgekehrte Warnung zu geben sein mag) nicht in der Noth, welche ein beginnender Haushalt zuweilen mit sich führt, sich eines zu eifrigen Haschens nach dem Einkommen schuldig machen. Nicht Vernachlässigung oder Verschleuderung des Eigenthums, nicht allzugroße Genauigkeit möge einem Pfarrer zur Last gelegt werden können. Ein heiliger, schlichter, redlicher Sinn kann durch den Rath der Erfahrenen leicht das finden, was zum heiligen Amte paßt. Die Erfahrung beweist, daß gerade die offen und gerade hindurchgehende Schlichtheit oft gesegnet ist. Es geht auch hier nach der Regel: „Die Gottseeligkeit ist zu allen Dingen nütze und hat die Verheißung dieses und des zukünftigen Lebens.“

95. Das Zweite, das wir noch zu erwähnen übrig haben, ist Folgendes. Unsere Zeit folgte auf eine glaubensleere Zeit, in welcher die meisten Pfarrer überhaupt und insonderheit durch Verbreitung vergifteter Erbauungsbücher das Gute vollends auszurotten suchten, welches etwa aus früheren Zeiten noch übrig war. So findet denn der neueintretende Pfarrer in seiner Gemeinde böse Bücher auszumerzen und natürlich auch die Nothwendigkeit, bessere an die Stelle zu setzen. Ist er nun überhaupt gleich anfangs in so weit gesegnet, daß er Aufmerksamkeit und Wohlgefallen seiner Gemeinde an sich zieht; so wird seine Empfehlung guter Bücher gerne angenommen und er bekommt
\quellseite{192}
leicht Anlaß, sie für seine Gemeinde selbst zu besorgen. Jeder nimmt sie gerne aus seiner Hand. Man will einen Vers oder Spruch, den Namen, den Ort, das Jahr von seiner Hand geschrieben im Buche haben \&c. Geht er nun hiebei recht zu Werk, so thut er nichts, als was sich für ihn ziemt. Die Besorgung der Bücher kann er wohl übernehmen; denn er wird dadurch in den Stand gesetzt, einen wohlfeileren Preis oder für Arme Gratisexemplare zu bekommen. Ein Buchbinder wird auch nicht dadurch beeinträchtigt, denn gebunden werden die Bücher doch, und sie wohlfeiler zu verkaufen, als es in den Buchläden von Buchhändlern und Buchbindern geschieht, ist unsträflich, wenn man auf eine unsträfliche Weise den geringeren Preis erlangt. Lassen sich Buchbinder selbst ins Interesse einer wohlfeileren Verbreitung guter Schriften ziehen, desto besser! — Aber das ist schlimm, wenn sich ein Pfarrer auf Borg einläßt oder sonst bei solchem Büchervertrieb nicht Ordnung hält. Beides bringt ihn zu Schaden, und zwar zu einem oft höchst empfindlichen Schaden, den er schwer trägt, ohne daß jemand eine Wohlthat geschieht. Es gibt viele Leichtsinnige, die ohne Absicht des Betrugs ein gutes Buch wohl borgen, das sie baar zu kaufen keine Lust bezeigen würden, und manchen Boshaften, welcher es mit Absicht des Betrugs thut. Der Boshafte schadet mit lachendem Munde; der leichtsinnige Borger aber wird eben dadurch, daß man ihm borgte, ohne ihn auf seine Zahlungsfähigkeit oder Unfähigkeit aufmerksam zu machen, vom Pfarrer entfernt. Aus Furcht, gemahnt zu
\quellseite{193}
werden, während doch inwendig ein strafendes Bewußtsein lebt, vermeidet er den Pfarrer, weicht ihm aus, versäumt wohl gar die Gelegenheiten, wo er ihm begegnen könnte, den Kirchweg 2c., — und damit wird Ohr und Herz vom Wort entfremdet, nicht bloß entfernt. Ja, es gibt Beispiele, daß dergleichen Bürger (das gilt natürlich auch von andern Borg, als dem der Bücher) einen Haß auf den Pfarrer werfen und auf sein Wort, bloß weil sie ein Paar Gulden oder Kreuzer nicht zahlen können, die sie schuldig sind. — Man frage nur Erfahrene, ob es nicht so ist! — Wir führen dies hier nur zum Beweis an, daß mit Borg in solchen Fällen niemand genützt, wohl aber der eigenen Kasse und fremden Seelen oft empfindlicher Schaden zugefügt wird. — Von jenem Schaden gar nicht zu reden, welcher durch Verschleuderung und maßloses Verschenken der Bücher erfolgt. Die Bücher werden verachtet, wenn sie zu gerne gegeben werden, und ein Pfarrer kann sich und seinem Haushalt ganz unbemerkt und ohne alle Noth bedeutenden Verlust bringen.

96. Am besten ist es, der Pfarrer macht sich seine Aufgabe klar. Die bösen Bücher müssen weg, das ist gewiß. Eben so gewiß ist es, daß bessere an die Stelle müssen. Aber 1. nicht viele statt weniger, 2. keine geringen oder mittelmäßigen. Mittelmäßige Bücher helfen dazu, des Volkes Sinn und Geschmack zu verderben, heben nicht, helfen nicht. Es dürfen gewaltige Bücher sein, welche das Volk aus dem Traume der zeitlichen Lust, Noth und Trägheit erheben sollen! — Viele Bücher werden eines dem andern hinderlich. Vielleserei ist dem
\quellseite{194}
Volke schädlich, und es ist gut, daß es sich damit in der Regel nicht einläßt, nicht einlassen kann. Die Volksbibliothek ist kurz beisammen. Die Bibel, die Concordia, eine Postille, ein Gesangbuch mit Psalter zum Singen, ein Beichtbuch, ein Betbuch für gesunde und kranke Tage, der Katechismus mit Spruchbuch, der Kalender mit Festbuch, — das reicht für die tägliche Speisung der Seelen hin. Um heilsame Bewegung zu erhalten und die Gemeinde im Zusammenhang des Ganzen zu erhalten, wird es gut sein, wenn ein Haushalter alljährlich ein christliches Buch ankauft, die wichtigsten Zeittractate anschafft, eine oder etliche christliche Zeitschriften mitliest. Dahin suche ein Pfarrer seine Gemeinde zu leiten, das sei sein Ziel. So weit kommt er auch vielleicht, so weit folgt ihm der bessere Theil des Volkes gern. Gut ists, wenn er allen zwar nicht ganz einerlei und genau dasselbe, aber nur nach kleiner Abwechselung und Wahl räth. Einerlei verursacht, daß die Erkenntniß der Gemeinde zu einförmig wird; es muß einiger Wechsel der Bücher, einige Vergleichung verschiedener Stimmen, einige Bewegung der Gedanken möglich sein. Doch darf der Kreis der empfohlenen Bücher nicht zu groß sein. Man muß sich in der Predigt, Katechese und Seelsorge auf die Volksbibliothek beziehen, sie lebendig und wirksam machen können, — was sich von selbst verbietet, wenn der Bücher zu viel sind, die es in der Gemeinde gibt, und daher immer der größte Theil die Beziehung, in welche der Pfarrer seine mündliche Belehrung und Ermahnung setzen will, nicht verstehen. — Er wähle also weniges, aber Treffliches, empfehle seinem
\quellseite{195}
Volke das Gewählte, mache es bekannt, — verrathe jedoch keine zu große Angelegentlichkeit bei der Anpreisung, denn da vermuthet das Volk Eigennutz; — sei nicht zu hastig und eifrig, verfahre nicht, wie wenn er einem Buchladen vorstünde, nicht als läge an dem Ankauf von Büchern das Heil, — er leihe neue, zum Verkauf angeschaffte Bücher nicht und borge nicht, — er schenke nicht leicht, — und mache überhaupt dergleichen Geschäfte einfach und kurz ab. Dabei führe er in Einnahme und Ausgabe die strengste Ordnung und bestelle nicht leicht wieder, ehe er zahlen kann. — — Das alles ist leicht einzusehen; jeder weiß es. Aber wie viel Ungemach erspart sich der, der darnach thut, — und wie viel mehr Ansehen hat eine geordnete Bücherverbreitung bei dem Volk, wie viel mehr Eindruck macht es, als die liederliche Bücherwirthschaft eines unbesonnenen Menschen.

\section{VIII. Gastfreundschaft}

97. „Seid gastfrei ohne Murmeln,“\footnote{Wiewohl die φιλοξενία, zu welcher die heiligen Apostel vermahnen (s. außer den oben angeführten Stellen noch 1. Tim. 3, 2. 5. 10. Tit. 1, 8.), etwas anderes ist, als was wir häufig unter Gastfreundschaft verstehen. Der christliche Pilger, der seine Glaubensgenossen aufsucht, um mit ihnen sich in der Liebe und heiligsten Gemeinschaft zu stärken, — oder der seinem Berufe in der Fremde obliegt und dabei die heimatliche Liebe seiner Glaubensbrüder sucht, ist meistens bei uns nicht der Gast, der oft kommt und das Gebot der Gastfreundschaft zu einem wichtigen und schwierigeren macht. Unsre Gäste und die Absicht ihrer Besuche sind oft von der Art, daß sich’s fragt, ob sie in den Bereich der φιλοξενία eingerechnet werden können.} sagt der Apostel (1. Petr. 4, 9. Ebr. 13, 2.) zu allen Menschen. Ohne
\quellseite{196}
Zweifel haben darum Pfarrer, das apostolische Gebot hauptsächlich sich selber zuzueignen. Indeß hat jedes Gebot seine besondere Anwendung auf verschiedene Lagen. Gastfreundschaft gehört im allgemeinen zum Geben. Das Geben aber hängt jedenfalls vom Empfangen ab. Wie man also im Allgemeinen nicht geben kann, was man nicht hat; so kann auch keiner die Gastfreundschaft weiter ausdehnen, als es seine Verhältnisse und seine Vermögensumstände erlauben. Der Unerfahrene weiß es nicht, was für große Ausgaben eine ausgedehnte Gastfreundschaft erfordert; aber erfahrene Hausväter kennen das, und eben deshalb ist eine mäßige Uebung der Gastfreundschaft von Seiten solcher werth- und ehrenvoller, als von Seiten des unbedachten, sorglosen Anfängers eine ausgedehnte. Wohl dem Anfänger, welcher sich gleich vornherein so stellt, daß er nicht hernachmals aufhören oder einziehen muß! Nur zu leicht verführt das anfängliche Behagen eigenen Haushaltens zur Unmäßigkeit in Ausübung der Gastfreundschaft. Und Mancher merkt es erst, wenn er sich geschadet hat. Wer gleich anfangs Maß hält, wird gewiß nicht leicht mißverstanden. Dagegen ist dem Urtheil der Menschen sehr ausgesetzt, wer, nachdem er ein Haus gemacht hat, anfangen muß, weniger zu thun. Das Maß kann übrigens nicht bloß dadurch überschritten werden, daß man zu viele
\quellseite{197}
aufnimmt, sondern auch dadurch, daß man in der Beurtheilung von der Einfalt weicht. Ist der Pfarrer überhaupt, auch wenn er in guten Umständen lebt, in äußerlichen Dingen nur den mittleren Ständen zuzuzählen; wie viel mehr wird er das festzuhalten haben, wenn er in Anfangszuständen lebt, wenn ihm Einkommen und Besoldung spärlich zugemessen ist. Ist auch die Sitte des Engländers, welcher an seine Gäste den Anspruch stellt, sich in seine Hausordnung genau zu fügen, nicht in englischer Schärfe festzuhalten, so könnte sie doch für einen Pfarrer die Regel geben. Ausnahmen werden leicht erkannt. — Einfalt in allen Dingen, Einfalt in der Gastfreundschaft ziemt dem Pfarrer. Er kann nichts anderes anstreben — Vermögens wegen! Er soll nichts anderes anstreben — von Standes wegen!

98. Bei der Gastfreundschaft ist übrigens die Rücksicht auf das Vermögen des Pfarrers zwar eine vorzügliche, aber nicht die alleinige. Zwei Rücksichten sind noch insonderheit hervorzuheben; die aufs Amt und die auf die Kinder.

Um in seinem Amte recht zu wirken, ist dem Pfarrer eine Abwechselung zwischen Einsamkeit und Oeffentlichkeit ersprießlich und nöthig. Diese Einsamkeit verschließt ihn in sein Haus. Sie erfordert nicht allein Stille der Anwesenheit, sondern nach dem Bedürfniß der meisten Pfarrer Abwesenheit aller, Alleinsein, wirkliche Einsamkeit. Das Bewußtsein, in ungewöhnlicher Umgebung zu sein, ist oft sehr störend, wenn auch alles ringsumher stille ist. Ist doch vielen schon das Bewußtsein, nicht völlig einsam zu sein,
\quellseite{198}
selbst eine geliebte Person in der Nähe zu haben, störend und hinderlich. Ist aber das, wie soll dann eine Gastfreundschaft, bei welcher jeder Besuch zum Mittelpunkt aller Hausbewohner wird, um den sich alles, also auch der Pfarrer, drehen und nach ihm sich richten muß, mit dieser Stille und Einsamkeit stimmen? In einem Gasthause ist der eintretende Gast Herr — der Hausherr tritt dienend zurück. So darf es bei der Gastfreundschaft eines Pfarrers durchaus nicht sein. Das heilige Amt, welches er verwaltet, erheischt die Rücksicht aller Besuchenden, — des Pfarrers Werk und darum auch seine Seelenruhe muß geschont werden. Jede innere Störung, jede Mißstimmung, jede Ermüdung bekommt die Gemeinde zu empfinden. Wer drum ein Pfarrhaus betritt, muß sich des Anspruchs begeben, den Pfarrer für sich zu haben und ohne Unterlaß von ihm beachtet, berücksichtigt, begleitet zu sein. Da nun aber viel Besuche so anspruchsvoll zu sein pflegen, so hat ein Pfarrer desto mehr zu wachen, daß nicht die Uebung der Gastfreundschaft ihm und seinem Amte zum Schaden gereiche. Je mehr er durch Besuche aufgeregt, aus der Stille der Seele gerissen, unruhig, unklar, abgespannt, ermüdet wird, desto weniger kann er sein Haus den Gästen öffnen. Seine Gemeinde und das Amt in ihr gehe vor allem. Alles muß sich seinem Amte unterordnen.

99. Die andere Rücksicht, welche ein Pfarrer bei Ausübung der Gastfreundschaft zu nehmen hat, sind seine Kinder. Wenn er auch Stärke der Seele und Klarheit des Geistes genug hat, sich durch Besuche in seinem heiligen Amte durchaus nicht stören zu lassen; so fragt es sich,
\quellseite{199}
ob er denn auch seine Pflicht gegen seine Kinder — eigene oder fremde, die er etwa im Hause hat, — ganz und völlig im Auge behalten kann? Eben weil einem Pfarrhause die Studierstube, einem Pfarrer selbst die Abgeschlossenheit und Einsamkeit so nöthig ist; eben weil sein Amt ihn so oft hinausführt in Kirche, Schule, Häuser; eben deswegen entbehren Pfarrerskinder, namentlich jüngere, den Vater ohnehin eher zu viel, als zu wenig. Wenn nun die Rücksicht auf Besuche des Vaters Herz und Auge, seine Kraft und Zeit den Kindern auch noch schmälert und entzieht, — wie bedauernswerth sind dann die armen Kinder. Sie fühlen sich unbeaufsichtigt, nicht beachtet, — so benehmen sie sich, wie es ihnen kommt. Ihre Unart wird vom Vater entweder bemerkt oder nicht bemerkt. Ist das erstere, so müssen die Kinder in Gegenwart des Besuchs gestraft werden, oder sie werden in der Bosheit gestärkt. Strafen und nicht strafen in Gegenwart des Besuchs hat beides sein Misliches. Wird die Unart gar nicht bemerkt, so ist es noch schlimmer; der Besuch bemerkt sie desto eher, bespricht aber in der Regel den Fehler vor den Ohren des Vaters doch nicht. Sind die Kinder über das Alter hinaus, in welchem man sie vor Fremden strafen kann, so tritt an die Stelle jugendlichen Uebermuths oft wirkliche unbändige Flegelei, und der Vater faßt dann leicht gegen die Kinder eine bitterere Stimmung, als gut ist, zumal er selbst nicht ohne Schuld am Uebel ist. Ist aber das alles nicht der Fall, so lernen doch Kinder, welche viele Fremde im Hause sehen, eine allzugroße Freiheit des Benehmens, und jene Schüchternheit, welche dem Kinde gegenüber dem Aelteren
\quellseite{200}
zusteht, wird allzufrüh verwischt. — Kurz, wenn der Vater den Kindern durch Besuche zu sehr entzogen wird, liegt es sehr nahe, daß die Kinder großen Schaden nehmen. Selbst wenn es der Aufsicht und des Raumes halber möglich ist, die Kinder von den Besuchen fern zu halten, also selbst im günstigsten Falle, bleibt der Nachtheil nicht aus, und es bleibt wahr, daß den jungen Kindern Besuche eben so hinderlich in ihrem Gedeihen werden können, als sie erwachsenen Söhnen und Töchtern förderlich werden können.

100. Was soll nun ein Pfarrer thun? Selbst wenn er dürfte und wollte, würde er sich der Uebung der Gastfreundschaft nicht völlig entziehen können. Ja, je mehr es ihm Ernst ist, allewege gutes Beispiel zu geben, desto schwerer würde er sich dem allgemeinen Gebot des Apostels entziehen können. Es bleibt daher zunächst nichts übrig, als das seinem Amte, seiner Familie ersprießliche, seinem Vermögen zupassende Maß — und die Weise der Einfalt zu finden und festzuhalten. Jeder Gast sei des Hausvaters Bruder oder Sohn und werde so behandelt. Die erwachsenen Kinder mögen also dem Gast wie einem Vater oder Bruder begegnen. Die kleineren aber, für die man die nöthige Beaufsichtigung findet, mögen ihren Aufenthalt in der Kinderstube finden, so lange es immer angeht. Müssen sie ins Wohnzimmer kommen, so erlaube man nicht, daß sie sich irgend bemerklich machen oder vom Gaste aus ihrem anspruchslosen Dasein hervorgezogen werden. Man mache es Gästen, die das nicht von selbst wissen, deutlich merkbar, was des Hauses Ton in diesem Falle sei. Den Kindern aber erlaube man nicht, sich allzuvertraulich
\quellseite{201}
an die Fremden zu hängen und mit ihnen umzugehen. Was von Seiten des Vaters, der Hausmutter als freundliche Achtung gegen den Gast erscheint, das präge sich beim Kind als Ehrerbietung und Hochachtung aus. Ein Vater belehre auch besonders seine Kinder zur Zeit, wo keine Gäste da sind, über die Achtung, die man dem Fremdling nach Gottes Wort schuldig ist.

Daß natürlich Verwandte, die als Gäste einsprechen, auch zu den Kindern des Hauses sich nicht wie pure Fremdlinge verhalten, versteht sich.

Das Haus eines christlichen Pfarrers sei demnach, so weit es ihm seine Verhältnisse erlauben, ein ξενοδοχεῖον des Herrn. Wohl und schön, wenn nicht bloß der selbst geistliche, gelehrte, gleichgestellte Fremdling, sondern auch der Arme, der empfohlen ist und anspruchslos mit dem zufrieden ist, was ihm werden kann, Aufnahme findet. Selbst nicht empfohlenen Fremdlingen darf ein Pfarrer mit Vorsicht und Weisheit ein Auge und Herz voll Barmherzigkeit und Liebe zuwenden. Und je weniger ein Fremdling Anspruch macht und machen kann, desto eher kann an ihm die heilige Pflicht erfüllt werden, desto weniger Einspruch machen Vermögen, Amt und Kinder.

\quellseite{202}
\section{IX. Landsläufer und Heckenleute}

101. Gegenwärtig gibt es allenthalben wieder erweckte Leute, und wo es diese gibt, gibt es auch deren Afterbilder. Das sind jene verzweifelten Menschen, welche sich ein Geschäft daraus machen, das Land zu durchstreichen, alle „christlichen Freunde“ und namentlich alle als christlich renommierten Pfarrer aufzusuchen. Sie verwechseln die Gemeinschaft der Heiligen mit Kameradschaft und setzen ein Gewaltiges darein, alle entschiedenen Christenleute zu kennen, d.\,i. einmal oder öfter gesehen zu haben. Ihre Themata, welche sie mit allen Christen abzuhandeln pflegen, sind die sogenannten Adiaphora und die Zukunft, der Welt Ende, die Offenbarung Johannis, allenfalls auch die Mission und insbesondere „Basel oder Leipzig“ etc. Sie lieben gefühliges Erregen neben scharfem Absprechen, Erweckungspredigten in methodistischer Gewalt, Thränen, Seufzer und Grimassen, — kurz alle Lüfte religiöser Thoren findet man aus ihnen sprechend und wirkend. Allenfalls erwählen sie sich, um desto leichter das Gewissen zu stillen, das Geschäft des Hausierers und bekommen so völlige Berechtigung, ihre geistlichen Gaben mit hausieren zu tragen. Sie haben zuweilen in ihrer Heimath einen ziemlichen Anhang, herrschen oft über die Erweckten daselbst, und ihre Stimme gilt wohl auch in weiteren Kreisen bei ihres Gleichen. Sie sprechen über Pfarrer und deren
\quellseite{203}
Wirken, über ihre Predigten, Katechesen und Gebete meisterlich ab, sprechen selig trotz dem Bischof von Rom und verdammen nicht minder fertig. Sie bringen die guten und bösen Gerüchte in Umlauf, durch welche St. Paulus und dessen Nachfolger allezeit gehen mußten. Mit der Zuversicht erwünschter Freunde treten sie zu jeder Tageszeit in die Pfarrhäuser ein, und wer sie niedersitzen heißt, ja auch nur stehend sie zu Worte kommen läßt, der sehe wohl zu, wie er ihrer wieder los wird. Sie hören sich gar zu gerne selber reden, zumal wenn sie außer den ihrigen noch andere willige Ohren finden. Im Tone und mit den Ausdrücken ihres Lieblingspfarrers, tragen sie ihre Urtheile, Lobsprüche, Warnungen, ihre apokalyptischen Schlüsse und apodiktischen Befehle in Sachen christlicher Freiheit vor. Finden sie Anklang, so werden sie wunders zuthunlich, voll Freundschaft und persönlicher Erbietung zu Hilfsleistungen aller Art. Werden sie gütig behandelt, so vertragen sie auch ein wenig Widerspruch und Zurechtweisung; die Freundschaft und Bekanntschaft wieder eines berühmten Pfarrers mehr ist nicht zu theuer erkauft, wenn man auch von dem geistlichen Herrn ein wenig zu leiden bekommt. Sie kommen immer wieder, und wenn sie können, so suchen sie sich irgendwie angenehm, wohl gar unentbehrlich zu machen. Zuweilen haben sie ganz aparte Sachen, z.\,B. die Wiederbringung aller Dinge, die erste Auferstehung \&c.; schont man sie darin, um sie nicht zu erzürnen, oder in Hoffnung, ihnen nach und nach auf eine sanfte Weise den Irrthum zu nehmen, so hangen sie einem an, auch wenn sie in Predigten \&c. die gegentheilige Lehre
\quellseite{204}
hören. Der gute Pfarrer wird dann mildiglich beurtheilt, er weiß es nur noch nicht besser oder „kommt schon herbei“ etc.

102. Gegen dieses leicht kenntliche Ungeziefer\footnote{Es sei übrigens ferne, hiermit denjenigen Pfarrern sanft zu thun, die an keinem Besuch und keinem Gespräche so großen Eifer haben, wie an dem erweckter Leute aus den geringeren Ständen. Die Art, wie solche Erweckte zu reden pflegen, ermangelt oft derjenigen Bildung, welche Mancher mit wahrer Erweckung und Religiosität völlig identifiziert. Sie sind unbeholfen, haben Sitten und Manieren, welche noch nicht verdienen, ein Ausspruch inneren göttlichen Lebens genannt zu werden, dennoch aber auch keine Ursache abgeben können, an der Redlichkeit und Aufrichtigkeit der Leute zu zweifeln. In so einer harten, ungefügigen, eckigen Schaale ist manchmal ein Kleinod, das ein Pfarrer zu finden und von seinen Hindernissen zu lösen Sinn, Liebe und Gabe haben sollte. Wehe, wenn ein Pfarrer keinen Unterschied machen könnte, sondern auch den einfältigen Besucher, das edle Gotteskind, wie Landläufer, Sectierer und Separatisten behandelte!} des Weinbergs Gottes, so wie gegen andere dergleichen Schwätzer innerhalb und außerhalb der Gemeinde ist nichts besser, als die schon bekannten Besucher dieser Art einfach abweisen, für andere aber ein für allemal eine Stunde festzusetzen, in welcher man jedermann zu sprechen sich bereit erklärt, während man sich zu Uebung der Seelsorge bei Tag und Nacht, in Ruhe und Arbeit Zeit zu nehmen verspricht. Die Kommenden lasse man getrost nach Wunsch und Begehr fragen und diejenigen, welche nicht in Sachen des Amtes oder der Seelsorge kommen, mit der offengestandenen Ursache des Zeitmangels abweisen. Schon das verdrießt oft die Läufer. Kommen sie aber in der Sprechstunde, so wissen sie zum voraus, daß der Pfarrer nach
\quellseite{205}
deren Ablauf abbrechen wird, sie können also nicht allzuweit ausholen; zuweilen treffen sie mit andern Leuten zusammen, die innerhalb derselben Stunde auch abgefertigt sein wollen und sollen. Dazu verweist nicht bloß die Anordnung einer solchen Stunde an und für sich auf kurze Fassung des zu Sagenden, sondern der Pfarrer hat an der Kürze seiner Sprechstunde guten Grund, auch ausdrücklich und ohne Beleidigung auf Abkürzung des Vortrags hinzuweisen, so wie ohne weiteres nach der Ursache des Besuches zu fragen. — Weist man die Leute von der obenbeschriebenen Gattung geradezu ab, so werden sie damit nicht gebessert, sondern ihr beleidigter Stolz treibt sie zu desto größerer Hartnäckigkeit und zur Verläumdung. Nimmt man sie in der angegebenen Weise an, so kann man ihnen mit der Wahrheit dienen und die ihnen heilsame Liebe üben, ohne daß man viel Zeit verliert, und ohne daß man durch sie in übles Gerücht kommt. Der Grundton aufrichtiger Liebe und der Ton ehrlicher Besorgnis um das Seelenheil jener Schwätzer hat doch hie und da schon so zusammengeklingelt, daß der bessere Theil in der Seele gestärkt und geläutert wurde.

Jeden Falles aber lasse mans in der Regel bei kurzer Begegnung. Man beherberge und hege solche Leute nicht, lade sie nicht zum Wiederkommen ein, weise ihre persönlichen Anerbietungen freundlich ab, bleibe etwas ferne von ihnen stehen, lasse sie nicht in Gemeinde- und Familienverhältnisse schauen, kurz man sei gegen sie ganz einfach, wie es einem liebreichen Pastor gegen einen irrenden und verkehrten Bruder geziemt.

\quellseite{206}
Das hieße Gastfreundschaft am verkehrten Orte üben, wenn man die aufnehmen wollte, welche aus dem Umgang mit Pfarrern Nahrung für den Stolz und für jene Anmaßung nehmen, welche sie gegen ihres Gleichen üben! Sünden Vorschub leisten ist der Pfarrer Sache nicht.

\section{X. Nachbarschaft}

103. Unter der Nachbarschaft verstehen wir hier nicht die Nachbarschaft im engern Sinn, nicht die Gemeinschaft mit denjenigen Gemeindegliedern, welchen der Pfarrer zunächst wohnt; sondern die Nachbarschaft im weiteren Sinn, die Gemeinschaft mit den benachbarten Pfarrersfamilien. Es ist zunächst von der Nachbarschaft der Landpfarrer die Sprache. Wer mit dem Leben der Landpfarrer bekannt ist, der weiß, wie sehr diese Nachbarschaft von vielen Pfarrern gepflegt wird. — Wer wollte nun bezweifeln, daß gerade dieses Zusammenleben benachbarter Landpfarrersfamilien recht heilsam und fördernd, nicht bloß recht wohlthuend und ergötzlich sein könne? Benachbarte Landpfarreien stehen in immerwährenden und sehr mannigfaltigen Beziehungen zu einander. Die Fluren, die Wiesen, die Wälder grenzen zusammen, die Gemeinsamkeit der Rechte, die Gewohnheiten und Bräuche und Sitten \&c. bringen viel Zusammenhang hervor. Die Gemeinden berühren sich, eben so die Pfarrer und ihre Familien. Auch sie haben gemeinsame Bedürfnisse, können einander vielfach aushelfen,
\quellseite{207}
berathen, — im leiblichen und geistlichen Leben aufs Schönste zusammengehen. Dazu macht es den besten Eindruck auf benachbarte Gemeinden, wenn die Hirten einig und Freunde sind. Gewiß, zu den edelsten Vorzügen des Landpfarrerlebens gehört dies Zusammenleben der Pfarrersfamilien. — Wer zumal schon einen bösen Amtsnachbar gehabt hat, weiß es zu rühmen, welch eine edle Gottesgabe „getreue Nachbarn“ sind.

104. Aber — und dies Aber werde nicht minder gewürdigt, die Nachbarschaft werde nicht zu einer abgeschmackten und langweiligen Geselligkeit benützt. Nicht der Zeitvertreib, sondern die Liebe, nicht bloß die Gemeinschaft zeitlicher Bedürfnisse, sondern insonderheit die Pflege geistlichen und geistigen Lebens führe die Pfarrer und ihre Familien zusammen. Da es dem Landpfarrerleben ohnehin gar leicht an Manchfaltigkeit, fehlt so führe man in der Unterhaltung nicht wieder die täglichen Vorkommnisse ins Unendliche aus, man thue kurz ab, was zeitlich ist, und gehe dann mit einander zu den höheren Interessen der Kirche Gottes vorwärts. Es gibt Gaue, in welchen die oft sehr nahe wohnenden Pfarrer ein wahres Phäakenleben führen. Jeder Tag hat da seinen Sammelplatz der benachbarten Geistlichkeit, hie und da buchstäblich jeder vom lieben Sonntag bis zum lieben Freitag, wenn nicht gar bis zum lieben Studiertag, dem Sonnabend. Was soll es da jeden Tag zu reden geben? Ja, was für eine Freude und Erholung kann da sein? Ohne Stillleben gibt es keine gesellschaftliche Erhebung, ohne Arbeit keine erquickliche Ruhe, ohne
\quellseite{208}
einsames Sinnen kein richtiges, interessantes Gespräch. Kurz, wer die Gesellschaftlichkeit alltäglich macht, nimmt ihr Reiz und Werth und ist überdies ein Faullenzer, welcher in dem Kreise guter Freunde ein böses Gewissen hat. Daher verweltlichen auch Pfarrersfamilien ganz und gar, welche allzunachbarlich leben und zu viel und oft beisammen sind. Die Unterhaltung, die sie führen, ist ein richtiges Wetterglas ihres Lebens. Oberflächlich im Amte, wie in der Gesellschaft sind solche Pfarrer. Nichts Edles, nichts Großes, keine gemeinsame Anstrengung fürs Heil der Gemeinde, keine Liebesabsicht beseelt sie. Bier, Tabak und Karte — wenn nicht gar noch ganz andere Dinge — sind die Bindemittel solcher Gesellschaften. Von ihnen sei mit Ekel und Indignation Abschied genommen.

105. Ein angehender Pfarrer sei selten, aber dann von ganzem Herzen bei seinem Amtsnachbar. Sein Amt, sein Tagewerk, seines Lebens Ziel, sein Studium, die väterliche Fürsorge für sein Haus halte ihn am Pfarrorte und in der Beobachtung seiner Residenzpflicht\footnote{Ein Ausdruck aus der römisch-katholischen Pastorallehre.}, welche eine oftmalige Entfernung vom Orte geradezu verbietet. Kommen sie dann, die Freunde, die sich gerne sehen und auf den Conferenztag sich freuen, zusammen, so geschehe es, wie beim Besuch der Leidenden und Kranken, daß man das Zeitliche vorweg kürzlich befrage, besage und ordne. Und nun wohlan vorwärts zum Werke des Amtes, dadurch die Heiligen zugerichtet werden. Allerdings müssen dann die Männer die Vorredner und Hauptredner sein.

\quellseite{209}
Den Frauen aber ist Frage und Antwort wohl zu gönnen. Ihnen wird es zumal wohl sein, dem Gespräche der Männer zu horchen und dadurch einmal der Haushaltungssphäre entrückt, gehoben, gefördert zu werden. Da bedarfs keiner Protokolle, keiner Formen, keines Zwanges. In freier Lust, innerhalb geliebter Grenzen des Schicklichen und Reinen, ergeht man sich unbefangen und harmlos. Man findet keine Freude daran, von „dem Hundertsten ins Tausendste“ zu kommen. Was man bespricht, wird abgehandelt nach Kraft und Möglichkeit. Nicht in der Oberflächlichkeit, sondern in der Gemeinsamkeit der Besprechung liegt der Reiz. Munter und fröhlich, Rede und Gegenrede vertragend, letztere erwartend und erbittend, ergießt man sich über die Gegenstände des Gesprächs. Widerspruch erweckt Ueberlegung, aber keine Bitterkeit, Lust, aber nicht Beleidigung, die innere Ruhe wird dadurch nicht gestört, sondern nur gewürzt und angenehm gemacht. Jeder darf alles sagen, was ihm aufstößt, — auch Unreifes und Unvollkommenes darf sich zu Tage wagen. Es ist vorherbereit ausgemacht, daß hie keine im Wort vollkommenen Männer beisammen sind. Drum braucht auch keiner durchaus Recht zu haben, keiner scheut sich, das falsche oder verkehrte Wort zurückzunehmen. Im Behaupten, Verständigen und Zurücknehmen mehrt sich Klarheit und Liebe; umsichtiger und voller, einsichtiger und gewichtiger wird Urtheil und Stimme, immer geistiger und bedeutender das Leben.

So ist es hie und da bei benachbarten Pfarrern und Pfarrfamilien, und gesegnet sind die Pfarrhäuser, in denen eine solche nachbarliche Gemeinschaft nicht allzuoft, nicht
\quellseite{210}
allzuselten gepflegt wird. Das Vorbild freier, heiliger Christenliebe zeigt sich da. Gott und seine Engel freuen sich solcher Versammlungen.

106. Bei diesem Zusammensein benachbarter Pfarrer und Pfarrfamilien herrsche doch ja kein Neid. Neid tödtet die Liebe. Liebe aber ist des Neides sicherer Tod. Darum überwinde man sich, den weisesten Mann der Versammlung zu lieben, statt ihn zu beneiden. Nicht Stimmenmehrheit, sondern Stimmengleichheit und Stimmeneinheit, und zwar nicht bloß eine äußere, sondern eine innere, ist das Ziel alles rechten Conferenzierens. Eine Stimmengleichheit aber, so gedacht, daß alle Anwesenden, jeder aus sich heraus, dasselbe mit derselben Klarheit, mit derselben Fülle und in derselben Begrenzung finde und sage, wie der andere, ist die unnatürlichste Forderung, die es gibt. Wo immer zwei Menschen zusammentreffen, ist einer der tüchtigere und fähigere, — treffen mehr zusammen, ist es natürlich nicht minder also. Da ist des Fähigeren Tugend demüthige Erweisung seiner Tüchtigkeit, — und der andern Tugend ist freie, einsichtsvolle Erwägung und Würdigung der besten Stimme, Prüfung alles dessen, was vorkommt, Behaltung des Guten. Es ist ja doch nur der Hochmuth der Schwachheit, wenn es anders ist, — und dieser schwache Hochmuth reizt den Hochmuth dessen, der die andern übersieht. Und eine armselige Zeit ist es, welche sich gegen die Austheilung der verschiedenen Gaben von Seiten des Schöpfers nicht anders zu rächen weiß, als durch die blähende Behauptung, daß einer dem andern an Gabe ebenbürtig sei, die Stimme des einen der Stimme des
\quellseite{211}
andern an Kraft und Inhalt gleichkomme. Daher dann die Erfahrung, daß so oft nichts ab- und auszumachen ist, weil man wider besseres Wissen und Gewissen nicht aufgeben mag, was offenbar besser ist. — Sie schenke Gott allen seinen Knechten die rechte Demuth!

\section{XI. Verhältnis zu Frauen}

107. Wir reden hier nicht vom Freien und Heirathen, das wir einem nachfolgenden Kapitel aufheben; sondern nur von der Freundschaft, welche zuweilen zwischen Männern und Frauen besteht. Je männlicher und männlich ausgezeichneter der Mann ist, desto mehr findet er Anerkennung bei Frauen, und zwar auch eine von geschlechtlichen Beziehungen völlig fremde. So natürlich es dem Manne, auch ohne eigennützige Absicht, ist, am Weibe Schönheit und Anmuth zu ehren und jede Frauentugend, wenn sie in diesem Lichte einhertritt, desto höher zu achten; so natürlich ist es dem Weibe, am Manne die innere Herrlichkeit eines männlich vollendeten Charakters zu suchen und zu ehren. Gerade hierin liegt eine Bürgschaft weiblichen Werthes. Eben daher kommt es, daß ausgezeichnetere Frauen oftmals Gemeinschaft mit männlichen Charakteren suchen. Und gerade im geistlichen Stande finden sie oft solche Charaktere am liebsten. Es geschieht daher auch sehr häufig, daß selbst sehr edle und hochgeborene Frauen besondere Freude an persönlicher und brieflicher Verbindung
\quellseite{212}
mit ausgezeichneten Geistlichen finden. Es fragt sich nun, wie sich in solchem Falle ein Geistlicher zu benehmen habe? — Fragen wir die Volksstimme, so wird sie immer ein Mißfallen daran haben, sie wird dergleichen im mildesten Falle als eine Verirrung bezeichnen oder zu etwas Lächerlichem stempeln. Fragen wir die eigene Stimme, so wird auch sie nicht eben etwas Natürliches drin finden. Es ist und bleibt ein eigenes Gefühl, freundschaftlichen Umgang mit Frauen persönlich oder brieflich pflegen zu sollen, — nemlich so, daß mehr als das gewöhnliche Interesse gesellschaftlichen Lebens genommen wird. Fragen wir die Alten, so werden auch sie ihr Urtheil dahin abgeben, daß man nicht zum Vergnügen und bloß aus Freundschaft dergleichen Verbindungen pflegen solle. Sie werden vor der Möglichkeit der Ausbildung eines verkehrten geschlechtlichen Verhältnisses zurückbeben. Sie werden einem Pfarrer die stille Ruhe seiner Seele wahren, ihn vor der üblen Nachrede böser Menschen sicher haben, die christliche Freiheit durch christliche Klugheit eingeschränkt wissen wollen. Und wer will sagen, daß sie Unrecht haben? Zumal wenn von unverheiratheten Frauen, gleichviel älteren oder jüngeren, die Rede ist, legt die Weisheit ihr veto ein und ein weiser Mann nimmt dieses veto mit Freuden an. Es ist den Frauen eben damit der sicherste Weg der Seelenruhe nicht minder angezeigt.

108. Es versteht sich übrigens von selbst, daß hier nur von berufslosen Verbindungen die Rede ist. Wo es der Beruf des Seelsorgers erfordert, Frauen näher zu treten, — schriftlich oder persönlich, mag ein Seelsorger
\quellseite{213}
die mögliche Anfechtung bedenken und auf seiner Hut sein; aber er wird deshalb keinen Schritt zurücktreten, noch zaudernd zu seinem Berufswerk schreiten. Der Gott, der seinen Knechten das Amt an alle Creaturen gibt, kann es ihnen auch leicht machen, es an allen Menschen auszuüben. Es ist nicht des Amtes, sondern ihre eigene Schuld, wenn sie da in irgend eine Verkehrtheit gerathen.

109. Was den Briefwechsel insonderheit betrifft, so ist er allewege zu zeitraubend, um von einem Diener des Evangeliums überhaupt ohne Anlaß geführt zu werden. Kaum werden alte Freunde, die im Amte leben, ohne Anlaß regelmäßig korrespondieren. Viel weniger wird ein freundschaftlicher Briefwechsel zwischen Pfarrern und Frauenspersonen Statt finden können. Dazu kommt noch, daß der Mensch in Briefen meistens ein anderer ist, als im gewöhnlichen Leben, — und zwar nicht bloß dem augenblicklichen Schwung und der Erhebung nach, sondern oft auch dem Wesen nach ein anderer. Nun ist das zwar möglich und geschieht auch wirklich oft, daß sich ein Mensch in Briefen wahrer gibt, als im Leben. Das Innerste ist oft so heilig, daß es öffentlich, vor vielen sich nicht darlegen läßt. So findet man oft wahrhaft zarte Seelen in ihrer Erscheinung ungeschickt und verletzend; oft ist die Liebe in eine abstoßende Gestalt verhüllt. Da kann es denn auch kommen, daß der verborgene Mensch des Herzens sich in einem Briefe richtiger gibt als im öffentlichen Wandel. Doch aber mag der umgekehrte Fall vielleicht der häufigere sein. Ja, es wird in Briefen, die man nur schreibt, um sich mitzutheilen, oft eine bewußte Verhüllung der Wahrheit,
\quellseite{214}
eine künstliche Gleisnerei und Heuchelei ausgeübt, — und für manche ist ein Brief, nur um sich mitzutheilen, geschrieben, eine wahre Versuchung zu selbstbeschönigender Lüge. Nun sind insbesondere die beiden Geschlechter geneigt, sich vor einander zu beschönigen, einander zu heucheln, — und bei einem Briefwechsel, wie der ist, von dem wir zunächst reden, kann deshalb doppelte Versuchung eintreten. Möge deshalb jeder Pfarrer auf seiner Hut sein und thun, was nach beiden Seiten gesund und heilsam ist, auch wenn es weniger gefällig erscheinen sollte. — Es heißt eben auch hier: Homo sum, nihil humani a me alienum puto. Wer sein und anderer sicher ist, ist unverbunden, diesen Rath zu befolgen.

\section{XII. Verhalten gegen die Armen}

110. Die Armen sind entweder völlig Verarmte und Bettler, oder solche, welche sich bei aller ihrer Dürftigkeit doch immer noch einen Theil ihres Lebensunterhalts selbst erwerben können. Gemäß den Bestimmungen der staatlichen Armenpflege sollen zwar durchaus keine Bettler geduldet werden, und der Pfarrer, welcher in der Regel an der Spitze der staatlichen Armenpflege steht, soll wohl gar in einzelnen Fällen diesen (gewiß unmöglichen) Grundsatz der Armenpflege in Ausübung bringen. Es steht in seiner Befugniß, wo nicht gar in seiner Pflicht, die Bettler greifen zu lassen, und hat er das nur erst ein oder etliche
\quellseite{215}
Male gethan, so kommt er schwerlich in den Fall, es noch einmal thun zu müssen. Allein so wenig der Bettel zu loben ist, so sehr man ihm in aller Weise steuern soll, so wenig vereinigt es sich doch mit dem geistlichen Charakter eines Pfarrers, gegen die Bettler die strengen Maßregeln des Staates auszuführen. Dem Bettler seine Gabe in einer Weise reichen, die ihm ein Antrieb zum Guten oder eine Bestärkung darin sein kann, ihm die leibliche Gabe in geistlicher Weise reichen, das ziemt dem Diener des Worts. Seine Stellung innerhalb der staatlichen Armenpflege darf ihn aber zur Ergreifung einer Maßregel vermögen, durch welche er nicht nur den Charakter des Seelenhirten mit dem eines weltlichen Polizeibeamten vertauschen, sondern auch dem Erzhirten Jesus unähnlich werden würde, der Arme selig preist und der kein Mittel weiß als die Liebe, wenn es gilt, die Unterschiede des Reichthums und der Armuth auszugleichen. Am besten ist es daher, wenn er seine Stellung innerhalb der staatlichen Armenpflege benützt, durch persönliche Einwirkung und Bemühung desto mächtiger dahin zu bringen, daß alle Maßregeln des Staats durch christliche Liebeswerke überflüssig werden. Wie er das anzufangen habe, das muß seiner Weisheit überlassen bleiben, da überall andere Verhältnisse obwalten, die auch andere Maßregeln und Mittel zum Zweck erfordern. — Es wird nie gelingen, den Bettel völlig aufzuheben; es gibt auch schlimmere und größere Uebel, aus denen er selber quillt, welche durch gewaltsame Unterdrückung, auch wenn diese möglich wäre, keineswegs gehoben werden würden. Es wäre z.\,B. viel besser,
\quellseite{216}
wenn die zahllosen Menschen, welche ihrem Mangel durch Diebstahl abzuhelfen suchen, zum Bitten und Betteln ihre Zuflucht nähmen. Betteln ist ehrlich, Stehlen ist ein Uebel, das vom Himmelreich ausschließt. Zwar würde man, wenn alle armen Diebe und Betrüger ehrliche Bettler würden, erst recht sehen, welch ein unabsehbares und unergründliches Meer des Elends und göttlicher Strafen unser Vaterland bedeckt; aber das Uebel wäre doch schon da, es würde damit nicht erst erzeugt oder größer werden, wohl aber könnte die Rückkehr zu einem sittlicheren Leben — denn das wäre doch der Schritt vom Diebstahl zum Bettel — auch mehr Kraft mitbringen, dem Jammer abzuhelfen, jeden Falls mehr Kraft zur Erduldung desselben verleihen. Es liegt in diesen Sätzen etwas auf die Spitze Stellendes, aber ohne Zweifel etwas Wahres, worin wir wünschten, nicht mißverstanden zu werden. Bettlerzünfte sind etwas Schreckliches, wider sie soll und muß man eifern; aber den Unterschied zwischen Bettlern und Bettlern sollte man nie vergessen; nicht Jeder, den man einen Bettler nennt, kann sich in seinen Umständen anders helfen und ist zu den zukünftigen Bettlern zu rechnen. Die Scheu, einem redlichen oder gar frommen Lazarus wehe zu thun, — die Lust zum gesegneten Becher kalten Wassers, Matth. 10, 42., von welchem Christus redet, und zu seinem großen Lohn sollte einen Pfarrer eifrig machen und ihm die Sinne schärfen; den Bettler vom Bettler zu scheiden und dem frommen Bettler zu dienen. Das geschieht nun freilich nicht dadurch, daß man kurzweg jedem Bettler eine Gabe reicht und sich dann getröstet, daß man auf diese Weise
\quellseite{217}
doch auch keinen frommen und redlichen Bettler versäumt habe. Das ist wohlfeile Barmherzigkeit, und wenn ein Pfarrer, der seine Stellung als Armenpflegvorstand nicht geltend zu machen weiß, weil er sie auch nicht benützt, auf diesem Wege in die Hände bettelspähender Gensdarmen geräth und als Uebertreter obrigkeitlicher Gebote angezeigt wird, so hat ers mindestens damit verdient, daß er den Bettlern nichts Besseres zu thun wußte oder zu träge war, es zu thun. Dagegen liegt es ganz in der Befugniß eines Pfarrers, als Vorstandes der örtlichen Armenpflege, sich der Armen und Bettler in jeder Weise ernstlich anzunehmen. Wider solche ist das Gesetz des Staates nicht; vielmehr stehen ihnen die Anstalten und Mittel des Staates zur Seite, und sie haben nur Anerkennung zu erwarten. — Wir reden hier nicht von der Seelsorge der Bettler, nicht von dem Bettel als einem pastoralen Zustande. Kein Mensch wird auch sagen, daß ein Pfarrer, welcher nach obigen Grundsätzen selbst gibt und seine Stellung als Vorstand der Armenpflege benützt, seinen Wirkungskreis überschreite, in fremdes Amt greife. Der Bettler fällt in die Behandlung heilender Liebe, ja auch der Seelsorge, wo er gerade bettelt, etwa wie der auswärtige Taglöhner, der in einer fremden Pfarrei krank wird, von dem Seelsorger der Pfarrei behandelt wird, in welcher ihn Gott krank werden ließ. Er wird von dem Pfarrer, in dessen Bereich er Gaben sammelt, für den Augenblick seelsorgerisch, weil liebend, behandelt — und dann seinem heimatlichen Pfarrer treulich wieder zugewiesen.

111. Ehedem, da die öffentliche Armenpflege noch
\quellseite{218}
ganz oder wenigstens theilweise Kirchensache war, wurde den Armen, auch den Bettlern, nach Anweisung des Pfarrers (man hatte dazu eigene Erlaubnißzettel) vom Kirchenpfleger eine Gabe gereicht, — natürlich aus dem Kirchensäckel. Dadurch war der Pfarrer in den Stand gesetzt, den Armen möglichst zu dienen; sie waren an ihn gewiesen; was er that, geschah ganz im Sinne des Amtes und im Namen einer christlichen Gemeinde. Diese schöne und gewissermaßen auch leichte Stellung hat gegenwärtig der Pfarrer nicht mehr. Es wird, da er nun aus eigenen Mitteln geben soll, zumal in manchen Gegenden, wo der Bettler viele sind, eine schwere Aufgabe sein, daß er mit dem wenigen, was er vermag, genüge.\footnote{Doch kann ja der Pfarrer, wenn er will, einen Verein gegen den Bettel gründen und im Namen desselben die Gabe abgeben. Das gewährt einigen Ersatz für den Entgang der früheren Stellung. Solche Vereine sind überall, wo Bettel einreißt, gut und nöthig, auch in Landpfarreien. Ein Pfarrer hielt Register über die Handwerksburschen, denen er Gaben verabreichte, und fand, daß es immer dieselben Handwerksburschen waren, welche zu bestimmter Frist wiederkamen. Ebenso war es mit den andern Reisenden und Bettlern. Die formale Behandlung der Sache brachte den Gewinn, daß man seine Leute und den Werth des bloßen Gebens kennen lernte. Man konnte sich hüten und den Bettlern die bessere Gabe geben, die sie nicht suchten, aber brauchten, Abweisung und Ermahnung.} Wiewohl die Aufgabe doch auch wieder nicht zu schwer zu nehmen ist! Manch frommes Pfarrkind wird ihm die Hände füllen, und umgekehrt, mancher Bettler wird sich gar nicht zum Pfarrer weisen lassen, wenn er weiß, daß er einem eingehenden Gespräche über seine innere und äußere Lage ausgesetzt sein wird. So schwer aber auch trotz dem die Aufgabe
\quellseite{219}
sein mag; sie ist es doch nur für den, welcher alle Augenblicke vergißt, daß er in dergleichen Dingen nie mehr soll, als er kann; sie ist leichter für den, der sich demüthig und verehrend auch im Geben in die Fügung Gottes und in seine Armuth findet, der großherzig den eignen Scherf zu schätzen mag, den er in JEsu Namen gibt, weil ihn JEsus achtete und achtet, und seine Seufzer über sein geringes Vermögen fröhlich in Fürbitten für die Armen verwandelt.

112. Bisher war von dem Benehmen eines Pfarrers gegen Bettler, d.\,i. gegen fremde Arme die Rede; denn die Bettelarmen und Bettler der eigenen Pfarrei hat er im Ganzen wie andere Arme der Pfarrei in gründlichere Behandlung und fortdauernde Pflege zu nehmen; bei ihnen wird noch mehr, als bei den fremden Bettlern die Quelle der Armuth aufzusuchen, der Bettel als einzelner Ausfluß des ganzen elenden Zustandes aufzufassen und eine Art radicaler Cur einzuleiten sein. Ein rechter Pfarrer wird in seinem Verhalten die Sprüche des göttlichen Wortes, nach denen kein Bettler in Israel sein sollte, mit jenem Worte Christi zu vereinigen wissen, da ER spricht: „Arme habt ihr allezeit bei euch.“ Er wird den Bettel zerstören und die Armen pflegen, sonderlich in der eigenen Pfarrei. Er wird auch bei den eigenen wie bei den fremden Armen durch Erforschung der Armuthsquellen zur Erforschung der religiösen und sittlichen Beschaffenheit der Armen und eben damit zur Unterscheidung der frommen Armen von den andern geleitet werden. Wie er dann als Seelsorger die frommen oder auch die gottlosen zu behandeln habe, das
\quellseite{220}
ist eine andere Sache. Hier ist eigentlich nur von dem allgemeinen Standpunkt und dem persönlichen Verhalten eines Pfarrers gegen seine Armen die Rede. Beharrlich gottlosen Armen kann sich der Pfarrer eben so wenig hingeben, als andern unbußfertigen Sündern. Was er allenfalls doch zu thun hat, ist casuell. Mit frommen Armen, d.\,i. mit solchen, welche die Armuth auf eine christliche Weise ertragen und ihr bitteres Loos durch eigenen Fleiß so viel als möglich zu lindern suchen, geht er aber auf das Freundlichste und Schonendste um. Er behandelt sie ganz nach der Lehre, welche er in seinem Amte von der Armuth vorträgt. Er stellt aber fromme Arme als Kleinodien der Gemeinden, mit Laurentius als ihre wahren Schatzkammern vor, — als Aufforderungen Gottes zu Lieb und Barmherzigkeit, — als Menschen, welche von Gott den theuern Beruf haben, die zumuthende Liebe auszuüben; — und das Wort des Herrn Joh. 12, 8: „Arme habt ihr allezeit bei euch“ wendet er gerne an, das Geschlecht der Armen als unsterblich für diese Welt darzustellen, als ein Geschlecht, in welchem Christus sich selbst ehren läßt und geehrt sein will, bis er sichtbar wiederkommt, Matth. 25, 40. Ertragung der Armuth erscheint als Nachfolge des armen Lebens Jesu. Andererseits wird das Geschlecht der Wohlhabenden und Reichen zu den Werken der Barmherzigkeit als zu seligem Thun vermahnt nach dem Wort des Herrn: „Geben ist seliger als nehmen;“ ihre ganze Lebensstellung wird ihnen im Lichte des himmlischen Berufes zu immerwährender Barmherzigkeit dargestellt. Wie der Arme den Einen Ton der Liebe, den der zumuthenden,
\quellseite{221}
gibt, so der Reiche den andern ergänzenden, den der entgegenkommenden und gewährenden Liebe. Arme und Reiche sind vom Herrn gewollte und bestellte Ausgangs- und Mittelpunkte thätiger Bruderliebe, ohne welche dem Leben die innigsten und seligsten Wechselbeziehungen fehlen würden. Einer hat den andern für nothwendig im Reiche Gottes anzusehen und zu schätzen; jeder den andern in seiner Weise, seiner Demuth zu ehren und zu fördern. Durch solche gewiß wahre Auffassung der Armuth und des Reichthums wird den Armen eine sehr würdige Stellung innerhalb der Gemeinde gegeben, und alle Tugenden der Armen wie der Vermögenden werden befördert, Barmherzigkeit und Dankbarkeit als Pflichten der Bruderliebe erkannt und so verklärt. So lehrt und danach handelt ein Pfarrer. — Sein Leben innerhalb der Gemeinde gibt ihm Gelegenheit, die Armen und ihre Bedürfnisse kennen zu lernen. Er kann ihnen rathen, nach Vermögen in aller Stille selbst helfen, ihr Fürsprecher bei Andern werden. Als Fürsprecher für die Armen findet er, wenn er anders die Lasten auszutheilen versteht, offene Ohren. Kein Mensch glaubt vorweg, wie leicht die Freiwilligkeit auch bei solchen Menschen zu wecken ist, welche sonst nicht etwa zum Guten eilen. Versteht ein Pfarrer, diesen Hebel zu gebrauchen; so wird er nicht leicht nöthig haben, viele Arme in die Armenconscriptionslisten des Staates aufzunehmen. Alle Bedürfnisse der Armuth werden durch freiwillige Liebe viel eher und lieber als auf Befehl des Staates gedeckt.

113. Freilich stößt man auch auf große Hindernisse. Eines wollen wir im Anschluß an Nro. 110 noch
\quellseite{222}
einmal erwähnen. Es gibt Gegenden, in denen Diebstahl von Alters her gemein ist. Wenn z.\,B. die örtliche Lage einer Gemeinde Armuth an Holz, Streu und Wieswachs mit sich bringt; so stehlen die Aermeren, ja sehr oft auch solche, die sich sehr leicht anders helfen könnten, in der Regel, was sie bedürfen. In solchen Gegenden fügt sich der Vermögende und Reiche in den altherkömmlichen Diebstahl, er läßt sich aber nicht dazu bewegen, zu geben, auch nicht so viel, als ihm nach sicherer Voraussicht durch Diebstahl abhanden kommt. Noch mehr! Bei dem allgemeinen Stehlen stößt der Arme selbst bei vielen Vermöglicheren an, wenn er nicht auch stiehlt, sondern die barmherzige Liebe in Anspruch nimmt: stiehlt er, so ist er etlichen beschwerlich, — der Anspruch auf Unterstützung aber beleidigt alle, selbst wenn er noch nicht ausgesprochen ist, sondern nur zu fürchten steht (und wie groß ist in diesen Stücken die Furcht des Landmanns!). Da vernimmt dann wohl der frömmere Arme, der sich gegen die allgemeine Diebssitte sträubt und mehr als andere darbt, nicht bloß von andern Armen, sondern sogar von Wohlhabenderen die Rede: „Stiehl dir was, so hast du was!“ Da ist schwer zu helfen, dem Armen, wie dem Reichen; jener läßt schwer vom Diebstahl, dieser gibt schwer. Habsucht und Geiz sind wie alter Rost, der nicht bloß auf dem Eisen sitzt, sondern das Eisen anfrißt; sie sind schwer auszufegen. Durch Geschlechter fortgepflanzte Dieberei verliert den Namen, wird schwer eingesehen; — wer sie läßt, kommt oft in übleren Ruf, als wer an ihr Theil hat; hier wird zum Schurken gestempelt, wer sich bessern will
\quellseite{223}
— und fast wird die Besserung Unmöglichkeit. Da ist Sieg und Triumph des Christenthums schwer zu erlangen. Selten kommt er zu Stande, und wenn, so freut man sich seiner kaum nach Würden, weil noch so viel Uebel zurückbleibt, so viel Uebel und Sünde derselben fast unüberwindlichen Art. — Ach, hier ist ein trüber Fleck des Pfarrerlebens in Gemeinden, wie sie bei uns zusammengesetzt sind! Diebsgesichter von allen Seiten — und wenig Gewissen, wenig Hilfe! Hier helfe Gott und sein heiliges Evangelium!

114. Ein besonderes Auge habe der Pfarrer auf arme Kranke. Den Kranken überhaupt, und namentlich den armen Kranken ist der Pfarrer eine liebe Person. Das Volk lebt in dem Glauben, daß Segen von ihm ausgehe. Eine einfache Suppe aus dem Pfarrhof wird wie eine Arznei voll Segens mit größtem Danke hingenommen. Der vom Arzte verordnete Trank, die von ihm gut geheißene Speise, wird, wenn sie der Pfarrer schickt, als eine doppelt theure und werthe Gabe hingenommen. Hier steht dem Pfarrer eine wohlfeile Art von Hilfsleistung offen. Dieselbe laße er sich empfohlen sein. Aber er hüte sich, auch wenn er, wie das nicht anders sein kann, immer erfahrener wird, den Arzt zu machen. Seine Küche werde nicht zur Apotheke. Er greife in keinen fremden Beruf; er sei zufrieden mit dem seinen und ehre den des
\quellseite{224}
andern. Er sei in medicinischen Dingen, so weit es seine Ueberzeugung erlaubt, der willige Diener des Arztes und verweise immer auf ihn. Eingriffe in den ärztlichen Beruf bestrafen sich oft durch beschämendes Mißlingen, immer aber durch Gewissensunruhe und Angst wegen des Gelingens. Auch wenn kein Schade geschieht, wird doch vielleicht die rechte Hülfe aufgehalten, welche eingetreten oder gesucht worden wäre, wenn niemand unbefugt eingegriffen hätte. Selbst bei homöopathischen Mitteln gilt das: mögen sie schaden können oder nicht, es geschieht doch oft Schade, weil der Kranke durch Pfuschereien unbefugter Menschen die rechte Zeit versäumt, zum Arzte zu gehen. Schaden sie nicht, so nützen sie doch auch nicht, was sie doch sollten. Aber auch abgesehen von Nutzen und Schaden, wird sich doch kein rechter Mann mit dem befassen, was er nicht versteht, und wozu er keinen Beruf hat. — Es ist Schade, daß die Pfarrer, insonderheit die Landpfarrer den Beruf des Arztes nicht mit dem ihrigen vereinigt führen können und dürfen. Der Pfarrer als Arzt — wie bequem wäre er für die Leidenden, wie viel wohlfeiler könnte er sein, wie förderlich wäre die Verbindung der beiden Berufe in Einer Person für die psychische und pastorale Pflege der Kranken! Aber was hilft es, dies zu meinen, da es nicht ist und nicht werden darf? Thue ein jeder, was seines Amtes ist und überlasse jedem andern seine Verantwortung.

— Etwas ganz anderes ist es, ob nicht ein Pfarrer so viel Kenntniß des menschlichen Leibes und seiner Krankheiten haben sollte, als nöthig ist, die psychischen Einwirkungen der Krankheiten zu kennen, durch welche der pastorale
\quellseite{225}
Zustand der Patienten maskiert wird. Man wird diese Frage wohl bejahen müssen. Bekümmerten sich die Geistlichen um so viele medicinische Bildung\footnote{Zu empfehlen ist Bering’s psychische Heilkunde. Leipzig, Barth. 1817.}, so würden sie für die Seelsorge der Kranken tüchtiger werden, ohne deshalb zum Arzneien mehr Befugnis oder Beruf — oder auch nur Tüchtigkeit zu bekommen. Sie würden aber auch innerhalb der ihnen gesteckten Grenzen an Krankenbetten genug Gelegenheit zu persönlichen Liebeserweisungen finden. Wie viel Barmherzigkeit kann z.\,B. geübt werden, wenn man sich der Pflege des Kranken, der geistlichen nicht bloß, sondern auch der leiblichen annimmt! Hier ist ein Gebiet, welches nicht bloß der Arzt zu beherrschen hat, da ja der Kranken viele sind, die keinen Arzt zu Rathe ziehen. Luft, Licht, Wärme, Wohnung, Speise etc. des Kranken, was gibt das alles zu denken! Wie oft muß da der Seelsorger, von der Noth der Kranken ergriffen, darein reden! Wie viele Sünde kann auf diesem Gebiete geschehen und vermieden werden! Wie hilfreich kann hier ein Pfarrer seinen Gemeindegliedern werden! Wie viel Verantwortung wird er geben müssen, wenn er nicht wenigstens versucht, zu rathen und zu helfen! — Bücher wie: „Die Pflege bei Kranken und Gesunden. Kurze Winke, den Frauen aller Stände gewidmet, von Florence Nightingale. Leipzig, F. A. Brockhaus. 1861“, könnten einen Seelsorger auf vieles aufmerksam machen, was er selbst vielleicht nicht sieht, aber für sein armes Volk brauchen kann, sofern er selbst aufmerksam gemacht ist.

\quellseite{226}
\section{XIII. Der Geistliche und seine Söhne}

115. „Diese sind es, die einst an unsere Stelle treten!“ Diese Worte, von einem alten Pfarrer mit Hinweisung auf die zur Schulprüfung versammelte Jugend der Gemeinde gesprochen, sprechen die ganze Wichtigkeit aus, welche die Jugend für das Alter hat. Aller Wechsel, aller Fortschritt, alle Besserung und Verschlimmerung der Zeit wird sich gestalten, je nachdem sich die Jugend artet. Wäre nun das jugendliche Geschlecht einem Stoffe gleich, aus dem man machen kann, was man will; so würde das reifere Alter die künftige Zeit nach Wohlgefallen bereiten können, — man würde in der Jugend die Zeit erziehen können. Aber so steht es nicht. Dies junge Volk ist wohl bildungsfähig und dem Unterricht, der Anweisung, dem Beispiel, dem Einfluß des abgehenden Geschlechtes ausgesetzt, und das letztere hat deshalb hohe Verantwortung in Betreff desselben; allein von diesen Knaben und Mädchen hat nicht bloß jedes einzelne seine angeborne Art, sondern auch seinen Willen, kraft dessen es dem gesammten Einfluß der Väter widerstehen kann und auch widersteht. Dazu hat jedes Geschlecht seinen Trieb, mit dessen Lüften es in’s Meer des Lebens fährt, ohne ihn zu verstehen, und ohne darin von dem Alter verstanden zu werden. Das abgehende Geschlecht versteht etwa den Trieb, den wundersamen, durch welchen es sich von seinem Vorgänger unterscheidet,
\quellseite{227}
aber es versteht nicht den neuen Trieb der Jugend und das Ziel, worauf hin der Geist der Zeiten strebt; und eben so wenig versteht die Jugend selbst den Trieb und Zug, von dem es beherrscht wird, und die Zeit, welche sich dadurch bildet. So kommt’s dann, daß das Alter die Jugend, die Jugend das Alter nicht versteht, und daß die Jugend sich dem Alter nicht einfach überläßt und zur Leitung übergibt. Das Geheimniß der nächsten Zukunft und viele Sünde begegnen sich hier und machen die obige Hinweisung des alten Pfarrers auf die Jugend zu einem Ausspruch hohen Bedenkens und großer Sorge. Dem Alter ist die Jugend selten die Hoffnung der kommenden Zeit. Sie kann es auch nicht sein, — am wenigsten in unsern Tagen. So mächtig schreitet der Geist der Welt voran, so unaufhaltsam entwickelt sich das Böse, und so empfänglich für diesen Fortschritt ist das junge Volk, daß wir auf die Jugend unsere Hoffnung nicht stellen können. Der unerfahrene, junge Geistliche, dem die Schwierigkeit der pastoralen Einwirkung auf die Erwachsenen sich schnell enthüllt, meint, er könne wenigstens die Jugend so heranbilden, daß eine Gemeinde nach Wunsch entstehe; über seinem Grabe, hofft er, wird es anders aussehen, als über dem seines Vorgängers. Allein — so wird’s nicht oder im seltensten Falle, wenn sich alle Factoren des Productes vereinigen. Gewöhnlich entwindet sich die Jugend dem geistlichen Einfluß sehr schnell, und es kann kommen, daß derselbe Pfarrer, der mit der Jugend ein Paradies bauen wollte, allmählich die Alten wieder zu Baugenossen erwählt und mit ihnen, wenn kein Paradies, doch eine
\quellseite{228}
Nachthütte Gottes im Kürbisgarten der Welt zu bauen versucht. — Jedenfalls darf man von der Jugend nicht sehr viel hoffen, so wichtig sie für die kommende Zeit ist.

116. Deshalb aber bleibt die Jugend dennoch ein Zielpunkt des treuesten Fleißes für einen Pfarrer. Es ist ein Glück, wenn er gleich vornherein frei von sanguinischen Hoffnungen in die Mitte der Kinder tritt, aber desto ernster tönt ihm das Wort „Weide meine Lämmer“ in seinen Ohren. Der Geistliche darf keinen Theil am Verderben der Jugend haben. Kann er aus ihr nicht machen, was er will, ist er vom Willen und Unwillen des jungen Volkes, sowie vom Geiste der kommenden Zeit gehindert; so bleibt ihm dennoch die Pflicht, die Jugend in Zucht und Vermahnung zum Herrn erziehen zu helfen. Er wählt den besten Samen, streut ihn zur besten Zeit, mit weihester Hand, begießt und pflegt seine Saat mit heiliger Sorgfalt und betet ohne Unterlaß für dieselbe. Sein Herz ist der Jugend unwandelbar zugewendet. Das Gedeihen ist in der Hand des Herrn, dafür geben wir keine Antwort, aber für unsere Arbeit sind wir verantwortlich. Daher muß sie mit höchstem Fleiße, ausdauernd bis an’s Ende und mit aller derjenigen Weisheit geschehen, die wir nach dem Maße unsrer Gabe und Erfahrung anwenden können. Das hat dann auch seine Verheißung und seinen Gnadenlohn.

117. Der Pfarrer sucht die jungen Kinder überall, in der Familie, in der Schule, in seinem eigenen Unterricht, in der Kirche. Ja er sucht die Kinder, ehe sie geboren sind, denn er sorgt für die neugeborenen Kinder,
\quellseite{229}
indem er die Eltern auf ihre heiligen Pflichten gegen sie aufmerksam macht und für die Ungeborenen betet. Er lehrt die jungen Eltern geistliche Fürsorge für die Neugeborenen. Er lehrt sie beten mit den Kindern und für sie. Er segnet die Säuglinge. Er naht schon den Unmündigen, um mit ihnen bekannt zu werden. Er lernt alle Kindlein kennen und sorgt dafür, daß er selbst, der Pfarrer, ihnen von Kindesbeinen an bekannt werde. Sein Auge wacht über ihnen, über ihren Spielplätzen, ihren Spielen, wie es eben sein kann. Er geleitet sie in die Schule und besucht sie in der Schule und zwar immer in Einer Eigenschaft, in der besten, welche er hat, als Pfarrer, als Seelsorger. So weit der Schulinspector von dem Pfarrer verschieden ist, ist er der Jugend wohl eine ziemlich unfaßliche Person, — als Pfarrer aber versteht ihn sein Volk, auch sein junges, am ersten und besten. Er ist und bleibt auch in allen Fällen als Pfarrer am gesegnetsten, vielleicht sogar von denen am wohlgelittensten, die seinen Einfluß auf die Schule nicht gerade wollen. Das ist ja so ganz in der Ordnung, daß ein Pfarrer die Gemeinde, also auch die Lämmer weide. Diese seine Stellung benützt denn auch ein weiser Pfarrer. Er geht in die Schule, ermahnt zu allem Guten, betet mit den und für die Kinder, segnet sie, lehrt sie, läßt seinen andern Religionsunterricht geben, ohne genau zu prüfen, was und wie gelehrt wird, — duldet nichts Verderbliches, — legt sich mit aller Macht seines Amtes und seiner ganzen Stellung drein, wenn sich Verführung zeigt, und schirmt die Lämmer Jesu mit jeder geziemenden Waffe gegen Wölfe.

\quellseite{230}
Die Sprüche, der Catechismus, die Lieder, die biblische Geschichte und das Bibellesen, der geistliche Gesang und das Gebet, diese Hauptsachen aller Lehre, welche jede Schule, so lange sie gelehrt werden, zu einem Vorhof der Kirche, zu einer Werkstatt der Kirche, zu einem Theil der Kirche machen und ihr den Stempel Jesu aufdrücken, werden von dem treuen Pfarrer unter allen Umständen beachtet, gepflegt, überwacht, auf dem Herzen getragen. Der gesammte geistliche Unterricht aber wird als Weg zur Confirmation und zum Sacramente angesehen und dargestellt. Weg und Ziel, Schule und Alter werden in ihrer Verbindung den Kindern immer und immer gezeigt. Die ganze Jugendzeit wird von dem Geistlichen als Vorbereitung zum Sacramente angesehen. Auch lehrt er die Kinder sie so betrachten und sucht in allen Kindern den Sinn für das Ziel und den Mittelpunkt alles geistlichen und kirchlichen Lebens auf Erden, die Sacramentsgemeinschaft, zu wecken und wach zu erhalten. So wird dann der Höhepunkt aller Beziehung des Geistlichen zur Jugend der eigentliche Confirmandenunterricht, der sacramentliche Unterricht. Auch wenn ein Pfarrer die Jugend sonst weniger lehrt, als lehren läßt, kann er in der vollen, oben in Grundzügen dargestellten Stellung zur Jugend sein; aber der sacramentliche Unterricht (in Franken oft der „Sechstwochenunterricht“ genannt) wird von ihm selbst gegeben, mit großem Ernst, mit Inbrunst, mit Gebet und Flehen, mit Erwecken aller seiner Gabe zu lehren, zu strafen, zu lenken, zu stärken, zu erziehen. Seine ganze Absicht und sein Fleiß geht dahin, Pfarrkinder, Beichtkinder, Brüder und Schwestern,
\quellseite{231}
Heilige Gottes zu erziehen, die, wenn die Zeit der Jugend vorüber ist, in Gemeinschaft mit dem Pfarrer und der Gemeinde zu Gott treten, mit ihm handeln, mit Gott wandeln und ihr ewiges Heil schaffen können. — Nichts ist gewisser, als daß wenige Wochen nach der Confirmation in gewöhnlicher weltlicher Umgebung hinreichen, alle Frucht einer ganzen seelsorgerischen Führung zu zerstören. Aber so oft dies auch erfahren werde, ein frommer Pfarrer bleibt in seinem Fleiße und sucht die Schafe Jesu, als gienge nie eines wieder verloren, als müßten wenigstens die unverloren bleiben, die er gerade in der Pflege hat. In diesem Stücke sucht er nur Recht zu thun, — überläßt alles Weitere Gott.

118. Das alles wird, so groß und schwer es ist, dennoch gelingen, — wenn Gott einem Pfarrer einen Schullehrer gibt, mit dem er nach einem Ziele streben kann, und der die heilige Gewalt über sich selbst hat, als der Zweite mit einem Ersten, mit seinem Pfarrer, zu gehen. Diese Gabe aber ist, namentlich in unserer Zeit, von so hoher, großer Bedeutung für Amt und Leben eines Pfarrers, daß sie des anhaltendsten und brünstigsten Flehens werth ist. Sie wird in der That nur selten gegeben. Der Schullehrer fühlt den Unterschied nicht der Bildung, aber des Bildungsganges, so wie der Stellung und der äußern Lebensverhältnisse, der nun einmal vorhanden ist, tief im Innern, oft mit Neid und Ingrimm. Sich unterordnen zu sollen, ist jedem Stolze schwer, wird aber dem Schullehrer, gegenüber dem Pfarrer, oft ganz besonders schwer, zumal wenn er etwa Gaben besitzt, vielleicht gar
\quellseite{232}
auf eine Seite hin, wo der Pfarrer keine hat, — oder wenn er in dem gegenwärtig herrschenden Wahne lebt, als sei eigentlich er und sein Stand der nützere und segensreichere. Wie viele Pfarrer wüßten hievon nicht zu reden. Wie viele aber sind auch der Aufgabe gewachsen, ihre Schullehrer mit Liebe, Geduld und Langmuth zu überwinden! O schwere Aufgabe und seltenes Gelingen! Die gegenwärtige Ausbildung der Schullehrer wird keiner loben, welcher die Ueberzeugung hat, daß die Schule nur dann recht gedeihen kann, wenn sie Organ und Theil der Kirche ist; aber so viele Hindernisse sie der Schule selbst, dem Verhältniß der Schule zur Kirche und der Kirche selbst in den Weg legt, — unüberwindlich sind die Hindernisse doch nicht. Daher sollte ein neu eintretender Pfarrer allen Fleiß anwenden, sie zu überwinden, den Schullehrer für die Kirche und sich in ihm einen Freund zu gewinnen. Ist denn etwas Schöneres nicht bloß, sondern etwas Segensreicheres, für Pfarrer und Schullehrer selbst aber Glücklicheres und Süßeres, als wenn Pfarrer und Schullehrer im richtigen Verhältniß zu einander stehen und dies Verhältniß als treue Freunde und Brüder aufrecht halten? — Möchte daher jeder junge Pfarrer weise und vorsichtig sein, sich dies Glück nicht zu verderben oder den Erwerb nicht gar zu schwierig durch das erste Benehmen zu machen! Vor allem erkenne der Geistliche Stand und Gabe des Schullehrers mit aller Aufrichtigkeit und Freudigkeit an — und zwar mit Worten und That. Sodann setze er sich mit demselben in das Verhältniß der anständigsten Freimüthigkeit und Offenheit, und mache sich völlig klar, wie
\quellseite{233}
er sich gegen den Schullehrer als Pfarrer und Schulinspector zu verhalten habe. Sein ganzes Verhalten lege er sich selbst, wie dem Lehrer offen dar. In der Aufrichtigkeit liegt Schlangenklugheit, die jederzeit den Weg vorwärts und rückwärts findet. Nach der Regel Jesu, daß der Vornehmste der Andern Diener sei („Wer Dir dient, der regieret“ — betet die Kirche), erkenne er sich als einen willigen Helfer und Diener des Lehrers in allem an, worin dieser selbst Christo und Seinem Worte dient. Je nachdem es dem Lehrer lieb ist, messe man ihm die Hülfe zu und drücke ihn nicht unnöthig durch Inspection und Ueberwachung. Da aber, wo das Heil der Jugend oder Gemeinde, oder des Lehrers selbst es verlangt, scheue sich ein Geistlicher auch nicht, stark vorwärts zu gehen wider des Lehrers Sinn. Der Lehrer muß wissen, daß der willige Dienst des Pfarrers nicht Schwachheit ist. In allem Falle bleibe man ruhig, leidenschaftslos, vor dem Angesicht Gottes, Herr seines Wortes, der theils sich nicht verredet, wenn aber je, auf der Stelle zurücknimmt, immer bei der Wahrheit bleibt, höher als Alles den Ruhm und die Ehre der Redlichkeit und Wahrhaftigkeit stellt. Dazu bete man ohne Unterlaß für den Lehrer, lasse die Sonne nicht untergehen, ohne den Groll gegen ihn wegzuwerfen, wenn man je Ursache hatte, Groll zu fassen. Man bewahre dem Lehrer allezeit ein treues, gütiges Herz, das nicht leicht die Hoffnung aufgibt. Solch treues priesterliches Benehmen segnet der Heilige und Allmächtige, oder, wenn es an dem Lehrer abprallt, macht es doch die Seele
\quellseite{234}
ruhig und der Herr wird an dem Glauben Wohlgefallen haben, der sich in so gutem Werke erweist.

\section{XIV. Dienstboten}

119. In der Regel bringen Eltern ihre Kinder gerne in Pfarrhäusern unter, wenigstens dürfte dies vom Landvolk gelten. Die ehrenvolle Stellung des Pfarrers innerhalb seiner Gemeinde geht zum Theil auch auf die Dienstboten über; die Dienste pflegen nicht sehr beschwerlich oder doch einträglicher als manch anderer zu sein; und was eine Hauptsache ist, die Eltern glauben ihre Kinder in Pfarrersfamilien in Betreff der Sittlichkeit sicherer als anderwärts aufgehoben. Auch die Dienenden selbst pflegen gern in Pfarrhäusern zu sein, wenn sie nicht schon gar zu selbstständig und eigensinnig oder in Sünden verhärtet sind. Ja, man hat Beispiele, daß der Eintritt in einen Pfarrersdienst im Leben eines Dienstboten ein Wendepunkt fürs ganze Leben wurde; einen so anziehenden und zugleich heilsamen Einfluß üben selbst in unserer Zeit, wo viele alte Bande sich lockern, noch immer Pfarrersfamilien auf die Gemeinden aus. — In der Regel hat daher auch ein Pfarrer in Betreff der Dienstboten größeres Glück als andere. Wenigstens ist es sehr häufig, daß man Pfarrersfamilien trifft, welche mit ihren Dienstboten, und umgekehrt Dienstboten, welche mit Pfarrersfamilien zufrieden sind. Um so mehr muß ein Pfarrer bedacht sein, den alten Ruhm der Pfarrhäuser in diesem Stücke zu bewahren.

\quellseite{235}
120. Ein Pfarrer braucht sich gegen seine Dienstboten nicht besonders annehmend zu stellen, da ohnehin das Amt, welches er führt, und die Verhältnisse, in welchen er lebt, annehmend und freundlich genug sind, um auch dem Dienstboten freudiges Behagen zu verschaffen. Daher kommt es auch, daß an einem Pfarrer selbst ein starker Mannesernst nicht verhindert, sein Herz und seinen Sinn zu finden und zu verstehen. Doch braucht es keinen Ernst, der mehr als ernste, sittliche Haltung wäre; es bedarf weder viel Befehlen, noch viel Schelten und Verbieten; die Predigt, die Katechese klingt im Hause nach und spricht auch da noch ans Herz. Männliche Haltung von Seiten des Pfarrers reicht aus, ein Pfarrhaus in Ordnung zu halten, Ehrerbietung einzuprägen, wenn nicht die häuslichen Verhältnisse allzusehr zerrüttet sind; denn im letzteren Fall kann es allerdings — man denke an die Einzelwirthschaft eines Pfarrers, sei er ledig oder verwitwet — schwierige Aufgaben zu lösen geben, zu denen nicht bloß eine männliche Haltung, sondern auch ein männliches Durchgreifen gehört, wenn nicht der ganze Haushalt zu Grunde gehen soll.

121. Besonderes Augenmerk, wohl mehr als in andern Diensten nöthig ist, hat ein Pfarrer auf Treue und Sittlichkeit seiner Dienstboten zu richten. Er belehre sie, daß ein gewisser Grad von Mißtrauen in jedem Menschenherzen und in seiner natürlichen Beschaffenheit Grund und Berechtigung genug finde, und erbitte sich’s von ihnen, nicht vorneherein trauen zu müssen, sondern im Gegentheil fragen und forschen zu dürfen. Das Dorfgehen, Besuchen der Märkte, der Tanzplätze und anderer
\quellseite{236}
Gelegenheiten, welche leicht zur Sünde verleiten, werden von vornherein ein für alle Mal verboten und der erste Dienstbote, welcher sich am hausväterlichen Verbote versündigt, unnachsichtlich entfernt. Strenge gleich von Anfang gewöhnt die Menschen. So weiß dann jeder neu eintretende Dienstbote gleich vornherein, was er zu erwarten hat; wer nicht Lust hat zu einem ordentlichen Leben, sieht sich vor und läßt sich mit einem Dienste nicht ein, der ihm nicht behagen würde. Je länger je mehr wird sich dann für den Pfarrer gar kein Dienstbote finden, der nicht Lust hätte am Worte Gottes und an einem stillen Familienleben. Es braucht dann keiner hervortretenden Strenge mehr: Liebe und Friede herrschen immer stiller und zugleich durchgreifender. Bei etwa doch nothwendig werdendem Wechsel heißt es dann: Similis simili gaudet („Gleich und gleich gesellt sich gern.“ — sagen wir Deutsche).

122. So wie unsre Gemeinden sind, fürchten sie von allen, die im Pfarrhause ab- und zugehen, Verleumdung, Verrath und Zuträgerei. Die meisten Gemeindeglieder haben beladene Gewissen und fühlen Unruhe, so wie nur der Gedanke kommt, es möchte jemand mit dem Pfarrer von ihnen reden. Es könnte ja, was sonst jedermann weiß, auch vor die Ohren des Mannes kommen, deß Wort und Auge man mehr als Gerechtigkeit und Gericht fürchtet. Gienge es, man würde das Pfarrhaus am liebsten durch Zaun und Mauer gegen allen andern als den nothwendigen amtlichen Verkehr abschließen, — man würde jeden, der nicht eine Taufe, Proklamation, Trauung oder Leiche anzusagen hat und im Sonntagsstaat, in Hut und
\quellseite{237}
Rock, zum Pfarrer geht, von Thor und Mauer abweisen laßen. Das geht nun nicht; daher ist das Mistrauen desto reger — und so oft der Pfarrer Treffendes öffentlich oder zum Einzelnen spricht, so oft weicht er vom Text, in dem ja von keinem jetzt lebenden Menschen steht, und mischt sich in Dinge, die ihn nichts angehen; es kann dann gar nicht anders sein, jemand muß ihm zugetragen haben, er muß Klatscherei und Verleumdung annehmen; denn Verleumdung nennt man auch eine ganz wahre Nachricht, so wie sie nur von fremdem Mund mitgetheilt wird und eine Sünde aufdeckt. Wer einige Erfahrung hat, der weiß auch, daß wir mit diesen voranstehenden Sätzen uns keiner Uebertreibung schuldig gemacht haben. Was könnte man hier alles klagen und erzählen, wenn etwas dran läge. Unter diesen Umständen wird man es ganz natürlich finden, daß der Verdacht der Zuträgerei besonders auf die Dienstboten im Pfarrhaus fällt. Einen Weg, sich selbst und seine Dienstboten völlig von jedem Vorwurf der Zuträgerei und Klatscherei frei zu erhalten, gibt es wohl nicht, so lange Gottes Wort trifft und der Pfarrer kein stummer Hund ist. Auch das ist rein unmöglich, jede Nachricht aus der Gemeinde vom Pfarrhaus fern zu halten, und ganz unstatthaft ist es, die Dienstboten, die ja hin- und hergehen und nothwendig dies und das hören müssen, zu völligem Schweigen zu verurtheilen. Das Einfache ist, auf brave Dienstboten zu sehen, sie mit den Vorwürfen, die einem Pfarrer wegen Klatscherei so gerne gemacht werden, bekannt zu machen und sie zu ermahnen, daß sie mit ihrem Hören und Reden nicht dem Amt hinderlich
\quellseite{238}
werden, daß sie langsam zum Reden und langsam zum Zorn seien, daß sie ihre Zunge schweigen, unnützes nicht sprechen, weder außer dem Hause, noch in demselben, — daß sie bei allem, was sie reden, zuvor Wahrheit, dann Barmherzigkeit und Mildigkeit sich regieren lassen, und es für eitel Unglück achten, wenn durch ihre Schuld ein Knecht Gottes zum Lotterbuben und Wäscher und sein väterliches Wort zur Klatscherei umgestempelt würde. Außerdem ahnde ein Pfarrer jedes Geschwätz seiner Dienstboten härter als anderes und predige den Seinen allen das achte Gebot mit Nachdruck und eingehender Sorgfalt, — und er selbst trachte mit unermüdlichem Fleiße nach dem, wenn schon nie völlig zu erreichenden Ruhm eines „vollkommenen“ Mannes, der im Wort nicht fehlt, die eigene Zunge im Zaume hat und alles, so lang es möglich ist, zum Besten kehrt. Er selbst und seine Tugend, sein aufrichtiges und edles Benehmen unter den elenden Vorwürfen seiner Pfarrkinder, die sich überall verrathen sehen, — das sind bei Weitem die Hauptstücke, auf welche es ankommt. Ein frommer, weiser Hirte weiß und kann die Kunst, sich mit Wort und That wenigstens vor verständigen und wohlwollenden Leuten von dem Vorwurf der Annahme von Klatschereien je länger je mehr zu befreien und frei zu erhalten.

\quellseite{239}
\section{XV. Bewirthschaftung der Pfarrgüter}

123. Bei Landpfarrstellen kommt es häufiger als bei städtischen Pfarrstellen vor, daß Zehnten und eigene Grundstücke zum Pfarrwitthum gehören und aus ihnen das Einkommen des Pfarrers gewonnen werden muß. Da handelt es sich nun um die rechte Art und Weise, aus Zehnten\footnote{Wie alt ist der Zehnten, und wie ganz werth einer Apologie ist er! Nichts Schöneres als freiwillige Zehnten! Während man bei uns unberechtigte Zehnten abschafft, kommt man hier und da, wo neue Kirchenbildungen entstehen, auf den Gedanken, sich zur Dotation der Pfarreien durch Zehnten und Grundbesitz frei zu vereinigen. Es ist für einen bayrischen Pfarrer zu spät, den Zehnten eine Vertheidigung aus göttlichen und menschlichen Rechten, aus Geschichte und Erfahrung zu schreiben, weil es doch nur ein Abschied wäre. Aber wahrlich, ganz abgesehen vom Nutzen, es ist doch Schade, daß durch Abschaffung der Zehnten ein aus göttlichem Gebot und dem Glauben hervorgegangenes uraltes edles Werk der Liebe, eines der allerschönsten, so roh und unverständig mit Füßen getreten wird!} und Grundstücken das Einkommen zu gewinnen. Entweder werden die Zehnten und Grundstücke verpachtet oder aber von den Pfarrern selbst bewirthschaftet und eingeheimst. Welche von beiden Arten nicht bloß die vortheilhafteste, sondern vor allem die einem Pfarrer geziemendere und seinem Amte zuträglichere sei, das ist je und je ein Gegenstand vielen Streites gewesen und wird es auch immer bleiben.

\quellseite{240}
124. Bequemer ist offenbar das Verpachten der Zehnten und Grundstücke. Da ist man aller Mühe los, wenn das Pachtprotokoll geschlossen ist; so scheint es wenigstens. Da man sich mit Landwirtschaft auf diese Weise nicht abzugeben braucht, so scheint auch alle Gefahr des Verderbens abgethan. Dagegen ist offenbar Selbsteinheimung des Zehntens mühevoller; eben so die eigene Bewirthschaftung des Pfarrwitthums. Der Sorge und Mühseligkeit ist so viel, daß die Gefahr, von der Höhe des geistlichen Lebens ins Irdische heruntergezogen zu werden, wirklich nahe liegt und groß ist. Hat man vollends, wie es gewöhnlich zu sein pflegt, keine treuen Leute, so hat man zur Sorge und Seelengefahr großen zeitlichen Schaden und muß in mancherlei Streit und Aerger leben. Das alles ist vermieden, wenn man verpachtet.

125. Die Bequemlichkeit, welche man beim Verpachten findet, scheint nicht bloß zum Verpachten zu rathen; sondern zu noch mehr, nämlich die Pfarrgüter auf möglichst lange Zeit zu verpachten, auf Dienstzeit etwa, und den Zehnten gar zu fixieren. Auf diese Weise hat man neben den übrigen Vortheilen auch den, für lange Zeit und beim Zehnten für immer das oft ärgerliche Geschäft des Verpachtens sich vom Halse zu schaffen.

126. Nun ist aber ein für alle Male gewiß, daß man durch Fixieren des Zehntens diejenigen beschränkt, welche über dergleichen Dinge zu reden eine eben so große Berechtigung haben, als es unmöglich ist, sie zu fragen. Wir meinen die Nachfolger. Wie bitter wird oft Jahrhunderte
\quellseite{241}
von diesen geklagt, wenn der Vorfahr irgend eine Einrichtung von schädlicher Nachwirkung gemacht hat. Wer das erfährt und erkennt, wird sich zum Fixieren nur gezwungen\footnote{Diese bereits mehrere Jahre vor 1848 geschriebenen Sätze sind für bayerische Pfarrer von keiner Wichtigkeit mehr, denn wir sind ja gezwungen worden, zu fixieren. Doch wollte man nichts ändern, weil hie und da, wo dies Büchlein vielleicht gelesen wird, noch kein Zwang stattgefunden hat.} und da auch nicht ohne Protest anschicken. — Dazu ist man, wenn man sein Recht auf den Zehnten in natura behält, von der Hand Gottes abhängig, durch welche uns fruchtbare und unfruchtbare Jahre zukommen. Diese göttliche Unsicherheit ist aber besser, als eine menschliche Gewißheit und Sicherheit. Denn so gewiß auch fixiert werde, allgemeine Unglücksfälle betreffen doch die Pfarrer härter, wenn sie ihr Einkommen aus dem Säckel anderer Unglücklicher nehmen müssen, als wenn es von Grund und Boden genommen wird. Kriege, Mißwachs, Wetterschläge etc. haben freilich keinen Einfluß auf fixierte Zehnten, so lange die Leute nicht zahlungsunfähig werden. Aber es kann dafür auch kommen, daß am Geld der Blutschweiß des armen Volkes hängt. Ist nicht fixiert, so hat der Pfarrer immer in gleichem Maße mit dem Volke, weint mit den Weinenden, freut sich mit den Fröhlichen, trägt alles, was Gott gibt, seis Glück, seis Unglück, mit der Gemeinde. Ist aber fixiert, so lassen sich Fälle genug denken, in welchen die Schlechtigkeit der Menschen den sorglos
\quellseite{242}
stehenden Pfarrern und ihrem Einkommen die Berechtigung abstreitet, es schmälert oder vernichtet. Denn warum sollen Rechte der Pfarrer weniger dem Widerstand und der Ungerechtigkeit der Menschen unterliegen, als Throne? Man könnte sagen: „Wenn es einmal so kommt, so ist auch das Recht auf Grund und Boden nicht sicher.“ Allein sicherer ist und bleibt es eben doch, als das Recht auf die Säcke der Gemeinden.

127. Was die Verpachtung des Zehntens und der Grundstücke anlangt, so ist allerdings für den keine andere Hilfe, welcher keine treuen Leute, kein Talent für den Landbau hat, den Sorgen oder der Last der Wirtschaft zu leicht erliegt und an sich spürt, daß jeden Falls sein Amt Schaden litte, wenn er sich selbst mit jenen Sachen befassen wollte. Es ist nun einmal so, daß nicht jeder gleich leicht und glücklich dieselbige Sache beschickt. Auch ist es in Wahrheit eine Seltenheit, daß ein Pfarrer treue und anhängliche Leute findet. Denn die Dienstboten, auch wenn sie einmal ihren Pfarrern treulich anhingen, reichen nicht aus, namentlich beim Zehnten: es werden mehr und andere Leute nöthig. Da finden sich denn gerne die Armen, die Nothigen ein, die sich Betrug verzeihen.

128. Jeden Falls ist es eine mißliche Sache, wenn der Pfarrer selbst oder die Pfarrerin dem Zehnten nachgehen müssen. Wie die Sachen stehen, geben die Leute nun einmal den Zehnten nicht gern; jeder vergißt, daß er schon bei Uebernahme seines Gutes des Pfarrers Gut mit übernommen hat; jeder thut, als wäre der zehntende
\quellseite{243}
Pfarrer nicht berechtigt. Alle mögliche List wird oft angewendet, um auch nur weniges zurückbehalten zu können. Wenn dann der Pfarrer oder die Pfarrerin selber beständig auf Diebereien und Betrug lauern müssen, so erwacht in den Herzen der Menschen eine Art von Haß und Ungunst, vermöge welcher man den nicht gern auf der Kanzel und am Altare sieht, der zu genau die einzelnen Gebrechen seiner Pfarrkinder mit Augen geschaut hat und sie zu schauen beflissen war. Auch erweckt das Mein und Dein in einem Pfarrer selbst, oft bei Gelegenheit kleiner Dinge, eine unwürdige Aufregung, die ihn zu Sünden verleitet. — Muß drum der Pfarrer selbst Hand ans Werk legen, so ist es für ihn allerdings gefährlich. Doch gibt es Charaktere, die auch in solchen Versuchungen auf das würdigste stehen können.

129. Findet ein Pfarrer zur Bewirthschaftung seiner Zehnten und Güter treue Leute, so scheint es allewege das Beste zu seyn, wenn er nicht verpachtet. Denn:

a) Vieles, was beim Verpachten nicht in seiner Macht steht, ist da in seine Hand gegeben. Der Eid, zur Aufnahme des Pfarrguts sein Möglichstes zu thun, wird so am redlichsten erfüllt. Die Güter können am besten bewirthschaftet, gehoben — und so dem Nachfolger das möglichst reichliche Einkommen überliefert werden. Dagegen ist das Verpachten den Gütern ingemein schädlich. Nicht bloß zieht jeder aus dem gepachteten Felde das Mögliche bei geringstem Aufwand; nicht bloß werden die Güter gegen Ende längerer Pachtzeiten auffallend vernachlässigt
\quellseite{244}
; nicht bloß kommen sie so immer mehr herunter und werfen je länger, je weniger Pachtschilling ab; sondern es werden bei dergleichen Gelegenheit manchmal die Grenzen geändert; oft wird das Gut des Nachbars, der gerne Pächter von Pfarrgrundstücken ist, unbillig vergrößert — und durch alles das ein Pfarrgut von bedeutendem Belang herunter gebracht. Ganze Stücke Landes wurden zuweilen auf diese Weise entwendet. Eben so ist es mit dem Zehntlande. Der Pächter läßt irgendwie Unordnung einschleichen, zehntet da und dort nicht, Jahre lang nicht, und hilft so für spätere Zeiten die Facta herbeischaffen, aus denen hernachmals bewiesen wird, daß einem Pfarrer ein Recht auf dies oder jenes Grundstück nicht zustehe, das ihm doch zusteht. — Alle Gefahren dieser Art fallen weg, wenn treue Leute unter des Pfarrers eigener Aufsicht das Werk thun.

b) Die Einnahme ist so die dem Ertrage entsprechendste, ehrlichste, und in vielen Fällen überdies die reichlichste. Beim Verpachten läuft viel List unter, viel schlaue Verabredung der Pachtlustigen. Andere — und zwar oftmals nicht die Armen, sondern reiche Pächter nehmen den ehrlichen Vortheil des Pfarrers lachend hin, versündigen sich und lassen ihm den Schaden. Anders bei eigener Bewirthschaftung. Wächst viel, so hat der Pfarrer, was ihm gebührt. Wächst wenig, so hat er auch wenig; er wird dann nicht mit Bitten um Nachlaß und Aufschub belästigt, sondern zu Gottes Stunde nimmt er aus Seiner Hand sein bescheidenes Theil. Ist auch im
\quellseite{245}
Pachtprotokoll ausdrücklich gesagt, daß kein Nachlaß, noch Aufschub gegeben werde, so entschuldigt ein solches Protokoll doch nur vor menschlichen Gerichten, wenn er in Unglücksfällen auf sein Eigenthum sieht. Sein Gewissen aber nöthigt ihn doch zur Barmherzigkeit. Und zwar muß er alsdann in einem Maße barmherzig sein, daß man ihm nicht Geiz nachredet, — und das geht dann armen Pfarrern schwer, zumal das Landvolk oftmals gegen Pfarrer sehr frech und unverschämt ist. Unter Schimpf und Schande seinerseits wird er dann doch oft noch betrogen und erleidet einen Schaden, der für seine Umstände sehr empfindlich ist. — Das Alles ist vermieden, wenn man selbst einheimsen kann.

c) Wenn der Pfarrer selbst einheimsen kann, lernt er nicht allein seine Pfarrkinder in Verhältnissen kennen, wo ihre Sünde und Tugend am öftesten und klarsten heraustritt, sondern er selbst kann hier die Pflicht der Gerechtigkeit und Billigkeit am schönsten üben, im schlimmsten Verhältniß das beste Beispiel geben. Je mehr er die Sünde seiner Pfarrkinder erkennt, je untadeliger er ihnen in gleichen Verhältnissen gegenübersteht, desto kräftiger vermag er zur Erkenntniß der Sünde zu führen und zur Buße zu rufen. — Es ist eine gefährliche Stellung für einen unlautern Pfarrer; aber ein Triumph des rechten Pfarrers.

d) Wenn er selbst einheimst und wirtschaftet, lernt er nicht allein die Verhältnisse des Landmanns am besten aus eigener Erfahrung kennen; sondern er kann auch
\quellseite{246}
erkennen, wo des Landmanns größte Nöthen sind.

e) Eben die eigene Wirtschaft setzt auch den Pfarrer in den Stand, manchen Armen zu versorgen — nämlich durch Arbeit, was am Ende doch die beste Art ist.

f) Es tritt ein Pfarrer in seiner patriarchalischen Würde nicht und niemals schöner hervor, als wenn er selber wirtschaftet. Die Beziehungen des Hausvaters zu seiner Familie, zu seinen Dienstboten, Taglöhnern und Arbeitern, je reichlicher sie sind, um so mehr Gelegenheit geben sie dem Pfarrer, unter seinem Volke recht geistlich schön zu stehen. Sein Ansehen wächst, sein Urtheil wird geachteter. Ein wahrhaft geistlicher Mann, der in den rein menschlichen Verhältnissen unsträflich wandelt, erweckt auch für sein Amt ein um so größeres Vertrauen. Wenn der Pfarrer außerhalb ländlicher Beziehungen lebt, heißt es oft: „Er würde anders reden, wenn er unsere Lage besser kennete.“ Kennt er sie nun und bleibt doch dem heiligen Wort getreu, so fällt sein Wort mit einem menschlichen Gewichte mehr ins Herz.

Was das Verbauern anlangt, so wollen wir nicht in Abrede stellen, daß mancher Pfarrer der Versuchung nicht gewachsen ist. Aber viele edle Männer haben sich herrlich bewährt! Viele haben ausgezeichnete Amtsthätigkeit entwickelt, unablässig sich selber fortgebildet, ihre eigenen Kinder unterrichtet und dazu ihr Pfarrgut bewirthschaftet. Wer will diesen die Palme verweigern? — Wollte Gott, es gäbe mehr Pfarrer dieser Art! Aber freilich,
\quellseite{247}
schon das immerwährende Trachten nach neuen Pfarreien verhindert ein Einwurzeln in der Gemeinde, wie es ein solches Thun zu Wege bringt und verlangt! Es fragt sich dann freilich, wo mehr Mietlingssinn ist, bei dem, der nie die rechte Gemeinde findet, oder bei dem, der, unter Versuchung zwar, bei Einer bleibt und lebt mit treuer Liebe, und sein Auskommen durch bessere Benutzung seines Bodens sucht?

\quellseite{248}

\chapter{Von der Ehe des Geistlichen}

\quellseite{249}

\quellseite{250}
\section{I. Cölibat oder Ehe?}

130. Daß der Cölibat den protestantischen Geistlichen nie verboten worden ist, weiß jedermann; es bedarf der Erwähnung nicht. Wie könnte man auch die Ehelosigkeit oder, mit den Alten zu reden, das „jungfräuliche Leben“ denen verbieten, welche dem Volke in allerlei Lob und Tugend voranzugehen berufen sind, zumal sich mehr als eine Stelle des göttlichen Wortes findet, in welchen dies Leben mit Lob und Ehre geschmückt wird. Matth. 19, 11. 12. und 1. Cor. 7, 26 ff. können nicht überhört werden und könnten wohl das Bedenken erregen, ob nicht das jungfräuliche Leben in der protestantischen Kirche gar zu gleichgültig angesehen zu werden pflege. Die Uebertreibungen der Alten und das gleichgültige Schweigen der Protestanten sind beides Extreme, innerhalb welcher Christi und St. Pauli wohlbegrenztes (s. 1. Cor. 7, 4. 5. 9.) Urtheil liegt. Lassen wir gelten, was Gott der Herr gelten läßt, lobt und ziert, und was Seine alleinige Gabe genannt werden muß. Es gibt durch Gottes Gnade einen
\quellseite{251}
reinen Cölibat, welcher in Zeiten der Verfolgung oder bei besonders schwierigen Aufgaben des Amtslebens sehr förderlich sein kann. Der Ehelose ist beweglicher, wagt und besteht Gefahren leichter, da er sie ja nur für sich besteht und durch seinen zeitlichen Tod und Hingang weiter niemand großen Verlust oder Schaden erleidet. Die Kühnheit der römisch-katholischen Heidenmissionen kommt gewiß mit auf Rechnung des Cölibats und kann zu 1. Cor. 7. manchen Beleg liefern.\footnote{Wer z.\,B. die römischen Missionsstationen in Nordamerika betrachtet, der findet, daß die römische Heidenmission der Colonisation noch voraneilt und sogar die Rührigkeit der Methodisten übertrifft, deren Grundsätze doch sehr zum Vorwärts blasen. Wo wir nicht hin dachten, steht hie und da schon das Kreuz der römisch-katholischen Missionäre.} Möchten nur auch alle jungen noch unverheiratheten protestantischen Geistlichen, welche im Stande der Verfolgung oder unter schweren Aufgaben des amtlichen Lebens seufzen, aus 1. Cor. 7. und dem Beispiel der Römischen lernen, was zu lernen ist, ihre Gabe erforschen, und falls sie die Gabe, ehelos zu leben, in sich finden, sich gesagt sein lassen, was z.\,B. Sarcerius in seinem trefflichen „Hirtenbuche“ (S. 137) sagt: „Wer die hohe Gabe hat, also keusch zu bleiben, der brauche derselben; denn die Zeit ist kurz, spricht der Apostel, und das Wesen dieser Welt vergeht.“ (Παράγει γὰρ τὸ σχῆμα τοῦ κόσμου τούτου. 1. Cor. 7, 31.)

131. So wie aber der Cölibat den Geistlichen nicht verboten ist, so ist er ihnen auch nicht geboten. Unserer Kirche ist je und je der römische Begriff des Cölibats und
\quellseite{252}
die unreine Ansicht von der Ehe, welche sich schon früh im Alterthum findet, fremd und ferne geblieben; sie hat auf Grund von apostolischem Wort und Beispiel, Matth. 8, 14., Gesch. 21, 8. 9., 1. Cor. 9, 5 ff.,\footnote{S. Baumgartens Kasuist. Pastoraltheologie S. 224.} 1. Tim. 3, 2., zeug mancher Kirchenordnungen des sechzehnten Jahrhunderts die Priesterehe empfohlen, wenngleich sie nicht mit der griechischen Kirche in Mißdeutung und Uebertreibung der apostolischen Stelle der Pfarrgeistlichkeit die Ehe zum Gebote macht. (Vgl. Guerike’s allgem. Christliche Symbolik. Leipzig 1839. p. 533 f.) Gefährlichen Irrthümern vorzubeugen, der Ueberschätzung des Cölibats zu begegnen, welcher aus der Geschichte so vieler Jahrhunderte des Lobes leider wenig genug davon gebracht hatte, durch die That Zeugniß gegen unreinen Cölibat abzulegen, wurde bei uns die Priesterehe gepriesen und empfohlen. Hatte man doch in den Stellen 1. Tim. 3, 2. Tit. 1, 6. dazu festen und guten Grund (vgl. auch 1. Tim. 4, 3.); enthält doch weder das alte, noch das neue Testament etwas, daraus man nachweisen könnte, daß die Priesterehe bei Gott in üblem Gerüchte stehe! Gab es doch in den ersten Jahrhunderten auch Zeugen von großem Namen, welche der falsch-ascetischen Auslegung oder vielmehr Verfälschung jener apostolischen Stellen die Wahrheit entgegenstellten! Ist doch die Zwiespältigkeit des Abend- und Morgenlandes Aufforderung genug gewesen, die heilige Schrift und ihren einfachen Wortlaut zwischen beiden entscheiden zu lassen! Und sind doch die Gegner genöthigt,
\quellseite{253}
den Ruhm des Cölibats zum Theil von heidnischen und muhamedanischen Gewährsmännern zu borgen!\footnote{Wohl möchte hier auf ein sehr reichhaltiges, im protestantischen Sinne geschriebenes, wenn gleich vergessenes Buch hinzuweisen sein, nemlich auf: Georgii Calixti De Conjugio Clericorum Liber. etc. Emendatius edidit Herr. Phil. Conr. Henke. Helmstädt 1783. P. I. S. daselbst namentlich P. I. Cap. II. Utrum coelibatus clericorum sit apostolicae institutionis? p. 9 ff.}

132. Wenn denn Ehelosigkeit und Ehe in der heiligen Schrift allen Menschen, und darum auch den Dienern der Kirche frei gegeben ist; so hat ein jeder Geistlicher für sich selbst zu entscheiden, ob jene oder diese nach seiner Beschaffenheit ihm und seinem Wirken förderlich sei.\footnote{„Die Ehe ist durch die Reformation zu Ehren gekommen und jener schmutzige Cölibat der Mönche und Nonnen ist zu Schanden geworden und in’s Abwesen gekommen. Aber es hat sich auch etwas breit gemacht, was die Reformatoren nicht wollten, nemlich ein zu geringschätziges, fast indifferentes Urtheil über den ehelosen Stand oder die Jungfrauschaft. Und doch werden beide Stände nur durch einander gehalten, und wenn der eine beschimpft wird, geschieht es auch dem andern. In der heiligen Schrift, namentlich 1. Cor. 7., finden wir beide Stände geehrt. Es ist richtig, daß St. Paulus zunächst um der schweren, Verfolgungs- und trübselvollen Zeit willen, in welcher er lebte, den ehelosen Stand anrieth; aber er nennt es doch ein καλόν (etwas Edles und Wohlgethanes), so zu bleiben, und hinter den Ehestand zurück kann darum der jungfräuliche Stand nicht treten. Beide sollen in Würden bleiben.\par Der Ehestand soll in Würden bleiben. Der zeitliche Zweck desselben, Kinder zu erzeugen und durch die von Gott verordnete leibliche Vereinigung das Herz vor Versuchung der Unzucht zu behüten, hindert ausgesprochener Maßen (1. Tim. 2, 15., 1. Cor. 7, 1 ff.) den Hauptzweck des Lebens, das Streben nach Seligkeit und Heiligkeit nicht. Vielmehr heiligt das gläubige Weib oder der gläubige Mann selbst in gemischter Ehe die Kinder (1. Cor. 7, 14.), zum Beweis daß die Ehe selbst müsse heilig sein können.\par Es ist deswegen die Ehe ein Stand des guten Gewissens, wer es nur bewahren will. Auch ist die Heiligkeit der Ehe etwas so Gewisses, daß sie sogar als unauflöslich von Seiten des Christen Matth. 19, 1 ff. und 1. Cor. 7. dargestellt wird. Nur durch die Schuld des ungläubigen Theils kann eine Ehe aufgelöst und getrennt werden, während der Gläubige, es gefalle ihm die Verbindung selbst mit dem Ungläubigen oder sie gefalle ihm nicht, unter keiner Bedingung den Anlaß zur Scheidung geben, und ihn nur in dem einzigen Fall der Hurerei nehmen darf. Denn im Falle von 1. Cor. 7, 15. nimmt er ihn nicht, sondern er wird ihm aufgedrungen. Auch das ist ein unüberdenklicher Beweis von der Heiligkeit der Ehe, daß sie Eph. 5, 22—33. als das Vorbild der allerheiligsten Vereinigung Gottes und seiner Kirche dargestellt wird. Die leibliche Vereinigung Mannes und Weibes in der heiligen Ehe ist ein von Gott geschaffenes Bild der Einigung Christi mit seiner Kirche, wir wiederholen es: sie stellt, selbst geheimnißvoll, das große Geheimniß der Wiedervereinigung der gefallenen Menschheit mit Gott in Christo Jesu dar. Im Manne wird der HErr, im Weibe die heilige Gemeinde dargestellt, wie sie eins werden, — und nach Maßgabe der Vereinigung Christi und seiner Kirche fließt dem Manne und dem Weibe, einem jeden sein besonderer Auftrag zu. 1. Cor. 11, 3 ff., 1. Tim. 2, 11 ff., 1. Petr. 3, 1. ff. — Da es nun mit der heiligen Ehe so steht, so ist sie von jedem Christen hoch zu ehren, und wehe dem, der sie mit leichtfertigem, unehrerbietigem Auge ansieht.\par Aus der rechten Betrachtung der Ehe ergibt sich:\par a) daß sie ein Gemeingut aller Menschen von Adam her ist, daß ein erzwungener Cölibat, sei er durch kirchliche oder weltliche Gesetze geboten (letzteres ist bei dem erzwungenen Cölibat der Armen der Fall), verwerflich, eine Correction göttlicher Ordnung, ein Beweis von Unglauben und Mangel an Vertrauen in Gottes väterliche Treue ist; b) daß sie ehrlich, unbefleckt, unauflöslich sein soll, daß Ehescheidung ein Uebel und innerhalb der heiligen Kirche, wenn beide Gatten rechte Christen sind, nicht denkbar ist; c) daß sie christlich ist, der Gottseligkeit nicht hinderlich, wofern sie nur christlich ist und so, wie Gottes Wort einem jeden ehelichen Menschen Befehl und Anweisung gibt.\par Daneben lehrt aber auch St. Pauli Wort und Beispiel, daß der ehelose Stand, wie bereits gesagt, etwas Edles und Wohlgefälliges (ein καλόν) und unter Umständen, besonders in Zeiten der Verfolgung im Vergleich mit dem Ehestande etwas Vorzüglicheres (ein περισσόν) sei. Wer sich dem Reiche Gottes, namentlich in böser Zeit, ungehindert widmen und alleine sorgen will, was des Herrn ist, wer den Geschäften des geistlichen Amtes in besonders schwierigen Verhältnissen obliegen will, ohne von Sorgen für Weib und Kind, von des Ehestandes Weh und Noth immer aufs Neue abgezogen zu werden, den verweist St. Paulus zum ehelosen Stande, welchen der Herr selbst in diesem Fall ein Verschnittensein um des Himmelreichs willen heißt. Das ist einfach und jedem verständlich.\par Es ist aber das eine, wie das andere von nicht geringer Wichtigkeit für alle Christen, sowohl zur Beurtheilung eigener, als fremder Fälle. Einerseits kann die rechte Lehre von der Ehe das gute Gewissen derjenigen stärken, die das ehrliche Verlangen tragen, in die Ehe zu treten, — auch jeden Christen bewegen, die Verehelichung junger Christen, welche selbst also wollen, zu fördern. Andererseits werden diejenigen, welche Hindernisse der Verehelichung finden, die sie nicht überwinden können, durch die Lehre von der Jungfräulichkeit angeleitet werden, ihren Beruf zur Ehelosigkeit zu erkennen. Wäre dieser göttliche Beruf zum jungfräulichen Stande in unsrer Kirche nicht so gar vergessen und in den Hintergrund getreten, wäre die Herrlichkeit und der Segen des jungfräulichen Standes mehr erkannt und zur Anerkennung gebracht worden; so würde dies mehr, als alles andere, haben beitragen können, die Jugend vor Sünden zu bewahren. Aus der heiligen Jungfrauschaft kann man jeden Falls die besten Gründe nehmen, um denen zu Hülfe zu kommen, die rücksichtlich ihrer Keuschheit angefochten sind. Das Gebot der Enthaltsamkeit wirkt weniger, als die Lehre von der Gottverlobtheit und Hingebung in Christi Dienst. Man lehre, lobe, vermahne und übe! Ein heiliger Gedanke wird, wenn Gott ihn schenkt, zu einer Macht über das Fleisch, und die Anfechtung zerrinnt, wenn der Geist Gottes Fügung und einen himmlischen Beruf erkennt. Wir reden jedoch von Christen, die wider ihr eigenes Fleisch Hilfe suchen, nicht von der Menge, die das Himmlische verhöhnt und nur begreift, was irdisch ist. — Der züchtigenden Liebe sind durch die vorgetragenen Sätze gewiß nicht geringe Mittel an die Hand gegeben, fromme Seelen zur Heiligung zu führen.“ (Vorschlag zur Vereinigung lutherischer Christen für apostolisches Leben §. 33 p. 70–75.)} Es
\quellseite{254}
ist vorauszusehen, daß nach richtiger Selbsterkenntniß die Mehrzahl zur Ehe wird greifen müssen. Ohne Anfechtung des Fleisches ehelos zu leben, ist eine besondere Gabe Gottes, und zwar eine ziemlich seltene, so gewiß es auch sein
\quellseite{255}
mag, daß mancher die vorhandene Gabe muthwilliger Weise verkennt oder verleugnet. Die Gabe ist öfter da, als es scheint. Muthwillige Neugier hindert oft, sie zu sehen. Viel eheliches Unglück solcher, die bisher ehelos geblieben
\quellseite{256}
wären ist Strafe des muthwilligen Eintritts in die Ehe. Ist nun Anfechtung vorhanden und mangelt die edle Gabe; so ist jedenfalls die Ehe das von Gott verordnete Heilmittel für Leib und Seele. Es ist viel besser, freien und so der Anfechtung des Fleisches los zu werden, als immerzu Klagen zu führen, wie Antonius, Hilarion, Pachomius und Andere. Wer deshalb außer den Zeiten der Verfolgung und anderer schwerer Arbeitslast in sich die Weisung und Mahnung zur heiligen Ehe findet und ihr nicht folgt, begibt sich in Seelengefahren, welchen zu unterliegen die schwere Verantwortung derjenigen nach sich zieht, die ihres Herrn Willen wissen und nicht thun.

133. Die innerlichen Anfechtungen und äußerlichen Versuchungen, welche ein mit der Gabe der Enthaltsamkeit nicht beschenkter Geistlicher zu überwinden hat, erweisen sich
\quellseite{257}
dem Amte viel hinderlicher und nachtheiliger, lähmen auch die Kraft des Mannes und berauben ihn der ungetheilten Aufmerksamkeit auf seinen heiligen Beruf viel mehr, als alle Nachtheile, welche man der Priesterehe nachsagt. Dagegen lernt man in der Ehe und durch sie nicht bloß das eheliche Leben und Wesen, sondern auch das eigene Herz, die eigene Kraft und Schwachheit von mancher Seite kennen, welche dem nie verheiratheten Geistlichen ein verhülltes Geheimniß bleibt. Die Seelsorge der Eheleute bewegt sich großen Theils auf einem Lebensgebiete, welches der am besten verstehen sollte, welcher es aus eigener Erfahrung kennt. Der Unverheirathete oder bisher nie verehelicht Gewesene kann hier kein Rathgeber sein wie der Erfahrene; dieser wird zuweilen mit wenig Worten Rath geben und mit Einfalt die vielen angefochtenen, namentlich jungen Eheleute bescheiden können, ohne daß er den heiligen Euphemismus und edlen Sprachgebrauch verletzt, welcher über das eheliche Geheimniß unter Zustimmung aller frommen Eheleute verbreitet ist. Unverheirathete Seelsorger werden dem weiblichen Theile ihrer Pfarr- oder Beichtkinder oft verlegen gegenüberstehen, werden auch selbst oft von jenem krankhaften und empfindsamen Süchteln kaum frei werden, das sehr vielen anhängt, die nie ein eigen Weib gehabt haben; der verheirathete Seelsorger dagegen sollte sich seinen weiblichen Pflegebefohlenen gegenüber unschuldiger, freier, zuversichtlicher fühlen; was gehen den treuen Ehemann alle fremden Frauenpersonen an?

134. Und nicht das allein; nicht bloß von Seite der weiblichen Seelsorge empfiehlt sich die Ehe; sondern
\quellseite{258}
— die Ehe setzt auch den Geistlichen erst in den Stand, nach allen Seiten hin das Beispiel zu geben, welches ihm geziemt. Die Gnade der Enthaltsamung vorausgesetzt, ist es ein Leichtes, allein durchs Leben zu wandern; da ist man von Adams Plage, von Kummer und Schweiß des Angesichts so ziemlich frei. Aber eine ganz andere Kraft und Höhe der Gesinnung verlangt es, selbander und unter einer heranwachsenden Familie zu leben. Da muß man sich in eine zweite Seele, in alle ihre Eigenheiten und Besonderheiten, Tugenden und Schwachheiten mit zarter Liebe finden; sich bessern lassen und alles tragen zum Zweck der Besserung der Liebsten. Da muß man unter den Wechselfällen der Ehe sich in immer gleicher Amtstüchtigkeit bewähren, in Tagen des Glücks und des Unglücks, an Kranken- und Sterbebetten der eigenen Familie, in Armuth und Fülle — ach, in wie vielen und verschiedenen Lagen gilt es da, immer Einer zu sein, des Herrn getreuer Knecht! So angesehen steht ein unverheiratheter Pfarrer hinter dem verheiratheten, wenn sonst auch alles gleich ist, doch merklich in der Aufgabe zurück. Weit entfernt, daß jener an Kraft und Tugend reicher wäre, ist vielmehr dieser in einer viel ernsteren und eingreifenderen Schule der Heiligung, — und welch ein heiliger Bekenner ist er, wenn er seine Aufgabe faßt und löst!

135. Schon aus der vorigen Nummer und dort gethaner beiläufiger Aeußerung läßt sich die Frage beantworten, ob das Leben im Cölibat oder in der Ehe das bequemere sei? — Der Geistliche hat ein Leben, welches sich durch Vielseitigkeit der Beziehungen auszeichnet und
\quellseite{259}
durch eine große Anzahl der verschiedensten Ansprüche von außen her. Es ist — oft bei scheinbarer Ruhe — eine große Unruhe vorhanden, eine reichliche Lebensplage. Da bedarf er auch mehr als manch anderer Mann eine Aufrichtung, Hilfe und Pflege. Auch ist dem kranken, müden, alten und sterbenden Pilger wohl eine treue Hand zu gönnen, die ihm Erquickung und Arznei reicht, den Schweiß vom Angesicht trocknet und das Auge schließt. Es ist so doch besser, als wenn Hieronymus’ Rath befolgt wird:

\begin{quote}
„Aegrotanti tibi quilibet sanctus frater assistat, et germana vel mater, aut probatae quaelibet omnes fidei.“
\end{quote}

(Ad Nepot.) Niemand schicklicher, niemand unanstößiger dient einem Diener Gottes, als die Gehilfin, die Gott schuf, daß sie um ihn sei, ein frommes Weib! Rechnet man aber diese Ausnahmen einer leidenden und mankenden Gesundheit, des Erkrankens und Sterbens ab; so wird man zugestehen müssen, daß der unverheirathete Mann das ruhigste, bequemste Leben führen, ganz seiner Meinung, seinem Lebensgeschmack sich ergeben und fröhnen kann, während der beweibte Mann sich immerzu fügen und beschränken muß, so selbständig er sei; und auch die Hilf und Pflege müder Stunden und kranker Tage reichlich durch all den Jammer und das Ungemach aufgehoben wird, welches er aus Liebe und Treue gegen Weib und Kind erträgt.

136. Die Verheirathung steht dem Geistlichen frei; eben so die zweite Heirath nach dem Tode der ersten Frau. Die Stelle 1. Tim. 3, 2. „Eines Weibes Mann“ (μιᾶς γυναικὸς ἄνδρα) redet nicht von successiver, sondern
\quellseite{260}
von gleichzeitiger Polygamie, — wenn es nemlich recht ist, die successive überhaupt eine Polygamie zu nennen.\footnote{Qui uxorem olim habuit, nunc autem non habet, is non est jam unius uxoris maritus, sed nullius, atque adeo viduus. Relata enim se mutuo ponunt et tollunt. Sublata igitur uxore tollitur correlatum, quod est maritum esse. G. Calixt. l. c. p. 41.} Sie verbietet auch einen Mann zum Bischof oder Presbyter zu machen, welcher die Regel Christi Matth. 19, 9. und die dem Geistlichen nothwendige σεμνότης durch willkürliche Auflösung (Scheidung) der Ehe und lüsternes Schließen einer zweiten etc. Ehe mit Füßen getreten hat, wie das unter Juden und Griechen gewöhnlich war. Die successiven zweiten Ehen aber hat man anfangs im apostolischen Worte mit nichten verboten gesehen. Hieronymus schreibt ad Oceanum: „Nosti problema tuum. Carterius Hispaniae episcopus, homo et aetate vetus et sacerdotio, unam antequam baptizaretur, alteram post lavacrum, priore mortua, duxit uxorem et arbitraris, eum contra apostoli fecisse sententiam, qui episcopum unius uxoris virum praeceperit ordinandum. Miror autem, te unum protraxisse in medium, quum omnis mundus his ordinationibus plenus sit: non dico de Presbyteris, non de inferiori gradu; ad Episcopos venio, quos si sigillatim voluero nominare, tantus numerus congregabitur, ut Ariminiensis synodi (quae numerosissima fuit) multitudo superetur.“ Solche historische Thatsachen beweisen noch mehr, als die von den Thatsachen absehenden Urtheile ausgezeichneter Väter, wie man sie zum Beispiel
\quellseite{261}
in dem oben angemerckten Buche Geo. Calixti finden kann, daß das früheste Alterthum wie gegen die Ehe der Geistlichen überhaupt, so gegen die zweite Ehe derselben nicht entrüstet war, daß nicht apostolische Befehle, sondern kirchenordnungsmäßige Beschlüsse späterer Zeiten den Geistlichen auch in diesem Stücke anders als andere Männer behandeln. Wenn die jetzige griechische Kirche den Pfarrgeistlichen die erste Ehe gebietet, die zweite aber sogar in dem Maße verbietet, daß sie denselben gebietet, als ἰσομόναχοι ins Kloster zu gehen; so ist ihr Verfahren auf eine doppelte, sei es auch alte Mißdeutung der apostolischen Stelle 1. Tim. 3, 2. gegründet, und es mag wohl für die einfache Auslegung der berühmten Stelle abermal ein Wort von Hieronymus ihr gegenüber angeführt werden: „Quod ait Apostolus „unius uxoris vir“, sic intellegere debemus, ut non omnem monogamum digamo putemus esse meliorem: sed quo is possit ad monogamiam et continentiam cohortari, qui sui exemplum praeferat in docendo. Esto quippe, aliquem adolescentulum conjugem perdidisse et carnis necessitate superatum accepisse uxorem secundam, quam et ipsam statim amiserit et deinceps vixerit continenter; alium vero usque ad senectam habuisse matrimonium et uxoris usum et, ut plerique existimant, felicitatem, et nunquam a carnis opere cessasse: quis vobis e duobus videtur esse melior, pudicior, continentior? Utique ille, qui infelix etiam in secundo matrimonio fuit et postea pudice et sancte conversatus est; et non is, qui ab uxoris amplexu nec senili separatus est aetate. Non sibi ergo applaudat, quicunque
\quellseite{262}
quasi monogamus eligitur, quod omni digamo sit melior, quum in eo magis sit electa felicitas quam voluntas. Quidam de hoc loco ita sentiunt: Judaicae, inquiunt, consuetudinis fuit, vel binas habere uxores, vel plures, quod etiam in veteri lege de Abraham et Jacob legimus: et hoc nunc volunt esse praeceptum, ne is, qui episcopus eligendus est, uno tempore duas pariter habeat uxores. Multi superstitiosius magis quam verius etiam eos, qui, quum gentiles fuerint, et unam uxorem habuerint, qua amissa post baptismum Christi alteram duxerint, putant in sacerdotio non legendos; quum utique, si hoc observandum sit, illi magis ab episcopatu arceri debeant, qui vagam per meretrices ante exercentes libidinem, unam regenerati uxorem acceperint, et multo detestabilius sit fornicatum esse cum pluribus, quam digamum reperiri, quia in alio infelicitas matrimonii est, in alio ad peccandum prona lascivia.“ (S. Calixt. l. c. P. I. p. 43 f.) — Aehnlich Theodoret: „Illud unius uxoris virum mihi videntur quidam recte exposuisse. Olim enim consueverant et Graeci et Judaei cum duabus et tribus et pluribus etiam uxoribus simul jungi matrimonio. Quin etiam nunc, quum leges imperatoriae prohibent duas simul uxores ducere, rem habent cum concubinis et meretricibus. Dixerunt itaque, divinum docuisse Apostolum, eum qui cum una sola uxore habitat, dignum esse, qui ordinetur Episcopus. Neque enim secundum rejicit matrimonium, qui illud saepe fieri. (Hierauf Anführung apostolischer Beweisstellen, dann weiter:) Si quis expulsa priore uxore alteri esset conjunctus
\quellseite{263}
, dignus foret, qui reprehenderetur, omni jure accusationi obnoxius. Si autem vis mortis priorem disjunxerit, et urgens natura, ut secundae conjungeretur, coegerit, non ex voluntate, sed ex casu processit secundum matrimonium. Haec considerans admitto eorum interpretationem, qui sic intellexerunt.“ (l. c. p. 42 f.)

An den apostolischen Stellen und deren richtiger Auslegung liegt hier alles.\footnote{Für den nemlich, der sie nicht für antiquierte Kirchenordnungen hält, denen die gegenwärtige Zeit entwachsen sei. Allerdings aber ist die Auslegung schwer. Daher auch in den ersten Jahrhunderten die Ungleichheit der Ansicht. Mildigkeit der That, wo das Wort nicht völlig deutlich.} Erweist sich nun schon aus dem Voraufstehenden, wie die einfache, mit dem übrigen Gottesworte harmonische Auslegung der Stelle 1. Tim. 3, 2. so alt und älter, wie die ihr gegenüberstehende ist, daß also Schrift und frühestes Alterthum im Einklang stehen; so ist nicht abzusehen, warum das Gewissen der protestantischen Kirche durch die römische und griechische Praxis solle des Weges ungewiß und in Unruhe versetzt werden.

\section{II. Brautwahl}

137. Ein interessantes Kapitel für alle Geistliche, sie mögen ledig, verheirathet oder Witwer sein. — Ein Weib wird vor der Ehe kaum jemals recht erkannt; erst durch die Ehe kommt an den Tag, was in ihr ist. Bis zur
\quellseite{264}
Unkenntlichkeit verschieden von der Jungfrau ist oftmal dieselbe Person als Frau. Das macht, ein Weib ist für den Mann geschaffen und vollendet sich in rechter harmonischer Weise erst dann, wenn sie in Verbindung mit demjenigen kommt, der ihr den Mangel und die Schwachheit ihres Geschlechts erstatten kann und wirklich erstattet. Ähnlich ist es mit dem Manne, wenn sich gleich das Gesagte auf ihn nicht so ausschließend und nicht in gleichem Maße anwenden läßt, wie auf das Weib. Des Mannes Lebensberuf vollendet sich nicht in und durch die Ehe, wie es bei dem Weibe der Fall ist; aber die Ehe wirkt mächtig auf seinen ganzen Lebensberuf ein. Je nach dem Berufe gestaltet sich der Einfluß der Ehe, aber kein Beruf entgeht ihrem Einfluß ganz und gar. Es treibt den jungen Mann insgemein zum ehelichen Leben, aber es ist in diesem Triebe mehr, als man bei dem heillosen Mißbrauch desselben von Seiten so vieler jungen Leute zu denken geneigt ist, nemlich ein Drang zur Vollendung. Wüßten junge Männer, was es mit der Ehe ist, so würden sie sich gewiß nicht übereilen, hineinzukommen; so viel nimmt sie, bei allem, was sie gibt — und so viel läßt sie auf. Aber es ist vom HErrn, daß niemand das öffentliche Geheimniß der Ehe anders als durch eigene Erfahrung erfassen kann, denn gerade diese Erfahrungen sind es, unter denen wahrhaft männliche Gemüther reifen. Nicht die Ehe vollendet einen Mann, aber werden wohl ohne sie alle Saiten seines Wesens zu einer heiligen, Gott wohlgefälligen Harmonie gestimmt? — Bei aller Verehrung für so viele heilige und große Seelen, die im jungfräulichen Leben dem HErrn
\quellseite{265}
und seinen Kindern dienten, — bei aller Demuth ihnen gegenüber sei doch diese Frage erlaubt, und unverhohlen sei die Antwort.

138. „Niemand hat jemals sein eigen Fleisch gehaßet, sondern er nähret es und pflegt sein, gleichwie der Herr die Gemeinde.“ Ephes. 5, 29. Es ist Mannesehre, das Weib zu nähren und zu pflegen; bevor er das kann, soll er die Hand zum Weibe nicht ausstrecken. Nun aber wird ein Mann des zeitlichen, ihn und die Seinigen nährenden Berufes viel später mächtig, als er in sich Lust und Drang zum ehelichen Leben zu finden pflegt; ach, wie wenig Bürgschaft vermag oftmals ein Freier dem erwählten Weibe zu bieten, daß sie bei ihm wohl geborgen sei, daß er sie nähren und pflegen könne! Man mißverstehe diese Worte nicht. Es soll keinerlei Geldheirathen Lob gesprochen, die Frage nach Speis und Kleidung nicht zur ersten Frage eines Freiers oder einer Jungfrau, die man zur Ehe begehrt, gemacht werden. Auch meinen wir nicht, daß ein Mann seiner Auserwählten bereits ein Leben in Reichthum und Fülle solle bieten können: Schweiß des Angesichts und Kummer der Ernährung gehören zur Ehe und sind auch gar nicht zu umgehen, sie machen die Ehe gesund und glücklich, so wunderlich das lautet. Aber der Mann muß doch dem Weibe eine Aussicht fürs Zeitliche bieten können; es muß doch Kraft, Grund und Hoffnung da sein für das Gelingen, und es sei das Gelingen auch nur ein Durchkommen, bei welchem man Nahrung und Kleider hat. Wer seiner Auserwählten nicht so viel bieten kann, wenn er um sie wirbt, liebt sie nicht, indem er um sie wirbt;
\quellseite{266}
er handelt eines frommen Christenmannes unwürdig und erweckt den Verdacht der Unbesonnenheit, wenn nicht einen schlimmeren. — Hier nun liegen für den Jüngling, auch für den jungen Geistlichen Kämpfe. Während einerseits die Natur oft zur Verehelichung drängt, wehrt die Weisheit und so geistliche, wie weltliche Pfleger und Vormünder sprechen ihr gewichtiges Veto gegen unbesonnene, übereilte Schritte. Es entsteht ein unfreiwilliger Cölibat, der in der Regel nicht reiner zu sein pflegt, als jener verrufene der Mönche und Nonnen. Dieser Cölibat und der oft lange und heiße Kampf, unter welchem er aufrecht erhalten, gereinigt und geheiligt werden soll, ist die ernste Aufgabe des jungen Mannes. Sich hier in Heiligung und Ehren, in Unschuld und fröhlichem Muthe zu bewähren, darin besteht des christlichen Jünglings größte Ehrenhaftigkeit, und wohl denen, welchen in dem Betracht aus 1. Joh. 2, 13. gelesen werden kann: „Ich schrieb euch Jünglingen, denn ihr habt den Bösewicht überwunden.“ Sie ist Geduld der Heiligen und jungfräulicher Seelen!

139. Es fragt sich nun, ob nicht vielleicht eine frühzeitige Verlobung, eine Art Brautstand, in welchem man sich vornherein ins Harren und Warten ergiebt, ein Mittel zur Erleichterung des Kampfes sei. Nicht bloß in den weltlichen Ständen, sondern auch im geistlichen scheinen unzählige Beispiele aus alter und neuer Zeit die Frage zu bejahen. Unter diesen Beispielen sind Namen von Glanz und Würde. Man denke an einen H. Bullinger, an J. J. Moser und andere. In der Gegenwart braucht man nicht lange zu suchen; es wird schwer ausfindig zu machen
\quellseite{267}
sein, ob mehr Studenten und Candidaten bereits „Bräute“ haben oder keine. — Es scheint hierbei ganz zu vermuthen, daß eine frühe Verlobung das Gemüth des Jünglings aus dem Sehnen und Suchen reißen, das Auge fixieren, dem Gang zu bestimmterer Richtung verhelfen könne, und es soll auch nicht in Abrede gestellt werden, daß es gewissen Charakteren bei früher Wahl gelingen sei. Wenn man sich aber in den Fall setzt, nicht zu geschehenen Verlobungen das Beste zu reden, sondern rathen zu sollen; so wird man doch schwerlich oder doch nur selten einmal Zuversicht gewinnen, zuzurathen. Aus dem Suchen ist der Jüngling freilich gerissen; er kennt sein zukünftiges Weib mit Namen, geht mit ihr um, lernt sie vor der Ehe zur Genüge (wenn nicht gar zum allzufrühen Ueberdruß) kennen, übt auf sie Einfluß und erleidet Einfluß etc. etc. Aber es bringt dieser Zwischenstand, besonders für gewisse Naturen, auch seine besondern und großen Versuchungen mit sich, — und da das Band nicht allzufest geschlungen ist, wenigstens noch nicht so fest, daß nicht eine Lösung, sei’s auch unter beiderseitiger Einwilligung, möglich erschiene; so ist nicht einmal die Fixation der Gedanken vollkommen, nicht einmal die Richtung entschieden. Es ist ein ungewisses, versuchungsvolles, oft sündenvolles Ding um einen langen Brautstand; man kann ihn im allgemeinen nicht rathen, sondern nur beanstanden. Ueberdies ist es auch ein ärgerliches Ding, das nichts Gutes stiftet. Der Bürger, der Landmann weiß zwischen dem Liebesverhältniß eines jungen Geistlichen (auch einer jungen Pfarrerstochter) und dem seines Sohnes (seiner Tochter)
\quellseite{268}
keinen rechten Unterschied zu machen, wie denn auch für die meisten Fälle zu bezweifeln steht, ob einer sei. Er sieht in beiden eine ganz gewöhnliche Liebschaft. Ist die eine erlaubt, so ists die andere auch; man kann von beiden Gutes und Böses denken. Warum predigt der junge Geistliche gegen Liebschaften des Landvolks und der bürgerlichen Jugend, wenn er selbst eine Liebschaft hat? Oder predigt er nicht dagegen, weil er gegen sich predigen müßte? Predigt er vielleicht auch gegen manch anderes nicht, weil seine Seele die Schärfe und Schneide durch das Bewußtsein eigenen Leichtsinns und Vergehens verloren hat? In einem, wie im andern Falle sitzt Schaden, wenn ein junger Geistlicher sich frühzeitig verlobt. Es wäre immerhin besser, wenn er, so lang er keine Frau nehmen kann, auch keine Geliebte und keine Verlobte und keine Braut hätte.

140. Zwar ist es nun auch im Gegentheil wahr und die Erfahrung belegt es mit zahlreichen Beispielen, daß Aufschub der Brautwahl bis zu der Zeit, wo man der Braut ein gewisses zeitliches Loos in Aussicht stellen und darbieten kann, auch wieder in Verlegenheit setzen kann. Meist ist man dann über die Jahre hinaus, in denen man fest zu wählen pflegt, weil man noch nicht weiß, wie folgenreich der Schritt ist, den man thut; aus dem Wählen wird dann oft ein Quälen. Je würdiger und eingezogener man gelebt hat, desto weniger Bekanntschaft hat man vielleicht unter dem weiblichen Geschlecht; der Kreis, unter dem man wählen kann, ist nur klein — und am Ende wählt man vielleicht gerade verkehrt. Das alles ist gewiß
\quellseite{269}
mäßig genug. Allein das kann so kommen, und kann auch anders und viel besser kommen und kommt auch recht oft viel besser. Gott wirds walten! Es ist nun einmal im Ganzen doch besser, nicht thun zu wollen, was man nicht kann. Es ist das Sichere, Ehrlichere, Edlere, erst durch eine Anstellung völlig Mann zu werden und das Weib, um die man freien will, samt ihrem Vater in den Stand zu setzen, daß sie den Schritt verstehen, den sie thun sollen, daß sie Ja und Nein mit Besonnenheit sprechen können. Was dann kommt, kommt im Frieden und zum Frieden, und es möchten sich unter den ehelichen Verbindungen, welche sich bewähren, vielleicht mehr finden, die zur rechten Zeit, als die zur Unzeit, d.\,i. verfrühten Zeit geschlossen sind. Ob aber auch nicht, für die meisten Menschen wird es doch besser sein, zu allem Ding die Zeit, die rechte Zeit zu erwarten.

141. Es mag nun aber die Wahl früher oder später vor sich gehen, so heimlich thut hierin kein frommer Jüngling, daß er ohne Beirath und Segen seiner Eltern zum Ziele seines ehrlichen Verlangens eilte. Es ist der schändlichste Undank, welcher sich denken läßt, in Berathung des Schrittes, welcher fürs eigene Leben der wichtigste und folgenreichste ist, diejenigen Menschen auszuschließen, welchen man nicht bloß alles verdankt, was man ist und hat und nun vergeben will, sondern welche auch unter allen Rathgebern die treuesten Herzen und eine über jeden Zweifel erhabene liebevolle Gesinnung haben. Dazu möchten spätere Leiden, die ja nicht ausbleiben werden, sich dem Gewissen als göttliche Strafen für die Eigenmächtigkeit
\quellseite{270}
der Wahl darstellen und erst dadurch recht schwer und lastend werden. Ein Jüngling besserer Art überlegt, daß er bei seiner Verheirathung ohnehin Vater und Mutter verläßt, um an seinem Weibe zu hangen, daß es eine Trennung gibt, es werde, wie es will, — und eben deshalb freut er sich, mit den geliebten Eltern wenigstens in dem Schritt recht einig zu sein, welcher zur Trennung führt.

124. Bei der Wahl der Braut wird ein Jüngling wohl thun, hohe Forderungen an die persönliche Trefflichkeit zu stellen. Das Herz läuft ohnehin in diesen Angelegenheiten oft genug der Ueberlegung und das Fleisch dem Geiste voran. Kaum wird selbst der, welcher geneigt ist, hohe Forderungen zu stellen, im eigenen Falle competent und unbestochen genug sein. Gut, wenn dann wenigstens die innere Aufforderung, Treffliches zu wählen, Mißtrauen in’s eigene Urtheil bewirkt und vor schnellem Zufahren behütet. Auch ist ja gar keine Frage, daß ein Mann diejenige zur Gehülfin erwählen soll, welche ihm vermöge ihrer Eigenschaften am meisten zu helfen vermag. Ein Streben nach dem Trefflichen ist hier Tugend, und das um so viel mehr, wenn es Licht und Verstand genug hat, nicht ungemeßene Forderungen zu stellen, sondern, was trefflich sei, nach der eigenen Beschaffenheit, dem eigenen Verstande und Berufe abzumessen. Demnach hätte man unter dem Trefflichen weder Jugend, noch Schönheit, noch Stand, noch Reichthum an und für sich zu verstehen. Es bleibt von allen diesen Dingen gewiß ein jedes in seinem Werthe; aber fänden sie sich auch alle in Einer Person
\quellseite{271}
zusammen, es wäre aber sonst nichts dabei, so wäre der arme Bühler dennoch sein Leben lang ein geschlagener Mann. Mehr als alles das ist, was man Bildung heißt. Aber wenn sich auch dieser Vorzug zu den schon genannten vier andern fände, so wäre dem Manne trotz dem noch nicht geholfen. Bei seinem Weibe hat ein Pfarrer Gottes Furcht und Liebe zum ersten, weibliche Bildung und Tüchtigkeit zum zweiten, Sinn für seinen Stand und Beruf zum dritten, herzliche Neigung aber zu seiner Person keineswegs zum letzten Augenmerk zu setzen; denn was ist Ehe ohne Neigung!

143. Von Neigung ist die Rede gewesen. Laßt uns dies Erforderniß gesondert wiederholen. Neigung ist blind, ist Thorheit, wenn sie bei der Wahl anfängt, mittelt und endet. Sie allein entscheidet nicht, hingegen gebührt ihr der rechte, volle Ausschlag unter den Personen, bei denen man unter Beirath der Eltern und treuer Freunde die rechten Eigenschaften gefunden hat. Eine Ehe ohne Neigung ist, wir betonen und wiederholen, eine grauenvolle, fast möchte man sagen, eine ruchlose Sache. Eine Neigung zu Unwürdigen aber ist eine Täuschung, welche sich bald genug in Jammer, wohl gar in Haß verkehren wird. Und wehe dann dem enttäuschten Manne und dem getäuschten Weibe!

\quellseite{272}
\section{III. Die Pfarrerin}

144. Der oft gebrauchte Name „Pfarrfrau“ muß als unschicklich beanstandet werden, wenn man das Wort nach seiner Bedeutung auffaßt. Man sollte billig den bescheideneren Namen Pfarrerin oder Pfarrersfrau gebrauchen. Jener Name drückt zum Unterschied von diesen den ganzen Irrtum und die ganze kokettierende Ueberschätzung aus, die man in unsern Tagen in der Ansicht von der Pfarrerin so häufig findet. Man will nemlich der Pfarrerin gerne eine einflußreichere und bedeutungsvollere Stellung in der Gemeinde ihres Mannes zuweisen, als früherhin. Etliche wollen sogar behaupten, der Locus von der „Pfarrfrau“ dürfe in keinem Pastorale fehlen. So fehle er denn auch in diesem Büchlein nicht, wenn auch nur, um die Pfarrerin aus dem Pastorale lieber wieder hinaus und in die anspruchslose Stellung ihres hausmütterlichen Berufes zu verweisen.

145. Es ist eine willkürliche, auf nichts beruhende Behauptung der neueren Zeit, daß die Pfarrerin in der Gemeinde ihres Mannes Diaconissin im biblischen Sinne sein soll. Ob 1. Tim. 3, 11. allein auf die Frauen der Diaconen oder auch auf die der Presbyter gehe, darüber mag man streiten; jeden Falls aber geht aus der Stelle hervor, daß nicht einmal die Frauen der Diaconen Diaconissinnen sein, sondern allein des Amtes ihrer Männer würdig wandeln sollen; und doch würde sich die Ehfrau des Diaconus
\quellseite{273}
immerhin noch natürlicher zur Diaconissin empfehlen, als die des Presbyters. Diaconissen sollen kinderlose Witwen sein, welche den Ueberrest ihrer Lebenszeit und die gesammelte Lebenserfahrung zum speciellen Dienste der Armen und Kranken anwenden wollen. S. 1. Tim. 5. Witwen sollen es sein, ganz natürlich, weil der Beruf der Armen- und Krankenpflege mit dem der Ehe- und Hausfrau sich nicht vereinigen läßt. Eine Ehe- und Hausfrau kann das liebevollste Herz für Arme und Kranke haben, aber zum Lebensberuf kann sie sich Armen- und Krankenpflege nicht machen, weil sie schon einen andern Beruf hat. So viel sie thue und zum Wohl der leidenden Menschen wirke, es geschieht doch Alles nur nebenher, es ordnet sich ihrem hausmütterlichen und ehelichen Berufe unter. Armenpflege und Krankendienst legen Anstrengungen und Aufopferungen auf die Diaconissin, welche oftmals die Berufsarbeit der Hausmutter geradezu aufheben und unmöglich machen würden.

146. Die Pfarrerin ist Ehefrau des Pfarrers, Mutter und Erzieherin seiner Kinder, seine Gehülfin zur Erreichung und Erfüllung der apostolischen Forderung, daß er seinem Hause wohl vorstehe und gehorsame Kinder habe. 1. Tim. 3, 4. Je vollkommener sie das ist und dabei auf die apostolische Vermahnung 1. Tim. 3, 11. sieht („Ihre Weiber sollen ehrbar sein, nicht Lästerinnen, nüchtern, treu in allen Dingen“ — σεμνὰς, μὴ διαβόλους, νηφαλίους, πιστὰς ἐν πᾶσι); desto völliger ist sie, was sie sein soll. Sich darin bewähren, darin sich vervollkommnen, das ist ihr Beruf und Lebensziel. Die größte Achtung
\quellseite{274}
der Gemeinde und Anderer, die sie kennen lernen, wird ihr nachfolgen, wenn sie dies scheinbar geringe, dies bescheidene Ziel getreulich verfolgt und erreicht hat. Es ist ein bescheidenes Ziel im Vergleich mit dem, was der Pfarrerin von verständigen Gönnern zugemessen wird; aber es ist dennoch ein hohes Ziel und schließt eine große Summe guter Werke ein, durch welche sie sich und ihren Glauben zieren kann bis in jene Welt hinein. Ihr stiller Wandel, der ihr Haus verklärt und zu einer angenehmen Friedenshütte macht, wird je länger je mehr auch von Andern erkannt werden und, ohne daß sie dran denkt (denn so Etwas läßt sich nicht erreichen, so wie es absichtsvoll angestrebt wird), einen tiefen Eindruck sogar auf die machen, welche dem Worte ihres Pfarrers nicht glauben. Der Eindruck ihres stillen Wandels wird um so tiefer gehen, je mehr sie selbst von der Herrlichkeit des Berufes erfüllt ist, den ihr Liebster trägt; ihr ganzes Thun und Lassen wird den Abglanz seines heiligen Amtes tragen; man wird ihr’s allenthalben abmerken, daß sie eine Pfarrersfrau ist. Ihr Mann wird des Ehre und Freude haben. Gleichwie er des priesterlichen Volkes Vorbild sein soll, so wird auch sie ein Vorbild und gutes Beispiel der Frauen sein. Sie wird es sein, nicht weil sie es mit Anstrengung affektiert, sondern weil sie es, abgesehen von aller Leute Augen, Anerkennung und Bewunderung, mit Lust und Liebe ist und gar nicht anders kann. So ist ihr, so gibt sie sich auch; sie dient im schönsten Loos der Frauen ihrem Herrn, ihr ist ein schön Erbtheil geworden; ihr Joch ist sanft und ihre Last ist leicht; auch wenn es hart hergeht,
\quellseite{275}
hat sie Freude und Stärke zum Guten. Sie hat ihr Nest nah an den Altären Jesu gefunden und grünt im Vorhof des Heiligthums.

147. So eingezogen und still auch eine Pfarrerin lebe, sie wird doch den Pfarrkindern ihres Mannes je länger je mehr persönlich bekannt werden. Das kann nicht ausbleiben. Sie lernt, ohne daß sie darnach forscht, die häuslichen Verhältnisse der einzelnen Familien kennen; ihr weiblich theilnehmendes Auge wird schnell jegliche Noth und jeden Jammer gewahr. Holdselige Worte, die aus ihrem Herzen quillen, weil Gottes Wort drin quillt, — liebliche Gebärde, welche mit dem Sinn des Herzens harmoniert und ihn mit einer Art von Natürlichkeit und Unmittelbarkeit widerspiegelt, prunkloses Helfen, daran sie selbst die unschuldigste Freude hat, — ein Auge voll Lieb und Freundlichkeit, ein Herz voll Gebet, wenn helfen nicht möglich ist: solche Eigenschaften machen die fromme Jüngerin mehr und mehr zum Liebling der Gemeinde. Bald wird sie gesucht werden, so wie man sie erkannte; bald wird jede Noth, die einem Pfarrkind zukommt, ein Fingerzeig auf die gute Helferin und ihr Name eine Erinnerung an allerlei gute Werke werden. — Ein rechtes Christenweib, in der Stellung einer Pfarrerin, braucht sich nicht bemerklich zu machen. Sie wird ohnehin bemerkt.

148. Die Amtsgeschäfte eines Pfarrers sind populärer, als die eines jeden andern Mannes; Jedermanns Augen sehen darauf, Jedermann nimmt Theil daran. Wie sollte denn die Pfarrerin nicht Theil nehmen? Ist sie doch auch ein Gemeindeglied und steht gewissermaßen an
\quellseite{276}
der Spitze aller weiblichen Gemeindeglieder; ihr eheliches Verhältnis stellt sie ihrem Hirten nicht so nahe, daß sie vergesse, daß er ihr Hirte ist; nur ihre Theilnahme an Allem, was die Gemeinde und das heilige Amt betrifft, wird durch die innige Verbindung mit ihrem Pfarrer erhöht. Hier, im Verhältnis der Pfarrerin gegenüber und zu dem Amte ihres Mannes, liegt die schönste Bewährung. Die Grenze, die zwischen der Nähe zu ihrem Mann und der Ehrerbietung vor ihrem Hirten und seinem Amte sich hinzieht, — das ist die Linie des Schönen und Schicklichen für sie. Dies Maß und Fern berechnen, das ist ihre Kunst. Das Maß, bis zu welchem sie es in dieser Kunst bringt, das ist auch das Maß, womit man eine Pfarrerin als solche mißt. Hier liegt ihre Würde und ihre Gemeinheit, ihre Hoheit und ihre Niedrigkeit, ihre Tugend und Untugend. Keine häusliche Tüchtigkeit, keine Bildung, kein sonstiger Vorzug kann ein Ersatz für einen Mangel in der Erkenntniß dieser Grenzlinie des Eheweibes und des Pfarrkindes werden, kann eine rohe Auffassung ihres Verhältnisses zu ihrem Hirten verdecken. Nimmt eine Pfarrerin an dem Werke ihres Eheherrn den innigsten Antheil, ohne sich zu einer Mittelsperson zwischen ihm und der Gemeinde zu erheben; ist sie ganz die Seine, sein Weib und seine Freundin, und doch auch wieder sein ehrerbietig Kind, seine Tochter; nahe wie Niemand, als Pfarrkind nicht näher als Alle; fürchtet sie sich vor der Möglichkeit, ihren Mann in amtlichen Dingen zu bestimmen; vernimmt sie immer des Herrn heiliges: „Weib, was hab ich mit dir zu schaffen?“ mit verständiger Freude und ist es ihre volle Angelegenheit,
\quellseite{277}
nur ihres Mannes Weib, seine Hausfrau, seine Freundin, nur sein Pfarr- und Beichtkind zu sein; und ist sie das und nicht mehr, gelingt ihr das; — so ist sie, was sie soll, und der Einfluß der Unschuld und Einfalt wird ihr selbst auf ihren Mann nicht entgehen. Er wird sich um so mehr scheuen, in seinem ehelichen Verhältnis, geschweige in seinem amtlichen zu vergeben, daß er Hirte ist, weil ein liebes Auge voll Lieb und Ehrerbietung seine Werke prüft und Bestätigung der Lieb und Ehrerbietung, eben damit weibliche Befriedigung und wahres eheliches Glück in seinem trefflichen Verhalten sucht. Nimmt hingegen die Pfarrerin Partie für ihren Mann: ist’s nicht das Amt, das sie an und in ihm ehrt, sondern er, um dessen willen sie das Amt ehrt; ist’s vielleicht gar gemeine Eigenliebe, was sie für das amtliche Thun ihres Mannes eifern macht; hat er Recht, weil er der Ihre ist; sind seine Predigten, Kinderlehren, Reden etc. die besten und die schönsten, weil er das Glück oder Unglück hat, ihr Gemahl zu sein; ist sie seine Vertreterin gegen Jedermann, die Mittelsperson zwischen Jedermann und ihm; kann man sich an sie wenden, sie gewinnen, wird sie dann seine Heva, die Adam verführt zu Thun und Lassen: — ach, was für ein gemeines Weib ist sie dann! Parteitische Weiber sind doch unter allen die unleidlichsten, — und Pfarrersweiber, die ihrer Männer Partie sind, können doch nur wieder die unleidlichsten unter allen unleidlichen sein. Wer wird sie achten? Die sie bedürfen, werden sie suchen und benützen; Andere werden sie scheuen; sie sind wahrer Achtung ohne. — Es soll damit gar nicht gesagt werden, daß die Pfarrerin von ihrem
\quellseite{278}
Manne gar niemals einen Auftrag bekommen könne, gar niemals etwas thun dürfe, sein Amt zu fördern. Sie ist die erste Schülerin ihres Mannes, sie wächst an seiner Seite in Erfahrung und Erkenntniß, und je mehr sie zunimmt und je mehr sie die oben angedeutete Grenzlinie des Schicklichen findet und bewahrt, desto untadelicher wird es sein, wenn sie zuweilen, namentlich in der Seelsorge der Frauen, einen vorübergehenden Auftrag ausführt oder eine dauerndere Aufgabe löst. Manche Dinge bringt kein Mann gegen andere Frauen über die Lippen, da kann und soll die Einzige, mit der er sie besprechen kann, eine Dolmetscherin der Hirtenliebe werden. Dann ist sie wahrlich auch Diaconissin, wenn auch nicht im biblischen Sinn. Bei aller Freude an solchem Dienen verlerne sie nur niemals die Lust und Liebe zur Demuth und Anspruchslosigkeit ihrer natürlichen Stellung. Ein frommer Mann leistet hierin Hülfe und weiß auch sein frommes Weib in dem zu erhalten, was ihr geziemt.

149. Von denjenigen Pfarrerinnen, welche ihre Männer und Brotschaffer hätscheln, keine Sorge wissen, als daß nur Gesundheit und Leben des Familienhauptes erhalten bleibe, die vor lauter fleischlicher Liebe und Angst um ihr eigenes Schicksal der Männer Schritte lähmen, sie hindern, ihrem Amte mit Aufopferung nachzugehen und ihre Schafe zu weiden, die wie Delila Simsons Stärke wegnehmen und es für Gewinn achten, die Starken zu bändigen und weiblich und weichlich zu machen: von denen wird der Herr eine Rechenschaft fordern, die sie nicht geben können. Möchte ein schreckendes Wehe in ihre elenden Herzen
\quellseite{279}
schallen, damit sie ablaßen, ihr und ihrer Männer Leben so zu lieben, daß sie es verlieren! — Im Gegenteil, ist ein Weib, so unglücklich, einen faulen Pfarrer zu haben, so werde sie seine Prophetin und rufe ihm ein Vorwärts und eine Ermunterung zur Arbeit zu! Was für eine Schmach, einem Faulen, und nun gar einem faulen Pfarrer zur Seite zu stehen! Noch vom sterbenden Munde der Pfarrerin vernehme der Pfarrer die Ermunterung, die Ermahnung zur Treue, zum Fleiß, zur Aufopferung für ihren Herrn und seine Schafe!\footnote{S. im Anhang p. 277 die „Erinnerungen für Pfarrerfrauen.“}

\section{IV. Kinderzucht}

150. Von der Kinderzucht in einem Pfarrhause sollte wohl nichts Besonderes zu erwähnen sein, als allenfalls, daß ein Pfarrer, eine Pfarrerin die allgemeine Pflicht der Eltern, ihre Kinder wohl zu erziehen, schon um des Wortes 1. Tim. 3. 4. willen vor Andern im Auge zu behalten und zu erfüllen haben. Denn es hat kein Stand der Hauptsache nach bei der Erziehung etwas Anderes als die übrigen zu erstreben — und auch die Art und Weise der Erziehung ist in den Mittelständen, zu welchen sich der Pfarrer rücksichtlich seiner äußern Stellung am besten zählen dürfte, so ziemlich eine und dieselbe, soweit sich nicht individuelle Besonderheiten geltend machen.

\quellseite{280}
151. Es sind vornehmlich etliche Fehler, zu welchen die Kinderzucht der Pfarrer leicht ausschlägt, um deren willen hier Etwas dieses Titels vorzubringen ist. Jedoch reden wir auch hier, wie meist in diesen Blättern, weniger von den Verhältnissen, Tugenden oder Fehlern städtischer Pfarrer, als von denen der Landpfarrer. Der Landpfarrer selbst steht in seiner Gemeinde ohne Zweifel als Augenmerk Aller. Jede Jugend leuchtet an ihm besonders, aber auch jeder Fehler greift augenfälliger. Scharf bemerkt, hoch übel genommen, oft und über’s Maß getadelt, aber nichts desto weniger äußerst selten vor seinen Ohren gerügt wird Alles, was ein Landpfarrer fehlt. Das gilt denn auch von den Fehlern der Kinderzucht.

152. Die Kinder des Landpfarrers müssen sich bald über die Kinder der andern Gemeindeglieder erhaben fühlen. Des Vaters Stellung und das Leben im Hause machen scharfen Gegensatz gegen die Landleute. Die Kinder aus der Gemeinde sehen auch von Anfang an auf die Pfarrerskinder hinauf, räumen ihnen, wenn sie sich des nicht in besonderer Weise unwürdig machen, einen Vorzug und das Recht der ersten Stimme ein. Leben nun Pfarrerskinder mit den andern Kindern des Orts zusammen; so herrschen sie meistens über dieselben, leiten sie, nehmen aber auch umgekehrt von ihnen an. Des Vaters Bildung findet sich deshalb bei Landpfarrerskindern nicht sehr oft wieder, desto häufiger aber eine häßliche, widerwärtige Mischung von bäuerischem und städtischem, herrischem und rohem, gemeinem Wesen, für welche es keinen Namen gibt. Es müssen geistig sehr gehobene Pfarrersfamilien sein, bei
\quellseite{281}
denen sich diese Bemerkungen nicht irgendwie geltend machen. — Vor dieser Mischung der Gemeinheit und Vornehmheit seine Kinder zu bewahren, ihnen einerseits eine rechte Bildung, andererseits aber Einfalt und Demuth, Bescheidenheit und Freundlichkeit gegen den Landmann zu geben und zu bewahren, dies wird die Aufgabe des Landpfarrers bei der Erziehung seiner Kinder sein. Eine schwere Aufgabe, zumal es zu einer völligen Abgeschlossenheit der Pfarrerskinder von der Dorfjugend nicht leicht kommt und kommen kann, und es wenige Eltern gibt, welche die Vortheile einer ländlichen Erziehung so zu benützen wissen, daß ihren Kindern zugleich der Sinn für’s Niedrige und für’s Hohe, für’s Kleine und für’s Große geweckt und gestärkt werden und sie in einem äußerlich geringen Leben der Herrlichkeit einer christlichen, zum Himmel reisenden Bildung nachjagen können.

153. Für junge Kinder in den ersten Jahren, ja bis zur Confirmation hin möchten sich das Leben und die Verhältnisse des Landpfarrers ganz insonderheit nützlich und segensreich erweisen. Wachsen die Söhne mehr heran und vermögen ihnen Vater und Mutter nicht mehr Alles in Allem zu sein, Alles zu ersetzen, was Stadtkinder an Bildungsmitteln voraus haben — und das vermögen wohl nur wenige Väter und Mütter; so bleibt nichts Anderes übrig, als sie auf eine städtische Schule zu schicken. So viel Nachtheiliges auch das Leben auf der Schule haben mag; so ist doch für manchen Knaben schon frühzeitig das Dorfleben, das Zusammengeben unter der heimatlichen Jugend und die Ansteckung der Dorflaster nicht minder
\quellseite{282}
gefährlich. Man achte seines Sohnes und entferne ihn, sobald es seine Individualität erheischt. Es muß ja ohnehin Alles, was wieder Pfarrer werden oder dereinst ein Staatsamt bekleiden soll, durch die Schulen gehen und diejenige Art von Bildung erringen, welche nun einmal eine stereotype Forderung für den genannten Zweck ist. Die Gefahren der Schule zu überstehen, durch ihre Anfechtung hindurchzugehen, ist Beruf des Knaben. Das mache man ihm klar und helfe ihm durch. Der heimathliche Sinn wird ja doch, ist nur ein rechtes Familienleben im Vaterhause vorhanden, durch die Entfernung viel mehr gestärkt und durch die daheim zugebrachten Ferien erfrischt, ja neu entzündet. Auch wächst in der Fremde die Gewandtheit des Lebens und die freiere Bewegung in zusammengesetzteren Lebensverhältnissen.

154. Was die Töchter anlangt, so sind sie leichter als die Söhne vom Umgang mit der Dorfjugend abzuhalten, und es wird ihnen, vorausgesetzt, daß Vater und Mutter irgend sind, was sie sollen, das Leben und Aufwachsen in ländlichen Verhältnissen sehr gedeihlich sein können. Frauenbildung wächst aus dem Familienleben heraus, so gewiß die Frauenpersonen für das Familienleben gebildet werden sollen; und das Familienleben des Landpfarrers ist fast das schönste, das es geben kann. Ueberdies hat jede Landpfarrersfamilie so viel Zusammenhang mit städtischen Familien, daß durch richtige Benützung desselben die ländliche Einseitigkeit der Töchter vermieden werden kann, welche man zu fürchten hat. Wächst das Landmädchen im Pfarrhause heran, so streckt sie sich, ist sie nur rechter Art, von
\quellseite{283}
selbst nach einer größeren Wirthschaft aus, und thut sie das, so ist ihr der Vortheil der vielseitigeren weiblichen Ausbildung, weil das Stadtmädchen keine Anreizung hat, ländliche Wirthschaft zu lernen, wohl aber die Landpfarrerstochter Anlaß genug, ihre größere Tüchtigkeit durch Kenntnißnahme der städtischen Ausbildung zu ergänzen und zu erhöhen. Auch was man höhere Bildung nennt, so weit es Frauen zu wünschen ist, erlangt die Landpfarrerstochter nicht so schwer, als es scheint, zumal wenn es die Eltern wagen, das vornehm fahren zu lassen, was auch den Stadtmädchen nur als eitle unbrauchbare Zier angehängt wird. Zwar ist es wahr, daß man unter den Pfarrerstöchtern auf dem Lande viele findet, denen Ton und feine Bildung fehlt; wer hat aber die Zahl derselben mit der Zahl der Stadttöchter verglichen, denen es an demselben Punkte fehlt? Und umgekehrt, wie viele Landpfarrerstöchter haben später als Hausmütter die angesehensten Familien geziert, — und wie viele Männer der edelsten Bildung holen nicht auch jetzt noch die Gattin gern und ohne Vorwurf aus der Stille eines ländlichen Pfarrhauses! Alles in Alles gerechnet, könnte man die Lage einer Pfarrerstochter, auf dem Lande, welche fromme und gebildete Eltern zur Leitung ihrer Jugend besitzt, eher beneidens- als bedauernswerth finden.

\quellseite{284}
\chapter{Schluß}

\section{Emeriti}

155. Man hat manchmal die Behauptung aufgestellt, es sei nichts ehrwürdiger, als ein frommer Pfarrer im Greisenalter. Seine Pfarrkinder sind, blieb er lange an Einer Stelle, von ihm getauft, unterrichtet, confirmiert und getraut; sie sind gewissermaßen alle seine Kinder. Die Gewöhnung ist im höchsten Grade zur Achtung und Ehrerbietung geworden. Und was man alles Wahres, Halbwahres und Falsches für die Ehrwürdigkeit und einflußreiche, segensvolle Stellung eines alten Pfarrers sagen kann. Denn es paßt am Ende das Meiste doch nur auf einen tüchtigen Mann, der bei einer und derselbigen Gemeinde lange aushielt; das aber ist grade, zumal wenn auch das Wörtchen „tüchtig“ ein wenig genauer genommen wird, kein gar oft wiederkehrender Fall. — Schon bei anderen Berufsarten, bei welchen die Fertigkeit das Größte ist, macht sich das Alter als ein großes Hinderniß geltend; wie erst bei dem Pfarrersberuf, in welchem Fertigkeit das Geringste ist, ein frischer, reger, immer wacher Sinn aber um so viel höher angeschlagen werden muß. Alle Kräfte reichen nun einmal selten aus, das Amt zu versehen, und es dürfte daher wohl die Krone eines alten Pfarrers sein, wenn er nicht darauf besteht, erst am Grabe den Hirtenstab niederzulegen, wenn er sich nicht selbst überleben mag, sondern weiß, wo es ihm geziemt und ehrenhaft ist, stille zu werden, zum Gebet für die Gemeinde, zur Liturgie sich zu halten, zurückzutreten von Allem, was ohne männliche Kraft und Frische
\quellseite{285}
nicht geschehen kann, und der Kraft den Platz zu räumen, welchen sie allein würdig einnehmen kann. In kraftlosen Tagen noch im geistlichen Amte arbeiten wollen, erweckt bei dem Volke nicht Mitleid, nicht Entschuldigung; das Volk fühlt sich durch den Eigensinn der Schwachheit — und gewiß mit Recht — höchst beschwert. Dagegen gönnt es dem Greise gern seine Ruhe und krönt ihm dieselbe mit Ehrerbietung und Liebe.

156. Vielleicht ist es ebenso oft die Lust am Leben und die Scheu vor der Erinnerung an den Tod, welchem man durch Zurücktreten vom Amte näher zu kommen scheint, als Armuth oder Geiz, was alten Pfarrern mißräth, ihr Amt in jüngere Hände niederzulegen. Deshalb sollten alte Pfarrer, zumal wenn ihre Schwachheit offenkundig groß ist, gehalten sein, einen Vicar oder Adjunct anzunehmen, und man sollte sie dazu, wie es auch hie und da geschieht, unterstützen. Wüßten die Pfarrer nicht anders, als daß sie einen Vicar anzunehmen haben, so würden sich voraus die Leidenschaften legen, welche dagegen streiten; das Herz würde sich drein finden und bei Zeiten darnach richten. Mancher Greis würde allgemach die schöne Seite der Sache herausfinden. Wie schön, die eigene Jugend noch einmal im Anschauen der Amtsverwaltung eines jugendlichen Vicarius zu feiern! Wie schön, so nach und nach dem Amt Valet zu sagen, das man ja nicht für die Ewigkeit auf sich nahm! Wie heilsam, so nach und nach zu sterben, ehe man stirbt, auf daß man lebe, wenn’s sterben gilt! — Es sollte sich’s schon jeder jüngere Pfarrer zur Aufgabe und zum festen Vorsatz machen, aufzuhören zu rechter
\quellseite{286}
Zeit, keine Last und Schwachheit seines Alters den Gemeinden zum Mitleiden aufzulegen.

157. Oft sind es andre Gründe, als Altersschwäche, welche einen Diener des Evangeliums nöthigen können, sich für einen Emeritus anzusehen und auf das Amt zu verzichten. Oft ist es anhaltende Krankheit und Siechthum, Abnahme des Gedächtnisses oder die Unmöglichkeit, sich von der Stelle ferner zu nähren etc. In diesen Fällen ist nichts zu tadeln an der Resignation. Zuweilen aber ist es auch fuga laborum, taedium adversitatum, amor status pinguioris. Dann freilich ist es eine Mietlingsseele, die vom Amte und der Gemeinde weicht. Hat sich doch Niemand selbst berufen; so kann sich auch Niemand selbst abrufen. Jeder muß dem Herrn sein Weh und Leid befehlen und seine Last tragen, so lang es Gott gefällt.

158. In einem andern Sinne emeritus, ausgedient und zwar ohne Lob und Verdienst, ist der, welcher falsche Lehre predigt ohne sich belehren und bekehren zu lassen, oder in sündlichem Leben ohne Buße und Besserung hartnäckig verharrt, sein Amt nachlässig versieht, Schisma anrichtet etc. Ein solcher Pfarrer ist ein Schandfleck und ein Unglück der Kirche. Hier ist nicht bloß die eigene Seele eines solchen Menschen, sondern es sind die Seelen ganzer Gemeinden in Gefahr. Zwar ist es wohlgethan, nach des Apostels Weisung 1. Tim. 5, 19., gegen einen Aeltesten Klage nur aus dem Munde zweier oder dreier Zeugen anzunehmen, langsam zu verfahren, weil sich an die Träger des geistlichen Amtes Lüge und Verleumdung, also auch ungerechte Anklage mehr als an andere anzuhängen pflegt.
\quellseite{287}
Doch allzulangen Beweis, allzulangsamen Prozeß, allzuspäte Sentenz sollte sich die Kirche auch nicht zu Schulden kommen lassen. Es steht gar zu viel auf dem Spiel, wenn ein Wolf in einer Gemeinde das Recht des Amtes in Anspruch nimmt. Dagegen sollten alle Pfarrer, die besseres Gewissen im Busen tragen, einmüthig und ernstlich Verwahrung einlegen. Ja, alle Besseren überhaupt sollten mit vereinter Macht des Zeugnisses sich gegen die Amtsverwaltung unverbesserlicher Irrlehrer und offenbarer, unbußfertiger Sünder aufmachen. Das sollte ein unverbrüchlicher Grundsatz sein: „Ausgedient hat, wer ein unverbesserlicher Irrlehrer, ein unbußfertiger, öffentlicher Sünder ist.“ — Es kommt oftmals vor, daß Pfarrer, welche alles Gute unterlassen und viel Böses thun, tief im Boden wurzeln und nicht zu beseitigen sind, während diejenigen, welche nur durch Uebermaß des Guten und Formenfehler sich das Mißfallen zuziehen, schnell weggeschafft sind, daß sie ihre Stätte nicht mehr kennt. Die Pfarrer, welche mit dem Regiment der Kirche betraut sind, entschuldigen sich mit dem Mangel an vernehmlicher Klage und Anklage gegen die ersteren und suchen im zweiten Fall Rechtfertigung in den dringenden, wirklich vorhandenen Klagen gegen diese. Es sollte wohl anders sein; denn die, welche im bischöflichen Amte sitzen, haben nicht auf Klage zu warten, sondern das eigene Auge aufzuthun; — und weil sie Aufseher, Meister sind, sollten sie auch wachen und Aufsicht haben und nach den Uebeln in der Kirche spähen. Wo so viel himmelschreiende Sünden herausfordern, die arme Heerde Christi von Wölfen zu entledigen und gegen
\quellseite{288}
sie zu verwahren; da sollten die „Augen der Kirche,“ d.\,i. ihre Hüter, mit scharfem Blicke die Uebel finden und richten. Aber wenn das nun nicht geschieht, sollten nicht fromme, tapfere Männer sich finden, die den Muth haben, Klage zu stellen? die im tiefen Frieden, in fester Ruh und Zuversicht der Seelen die Wölfe fassen, welche die arme Heerde angreifen und verzehren? Ist’s denn edel und großmüthig, den Frevler und sein Haus schonen, während das Haus Gottes, die Gemeinde Jesu, Schaden leidet? Gibt’s Keine mehr, die ein „emeritus“ rufen, wenn Hirten ihre Heerden nicht bloß verderben lassen, sondern selbst verderben? — — Was hilft alles Pastorale, wenn es durch Wort und Wandel reißender Wölfe ungestraft verhöhnt und damit für Tausende seiner Frucht beraubt wird? Das treue Thun frommer Knechte Gottes ist zum großen Theil wirkungslos, wenn die Wölfe straflos im Amte bleiben!

\quellseite{289}
Wie wir handeln und wandeln sollen im Hause Gottes, in der Gemeinde des lebendigen Gottes, das wissen wir nun wohl. — Du aber, o Herr, erbarme Dich unser und vergib uns Allen unsre Sünde und Missethat, damit wir wider Dich und Dein heiliges Wort und Amt gesündigt haben. Gib aber uns, Deinen Knechten allen, große Kraft und Stärk in uns Dein Leben, daß wir Alle, fertig und an Beinen gestiefelt, zu treiben das Evangelium des Friedens, dahin gehen in Geduld mit unserm Samenwurf und arbeiten und warten auf die köstliche Frucht der Erden!

Amen.

\quellseite{290}
\chapter{Anhang}

\quellseite{291}

\quellseite{292}
\section{Erinnerungen an Pfarrerfraun}

\begin{center}
(Einer jungen Pfarrerin im Jahre 1837 geschrieben.)
\end{center}

\subsection{1.}

Eine Pfarrersfrau erkenne ihren Beruf richtig. Nicht das Amt theilt sie mit ihrem Manne, sondern das Haus, wie jede andere Frau mit ihrem Manne; nicht zunächst im Amte hat sie zu helfen, sondern ihre Sorge sei, daß es im Hause und in der Haushaltung allezeit fein priesterlich hergehe; nicht mit Wort und Predigt, sondern mit heiligem Wandel hat sie ihren Mann in seinem Amte zu unterstützen. Es steht geschrieben 1. Petr. 3, 1. 2., daß die, welche nicht glauben an das Wort, durch der Weiber Wandel ohne Wort gewonnen werden sollen. (Vergl. 1. Tim. 2, 11. 12.)

\subsection{2.}

Es ist aber des Weibes Wandel, wenn er anders nach dem Sinn des HErrn sein soll, kein geringes und leichtes Werk.

— Ein Weib hat zunächst irdische Besorgungen im Hause, und eben darum ist es für sie schwer, durch’s Zeitliche so zu wandeln, daß sie das Ewige nicht verliere. Sie hat überdieß
\quellseite{293}
ihre Kinder zu versehen, und dabei selig zu werden nach 1. Tim. 2, 15. ist eine unmögliche Sache, wenn nicht der HErr hilft. Summa: ein Weib, das ihren Beruf erfüllt, ist ehrwürdig, wie der Mann, der seinen Beruf erfüllt; sie braucht also noch des Mannes Beruf nicht zu greifen.

\subsection{3.}

Ein Weib, welches in den Beruf des Mannes sich eindrängt, hat die Krone der Weiblichkeit verloren; denn ein rechtes Weib kann nicht männlich sein und mag es nicht. Weiblichkeit und Demuth gehen zusammen; es ist der größte Muth einer Frau, wenn sie Demuth, das ist doch auch „Muth zu dienen“ hat, — und welche in solchem Muth geduldig ist, scheint eine Gabe aus einem Reiche empfangen zu haben, in dem der vornehmste, der ewige König, ein Diener Aller geworden ist bis in den Tod. Eine Frau muß sich gerne regieren lassen, und unter dem Regimente ihres frommen Mannes zu stehen und zu gehen, muß sie für Glück halten, wie es die Kirche für ihre Seligkeit hält, unter dem Regimente und Geiste ihres ewigen Bräutigams zu leben. (Vergl. 1. Tim. 2, 12.)

\subsection{4.}

Ein Weib, welches sich über den Mann stellen und anstatt seiner seinen Beruf üben will (denn das heißt sich über den Mann stellen), verderbt, so viel an ihr liegt, die Seele, welcher sie als eine Gehilfin zum ewigen Leben geschaffen ist, nämlich die Seele ihres Mannes. Denn entweder weiß der Mann, daß es seine Demuth ist, im Namen des HErrn selber Mann zu sein; so ist ihm ein derartiges Thun seines Weibes
\quellseite{294}
eine böse Anfechtung und eine Ursache vieler Versündigungen gegen sie: — oder er beugt sich unter das Weib, dann ist das Weib schuldig an seiner großen Sünde, daß er nicht mehr des Herrn Bild ist, gleichwie das Weib der Gemeinde; man soll sie Delila heißen, weil sie ihres Simsons Seele matt gemacht und ihm sein Nasiräat gestohlen hat!

\subsection{5.}

Eine Pfarrersfrau ist ein heiliges Weib: sie wandelt in Demuth ihrem Manne nach — in Reinigkeit des Herzens und heiliger Ordnung des innwendigen Lebens. Wo kein reines Herz ist, da ist Gemeinheit — und wo ein Mann ein gemeines Weib hat, ist’s eben, als wenn man dem Vogel seine Flügel beschneidet. Und wo keine Ordnung im innwendigen Leben ist, wo man nicht fleißig in der Heilsordnung Gottes wandelt, da gibt’s bald Oede und Leere, eilendes Altern der Seelen, trostloses Seufzen, weil man keine dauernde Befriedigung gefunden und wohl gar seines Lebens Zweck verfehlt zu haben glaubt. Hat ein Pfarrer ein solches Weib, so hat er ein unglückliches Weib, und des Mannes größtes Unglück auf Erden scheint fast, wenn er ein unglückliches Weib hat.

\subsection{6.}

Eine Pfarrersfrau muß in ihrem eigenen Aeußern und in ihrem Haushalt Beweis liefern, daß innwendig Reinigkeit und Ordnung herrscht. Das ist Probe zum Exempel. Man muß ihr und ihrem Hause, man komme zu welcher Zeit oder Stunde man wolle, ansehen können, daß hier ein Diener einer reinen Welt wohne und das Kodesch la Jehova
\quellseite{295}
muß über dem Haushalt, aber nicht mit Buchstaben, leserlich geschrieben sein. Wenn die Pfarrersfrauen bedächten, wie viel Achtung ein priesterliches Wesen dem Pfarrer gewinnt, sie würden sich schon um dieses Grundes willen ein reines Herz voll himmlischer Ordnung erbitten.

\subsection{7.}

Eine Pfarrersfrau hat (soll haben) ein stilles Herz und einen schweigsamen Mund. Der Mann ist unglücklich, welcher ein geschwätziges und unruhiges Weib hat (1. Tim. 5, 13.); doppelt unglücklich, wenn sie jene verderblichen Laster an sich hat und es nicht weiß, noch erkennt! Ein solcher hat im Weibe seinem Amte einen Mühlstein angehängt, der es in ein Meer von Hindernissen und Verlegenheiten hinabzieht, aus dem es schwer auftauchen wird.

Bekanntlich trifft Gottes Wort die Herzen manchmal recht scharf und tief. Der Zuhörer, der getroffen wird, weiß nicht, wer ihn schlägt; von dem Schwerte Ebr. 4, 12. 13. weiß er nichts; er sieht bloß den Pfarrer um die Wunde an, die er im Gewissen empfangen hat — und urteilt: „Es muß ihm Jemand meine Sache verrathen haben!“ Je ernster und würdiger dann des Pfarrers Benehmen ist, so daß man ihm nicht zutraut, daß er die Schwätzer frage; desto mehr wird die Schuld auf die Pfarrersfrau geworfen, als von der er Etwas gehört haben könne, was sie von Schwätzerinnen empfangen. Wenn nun nicht der Pfarrerin stiller Wandel und schweigendes Wesen dagegen kämpft, oder sie gar einigen Anlaß für solche Geschwätze und Ausreden getroffener Seelen gibt; so ist sie schuldig, daß der Pfarrer als ein Weibbediener und Gottes zweischneidiges
\quellseite{296}
Schwert für Schwert, Spieß und Nägel des Verläumders angesehen wird, somit einen guten Theil seiner Kraft verliert. Eine Pfarrerin habe deswegen immer auf sich Acht, daß sie gewissermaßen über der Gemeinde und von ihr gesondert stehe, und betrage sich so, daß zwar Jedermann Lust zu ihr habe, namentlich die Elenden, daß aber Niemand auf den Gedanken kommen kann, sie habe in unwürdiger Vertraulichkeit mit irgend einem Gemeindebeamte, sei es nun Schullehrerin, Amme oder wer sonst verkehrt. — Eine Pfarrersfrau muß kurzum gewöhnliche Geselligkeit und Unterhaltung mehr, als andere Frauen entbehren können!

\subsection{8.}

Eine Pfarrersfrau darf sich in ihrem Urtheil über die Gemeindeglieder in keiner Weise bestechen lassen und keine Günstlinge haben. Wer etwa ihren Mann ihr in’s Angesicht, das ist in sein eigenes Angesicht hineinlobt, seine Predigten und Amtsverwaltung rühmt, seine Fehler zudeckt, Uebles über den Vorfahr und dessen Hausgenossen, über benachbarte Pfarrer u.\,s.\,w. redet, unmäßige Hoffnungen von ihres Pfarrers Amtswirksamkeit vorbringt; den sehe sie als verdächtig an und schaue auf, ob er nicht Etwas erschleichen wolle, was er ehrlicher Weise sich nicht getraut zu gewinnen. Einem solchen begegne sie mit allem Ernste und sage ihm, wiewohl mit aller Sanftmuth, doch rund und nett, daß ihr Mann bei ihr des Lobes nicht bedürfe, daß sie in seinem Namen sich schäme, sein Lob zu hören, daß sie sich vor Leuten der Art fürchte und sich ihr Geschwätze verbitten müsse.

\quellseite{297}
\subsection{9.}

Man hat oft gesehen, daß Pfarrersfrauen auf die Bequemlichkeit und Körperpflege ihrer Männer mehr, als auf die treue Ausübung ihrer amtlichen Pflichten gesehen haben. Da unterscheidet sich nun eine treue Magd des Herrn von der selbstsüchtigen Besitzerin eines Mannes dadurch, daß sie allezeit Wohlgefallen an seiner treuen Amtsverrichtung haben und äußern kann, ihn bei ihm anwandelnder Trägheit mahnt und zum Fleiße reizt und allewege mehr auf seine ewige Krone, als auf ein langes und dabei verwerfliches Leben sieht. Es muß ihr mehr werth sein, ihn dereinst leuchten zu sehen in des Himmels Glanz, wie die Sterne immer und ewiglich, als ihn unaufhörlich und nur recht lange an der Seite zu haben. Deshalb braucht aber ein Weib nicht ihres Mannes leibliche Erleichterung zu vernachlässigen; aber sie, wie alle Weiber soll einen Mann haben, welcher Erleichterung bedarf, weil er im Schweiße seines Angesichtes seinen Acker baut. Den Einfältigen gibt der Herr Gnade, hierin klar zu sehen und das rechte Maß zu halten — in dem Namen „Gehilfin“ liegt Beides: reizen zu guten Werken und Erleichterung schaffen.

\subsection{10.}

Eine Pfarrersfrau ist fleißig in ihrem Berufe, ist häuslich, schämt sich der Arbeit nicht, auch nicht der geringen. Sprüchw. 31, 10–34. ist ihr ein Spiegel, darin sie sich gern beschaut. Von ihrem Manne nach Gottes Befehl v. 28. 30. gelobt zu werden, ist ihr bei aller Demuth lieblich und vergönnt.
\quellseite{298}
Eine träge Pfarrersfrau schändet ihres Mannes Amt und macht ihn mehr als einmal auf und unter der Kanzel verstummen und erröthen. Ihr stiller Fleiß aber macht sein Angesicht heiter und erwirbt Dank; stärkt ihn auch, öffentlich und in den Häusern den Fleiß der Weiber zu preisen. — Eine Pfarrersfrau soll auch hierin dem nachleben, was ihr Mann predigt, und wenn er weibliche Tugend rühmt, sollen die Blicke der Gemeinde sich unwillkürlich auf die Pfarrerin mit Wohlgefallen richten, oder doch richten können.

\subsection{11.}

Eine Pfarrerin arbeitet gerne, thut gern auch das Niedrigste und Geringste, wenn es sein muß, aber ihr Sinn ist nicht niedrig; arbeitet sie, so thut sie es in Demuth und in freier Liebe, nicht wie gedungene Mägde, sondern daß man das adelige Wesen der Hausfrau eines Pfarrers erkenne. Sie arbeitet, aber sie läßt auch gerne arbeiten, daß nicht ihr Arbeiten für Geiz gehalten werde, sondern für Tugend, was es ja auch sein soll.

\subsection{12.}

Eine Pfarrerin ist gastfrei nach des Apostels Gebot 1. Petr. 4, 9. Ebr. 13, 2. — und ist’s eben danach ohne Murren; aber sie ist es nur in soweit, als des Ehegemahls Umstände es zulassen. Sie will mit ihrer Gastfreiheit nicht prunken, sondern die Einfalt, welche alle Tage in einem Pfarrhause herrscht, erhält sie auch, wenn Gäste kommen. Sie ist nicht gastfrei, um gesellschaftlich zu sein, sondern um Liebe zu üben. Sie ist es darum am liebsten gegen Dürftige und gegen
\quellseite{299}
Solche, welchen Andere keine Liebe erweisen. So thut sie und hat sammt ihrem Manne Freiheit der Seele genug, um über die hinweg zu schauen, die aus der Gastfreundschaft gerne ein weltliches Hin- und Herbesuchen machen möchten.

\subsection{13.}

Eine Pfarrerin ist sparsam in ihrem Haushalt und hält zu Rathe, was nur immer möglich ist; aber es ist ihr nicht darum zu thun, für den Kasten zu sparen; sondern sie spart, um nicht mehr auszugeben, als sie einnimmt, — und damit sie habe zu geben den Dürftigen.

Ein Pfarrer hat viel Versuchung geizig zu sein und zu werden; viele Pfarrer werden es auch, wie die Erfahrung beweist, und doch sollte es Niemand weniger werden, weil es Niemand für sich und sein Amt so ganz zum Schaden wird, als eben er. Leider sagt auch die Erfahrung, daß die Frauen in der Regel viel Anlaß, ja auch Anreizung zum Geiz der Männer geben. Ein frommes Weib aber kann nach Anleitung von Jacobi 1, 5. leicht weise werden — und ein weises Weib kann Beides, durch Sparen und Geben, für ihren Mann und sein Amt sehr wirksam werden.

\subsection{14.}

Eine fromme Pfarrersfrau ist eine dem Pfarrer durch Gottes Fügung zugewiesene Helferin für besondere Fälle in der Seelsorge der Weiber und Mädchen, so weit als es ihr Hauptberuf, die Verwaltung ihres Hauses zuläßt. Sie betet mit ihrem Manne täglich für die Gemeinde und thut hierin nicht weniger als ihr Mann selber — hilft auf diese stille verborgene Weise
\quellseite{300}
mächtig zum Amt. Sie begleitet auch, wenn ihr Mann zum Amte geht, in der Stille ihres Herzens seinen Gang mit Gebeten; sie bittet, daß sein Amt gelinge, — wenn es gelingt, bittet sie um Demuth für ihren Liebsten; wenn’s nicht gelingen will und der Herr seinen Segen zurückhält, so betet sie, daß ihr Mann nur treu erfunden werde und Gott zum Troste habe. Denn das Amt und der Segen ist des Herrn. Er kann segnen, welchen Er will — und prüft väterlich den Knecht, der keinen Segen schauen darf, auf daß er desto gewisser halte an der ewigen Gnade, die im Worte offenbart ist. — Die fromme Pfarrerin hilft aber auch sonst noch thätig zum Amte des Mannes. So wenig sie in den Häusern der Gemeindeglieder ohne Ursache, bloß der Unterhaltung wegen herumgeht, so gerne und fleißig sucht sie die Kranken, die Armen, die Witwen, die Waisen, die Betrübten, die Alten auf, — gießt in die Wunden Oel und Wein und ist eine Trostquelle der Elenden und ein Engel Gottes in ihren Hütten. Bei diesen Gelegenheiten geht ihr sonst schweigsamer Mund auf, und ihre Rede wirkt Segen.

\quellseite{301}
\section[Eine Stimme aus dem grauen Alterthum]{Eine Stimme aus dem grauen Alterthum über Pfarrersfrauen und Pfarrerskinder}

\begin{quote}
Τὰς γυναῖκας ἡμῶν ἐπὶ τὸ ἀγαθὸν διορθωσώμεθα. Τὸ ἀξιαγάπητον ἦθος τῆς ἁγνείας ἐνδειξάσθωσαν. Τὸ ἀκέραιον τῆς πραΰτητος αὐτῶν βούλημα ἀποδειξάτωσαν. Τὸ ἐπιεικὲς τῆς γλώσσης αὐτῶν διὰ τῆς σιγῆς φανερὸν ποιησάτωσαν, ἀγάπην αὐτῶν μὴ κατὰ προσκλίσεις ἀλλὰ πᾶσι τοῖς φοβουμένοις τὸν θεὸν ὁσίως ἴσην παρεχέτωσαν. Τὰ τέκνα ἡμῶν τῆς ἐν Χριστῷ παιδείας μεταλαβέτωσαν, μαθέτωσαν τί ταπεινοφροσύνη παρὰ θεῷ ἰσχύει, τί ἀγάπη ἀγνὴ παρὰ θεῷ δύναται, πῶς ὁ φόβος τοῦ κυρίου καλὸς καὶ μέγας, σώζων πάντας τοὺς ἐν αὐτῷ ὁσίως ἀναστρεφομένους ἐν καθαρᾷ καρδίᾳ.
\end{quote}

\quelle{Clemens Rom. bei Clem. Alex.}

\clearpage
\quellseite{302}
\begin{sideways}
\begin{minipage}{0.92\textheight}
\noindent Anhang zu §. 48 p. 97.
\begin{center}
{\large\bfseries Auszug aus einem Familien- oder Seelen-Register.}\\[0.8em]
Hausnummer 24.\\[0.3em]
{\large \gesperrt{Büntel}\footnote{Fingirte Familien- und Ortsnamen.}}\\[0.6em]
\footnotesize
\begin{tabular}{|r|r|c|c|c|c|p{3.6cm}|}
\hline
8. & \zl{4. und 7.\\ Maria Magdalena.\\ fil.} & \zl{Dahier\\ 1. April 1847.\\ baptiz. eod. die.} & & & & \\ \hline
9. & \zl{dito 2.\\ Maria Barbara.\\ fil.} & \zl{Dahier\\ 26. April 1848,\\ ¾ 10 U. VM.\\ bapt. eod. d.} & — & — & \zl{†\\ 11. Mai 1848.} & \\ \hline
10. & \zl{dito 3.\\ Johann Konrad.\\ fil.} & \zl{Dahier\\ 13. März 1849,\\ früh 5 U.\\ bapt. eod. d.} & — & — & \zl{†\\ 25. März 1849.} & \\ \hline
11. & \zl{dito 4.\\ Maria Barbara.\\ fil.} & \zl{Dahier\\ 13. Febr. 1850,\\ A. ½ 6 U.\\ bapt. 14 ejusd.} & & & & \\ \hline
12. & \zl{dito 5.\\ Johann Leonhard.\\ fil.} & \zl{Dahier\\ 1. Jul. 1851,\\ N. ¾ 12 U.\\ bapt. 2. ejusd.} & & & & \\ \hline
\end{tabular}
\end{center}
\end{minipage}
\end{sideways}

\clearpage
\quellseite{303}
\begin{sideways}
\begin{minipage}{0.92\textheight}
\begin{center}
Hausnummer 73.\\[0.3em]
{\large \gesperrt{Füll.}}\\[0.6em]
\footnotesize
\begin{tabular}{|r|c|c|c|c|c|p{3.6cm}|}
\hline
Nro. & Taufnamen. & \zl{Geburts- und\\ Tauftag.} & Copulationstag. & Stand. & Todestag. & \centering Bemerkungen. \tabularnewline \hline\hline
1. & Johann Michael. & \zl{Helpersdorf bei\\ Dietenburg\\ 4. Oktbr. 1806.} & \zl{Dahier\\ 24. Oktbr. 1841.} & Gütler. & \zl{†\\ 11. Jan. 1848,\\ A. ¾ 4 U.\\ Typhus.} & \scriptsize Eine redliche Christenseele. Er hatte den Bund mit Christo gemacht, daß ER Hausherr, Füll Hausknecht sein und allen seinen Bauernberuf nach Christi Sinn erfüllen wollte. \\ \hline
2. & \zl{Anna Apollonia,\\ geb. Stegerin.\\ uxor.} & \zl{Dahier\\ 28. Jun. 1820.} & Desgl. & \zl{Ledige\\ Gütlerstochter.} & & \\ \hline
3. & \zl{Erstes Kind.\\ Maria Barbara.\\ fil.} & \zl{Dahier\\ 14. Jul. 1842.\\ bapt. d. seq.} & & & & \\ \hline
4. & \zl{dito 2.\\ Johann Bartholo-\\ mäus.\\ fil.} & \zl{Dahier\\ 10. Jun. 1845.\\ bapt. d. seq.} & — & — & \zl{†\\ 12. Juni 1845.} & \\ \hline
\end{tabular}
\end{center}
\end{minipage}
\end{sideways}

\clearpage
\quellseite{304}
\begin{sideways}
\begin{minipage}{0.92\textheight}
\begin{center}
{\large\bfseries Auszug aus einem Beichtregister.}\\[0.8em]
Hausnummer 24.\\[0.3em]
{\large \gesperrt{Büntel.}}\\[0.3em]
Beichtjahre und Tage.\\[0.6em]
\footnotesize
\begin{tabular}{|r|c|c|c|c|c|c|}
\hline
Nro. & Namen. & 1850. & 1851. & 1852. & 1853. & 1854. \\ \hline\hline
1. & \zl{Büntel,\\ Johann Michael,\\ geb. 17/7/11.} & \zl{Invoc.\\ D. Virid.\\ D. F. Ascens.\\ DD. F. Pentec.\\ p. Tr. XXIII.\\ D. F. Nat.} & \zl{Sept.\\ D. Virid.\\ DD. F. Pentec.\\ p. Tr. VI.\\ p. Tr. XXIII.} & \zl{Ocul.\\ Virid.\\ Ascens.\\ Pentec.} & & \\ \hline
2. & \zl{Ursula Barbara,\\ geb. Hollmännin.\\ geb. 13/6/21.\\ uxor.} & \zl{Invoc.\\ D. Virid.\\ D. F. Ascens.\\ —\\ p. Tr. XXIII.\\ D. F. Nat.} & \zl{Sept.\\ D. Virid.\\ DD. F. Pentec.\\ p. Tr. VI.\\ p. Tr. XXIII.} & \zl{Ocul.\\ Virid.\\ Ascens.\\ Pentec.} & & \\ \hline
3. & \zl{Maria Magdalena\\ Büntelin,\\ geb. 9/4/19.\\ soror.} & \zl{Laet.\\ DD. F. Pentec.\\ p. Tr. XXII.\\ D. F. Nat.} & \zl{Invoc.\\[1em] Nach N., Pfarrei N.,\\ abgezogen.} & & & \\ \hline
 & & & \zl{DD. F. Pasch.\\ M. B. Meierin,\\ Dienstmagd.\\ Leonh. Schroll,\\ Knecht.\\ NB. Beichtscheine\\ vorgezeigt.} & & & \\ \hline
\end{tabular}
\end{center}
\end{minipage}
\end{sideways}

\clearpage
\quellseite{305}
\begin{unterschrift}
Druck von G. Buchbinder in Neu-Ruppin.
\end{unterschrift}

\quellseite{306}

\quellseite{307}

\quellseite{308}

\quellseite{309}
